\documentclass[11pt,a4paper]{article}

% ============================================================
% Packages
% ============================================================

\usepackage[utf8]{inputenc}
\usepackage[T1]{fontenc}
\usepackage{lmodern}

\usepackage{amsmath}
\usepackage{amssymb}
\usepackage{mathtools}
\usepackage{bbm}

\usepackage{graphicx}
\usepackage{booktabs}
\usepackage{array}
\usepackage{enumitem}

\usepackage[
    margin=2.5cm
]{geometry}

\usepackage{setspace}
\onehalfspacing

\usepackage{hyperref}
\hypersetup{
    colorlinks=true,
    linkcolor=blue,
    citecolor=blue,
    urlcolor=blue
}

% ============================================================
% Document Information
% ============================================================

\title{\textbf{Assessing Quantitative Investment Strategies: The Backtester as a Scientific Protocol}\\
\large A Technical Framework for Quantitative Strategy Backtesting}

\author{Matéo Molinaro}

\date{August 2026}

% ============================================================
% Document
% ============================================================

\begin{document}

\maketitle

\tableofcontents

\newpage


% ============================================================
\part{Backtesting Framework and Research Design}
% ============================================================


% ============================================================
\section{Problem Definition and Backtest Conventions}
\label{sec:problem_definition}
% ============================================================

A backtest is a numerical simulation of an investment strategy using historical
data. Its objective is to estimate how a systematic investment process would
have behaved if it had been implemented historically under a precisely defined
set of information, portfolio construction, execution, and transaction-cost
assumptions.

A quantitative investment strategy can be represented as a sequence of
transformations:

\begin{equation}
\boxed{
\mathcal{D}_t
\longrightarrow
X_t
\longrightarrow
s_t
\longrightarrow
\widetilde{s}_t
\longrightarrow
w_t^{\mathrm{target}}
\longrightarrow
w_t
\longrightarrow
R_{p,t\rightarrow t+1}
}
\end{equation}

where:

\begin{itemize}
    \item $\mathcal{D}_t$ denotes the raw data available to the investor at time $t$;

    \item $X_t$ denotes the set of features constructed from these data;

    \item $s_t$ denotes the raw investment signal;

    \item $\widetilde{s}_t$ denotes the final processed signal after operations
    such as standardization, ranking, combination, smoothing, or neutralization;

    \item $w_t^{\mathrm{target}}$ denotes the target portfolio produced by the
    portfolio-construction process;

    \item $w_t$ denotes the portfolio that is actually implemented after
    accounting for constraints, execution assumptions, and market frictions;

    \item $R_{p,t\rightarrow t+1}$ denotes the subsequent realized portfolio
    return over the holding period.
\end{itemize}

The central principle underlying the entire backtesting framework is that every
investment decision made at time $t$ must depend exclusively on information
that would have been available to the investor at that time.

Let $\mathcal{I}_t$ denote the investor's information set at time $t$. The
portfolio decision must satisfy

\begin{equation}
\boxed{
w_t = f(\mathcal{I}_t)
}
\end{equation}

for some portfolio-construction function $f$.

More formally, every input used in the portfolio decision must be
$\mathcal{I}_t$-measurable. In particular,

\begin{equation}
X_t,
\qquad
s_t,
\qquad
\widetilde{s}_t,
\qquad
w_t
\end{equation}

must be functions exclusively of information contained in $\mathcal{I}_t$.

Future information must therefore not enter the decision rule:

\begin{equation}
\mathcal{I}_t
\cap
\left\{
\text{information revealed strictly after } t
\right\}
=
\varnothing.
\end{equation}

This apparently simple condition is one of the most important requirements of
a valid backtest. Violations generate look-ahead bias and can substantially
overstate the historical performance of an investment strategy.

Before constructing signals or portfolios, the investment problem and the
backtesting conventions must therefore be explicitly defined.


% ------------------------------------------------------------
\subsection{Investment Objective}
\label{subsec:investment_objective}
% ------------------------------------------------------------

The first step consists of defining what the investment strategy is designed
to achieve.

Depending on the application, the objective may be to:

\begin{itemize}
    \item maximize absolute returns;

    \item maximize returns relative to a benchmark;

    \item maximize risk-adjusted returns;

    \item generate market-neutral alpha;

    \item minimize portfolio risk;

    \item track a benchmark while generating incremental alpha;

    \item maximize expected utility;

    \item maximize expected returns subject to risk, exposure, liquidity,
    turnover, or implementation constraints.
\end{itemize}

At a general level, portfolio construction can be represented as the
optimization problem

\begin{equation}
w_t^*
=
\underset{w}{\arg\max}
\;
\mathcal{J}
\left(
w;
\mu_t,
\Sigma_t,
w_{t^-},
\mathcal{C}_t
\right),
\end{equation}

where:

\begin{itemize}
    \item $\mu_t \in \mathbb{R}^{N_t}$ is the vector of expected asset returns;

    \item $\Sigma_t \in \mathbb{R}^{N_t \times N_t}$ is the estimated covariance
    matrix of asset returns;

    \item $w_{t^-}$ denotes the portfolio immediately before rebalancing at
    time $t$;

    \item $\mathcal{C}_t$ represents portfolio constraints and market frictions;

    \item $N_t$ denotes the number of investable securities at time $t$.
\end{itemize}

A common mean--variance specification augmented with implementation costs is

\begin{equation}
\boxed{
w_t^*
=
\underset{w}{\arg\max}
\left[
\mu_t^\top w
-
\frac{\lambda}{2}
w^\top \Sigma_t w
-
C_t(w,w_{t^-})
\right]
}
\end{equation}

subject to

\begin{equation}
w \in \mathcal{W}_t,
\end{equation}

where $\lambda>0$ is a risk-aversion parameter,
$C_t(w,w_{t^-})$ represents expected implementation costs, and
$\mathcal{W}_t$ denotes the feasible set of portfolios.

Importantly, not every systematic investment strategy requires explicit
optimization. A strategy may instead directly transform cross-sectional signal
values into portfolio weights.

For example,

\begin{equation}
w_{i,t}
=
\frac{s_{i,t}}
{\sum_{j=1}^{N_t}|s_{j,t}|}.
\end{equation}

Portfolio optimization should therefore be viewed as one possible
signal-to-weight mapping rather than as a necessary component of every
backtesting framework.


% ------------------------------------------------------------
\subsection{Investment Universe and Asset Class}
\label{subsec:investment_universe}
% ------------------------------------------------------------

Let

\begin{equation}
\mathcal{U}_t
=
\left\{
1,\ldots,N_t
\right\}
\end{equation}

denote the investment universe at time $t$.

The universe may vary through time:

\begin{equation}
\mathcal{U}_t
\neq
\mathcal{U}_{t+1}.
\end{equation}

This point is fundamental. Using the current investment universe throughout
historical periods creates survivorship bias because securities that disappeared
through bankruptcy, delisting, acquisition, or other corporate events are
incorrectly excluded from historical observations.

The universe definition should specify at least:

\begin{itemize}
    \item asset class;

    \item geographical coverage;

    \item eligible exchanges;

    \item eligible security types;

    \item minimum market capitalization;

    \item minimum liquidity;

    \item minimum security price;

    \item minimum trading history;

    \item data-availability requirements;

    \item shortability requirements for short portfolios.
\end{itemize}

The eligibility of security $i$ at time $t$ can be represented by the indicator

\begin{equation}
I_{i,t}^{\mathrm{eligible}}
=
\begin{cases}
1,
&
\text{if security } i
\text{ satisfies all eligibility conditions at } t,
\\[0.2cm]
0,
&
\text{otherwise}.
\end{cases}
\end{equation}

The dynamic investment universe is therefore

\begin{equation}
\boxed{
\mathcal{U}_t
=
\left\{
i:
I_{i,t}^{\mathrm{eligible}}=1
\right\}.
}
\end{equation}

Importantly, the eligibility rule must itself respect the information set
$\mathcal{I}_t$. For example, a future delisting cannot be used retrospectively
to exclude a stock from today's investment universe.


% ------------------------------------------------------------
\subsection{Prediction Target and Investment Horizon}
\label{subsec:prediction_target}
% ------------------------------------------------------------

For predictive strategies, the target variable must be precisely defined.

Let $P_{i,t}$ denote the price of asset $i$ at time $t$. The simple return
between $t$ and $t+h$ is

\begin{equation}
R_{i,t\rightarrow t+h}
=
\frac{P_{i,t+h}}{P_{i,t}}-1.
\end{equation}

The corresponding logarithmic return is

\begin{equation}
r_{i,t\rightarrow t+h}
=
\log(P_{i,t+h})
-
\log(P_{i,t}).
\end{equation}

A predictive model may therefore estimate

\begin{equation}
\widehat{R}_{i,t\rightarrow t+h}
=
f(X_{i,t}),
\end{equation}

where $X_{i,t}$ contains exclusively information available at time $t$.

The target does not necessarily have to be the future raw return. Possible
targets include:

\begin{itemize}
    \item future absolute returns;

    \item future excess returns;

    \item benchmark-relative returns;

    \item market- or factor-adjusted returns;

    \item cross-sectional return ranks;

    \item return signs;

    \item probabilities of outperforming;

    \item future volatility;

    \item future risk-adjusted returns.
\end{itemize}

For example, a cross-sectional model may directly predict a relative return:

\begin{equation}
Y_{i,t+h}
=
R_{i,t\rightarrow t+h}
-
R_{B,t\rightarrow t+h},
\end{equation}

where $R_{B,t\rightarrow t+h}$ denotes the benchmark return over the same
horizon.

Alternatively, the target may be a cross-sectional rank:

\begin{equation}
Y_{i,t+h}
=
\operatorname{Rank}
\left(
R_{i,t\rightarrow t+h}
\right).
\end{equation}

The prediction horizon $h$ should be economically consistent with the expected
persistence of the underlying signal and the intended holding period.

However, the prediction horizon and the holding period need not be identical.
This distinction becomes particularly important when overlapping forecasts or
staggered portfolios are employed.


% ------------------------------------------------------------
\subsection{Rebalancing Frequency and Holding Period}
\label{subsec:rebalancing}
% ------------------------------------------------------------

Let

\begin{equation}
\mathcal{T}^{\mathrm{reb}}
=
\left\{
t_0,t_1,\ldots,t_K
\right\}
\end{equation}

denote the set of portfolio rebalancing dates.

Typical rebalancing frequencies include:

\begin{itemize}
    \item daily;

    \item weekly;

    \item monthly;

    \item quarterly.
\end{itemize}

At each rebalancing date $t_k$, the strategy determines a new target portfolio

\begin{equation}
w_{t_k}^{\mathrm{target}}.
\end{equation}

If no trading occurs between $t_k$ and $t_{k+1}$, the \emph{number of shares}
held remains unchanged. The portfolio weights, however, generally drift because
individual securities experience different returns.

If $w_{i,t_k}$ denotes the implemented weight immediately after rebalancing,
the weight immediately before the next rebalance satisfies

\begin{equation}
w_{i,t_{k+1}^{-}}
=
\frac{
w_{i,t_k}
\left(
1+R_{i,t_k\rightarrow t_{k+1}}
\right)
}{
1+
R_{p,t_k\rightarrow t_{k+1}}
},
\end{equation}

in the simplified case of a fully invested long-only portfolio without cash
flows or intermediate trading.

Thus,

\begin{equation}
w_{t_{k+1}^{-}}
\neq
w_{t_k}
\end{equation}

in general.

It is therefore important to distinguish:

\begin{equation}
\boxed{
w_{t_k}
\quad
\text{from}
\quad
w_{t_{k+1}^{-}}.
}
\end{equation}

The former denotes the portfolio immediately after rebalancing, while the
latter denotes the portfolio obtained after market-driven weight drift and
immediately before the next rebalance.

The rebalancing frequency and the prediction horizon are related but
conceptually distinct.

For example, a model may predict one-month returns while the portfolio is
rebalanced weekly. Successive predictions then correspond to overlapping
forecast horizons.

This feature has implications for:

\begin{itemize}
    \item portfolio construction;

    \item turnover;

    \item strategy capacity;

    \item autocorrelation of portfolio returns;

    \item statistical inference.
\end{itemize}


% ------------------------------------------------------------
\subsection{Long-Only and Long--Short Strategies}
\label{subsec:long_short}
% ------------------------------------------------------------

Let

\begin{equation}
w_t
=
\begin{bmatrix}
w_{1,t}
&
\cdots
&
w_{N_t,t}
\end{bmatrix}^{\top}
\in
\mathbb{R}^{N_t}
\end{equation}

denote the portfolio-weight vector.

For a long-only portfolio,

\begin{equation}
w_{i,t}\geq 0
\qquad
\forall i.
\end{equation}

A fully invested long-only portfolio generally satisfies

\begin{equation}
\sum_{i=1}^{N_t}
w_{i,t}
=
1.
\end{equation}

For a long--short portfolio,

\begin{equation}
w_{i,t}
\in
\mathbb{R},
\end{equation}

where positive weights represent long positions and negative weights represent
short positions.

Define the total long exposure as

\begin{equation}
L_t
=
\sum_{i=1}^{N_t}
\max(w_{i,t},0),
\end{equation}

and the absolute short exposure as

\begin{equation}
S_t
=
\sum_{i=1}^{N_t}
\max(-w_{i,t},0).
\end{equation}

The gross exposure is

\begin{equation}
\boxed{
G_t
=
L_t+S_t
=
\sum_{i=1}^{N_t}
|w_{i,t}|
}
\end{equation}

while the net exposure is

\begin{equation}
\boxed{
E_t^{\mathrm{net}}
=
L_t-S_t
=
\sum_{i=1}^{N_t}
w_{i,t}.
}
\end{equation}

The notation $E_t^{\mathrm{net}}$ is used rather than $N_t$ to avoid confusion
with $N_t$, the number of securities in the investment universe.

For example, a dollar-neutral portfolio with $100\%$ long exposure and
$100\%$ short exposure satisfies

\begin{equation}
L_t=1,
\qquad
S_t=1.
\end{equation}

Consequently,

\begin{equation}
G_t=2,
\qquad
E_t^{\mathrm{net}}=0.
\end{equation}

The portfolio therefore has $200\%$ gross exposure while remaining dollar
neutral.


% ------------------------------------------------------------
\subsection{Benchmark Definition}
\label{subsec:benchmark}
% ------------------------------------------------------------

Let $R_{B,t}$ denote the return of the strategy benchmark.

The appropriate benchmark depends on the investment objective.

For an absolute-return or market-neutral strategy, the relevant benchmark may
be cash or the risk-free rate.

For a benchmark-relative strategy, the active return is

\begin{equation}
R_t^{\mathrm{active}}
=
R_{p,t}
-
R_{B,t}.
\end{equation}

The active portfolio weights can similarly be defined as

\begin{equation}
w_t^{\mathrm{active}}
=
w_t-w_t^B,
\end{equation}

where $w_t^B$ denotes the benchmark portfolio weights.

For a fully invested portfolio and benchmark,

\begin{equation}
\mathbf{1}^{\top}w_t
=
1
\end{equation}

and

\begin{equation}
\mathbf{1}^{\top}w_t^B
=
1.
\end{equation}

Therefore,

\begin{equation}
\mathbf{1}^{\top}
w_t^{\mathrm{active}}
=
0.
\end{equation}

This distinction becomes particularly important when analyzing:

\begin{itemize}
    \item excess returns;

    \item tracking error;

    \item information ratios;

    \item active risk;

    \item active factor exposures;

    \item performance attribution.
\end{itemize}


% ------------------------------------------------------------
\subsection{Information Sets and Timing Conventions}
\label{subsec:timing}
% ------------------------------------------------------------

Timing conventions constitute one of the fundamental components of a valid
backtest.

Let

\begin{equation}
\mathcal{I}_t
\end{equation}

denote all information that could realistically have been observed by the
investor at decision time $t$.

Every quantity used to construct the portfolio must be measurable with respect
to $\mathcal{I}_t$.

Therefore,

\begin{equation}
X_t,
\qquad
s_t,
\qquad
\widetilde{s}_t,
\qquad
w_t
\end{equation}

must all be $\mathcal{I}_t$-measurable.

By contrast, future returns are not observable at the decision date:

\begin{equation}
R_{t\rightarrow t+1}
\notin
\mathcal{I}_t.
\end{equation}

Hence, the fundamental backtesting sequence is

\begin{equation}
\boxed{
\mathcal{I}_t
\longrightarrow
X_t
\longrightarrow
s_t
\longrightarrow
w_t
\longrightarrow
R_{t\rightarrow t+1}.
}
\end{equation}

More generally, it is useful to distinguish four timestamps:

\begin{equation}
t^{\mathrm{data}}
\leq
t^{\mathrm{signal}}
\leq
t^{\mathrm{decision}}
\leq
t^{\mathrm{execution}}.
\end{equation}

These correspond respectively to:

\begin{itemize}
    \item the time at which the latest required data become available;

    \item the time at which the signal can be computed;

    \item the time at which the portfolio decision is made;

    \item the time at which the corresponding trade is executed.
\end{itemize}

For example, suppose that a strategy uses the closing price on trading day $t$
to compute a momentum signal.

The signal cannot be known until the closing price itself has been observed.
Therefore, assuming both

\begin{equation}
s_t
=
f(P_t^{\mathrm{close}})
\end{equation}

and execution at exactly $P_t^{\mathrm{close}}$ may introduce look-ahead bias
unless the trading mechanism explicitly makes such execution feasible.

A conservative timing convention is

\begin{equation}
\boxed{
P_t^{\mathrm{close}}
\rightarrow
s_t
\rightarrow
w_t^{\mathrm{target}}
\rightarrow
\text{execution at } t+1.
}
\end{equation}

The exact convention depends on the trading process, data timestamp, order type,
market, and assumed execution mechanism.

The same principle applies to fundamental data.

Suppose a firm's annual accounting variable corresponds to the fiscal year
ending on December 31 but the financial statement is released on March 15 of
the following year.

The accounting information cannot enter $\mathcal{I}_t$ before March 15.

Thus,

\begin{equation}
\boxed{
\text{economic reference date}
\neq
\text{information availability date}.
}
\end{equation}

For backtesting purposes, the relevant date is generally the date at which the
information became observable by market participants.


% ------------------------------------------------------------
\subsection{Execution Assumptions}
\label{subsec:execution}
% ------------------------------------------------------------

A backtest must define the price at which hypothetical transactions are assumed
to occur.

Common execution assumptions include:

\begin{itemize}
    \item closing price;

    \item next-day opening price;

    \item Volume-Weighted Average Price (VWAP);

    \item Time-Weighted Average Price (TWAP);

    \item market mid-price plus an estimated execution cost.
\end{itemize}

Let

\begin{equation}
P_{i,t}^{\mathrm{exec}}
\end{equation}

denote the assumed execution price of security $i$.

If $\Delta q_{i,t}$ shares are traded at time $t$, the corresponding traded
notional is

\begin{equation}
V_{i,t}^{\mathrm{trade}}
=
\Delta q_{i,t}
P_{i,t}^{\mathrm{exec}}.
\end{equation}

Execution assumptions must be compatible with the information timing
convention.

In particular, an execution price that was already determined before the
signal could have been computed cannot generally be used.

The backtest should therefore distinguish explicitly between

\begin{equation}
\boxed{
\text{signal price}
\neq
\text{execution price}
\neq
\text{valuation price},
}
\end{equation}

unless the investment process provides a valid economic and operational reason
for these prices to coincide.

Execution assumptions become increasingly important as turnover increases and
the strategy's investment horizon decreases.


% ------------------------------------------------------------
\subsection{Gross Exposure, Net Exposure, and Leverage}
\label{subsec:gross_net_leverage}
% ------------------------------------------------------------

Gross and net exposures characterize different dimensions of portfolio
positioning.

Recall that

\begin{equation}
G_t
=
\sum_{i=1}^{N_t}
|w_{i,t}|
\end{equation}

and

\begin{equation}
E_t^{\mathrm{net}}
=
\sum_{i=1}^{N_t}
w_{i,t}.
\end{equation}

A portfolio can therefore have low net exposure while simultaneously carrying
substantial gross exposure.

For example,

\begin{equation}
L_t=1.5,
\qquad
S_t=1.5
\end{equation}

implies

\begin{equation}
G_t=3,
\qquad
E_t^{\mathrm{net}}=0.
\end{equation}

The portfolio is dollar neutral but carries $300\%$ gross exposure.

This illustrates the important distinction

\begin{equation}
\boxed{
\text{Dollar neutrality}
\neq
\text{Low risk}
\neq
\text{Beta neutrality}.
}
\end{equation}

Suppose

\begin{equation}
\beta_t
=
\begin{bmatrix}
\beta_{1,t}\\
\vdots\\
\beta_{N_t,t}
\end{bmatrix}
\end{equation}

denotes the vector of individual security betas relative to a market factor.

Under the standard linear exposure approximation, the portfolio beta is

\begin{equation}
\boxed{
\beta_{p,t}
=
w_t^\top\beta_t.
}
\end{equation}

Dollar neutrality requires

\begin{equation}
\mathbf{1}^\top w_t
=
0,
\end{equation}

whereas beta neutrality requires

\begin{equation}
\beta_t^\top w_t
=
0.
\end{equation}

In general,

\begin{equation}
\boxed{
\mathbf{1}^\top w_t=0
\;\not\Rightarrow\;
\beta_t^\top w_t=0.
}
\end{equation}

A dollar-neutral portfolio may therefore retain significant systematic market
risk.

Gross exposure, net exposure, beta exposure, factor exposures, and ex-ante
portfolio volatility should consequently be treated as distinct quantities
throughout portfolio construction and risk management.


% ------------------------------------------------------------
\subsection{Risk and Volatility Targets}
\label{subsec:volatility_target}
% ------------------------------------------------------------

Many systematic strategies target a predefined level of portfolio volatility.

Let

\begin{equation}
\widehat{\sigma}_{p,t}
\end{equation}

denote the ex-ante estimate of annualized portfolio volatility available at
time $t$, and let

\begin{equation}
\sigma^{\mathrm{target}}
\end{equation}

denote the desired volatility target.

A simple volatility-scaling coefficient is

\begin{equation}
\lambda_t
=
\frac{
\sigma^{\mathrm{target}}
}{
\widehat{\sigma}_{p,t}
}.
\end{equation}

The scaled portfolio is then

\begin{equation}
\boxed{
w_t^{\mathrm{scaled}}
=
\lambda_t w_t.
}
\end{equation}

If the portfolio volatility is estimated using a covariance matrix,

\begin{equation}
\widehat{\sigma}_{p,t}
=
\sqrt{
w_t^\top
\widehat{\Sigma}_t
w_t
},
\end{equation}

where $\widehat{\Sigma}_t$ is an ex-ante covariance estimate constructed using
information available at time $t$.

If $\widehat{\Sigma}_t$ represents a daily covariance matrix, annualized
volatility can be written as

\begin{equation}
\widehat{\sigma}_{p,t}^{\mathrm{ann}}
=
\sqrt{
252
w_t^\top
\widehat{\Sigma}_t
w_t
}.
\end{equation}

Alternatively, volatility may be estimated directly from historical portfolio
returns.

For example,

\begin{equation}
\widehat{\sigma}_{p,t}
=
g
\left(
R_{p,t},
R_{p,t-1},
\ldots,
R_{p,t-L+1}
\right).
\end{equation}

The estimator must satisfy the information constraint and must therefore not
use future portfolio returns.

In practice, leverage caps are frequently combined with volatility targeting:

\begin{equation}
\boxed{
\lambda_t
=
\min
\left(
\frac{
\sigma^{\mathrm{target}}
}{
\widehat{\sigma}_{p,t}
},
\lambda_{\max}
\right).
}
\end{equation}

This prevents the strategy from taking extreme leverage when estimated
volatility becomes unusually low.

Additional smoothing or bounds may also be imposed on changes in leverage in
order to avoid excessive turnover.


% ------------------------------------------------------------
\subsection{Transaction Cost Assumptions}
\label{subsec:transaction_costs}
% ------------------------------------------------------------

A realistic backtest must account for the costs associated with implementing
portfolio decisions. A strategy may exhibit attractive gross performance while
becoming economically unattractive once trading costs, market impact, financing,
and short-borrowing costs are taken into account.

It is useful to distinguish between two broad categories of implementation
costs:

\begin{enumerate}
    \item \textbf{Trading costs}, which are generated when positions are bought
    or sold;

    \item \textbf{Holding and financing costs}, which arise from maintaining
    positions through time.
\end{enumerate}

The distinction is important because trading costs are primarily related to
changes in portfolio positions, whereas financing and borrowing costs depend
primarily on the positions held and the duration for which they are maintained.


% ------------------------------------------------------------
\subsubsection{Portfolio Trades and Weight Changes}
% ------------------------------------------------------------

Let

\begin{equation}
w_{t^-}
\end{equation}

denote the vector of portfolio weights immediately before the rebalance at time
$t$, after accounting for the drift of the existing portfolio, and let

\begin{equation}
w_t
\end{equation}

denote the implemented portfolio weights immediately after the rebalance.

The required change in portfolio weights is therefore

\begin{equation}
\boxed{
\Delta w_t
=
w_t-w_{t^-}.
}
\end{equation}

At the security level,

\begin{equation}
\Delta w_{i,t}
=
w_{i,t}-w_{i,t^-}.
\end{equation}

A positive value,

\begin{equation}
\Delta w_{i,t}>0,
\end{equation}

corresponds to a net purchase of security $i$, whereas

\begin{equation}
\Delta w_{i,t}<0
\end{equation}

corresponds to a net sale.

Define the purchased and sold notionals, expressed as fractions of portfolio
net asset value, as

\begin{equation}
B_t
=
\sum_{i=1}^{N_t}
\max(\Delta w_{i,t},0),
\end{equation}

and

\begin{equation}
S_t^{\mathrm{trade}}
=
\sum_{i=1}^{N_t}
\max(-\Delta w_{i,t},0).
\end{equation}

The notation $S_t^{\mathrm{trade}}$ is used here to distinguish traded sales
from the portfolio's short exposure.

The total traded notional is therefore

\begin{equation}
\boxed{
T_t
=
B_t
+
S_t^{\mathrm{trade}}
=
\sum_{i=1}^{N_t}
|\Delta w_{i,t}|.
}
\end{equation}


% ------------------------------------------------------------
\subsubsection{One-Way Turnover and Total Traded Notional}
% ------------------------------------------------------------

Turnover is not uniquely defined in the literature or in industry practice.
Two related conventions are commonly encountered.

The first reports the \emph{total traded notional}:

\begin{equation}
\boxed{
TO_t^{\mathrm{gross}}
=
\sum_{i=1}^{N_t}
|\Delta w_{i,t}|.
}
\end{equation}

This quantity counts both purchases and sales.

A second convention reports \emph{one-way turnover}:

\begin{equation}
\boxed{
TO_t^{\mathrm{one\text{-}way}}
=
\frac{1}{2}
\sum_{i=1}^{N_t}
|\Delta w_{i,t}|.
}
\end{equation}

The origin of the factor $1/2$ can be seen directly for a fully invested
portfolio.

Suppose that both the pre-rebalance and post-rebalance portfolios satisfy

\begin{equation}
\mathbf{1}^{\top}w_{t^-}
=
\mathbf{1}^{\top}w_t
=
1.
\end{equation}

Then

\begin{equation}
\sum_{i=1}^{N_t}\Delta w_{i,t}
=
\sum_{i=1}^{N_t}
\left(
w_{i,t}-w_{i,t^-}
\right)
=
0.
\end{equation}

Consequently, total purchases must equal total sales:

\begin{equation}
B_t
=
S_t^{\mathrm{trade}}.
\end{equation}

Since

\begin{equation}
\sum_{i=1}^{N_t}
|\Delta w_{i,t}|
=
B_t+S_t^{\mathrm{trade}},
\end{equation}

it follows that

\begin{equation}
\sum_{i=1}^{N_t}
|\Delta w_{i,t}|
=
2B_t
=
2S_t^{\mathrm{trade}}.
\end{equation}

Therefore,

\begin{equation}
\boxed{
\frac{1}{2}
\sum_{i=1}^{N_t}
|\Delta w_{i,t}|
=
B_t
=
S_t^{\mathrm{trade}}.
}
\end{equation}

The one-way turnover convention can thus be interpreted as the fraction of the
portfolio that has been replaced, whereas total traded notional measures the
sum of all purchases and sales.

For example, suppose that $50\%$ of a portfolio is sold and the proceeds are
used to purchase new securities. Then

\begin{equation}
B_t=0.5,
\qquad
S_t^{\mathrm{trade}}=0.5.
\end{equation}

Hence,

\begin{equation}
TO_t^{\mathrm{gross}}
=
0.5+0.5
=
1,
\end{equation}

whereas

\begin{equation}
TO_t^{\mathrm{one\text{-}way}}
=
\frac{1}{2}(1)
=
0.5.
\end{equation}

The strategy can therefore be described as having replaced $50\%$ of its
portfolio while having traded an aggregate notional equal to $100\%$ of its
net asset value.

The factor $1/2$ is consequently a \emph{turnover reporting convention}. It
does not imply that only one side of a transaction generates implementation
costs.


% ------------------------------------------------------------
\subsubsection{Proportional Transaction Costs}
% ------------------------------------------------------------

Both purchases and sales generally generate transaction costs. Therefore, if
$c$ denotes the proportional cost per unit of traded notional and costs are
assumed to be symmetric between purchases and sales, transaction costs are

\begin{equation}
\boxed{
TC_t
=
c
\sum_{i=1}^{N_t}
|\Delta w_{i,t}|.
}
\end{equation}

Equivalently,

\begin{equation}
TC_t
=
c
\left(
B_t+S_t^{\mathrm{trade}}
\right).
\end{equation}

The factor $1/2$ should therefore \emph{not} be introduced into this expression
when $c$ represents the cost per unit of actual traded notional.

For example, suppose that a portfolio sells $50\%$ of its NAV and purchases
another $50\%$, such that

\begin{equation}
\sum_i |\Delta w_{i,t}|=1.
\end{equation}

If the assumed transaction cost is $10$ basis points per unit of traded
notional,

\begin{equation}
c=0.001,
\end{equation}

then

\begin{equation}
TC_t
=
0.001\times 1
=
0.001,
\end{equation}

corresponding to a $10$ basis-point reduction in portfolio NAV.

Using one-way turnover without adjusting the cost convention would instead
produce

\begin{equation}
0.001\times 0.5
=
0.0005,
\end{equation}

or only $5$ basis points, thereby understating the true assumed trading cost.

If transaction costs vary across securities, the model becomes

\begin{equation}
\boxed{
TC_t
=
\sum_{i=1}^{N_t}
c_{i,t}
|\Delta w_{i,t}|.
}
\end{equation}


% ------------------------------------------------------------
\subsubsection{Asymmetric Buy and Sell Costs}
% ------------------------------------------------------------

Transaction costs need not be symmetric between purchases and sales.

Let

\begin{equation}
c_{i,t}^{\mathrm{buy}}
\end{equation}

and

\begin{equation}
c_{i,t}^{\mathrm{sell}}
\end{equation}

denote the proportional costs associated respectively with buying and selling
security $i$.

A more general transaction-cost model is then

\begin{equation}
\boxed{
TC_t
=
\sum_{i=1}^{N_t}
\left[
c_{i,t}^{\mathrm{buy}}
\max(\Delta w_{i,t},0)
+
c_{i,t}^{\mathrm{sell}}
\max(-\Delta w_{i,t},0)
\right].
}
\end{equation}

If

\begin{equation}
c_{i,t}^{\mathrm{buy}}
=
c_{i,t}^{\mathrm{sell}}
=
c_{i,t},
\end{equation}

the expression reduces to

\begin{equation}
TC_t
=
\sum_{i=1}^{N_t}
c_{i,t}
|\Delta w_{i,t}|.
\end{equation}

Allowing asymmetric costs can be important when taxes, exchange fees, market
microstructure, or other institutional features affect purchases and sales
differently.


% ------------------------------------------------------------
\subsubsection{Components of Trading Costs}
% ------------------------------------------------------------

The proportional coefficient $c_{i,t}$ is itself a simplified representation
of several distinct implementation costs.

Trading costs can broadly be decomposed as

\begin{equation}
\boxed{
TC_t
=
TC_t^{\mathrm{commission}}
+
TC_t^{\mathrm{spread}}
+
TC_t^{\mathrm{slippage}}
+
TC_t^{\mathrm{impact}}
+
TC_t^{\mathrm{tax}}
+
TC_t^{\mathrm{fees}}.
}
\end{equation}

These components have different economic origins.

\paragraph{Commissions.}

Broker commissions and explicit execution fees may be charged as a fixed amount
per order, per share, or as a proportion of traded notional.

A simple proportional representation is

\begin{equation}
TC_t^{\mathrm{commission}}
=
\sum_i
c_{i,t}^{\mathrm{commission}}
|\Delta w_{i,t}|.
\end{equation}


\paragraph{Bid--Ask Spread.}

A buyer generally executes closer to the ask price, while a seller executes
closer to the bid price.

Let

\begin{equation}
P_{i,t}^{\mathrm{mid}}
=
\frac{
P_{i,t}^{\mathrm{ask}}
+
P_{i,t}^{\mathrm{bid}}
}{2}
\end{equation}

denote the mid-price and

\begin{equation}
\mathrm{Spread}_{i,t}
=
P_{i,t}^{\mathrm{ask}}
-
P_{i,t}^{\mathrm{bid}}
\end{equation}

the quoted bid--ask spread.

Under the simplified assumption that a market order executes at the best bid
or ask, the cost relative to the mid-price is approximately one half of the
spread:

\begin{equation}
c_{i,t}^{\mathrm{spread}}
\approx
\frac{
P_{i,t}^{\mathrm{ask}}
-
P_{i,t}^{\mathrm{bid}}
}{
2P_{i,t}^{\mathrm{mid}}
}.
\end{equation}

Both purchases and sales therefore incur spread costs relative to the
mid-price.


\paragraph{Slippage.}

Slippage captures the difference between the theoretical or expected execution
price and the price actually obtained.

It may result from price movements between signal formation and execution,
order delays, limited liquidity, or execution uncertainty.


\paragraph{Market Impact.}

Large orders may move market prices against the investor. Market impact
therefore generally increases nonlinearly with order size.

A generic market-impact model can be represented as

\begin{equation}
TC_{i,t}^{\mathrm{impact}}
=
f
\left(
|\Delta w_{i,t}|,
\mathrm{ADV}_{i,t},
\sigma_{i,t},
\ldots
\right),
\end{equation}

where $\mathrm{ADV}_{i,t}$ denotes Average Daily Volume and $\sigma_{i,t}$
denotes a measure of asset volatility.

A commonly used stylized specification relates impact to participation in
daily volume:

\begin{equation}
TC_{i,t}^{\mathrm{impact}}
\propto
\sigma_{i,t}
\left(
\frac{
Q_{i,t}
}{
\mathrm{ADV}_{i,t}
}
\right)^{\gamma},
\end{equation}

where $Q_{i,t}$ is the traded quantity and typically

\begin{equation}
0<\gamma\leq 1.
\end{equation}


\paragraph{Taxes and Exchange Fees.}

Certain markets impose transaction taxes, stamp duties, exchange fees, or
regulatory fees. These costs may be asymmetric and may apply only to purchases
or only to sales.

They should therefore be modeled separately when economically material.


% ------------------------------------------------------------
\subsubsection{Trading Costs vs.\ Holding Costs}
% ------------------------------------------------------------

Trading costs should be distinguished from costs generated by maintaining
positions.

Trading costs are primarily functions of portfolio changes:

\begin{equation}
TC_t
=
C(\Delta w_t).
\end{equation}

By contrast, holding and financing costs depend primarily on the positions
maintained through time:

\begin{equation}
FC_{t\rightarrow t+1}
=
F
\left(
w_t,
t\rightarrow t+1
\right).
\end{equation}

For a long--short strategy, relevant holding costs may include:

\begin{itemize}
    \item short-stock borrowing fees;

    \item financing costs associated with leverage;

    \item financing spreads;

    \item other position-specific carrying costs.
\end{itemize}

A short sale therefore involves two conceptually different types of costs.

First, initiating or modifying the short position generates a trading cost,
just as buying or selling a long position does.

Second, maintaining the short position may generate an ongoing stock-borrow
cost.

These costs should not be combined indiscriminately because they depend on
different economic quantities.


% ------------------------------------------------------------
\subsubsection{General Transaction-Cost Function}
% ------------------------------------------------------------

A realistic transaction-cost function may therefore depend on several
security-level and market-level variables:

\begin{equation}
\boxed{
TC_t
=
C
\left(
\Delta w_t,
\mathrm{Spread}_t,
\mathrm{ADV}_t,
\mathrm{Volatility}_t,
\mathrm{Liquidity}_t,
\mathrm{ParticipationRate}_t,
\ldots
\right).
}
\end{equation}

At the simplest level, a backtest may assume a constant proportional cost:

\begin{equation}
TC_t
=
c
\sum_i |\Delta w_{i,t}|.
\end{equation}

At a more sophisticated level, costs may vary across securities, dates,
directions, and trade sizes:

\begin{equation}
TC_t
=
\sum_i
C_{i,t}
\left(
\Delta w_{i,t},
\mathrm{Spread}_{i,t},
\mathrm{ADV}_{i,t},
\sigma_{i,t},
\ldots
\right).
\end{equation}

The appropriate level of complexity depends on the strategy. Transaction-cost
modeling is particularly important for high-turnover strategies, strategies
trading less-liquid securities, and portfolios whose size represents a
significant fraction of available market liquidity.


% ------------------------------------------------------------
\subsubsection{Gross and Net Portfolio Returns}
% ------------------------------------------------------------

If the portfolio is formed at time $t$ and subsequently earns returns over
$[t,t+1]$, gross portfolio performance is

\begin{equation}
R_{p,t\rightarrow t+1}^{\mathrm{gross}}
=
w_t^\top
R_{t\rightarrow t+1},
\end{equation}

under the simplified assumption of no intra-period trading.

The corresponding net return can be represented schematically as

\begin{equation}
\boxed{
R_{p,t\rightarrow t+1}^{\mathrm{net}}
=
R_{p,t\rightarrow t+1}^{\mathrm{gross}}
-
TC_t
-
FC_{t\rightarrow t+1}.
}
\end{equation}

Here,

\begin{itemize}
    \item $TC_t$ represents costs generated by the trades required to establish
    the portfolio at time $t$;

    \item $FC_{t\rightarrow t+1}$ represents financing, stock-borrowing, and
    other holding costs incurred while maintaining the portfolio over
    $[t,t+1]$.
\end{itemize}

This decomposition emphasizes the distinction

\begin{equation}
\boxed{
\underbrace{TC_t}_{\text{cost of changing positions}}
\qquad\text{vs.}\qquad
\underbrace{FC_{t\rightarrow t+1}}_{\text{cost of holding positions}}.
}
\end{equation}

The precise portfolio accounting, transaction-cost models, financing
conventions, and treatment of market impact will be developed in greater
detail in the portfolio implementation and backtest return construction
sections.

% ------------------------------------------------------------
\subsection{The Complete Backtesting Timeline}
\label{subsec:complete_timeline}
% ------------------------------------------------------------

The conventions introduced above can be summarized through the complete
investment timeline.

At time $t$, the investor observes the information set

\begin{equation}
\mathcal{I}_t.
\end{equation}

The raw point-in-time data available to the investor satisfy

\begin{equation}
\mathcal{D}_t
\subseteq
\mathcal{I}_t.
\end{equation}

These data are transformed into features:

\begin{equation}
X_t
=
g(\mathcal{D}_t).
\end{equation}

A systematic investment rule or predictive model generates a raw signal:

\begin{equation}
s_t
=
f(X_t).
\end{equation}

Depending on the strategy, this signal may originate from:

\begin{itemize}
    \item a directly constructed firm-level characteristic;

    \item several combined characteristics;

    \item an econometric model;

    \item a machine-learning model;

    \item a deep-learning model.
\end{itemize}

The raw signal is subsequently transformed into a final signal:

\begin{equation}
\widetilde{s}_t
=
h
\left(
s_t,
Z_t
\right),
\end{equation}

where $Z_t$ may include variables required for transformations such as
industry, beta, size, or factor neutralization.

Portfolio construction maps the final signal into target weights:

\begin{equation}
w_t^{\mathrm{target}}
=
\phi
\left(
\widetilde{s}_t,
\widehat{\Sigma}_t,
w_{t^-},
\mathcal{C}_t
\right).
\end{equation}

The target portfolio may subsequently be modified because of implementation
constraints, tradability, rounding, short availability, or other market
frictions, producing the implemented portfolio

\begin{equation}
w_t.
\end{equation}

The portfolio then experiences future asset returns

\begin{equation}
R_{t\rightarrow t+1}
=
\begin{bmatrix}
R_{1,t\rightarrow t+1}\\
\vdots\\
R_{N_t,t\rightarrow t+1}
\end{bmatrix}.
\end{equation}

Ignoring intra-period trading, portfolio cash flows, and weight-drift
complications for the moment, the gross portfolio return is

\begin{equation}
\boxed{
R_{p,t\rightarrow t+1}^{\mathrm{gross}}
=
w_t^\top
R_{t\rightarrow t+1}.
}
\end{equation}

After implementation costs,

\begin{equation}
\boxed{
R_{p,t\rightarrow t+1}^{\mathrm{net}}
=
w_t^\top
R_{t\rightarrow t+1}
-
TC_t
-
FC_{t\rightarrow t+1}.
}
\end{equation}

The complete backtesting pipeline can therefore be summarized as

\begin{equation}
\boxed{
\mathcal{I}_t
\rightarrow
\mathcal{D}_t
\rightarrow
X_t
\rightarrow
s_t
\rightarrow
\widetilde{s}_t
\rightarrow
w_t^{\mathrm{target}}
\rightarrow
w_t
\rightarrow
R_{t\rightarrow t+1}
\rightarrow
R_{p,t\rightarrow t+1}^{\mathrm{net}}.
}
\end{equation}

Equivalently, at a higher conceptual level,

\begin{equation}
\boxed{
\text{Information}
\rightarrow
\text{Data}
\rightarrow
\text{Features}
\rightarrow
\text{Signal}
\rightarrow
\text{Portfolio}
\rightarrow
\text{Execution}
\rightarrow
\text{Performance}.
}
\end{equation}

This sequence provides the conceptual backbone of the entire backtesting
framework developed in the following sections.

Each subsequent chapter studies one or several of these transformations in
detail while maintaining the fundamental constraint that every investment
decision made at time $t$ must depend exclusively on the information available
at time $t$.


% ============================================================
\section{Backtesting Biases and Research Integrity}
% ============================================================

A backtest is only meaningful if the simulated investment process could have been
implemented using the information and market conditions that were actually
available at each historical date.

Backtesting errors generally arise when the historical simulation benefits from
information, securities, parameter choices, or trading conditions that would
not have been available to an investor in real time.

These errors can be grouped into several broad categories:

\begin{itemize}
    \item information-timing biases;
    \item universe-construction biases;
    \item sample-selection biases;
    \item model-selection and multiple-testing biases;
    \item data-processing leakage;
    \item unrealistic implementation assumptions.
\end{itemize}

The fundamental principle remains the information-set condition introduced in
Section~\ref{sec:problem_definition}. If $\mathcal{I}_t$ denotes all information
available to an investor at time $t$, then every investment decision must satisfy

\begin{equation}
\boxed{
w_t=f(\mathcal{I}_t).
}
\end{equation}

Any backtest in which $w_t$ depends directly or indirectly on information
revealed after $t$ is contaminated by future information.

However, avoiding explicit future variables is not sufficient. Bias can also
enter through the investment universe, data preprocessing, hyperparameter
selection, model comparison, transaction-cost assumptions, or even through the
researcher's decision to continue or abandon a strategy after observing its
historical performance.

This section describes the principal sources of bias and establishes a set of
research-integrity principles that should be enforced throughout the complete
backtesting pipeline.


% ------------------------------------------------------------
\subsection{Look-Ahead Bias}
\label{subsec:look_ahead_bias}
% ------------------------------------------------------------

Look-ahead bias occurs whenever a portfolio decision at time $t$ uses
information that was not yet available at that time.

Formally, if the portfolio is

\begin{equation}
w_t
=
f(X_t),
\end{equation}

then every component of $X_t$ must be measurable with respect to
$\mathcal{I}_t$.

A look-ahead violation occurs if

\begin{equation}
X_t
=
g
\left(
\mathcal{I}_t,
Z_{t+h}
\right),
\qquad
h>0,
\end{equation}

where $Z_{t+h}$ denotes information that becomes observable only after time
$t$.

A trivial example would be constructing today's portfolio using tomorrow's
return:

\begin{equation}
w_{i,t}
=
f
\left(
R_{i,t\rightarrow t+1}
\right).
\end{equation}

Such an error is obvious. In practice, however, look-ahead bias often enters
the research pipeline in more subtle ways.


\subsubsection{Price Timing}

Suppose a signal is computed using the closing price on day $t$:

\begin{equation}
s_t
=
f
\left(
P_t^{\mathrm{close}}
\right).
\end{equation}

If the closing price is not fully known until the market close, the strategy
cannot generally assume that it simultaneously executes the resulting trade at
that exact same closing price.

A more defensible sequence is

\begin{equation}
P_t^{\mathrm{close}}
\rightarrow
s_t
\rightarrow
w_t
\rightarrow
P_{t+1}^{\mathrm{exec}}.
\end{equation}

Alternatively, same-day execution may be valid if the strategy uses information
available before the execution window. The critical requirement is that the
timing assumption matches the actual information and trading process.


\subsubsection{Fundamental Data}

Accounting variables illustrate another common source of look-ahead bias.

Suppose a firm's fiscal-year earnings correspond to December 31 but are
published on March 15 of the following year.

The accounting value cannot be used on January 1 simply because its fiscal
period ended on December 31.

Let

\begin{equation}
t_i^{\mathrm{period}}
\end{equation}

denote the economic reference date of observation $i$, and

\begin{equation}
t_i^{\mathrm{available}}
\end{equation}

denote the date on which it becomes observable.

A valid backtest must satisfy

\begin{equation}
\boxed{
t
\geq
t_i^{\mathrm{available}}
}
\end{equation}

before the observation can enter $\mathcal{I}_t$.

The relevant backtesting date is therefore the information-availability date,
not necessarily the fiscal or economic reference date.


\subsubsection{Rolling Statistics}

Look-ahead bias can also enter through incorrectly constructed rolling
statistics.

Suppose volatility is estimated as

\begin{equation}
\widehat{\sigma}_{i,t}
=
\operatorname{Std}
\left(
R_{i,t-L+1},
\ldots,
R_{i,t}
\right).
\end{equation}

This estimator is valid if $R_{i,t}$ is already observable when the portfolio
decision is made.

However, using

\begin{equation}
\operatorname{Std}
\left(
R_{i,t-L+2},
\ldots,
R_{i,t+1}
\right)
\end{equation}

to construct the portfolio at time $t$ introduces future information.

The same issue applies to:

\begin{itemize}
    \item moving averages;
    \item rolling betas;
    \item rolling correlations;
    \item volatility estimates;
    \item rolling z-scores;
    \item rolling normalization parameters.
\end{itemize}


\subsubsection{Prevention}

The strongest protection against look-ahead bias is to define the backtest as
an explicit chronological process:

\begin{equation}
\boxed{
\mathcal{I}_t
\rightarrow
\text{data processing}
\rightarrow
\text{signal}
\rightarrow
\text{portfolio}
\rightarrow
\text{execution}
\rightarrow
R_{t\rightarrow t+1}.
}
\end{equation}

At every stage of the pipeline, one should be able to answer the question:

\begin{quote}
Could this quantity have been computed using only information that was
available at this exact historical time?
\end{quote}


% ------------------------------------------------------------
\subsection{Survivorship Bias}
\label{subsec:survivorship_bias}
% ------------------------------------------------------------

Survivorship bias arises when the historical universe contains only securities
that survive until a later date.

Suppose the current investment universe is

\begin{equation}
\mathcal{U}_{T},
\end{equation}

where $T$ denotes the present date.

Using this same universe throughout the historical backtest,

\begin{equation}
\mathcal{U}_t
=
\mathcal{U}_T
\qquad
\forall t<T,
\end{equation}

implicitly assumes that the investor knew which firms would survive until $T$.

This systematically removes securities that subsequently:

\begin{itemize}
    \item went bankrupt;
    \item were delisted;
    \item were acquired;
    \item ceased trading;
    \item became financially distressed;
    \item left the relevant index or market.
\end{itemize}

The resulting sample tends to contain firms with better ex-post outcomes than
the true historical opportunity set.

A valid universe must therefore be reconstructed dynamically:

\begin{equation}
\boxed{
\mathcal{U}_t
=
\left\{
i:
I_{i,t}^{\mathrm{eligible}}=1
\right\},
}
\end{equation}

where eligibility is determined using only information available at time $t$.

Survivorship bias is particularly dangerous in equity strategies based on
financial distress, value, profitability, or momentum, since securities that
eventually disappear may have extreme characteristics before their exit.


% ------------------------------------------------------------
\subsubsection{Delisting Returns}
\label{subsubsec:delisting_returns}
% ------------------------------------------------------------

Simply retaining delisted securities in the historical database is not always
sufficient to eliminate survivorship bias. A security disappearing from the
database or investment universe does not imply that the corresponding
investment position disappears without an economic gain or loss.

Suppose security $i$ has its last regular trading observation at time
$T_i-1$ and is subsequently delisted at time $T_i$. The investor may receive a
terminal economic payoff resulting from bankruptcy, liquidation, acquisition,
merger, privatization, or another corporate event.

Let

\begin{equation}
P_{i,T_i-1}
\end{equation}

denote the last relevant market price before delisting and let

\begin{equation}
V_{i,T_i}^{\mathrm{terminal}}
\end{equation}

denote the terminal economic value received by the shareholder. A simplified
delisting return can then be defined as

\begin{equation}
\boxed{
R_{i,T_i}^{\mathrm{delist}}
=
\frac{
V_{i,T_i}^{\mathrm{terminal}}
}{
P_{i,T_i-1}
}
-1.
}
\end{equation}

The terminal value may correspond, depending on the event, to cash proceeds,
liquidation proceeds, the value of securities received in a merger, or another
form of economic compensation.


\paragraph{Bankruptcy Example.}

Suppose a security is purchased at a price of

\begin{equation}
P_{i,t}=10.
\end{equation}

Before being delisted, its last observed market price falls to

\begin{equation}
P_{i,T_i-1}=8.
\end{equation}

If the company subsequently goes bankrupt and shareholders receive nothing,
then

\begin{equation}
V_{i,T_i}^{\mathrm{terminal}}=0.
\end{equation}

The final delisting return is therefore

\begin{equation}
R_{i,T_i}^{\mathrm{delist}}
=
\frac{0}{8}-1
=
-100\%.
\end{equation}

If the database simply stops the return series at the last observed price of
$8$, the backtest records only the decline

\begin{equation}
\frac{8}{10}-1
=
-20\%.
\end{equation}

However, the investor's true cumulative economic return is

\begin{equation}
\frac{0}{10}-1
=
-100\%.
\end{equation}

The correct economic path is therefore

\begin{equation}
\boxed{
10
\longrightarrow
8
\longrightarrow
0,
}
\end{equation}

rather than

\begin{equation}
10
\longrightarrow
8
\longrightarrow
\mathrm{NA}.
\end{equation}

Treating the missing observation as the disappearance of the position would
therefore substantially overstate portfolio performance.


\paragraph{Acquisition Example.}

A delisting does not necessarily generate a negative return.

Suppose a company trades at

\begin{equation}
P_{i,T_i-1}=80
\end{equation}

immediately before being acquired for

\begin{equation}
V_{i,T_i}^{\mathrm{terminal}}=100
\end{equation}

per share.

The corresponding delisting return is

\begin{equation}
R_{i,T_i}^{\mathrm{delist}}
=
\frac{100}{80}-1
=
25\%.
\end{equation}

If the historical price series simply terminates at $80$, the backtest would
fail to capture the $25\%$ terminal gain received by the shareholder.

Delisting returns should therefore not be interpreted exclusively as
bankruptcy losses. More generally, they represent the final economic payoff
associated with the disappearance of a security from regular trading.


\paragraph{Portfolio-Level Consequences.}

The distinction becomes particularly important when translating security-level
returns into portfolio returns.

Suppose the portfolio holds security $i$ immediately before its delisting with
weight

\begin{equation}
w_{i,T_i-1}=2\%.
\end{equation}

If the delisting return is

\begin{equation}
R_{i,T_i}^{\mathrm{delist}}=-100\%,
\end{equation}

then the contribution of the security to the portfolio return over the
delisting event is approximately

\begin{equation}
w_{i,T_i-1}
R_{i,T_i}^{\mathrm{delist}}
=
0.02\times(-1)
=
-2\%.
\end{equation}

Simply removing the security from the portfolio and setting

\begin{equation}
w_{i,T_i}=0
\end{equation}

without first accounting for the delisting return would incorrectly eliminate
this loss. Economically, such a procedure would implicitly assume that the
position could be liquidated at its previous valuation immediately before
disappearing.

The correct sequence is instead

\begin{equation}
\boxed{
\text{Position at }T_i-1
\rightarrow
R_{i,T_i}^{\mathrm{delist}}
\rightarrow
V_{i,T_i}^{\mathrm{terminal}}
\rightarrow
\text{Position removed from the portfolio}.
}
\end{equation}

Thus, the removal of a security from the investment universe and the economic
liquidation of the corresponding position are two distinct operations.

A historical return series should conceptually contain

\begin{equation}
\boxed{
R_{i,1},
R_{i,2},
\ldots,
R_{i,T_i-1},
R_{i,T_i}^{\mathrm{delist}},
}
\end{equation}

rather than

\begin{equation}
R_{i,1},
R_{i,2},
\ldots,
R_{i,T_i-1},
\mathrm{NA}.
\end{equation}

Only after the terminal economic payoff has been incorporated should the
security disappear from the portfolio and subsequent investment universe.

This issue is particularly important for strategies exposed to distressed,
small-cap, value, or otherwise financially fragile securities, since the
probability of delisting may be systematically related to the investment
signal itself. Ignoring delisting returns can therefore introduce not only a
general upward bias in performance, but also a bias that is correlated with
the characteristics being studied.


% ------------------------------------------------------------
\subsection{Selection Bias}
\label{subsec:selection_bias}
% ------------------------------------------------------------

Selection bias occurs when the sample used for research is selected using
information related to subsequent performance.

Suppose the true historical population is

\begin{equation}
\mathcal{P}_t,
\end{equation}

but the researcher constructs a sample

\begin{equation}
\mathcal{S}_t
\subset
\mathcal{P}_t
\end{equation}

according to criteria that depend, directly or indirectly, on future outcomes.

Then the estimated relationship

\begin{equation}
\mathbb{E}
\left[
R_{t+1}
\mid
X_t,
i\in\mathcal{S}_t
\right]
\end{equation}

need not correspond to the relationship that would have existed in the true
investable population.

Selection bias may arise through:

\begin{itemize}
    \item choosing markets that historically performed well;
    \item selecting factors after inspecting historical results;
    \item restricting the sample to companies with complete future histories;
    \item excluding periods in which the strategy performed poorly;
    \item choosing a backtest start date after examining results;
    \item studying only strategies that survived internal research.
\end{itemize}

An important distinction is that an economically motivated universe filter is
not necessarily selection bias.

For example, imposing at time $t$

\begin{equation}
\mathrm{ADV}_{i,t}
>
\mathrm{ADV}_{\min}
\end{equation}

may be economically justified because the strategy cannot realistically trade
illiquid securities.

The problem arises when the filtering rule itself is chosen ex post because it
improves historical performance.

Research decisions should therefore be motivated by economic reasoning,
implementation constraints, or hypotheses defined independently of the
backtest outcome whenever possible.


% ------------------------------------------------------------
\subsection{Data Snooping and Multiple Testing}
\label{subsec:data_snooping}
% ------------------------------------------------------------

Quantitative research typically involves evaluating many candidate signals,
models, hyperparameters, universes, and portfolio-construction rules.

Suppose the researcher evaluates $M$ strategies:

\begin{equation}
S_1,
S_2,
\ldots,
S_M.
\end{equation}

Even if every strategy has zero true alpha,

\begin{equation}
\mathbb{E}[R^{S_j}]=0
\qquad
\forall j,
\end{equation}

sampling noise implies that some strategies will exhibit strong historical
performance purely by chance.

If the researcher selects

\begin{equation}
S^*
=
\underset{j=1,\ldots,M}{\arg\max}
\;
\widehat{\mathrm{Sharpe}}_j,
\end{equation}

then

\begin{equation}
\widehat{\mathrm{Sharpe}}_{S^*}
\end{equation}

is upward biased as an estimate of the strategy's true Sharpe ratio.

The bias generally increases with the number of experiments performed.


\subsubsection{Repeated Hypothesis Testing}

Suppose each independent hypothesis test is conducted at significance level
$\alpha$.

For a single null hypothesis,

\begin{equation}
P(\text{false rejection})=\alpha.
\end{equation}

If $M$ independent null hypotheses are tested, the probability of obtaining at
least one false positive is

\begin{equation}
\boxed{
P(\text{at least one false positive})
=
1-(1-\alpha)^M.
}
\end{equation}

For example, with

\begin{equation}
\alpha=5\%,
\qquad
M=20,
\end{equation}

the probability becomes

\begin{equation}
1-0.95^{20}
\approx
64\%.
\end{equation}

Thus, observing one apparently significant result among many experiments
provides substantially weaker evidence than observing the same result when only
one hypothesis was tested.


\subsubsection{Hidden Multiple Testing}

The number of effective tests is often much larger than the number of models
reported in the final paper or research note.

Researcher degrees of freedom may include:

\begin{itemize}
    \item alternative signal definitions;
    \item different winsorization thresholds;
    \item different lags;
    \item different holding periods;
    \item alternative universes;
    \item different neutralization procedures;
    \item different model architectures;
    \item different hyperparameters;
    \item alternative transaction-cost assumptions;
    \item different portfolio mappings.
\end{itemize}

All of these decisions constitute implicit experiments if they are adjusted
after observing backtest performance.


\subsubsection{Mitigation}

Possible safeguards include:

\begin{itemize}
    \item keeping a genuinely untouched test sample;
    \item walk-forward validation;
    \item predefining the research hypothesis;
    \item recording all experiments;
    \item correcting for multiple testing;
    \item evaluating economic rather than purely statistical significance;
    \item testing robustness across independent markets and periods;
    \item using techniques such as the Deflated Sharpe Ratio.
\end{itemize}

The purpose is not to eliminate experimentation. Experimentation is central to
quantitative research. The objective is instead to recognize that repeated
experimentation changes the statistical interpretation of the best observed
backtest.


% ------------------------------------------------------------
\subsection{Publication and Restatement Bias}
\label{subsec:publication_restatement_bias}
% ------------------------------------------------------------

Historical databases frequently contain the latest corrected or restated value
of an economic or accounting variable rather than the value that was originally
available to investors.

Suppose a firm initially reports

\begin{equation}
X_{i,t}^{(0)}
\end{equation}

and later restates the same accounting item as

\begin{equation}
X_{i,t}^{(1)}.
\end{equation}

If the backtest uses

\begin{equation}
X_{i,t}^{(1)}
\end{equation}

at a historical date when investors could observe only

\begin{equation}
X_{i,t}^{(0)},
\end{equation}

then future information has implicitly entered the strategy.

The same issue appears in macroeconomic data.

For example, GDP, inflation, employment, and other economic statistics may be
revised several times after their initial release.

The value stored today,

\begin{equation}
X_t^{\mathrm{final}},
\end{equation}

may therefore differ from the real-time value

\begin{equation}
X_t^{\mathrm{vintage}}.
\end{equation}

For a historically realistic backtest, one ideally requires point-in-time or
vintage data:

\begin{equation}
\boxed{
X_t^{\mathrm{backtest}}
=
X_t^{\mathrm{available\ at\ }t}.
}
\end{equation}

Using final revised data may materially overstate the predictability available
to a real-time investor.


% ------------------------------------------------------------
\subsection{Universe Reconstruction Bias}
\label{subsec:universe_reconstruction_bias}
% ------------------------------------------------------------

Universe reconstruction bias is related to, but broader than, survivorship
bias.

A historical investment universe should reproduce the securities that would
actually have satisfied the strategy's eligibility rules at each date.

Let the universe-selection function be

\begin{equation}
\mathcal{U}_t
=
g
\left(
\mathcal{I}_t
\right).
\end{equation}

The universe itself is therefore part of the investment decision process.

Errors can arise when historical membership is reconstructed using information
that was not available contemporaneously.

Examples include:

\begin{itemize}
    \item using today's index constituents historically;
    \item applying today's sector classification to past observations;
    \item using future market capitalization;
    \item using future liquidity measures;
    \item excluding securities that later delist;
    \item incorrectly handling IPO entry dates;
    \item applying filters using future data availability.
\end{itemize}


\subsubsection{Index Membership}

Suppose a strategy trades constituents of an index.

The correct universe is

\begin{equation}
\mathcal{U}_t
=
\text{index constituents known at }t,
\end{equation}

not

\begin{equation}
\mathcal{U}_t
=
\text{current constituents}.
\end{equation}

Moreover, index changes are often announced before their effective date. The
researcher must therefore specify whether the strategy trades using the
announcement date, the effective date, or another operational convention.


\subsubsection{Universe Filters}

Suppose a minimum market-capitalization threshold is imposed:

\begin{equation}
\mathrm{MCap}_{i,t}
>
M_{\min}.
\end{equation}

The market capitalization used must be based on information available at $t$.

Similarly, a liquidity filter such as

\begin{equation}
\mathrm{ADV}_{i,t}^{(L)}
>
ADV_{\min}
\end{equation}

must use trailing rather than future volume.

The universe-selection function must therefore obey exactly the same
information-set discipline as the signal itself.


% ------------------------------------------------------------
\subsection{Data Leakage}
\label{subsec:data_leakage}
% ------------------------------------------------------------

Data leakage occurs when information from outside the training information set
influences model estimation, preprocessing, feature construction, or
hyperparameter selection.

Leakage is particularly important for econometric, machine-learning, and
deep-learning strategies.


\subsubsection{Preprocessing Leakage}

Consider standardizing a feature:

\begin{equation}
Z_{i,t}
=
\frac{
X_{i,t}-\mu_X
}{
\sigma_X
}.
\end{equation}

If $\mu_X$ and $\sigma_X$ are estimated using the entire dataset, including
future observations, then the transformed training data contain information
about the future.

Instead, preprocessing parameters must be estimated from the training sample:

\begin{equation}
\widehat{\mu}_{X,\mathrm{train}}
=
\frac{1}{N_{\mathrm{train}}}
\sum_{j\in\mathrm{train}}
X_j,
\end{equation}

and

\begin{equation}
\widehat{\sigma}_{X,\mathrm{train}}
=
\operatorname{Std}
\left(
X_j:
j\in\mathrm{train}
\right).
\end{equation}

Validation and test observations are then transformed using these fixed
training parameters:

\begin{equation}
Z_j^{\mathrm{test}}
=
\frac{
X_j^{\mathrm{test}}
-
\widehat{\mu}_{X,\mathrm{train}}
}{
\widehat{\sigma}_{X,\mathrm{train}}
}.
\end{equation}


\subsubsection{Feature-Selection Leakage}

Suppose features are selected by computing their relationship with the target
over the entire dataset:

\begin{equation}
\rho_j
=
\operatorname{Corr}
\left(
X_j,
Y
\right).
\end{equation}

Selecting features with the largest $|\rho_j|$ before splitting the data causes
future target information to influence the model.

Feature selection must instead be performed entirely inside the training
procedure.


\subsubsection{Hyperparameter Leakage}

Suppose hyperparameters $\theta$ are chosen as

\begin{equation}
\theta^*
=
\underset{\theta}{\arg\max}
\;
\mathrm{Performance}_{\mathrm{test}}(\theta).
\end{equation}

The test sample is then no longer a true out-of-sample sample.

The proper structure is

\begin{equation}
\boxed{
\text{Train}
\rightarrow
\text{Validation}
\rightarrow
\text{Test}.
}
\end{equation}

Training data estimate model parameters. Validation data select models and
hyperparameters. Test data are used only for final evaluation.

For time-series applications, these sets must also respect chronology.


\subsubsection{Cross-Sectional vs.\ Temporal Leakage}

In panel datasets, leakage can occur along both the time and asset dimensions.

A random train-test split may place observations from the same economic period
in both training and testing samples.

For financial prediction, a chronological split is generally preferable:

\begin{equation}
\underbrace{
\{1,\ldots,T_{\mathrm{train}}\}
}_{\text{training}}
<
\underbrace{
\{T_{\mathrm{train}}+1,\ldots,T_{\mathrm{val}}\}
}_{\text{validation}}
<
\underbrace{
\{T_{\mathrm{val}}+1,\ldots,T_{\mathrm{test}}\}
}_{\text{testing}}.
\end{equation}

Walk-forward validation will be developed in detail in the model-training
section.


% ------------------------------------------------------------
\subsection{Overfitting}
\label{subsec:overfitting}
% ------------------------------------------------------------

Overfitting occurs when a model captures idiosyncratic patterns in the
historical sample rather than persistent economic relationships.

Suppose observed returns satisfy

\begin{equation}
Y
=
f^*(X)
+
\varepsilon,
\end{equation}

where $f^*$ denotes the true underlying relationship and $\varepsilon$ denotes
noise.

A highly flexible model may estimate

\begin{equation}
\widehat{f}
\end{equation}

that fits both

\begin{equation}
f^*(X)
\end{equation}

and part of the sample-specific noise $\varepsilon$.

This may produce excellent in-sample performance:

\begin{equation}
\mathrm{Error}_{\mathrm{IS}}
\downarrow,
\end{equation}

while out-of-sample performance deteriorates:

\begin{equation}
\mathrm{Error}_{\mathrm{OOS}}
\uparrow.
\end{equation}


\subsubsection{Backtest Overfitting}

Overfitting is not restricted to predictive models.

An entire investment strategy can be overfit through repeated choices of:

\begin{itemize}
    \item signal definitions;
    \item lookback windows;
    \item portfolio thresholds;
    \item neutralization factors;
    \item rebalancing frequencies;
    \item risk targets;
    \item transaction-cost assumptions;
    \item universes.
\end{itemize}

The final strategy may therefore be viewed as the result of a large implicit
optimization:

\begin{equation}
\theta^*
=
\underset{\theta\in\Theta}{\arg\max}
\;
\widehat{S}(\theta),
\end{equation}

where $\theta$ represents all research choices and $\widehat{S}$ denotes an
estimated performance statistic such as the Sharpe ratio.

The larger the effective parameter space $\Theta$, the greater the risk that

\begin{equation}
\widehat{S}(\theta^*)
\gg
S_{\mathrm{true}}(\theta^*).
\end{equation}


\subsubsection{Signs of Overfitting}

Potential warning signs include:

\begin{itemize}
    \item very high in-sample Sharpe ratios;
    \item narrow optimal parameter regions;
    \item strong sensitivity to small parameter changes;
    \item performance concentrated in a short subperiod;
    \item performance driven by a small number of securities;
    \item poor performance across related markets;
    \item large degradation from validation to test samples;
    \item economically implausible signal mechanisms.
\end{itemize}


\subsubsection{Mitigation}

Overfitting can be reduced through:

\begin{itemize}
    \item simple and economically interpretable hypotheses;
    \item regularization;
    \item walk-forward validation;
    \item strict separation of training, validation, and testing;
    \item parameter-stability analysis;
    \item robustness tests;
    \item independent datasets;
    \item realistic transaction costs;
    \item explicit accounting for the number of research trials.
\end{itemize}

The objective is not necessarily to choose the model with the highest
historical performance, but rather the model whose performance is most likely
to generalize beyond the historical sample.


% ------------------------------------------------------------
\subsection{Unrealistic Trading and Cost Assumptions}
\label{subsec:unrealistic_trading}
% ------------------------------------------------------------

Even a statistically valid signal can produce an unrealistic backtest if the
simulation assumes impossible or excessively favorable trading conditions.

The backtest must therefore distinguish between

\begin{equation}
\boxed{
\text{predictive performance}
}
\end{equation}

and

\begin{equation}
\boxed{
\text{economically implementable performance}.
}
\end{equation}


\subsubsection{Execution Prices}

A common unrealistic assumption is executing all trades at the closing price
used to compute the signal.

If

\begin{equation}
s_t
=
f(P_t^{\mathrm{close}}),
\end{equation}

then execution at the same closing price may not be feasible.

Execution should instead respect the chronological sequence established in the
backtest design.


\subsubsection{Bid--Ask Spread}

Assuming execution at the mid-price ignores the bid--ask spread.

For a security with bid and ask prices

\begin{equation}
P_{i,t}^{\mathrm{bid}}
\qquad\text{and}\qquad
P_{i,t}^{\mathrm{ask}},
\end{equation}

the mid-price is

\begin{equation}
P_{i,t}^{\mathrm{mid}}
=
\frac{
P_{i,t}^{\mathrm{ask}}
+
P_{i,t}^{\mathrm{bid}}
}{2}.
\end{equation}

A market buy generally occurs closer to the ask and a market sell closer to the
bid.

The one-way spread cost relative to the mid-price is approximately

\begin{equation}
c_{i,t}^{\mathrm{spread}}
\approx
\frac{
P_{i,t}^{\mathrm{ask}}
-
P_{i,t}^{\mathrm{bid}}
}{
2P_{i,t}^{\mathrm{mid}}
}.
\end{equation}


\subsubsection{Market Impact}

A backtest that assumes transaction costs are independent of portfolio size may
substantially overstate strategy capacity.

Let

\begin{equation}
Q_{i,t}
\end{equation}

denote traded quantity and

\begin{equation}
\mathrm{ADV}_{i,t}
\end{equation}

denote Average Daily Volume.

The participation rate is

\begin{equation}
\boxed{
PR_{i,t}
=
\frac{
Q_{i,t}
}{
\mathrm{ADV}_{i,t}
}.
}
\end{equation}

As $PR_{i,t}$ increases, market impact generally increases.

A stylized impact function may take the form

\begin{equation}
c_{i,t}^{\mathrm{impact}}
=
\eta
\sigma_{i,t}
PR_{i,t}^{\gamma},
\end{equation}

where $\eta>0$ and typically

\begin{equation}
0<\gamma\leq 1.
\end{equation}

A strategy that appears highly profitable at small capital may therefore become
unattractive at large scale.


\subsubsection{Short-Selling Constraints}

Backtests of long--short strategies frequently assume that every security can
be shorted without restriction.

In reality, a short position may face:

\begin{itemize}
    \item limited borrow availability;
    \item stock-loan fees;
    \item hard-to-borrow premiums;
    \item recalls;
    \item short-sale restrictions;
    \item position limits.
\end{itemize}

The cost and feasibility of shorting may also be correlated with the investment
signal itself.

For example, distressed or highly shorted securities may simultaneously appear
attractive to a short signal and be particularly expensive to borrow.


\subsubsection{Liquidity Constraints}

A realistic portfolio should constrain trading relative to available market
liquidity.

For example,

\begin{equation}
Q_{i,t}
\leq
\kappa
\,
\mathrm{ADV}_{i,t},
\end{equation}

where $\kappa$ denotes the maximum acceptable participation rate.

Ignoring such constraints can result in hypothetical portfolios that could not
have been executed in practice.


\subsubsection{Financing and Leverage}

Leveraged portfolios incur financing costs.

If leverage at time $t$ is denoted by

\begin{equation}
L_t,
\end{equation}

then the economic return of a leveraged strategy cannot generally be obtained
simply by multiplying the unlevered return by $L_t$.

Financing spreads, collateral requirements, borrowing constraints, and margin
rules must also be considered.


\subsubsection{Capacity}

The profitability of a strategy is generally a function of assets under
management:

\begin{equation}
\Pi
=
\Pi(A),
\end{equation}

where $A$ denotes strategy capital.

Because transaction costs and market impact generally increase with traded
notional,

\begin{equation}
\frac{\partial TC(A)}{\partial A}
>
0.
\end{equation}

Strategy capacity can therefore be interpreted as the level of assets at which
the incremental economic attractiveness of the strategy becomes insufficient.

Capacity analysis will be developed in greater detail in the performance and
implementation sections.


% ------------------------------------------------------------
\subsection{Research Integrity and Backtest Governance}
\label{subsec:research_integrity}
% ------------------------------------------------------------

The biases described above are not independent. A single research workflow may
simultaneously contain survivorship bias, preprocessing leakage, repeated
parameter tuning, and unrealistic implementation assumptions.

A robust quantitative research process should therefore impose systematic
controls throughout the entire pipeline.


\subsubsection{Chronological Integrity}

Every transformation should preserve the chronological ordering

\begin{equation}
\boxed{
\text{information}
\rightarrow
\text{decision}
\rightarrow
\text{execution}
\rightarrow
\text{realized outcome}.
}
\end{equation}

Future observations must never influence earlier investment decisions.


\subsubsection{Point-in-Time Integrity}

Data used at time $t$ should reproduce the data that would genuinely have been
available at that date:

\begin{equation}
\boxed{
\mathcal{D}_t
=
\mathcal{D}_t^{\mathrm{PIT}}.
}
\end{equation}

This includes:

\begin{itemize}
    \item security membership;
    \item fundamental releases;
    \item analyst forecasts;
    \item classifications;
    \item corporate actions;
    \item macroeconomic vintages.
\end{itemize}


\subsubsection{Out-of-Sample Integrity}

All decisions involving model selection should be separated from final
evaluation.

Conceptually,

\begin{equation}
\boxed{
\underbrace{\text{Training}}_{\text{parameter estimation}}
\rightarrow
\underbrace{\text{Validation}}_{\text{model selection}}
\rightarrow
\underbrace{\text{Testing}}_{\text{final evaluation}}.
}
\end{equation}

Repeatedly consulting the test sample gradually converts it into another
validation sample.


\subsubsection{Experiment Tracking}

Every material research experiment should ideally record:

\begin{itemize}
    \item data version;
    \item universe definition;
    \item signal definition;
    \item preprocessing parameters;
    \item model parameters;
    \item portfolio-construction parameters;
    \item transaction-cost assumptions;
    \item performance results;
    \item date of the experiment.
\end{itemize}

This helps quantify researcher degrees of freedom and prevents the final
strategy from appearing as if it were the only hypothesis ever tested.


\subsubsection{Economic Plausibility}

Statistical performance alone is insufficient evidence of a robust investment
strategy.

A credible strategy should ideally have an economic, behavioral, structural,
or institutional explanation for why the predictive relationship may persist.

The relevant question is not only

\begin{equation}
\text{``Did it work historically?''}
\end{equation}

but also

\begin{equation}
\boxed{
\text{``Why should this relationship survive out of sample?''}
}
\end{equation}


\subsubsection{Backtest Integrity Checklist}

Before interpreting historical performance, the researcher should verify at
least that:

\begin{itemize}
    \item no future information enters signal construction;

    \item preprocessing parameters are estimated using only admissible data;

    \item the historical investment universe is reconstructed point-in-time;

    \item delisted and disappeared securities are treated correctly;

    \item model selection is separated from final testing;

    \item the number of experiments and parameter searches is recognized;

    \item execution timing is feasible;

    \item transaction costs include both purchases and sales;

    \item spread, market impact, borrowing, and financing are modeled when
    economically relevant;

    \item portfolio size is consistent with market liquidity;

    \item robustness is evaluated across alternative periods and reasonable
    parameter choices.
\end{itemize}

A backtest should therefore not be interpreted as a direct estimate of future
performance. It is better understood as a controlled historical experiment
designed to determine whether an investment hypothesis remains economically and
statistically plausible after accounting for realistic information,
implementation, and research constraints.

% ============================================================
\part{Data Engineering}
% ============================================================

% ============================================================
\section{Raw Data Acquisition}
\label{sec:raw_data_acquisition}
% ============================================================

Data constitute the first operational input of a quantitative investment
process. Before signals can be constructed, the backtester must collect the
historical information required to reproduce the investment decisions that
could have been made at each point in time.

Let the complete collection of raw data available to the research system be
denoted by

\begin{equation}
\mathcal{D}.
\end{equation}

This collection may contain several heterogeneous data sources:

\begin{equation}
\boxed{
\mathcal{D}
=
\left\{
\mathcal{D}^{\mathrm{market}},
\mathcal{D}^{\mathrm{fund}},
\mathcal{D}^{\mathrm{est}},
\mathcal{D}^{\mathrm{alt}},
\mathcal{D}^{\mathrm{risk}},
\mathcal{D}^{\mathrm{class}},
\mathcal{D}^{\mathrm{corp}},
\mathcal{D}^{\mathrm{impl}}
\right\}.
}
\end{equation}

where the superscripts correspond respectively to market data, fundamental
data, analyst estimates, alternative data, risk and market variables,
classifications, corporate actions, and implementation-related data.

These sources may differ substantially in:

\begin{itemize}
    \item observation frequency;
    \item historical coverage;
    \item entity granularity;
    \item timestamp conventions;
    \item identifiers;
    \item revision policies;
    \item data quality;
    \item point-in-time availability.
\end{itemize}

Consequently, raw data acquisition should not be viewed as the construction of
a single homogeneous table. Rather, it consists of collecting a set of
heterogeneous historical sources that will subsequently be aligned into a
consistent point-in-time information set.

A useful distinction is

\begin{equation}
\boxed{
\text{Raw Data}
\neq
\text{Point-in-Time Data}
\neq
\text{Model-Ready Data}.
}
\end{equation}

Raw data correspond to observations obtained from the original data sources.
Point-in-time data reconstruct the information that was genuinely available at
each historical date. Model-ready data are subsequently cleaned, aligned,
transformed, and converted into features suitable for signal construction.

The present section focuses exclusively on the nature of the raw inputs.


% ------------------------------------------------------------
\subsection{Market Data}
\label{subsec:market_data}
% ------------------------------------------------------------

Market data describe the observable trading history of financial securities.

For equity security $i$, a standard daily record may contain

\begin{equation}
\boxed{
\left(
P_{i,t}^{\mathrm{open}},
P_{i,t}^{\mathrm{high}},
P_{i,t}^{\mathrm{low}},
P_{i,t}^{\mathrm{close}},
V_{i,t}
\right),
}
\end{equation}

where:

\begin{itemize}
    \item $P_{i,t}^{\mathrm{open}}$ denotes the opening price;
    \item $P_{i,t}^{\mathrm{high}}$ denotes the highest traded price;
    \item $P_{i,t}^{\mathrm{low}}$ denotes the lowest traded price;
    \item $P_{i,t}^{\mathrm{close}}$ denotes the closing price;
    \item $V_{i,t}$ denotes trading volume.
\end{itemize}

Depending on the strategy and data provider, market data may additionally
contain:

\begin{itemize}
    \item VWAP;
    \item number of trades;
    \item traded notional;
    \item intraday trades;
    \item quotes;
    \item order-book information;
    \item trading venue;
    \item exchange status.
\end{itemize}

Detailed liquidity measures such as bid--ask spreads, market depth, ADV, and
participation rates are treated separately in
Section~\ref{subsec:liquidity_borrowing_transaction_cost_data}.


\subsubsection{Prices and Returns}

Given a sequence of prices, the simple return of security $i$ between $t-1$ and
$t$ is

\begin{equation}
R_{i,t}
=
\frac{P_{i,t}}{P_{i,t-1}}-1.
\end{equation}

The corresponding logarithmic return is

\begin{equation}
r_{i,t}
=
\log(P_{i,t})
-
\log(P_{i,t-1}).
\end{equation}

The two quantities satisfy

\begin{equation}
r_{i,t}
=
\log(1+R_{i,t}),
\end{equation}

and

\begin{equation}
R_{i,t}
=
e^{r_{i,t}}-1.
\end{equation}

For small returns,

\begin{equation}
r_{i,t}
\approx
R_{i,t}.
\end{equation}

Price returns do not necessarily represent the complete economic return earned
by an investor.

If $D_{i,t}$ denotes cash distributions received during the period, a
simplified total return is

\begin{equation}
\boxed{
R_{i,t}^{\mathrm{total}}
=
\frac{
P_{i,t}
+
D_{i,t}
}{
P_{i,t-1}
}
-1.
}
\end{equation}

The treatment of dividends, splits, mergers, and other events is developed in
Section~\ref{subsec:corporate_actions}.


\subsubsection{Raw and Adjusted Prices}

Historical databases frequently provide both raw and adjusted prices.

Let

\begin{equation}
P_{i,t}^{\mathrm{raw}}
\end{equation}

denote the observed market price and

\begin{equation}
P_{i,t}^{\mathrm{adj}}
\end{equation}

an adjusted price.

A generic relationship is

\begin{equation}
P_{i,t}^{\mathrm{adj}}
=
A_{i,t}
P_{i,t}^{\mathrm{raw}},
\end{equation}

where $A_{i,t}$ denotes an adjustment factor.

Whenever possible, the research system should retain separately

\begin{equation}
\boxed{
P_{i,t}^{\mathrm{raw}},
\qquad
A_{i,t},
\qquad
P_{i,t}^{\mathrm{adj}}.
}
\end{equation}

This allows the historical transformation to remain auditable rather than
depending exclusively on vendor-preprocessed series.


\subsubsection{Data Frequency}

Market data may be observed at different frequencies:

\begin{equation}
\boxed{
\text{Tick}
\rightarrow
\text{Intraday}
\rightarrow
\text{Daily}
\rightarrow
\text{Weekly}
\rightarrow
\text{Monthly}.
}
\end{equation}

The required frequency depends on both signal horizon and execution horizon.

For example, a monthly equity strategy may still require daily prices to
construct signals or simulate execution, whereas an intraday strategy may
require individual transactions and quotes.


\subsubsection{Security Identifiers}

A security should not generally be identified exclusively by its ticker.

Tickers may:

\begin{itemize}
    \item change through time;
    \item be reused;
    \item differ across exchanges;
    \item refer to different share classes.
\end{itemize}

A robust architecture should therefore use persistent identifiers whenever
possible.

Conceptually,

\begin{equation}
\boxed{
\text{Company}
\longleftrightarrow
\text{Security}
\longleftrightarrow
\text{Listing}
\longleftrightarrow
\text{Ticker}.
}
\end{equation}

These objects are distinct. A company may issue multiple securities, and one
security may be associated with multiple listings or trading venues.


% ------------------------------------------------------------
\subsection{Fundamental Data}
\label{subsec:fundamental_data}
% ------------------------------------------------------------

Fundamental data describe the accounting and economic characteristics of
companies.

They are commonly extracted from:

\begin{itemize}
    \item income statements;
    \item balance sheets;
    \item cash-flow statements;
    \item regulatory filings;
    \item annual and interim reports;
    \item company disclosures.
\end{itemize}

For company $i$ and accounting period $\tau$, let

\begin{equation}
F_{i,\tau}
=
\begin{bmatrix}
F_{i,\tau}^{(1)}\\
\vdots\\
F_{i,\tau}^{(K)}
\end{bmatrix}
\in
\mathbb{R}^{K}
\end{equation}

denote a vector of $K$ fundamental variables.


\subsubsection{Main Accounting Variables}

Typical income-statement variables include:

\begin{itemize}
    \item revenue;
    \item gross profit;
    \item EBITDA;
    \item EBIT;
    \item operating income;
    \item interest expense;
    \item taxes;
    \item net income;
    \item earnings per share.
\end{itemize}

Typical balance-sheet variables include:

\begin{itemize}
    \item cash;
    \item inventories;
    \item receivables;
    \item total assets;
    \item short- and long-term debt;
    \item total liabilities;
    \item shareholders' equity;
    \item book value.
\end{itemize}

Typical cash-flow variables include:

\begin{itemize}
    \item operating cash flow;
    \item capital expenditures;
    \item investing cash flow;
    \item financing cash flow;
    \item free cash flow;
    \item dividends;
    \item share repurchases.
\end{itemize}


\subsubsection{Stock and Flow Variables}

Accounting variables differ according to whether they are measured at a date or
over a period.

A stock variable, such as total assets, is measured at a particular date:

\begin{equation}
\mathrm{Assets}_{i,\tau}.
\end{equation}

A flow variable, such as revenue, is accumulated over an interval:

\begin{equation}
\mathrm{Revenue}_{i,[\tau_0,\tau_1]}.
\end{equation}

This distinction matters when constructing financial ratios.

For example,

\begin{equation}
ROA_i
=
\frac{
\mathrm{NetIncome}_{i,[\tau_0,\tau_1]}
}{
\frac{1}{2}
\left(
\mathrm{Assets}_{i,\tau_0}
+
\mathrm{Assets}_{i,\tau_1}
\right)
}.
\end{equation}

The numerator is measured over a period, whereas the denominator approximates
the average balance-sheet stock over that period.


\subsubsection{Reporting Structure}

Fundamental data may be reported annually, semi-annually, or quarterly,
depending on the company, country, regulatory regime, and historical period.

Importantly, the accounting period to which an observation refers is not
necessarily the date on which that observation becomes known to investors.

The precise distinction between reporting date, publication date, vendor
timestamp, and information-availability date is deferred to the
Point-in-Time Data Construction section.


% ------------------------------------------------------------
\subsection{Analyst Estimates}
\label{subsec:analyst_estimates}
% ------------------------------------------------------------

Sell-side analysts publish forecasts for future company fundamentals and other
financial quantities.

For company $i$, analyst $a$, information date $t$, target fiscal period
$\tau$, and forecasted variable $k$, an individual estimate can be represented
as

\begin{equation}
E_{i,a,t}^{(\tau,k)}.
\end{equation}

Typical forecast variables include:

\begin{itemize}
    \item revenue;
    \item EBITDA;
    \item EBIT;
    \item earnings per share;
    \item net income;
    \item free cash flow;
    \item target price.
\end{itemize}


\subsubsection{Individual and Consensus Estimates}

Suppose $A_{i,t}$ analysts provide estimates for company $i$.

A simple consensus estimate is

\begin{equation}
\boxed{
\bar{E}_{i,t}
=
\frac{1}{A_{i,t}}
\sum_{a=1}^{A_{i,t}}
E_{i,a,t}.
}
\end{equation}

Alternative aggregation methods may use:

\begin{itemize}
    \item the median;
    \item trimmed means;
    \item recency weights;
    \item analyst-specific weights.
\end{itemize}

The number of contributing analysts,

\begin{equation}
A_{i,t},
\end{equation}

may itself be informative because coverage and forecast uncertainty often vary
across securities.


\subsubsection{Forecast Revisions and Dispersion}

A simple revision over horizon $h$ is

\begin{equation}
\boxed{
\Delta E_{i,t}^{(h)}
=
E_{i,t}
-
E_{i,t-h}.
}
\end{equation}

A normalized revision may be written as

\begin{equation}
\Delta E_{i,t}^{\mathrm{rel}}
=
\frac{
E_{i,t}-E_{i,t-h}
}{
|E_{i,t-h}|
}.
\end{equation}

If individual analyst estimates are available, disagreement may be measured as

\begin{equation}
\mathrm{Dispersion}_{i,t}
=
\sqrt{
\frac{1}{A_{i,t}-1}
\sum_{a=1}^{A_{i,t}}
\left(
E_{i,a,t}
-
\bar{E}_{i,t}
\right)^2
}.
\end{equation}

Revisions measure changes in expectations, whereas dispersion measures
cross-sectional disagreement among analysts.


\subsubsection{Actuals}

Once the corresponding realized accounting quantity becomes available, analyst
forecasts can be compared with the realized value.

The timing of forecasts, revisions, and actual releases is essential and will
be handled explicitly in the point-in-time construction stage rather than in
the raw-data description.


% ------------------------------------------------------------
\subsection{Alternative Data}
\label{subsec:alternative_data}
% ------------------------------------------------------------

Alternative data refer broadly to information sources outside conventional
market, accounting, and analyst-estimate databases.

Examples include:

\begin{itemize}
    \item news;
    \item earnings-call transcripts;
    \item regulatory text;
    \item social-media activity;
    \item web traffic;
    \item job postings;
    \item credit-card transactions;
    \item satellite imagery;
    \item geolocation information;
    \item supply-chain information;
    \item product reviews;
    \item search activity.
\end{itemize}

Unlike conventional financial data, these sources may be:

\begin{itemize}
    \item unstructured;
    \item high dimensional;
    \item irregularly sampled;
    \item event driven;
    \item noisy;
    \item difficult to associate with tradable securities.
\end{itemize}


\subsubsection{Raw Data Representations}

Alternative-data observations need not initially be numerical.

For example,

\begin{equation}
\text{Text Document},
\qquad
\text{Image},
\qquad
\text{Event},
\qquad
\text{Graph},
\qquad
\text{Transaction Record}
\end{equation}

may all constitute raw inputs.

Their conversion into numerical features belongs to the Feature Engineering
section rather than to raw data acquisition.


\subsubsection{Entity Mapping}

Alternative datasets frequently refer to economic entities rather than directly
to securities.

A generic mapping may therefore require

\begin{equation}
\boxed{
\text{Raw Entity}
\longrightarrow
\text{Company ID}
\longrightarrow
\text{Security ID}.
}
\end{equation}

Examples include mapping a product, brand, subsidiary, website, or executive to
the corresponding publicly traded company.

The historical validity and timing of these mappings will be addressed in the
point-in-time construction stage.


% ------------------------------------------------------------
\subsection{Risk Factors and Market Variables}
\label{subsec:risk_factors_market_variables}
% ------------------------------------------------------------

Quantitative strategies frequently require market-wide variables in addition to
security-specific information.

Let

\begin{equation}
F_t
=
\begin{bmatrix}
F_{1,t}\\
\vdots\\
F_{K,t}
\end{bmatrix}
\end{equation}

denote a vector of $K$ systematic factor returns.


\subsubsection{Systematic Risk Factors}

Possible factors include:

\begin{itemize}
    \item market;
    \item size;
    \item value;
    \item momentum;
    \item profitability;
    \item investment;
    \item quality;
    \item low volatility;
    \item industry;
    \item country.
\end{itemize}

A linear factor model may be represented as

\begin{equation}
R_{i,t}^{e}
=
\alpha_i
+
\beta_i^\top F_t
+
\varepsilon_{i,t},
\end{equation}

where $\beta_i$ denotes the vector of factor exposures.

Such factors may later be used for:

\begin{itemize}
    \item beta estimation;
    \item signal neutralization;
    \item portfolio constraints;
    \item risk modeling;
    \item performance attribution.
\end{itemize}


\subsubsection{Risk-Free Rates and Macroeconomic Variables}

Other relevant market-level inputs may include:

\begin{itemize}
    \item risk-free rates;
    \item government bond yields;
    \item yield-curve slopes;
    \item credit spreads;
    \item implied volatility;
    \item inflation;
    \item economic growth;
    \item monetary-policy variables;
    \item foreign-exchange rates;
    \item commodity prices.
\end{itemize}

If $R_{f,t}$ denotes the relevant risk-free return, an excess return is

\begin{equation}
\boxed{
R_{i,t}^{e}
=
R_{i,t}-R_{f,t}.
}
\end{equation}

Macroeconomic variables may subsequently require vintage reconstruction because
the values observed today may differ from the originally released values.


\subsubsection{Currency Data}

Foreign-exchange data are required when assets and portfolio NAV are expressed
in different currencies.

Suppose an asset earns return $R_{i,t}^{F}$ in foreign currency $F$, while the
foreign currency earns return $R_{FX,t}$ relative to domestic currency $D$.

Then

\begin{equation}
1+R_{i,t}^{D}
=
\left(
1+R_{i,t}^{F}
\right)
\left(
1+R_{FX,t}
\right),
\end{equation}

so that

\begin{equation}
\boxed{
R_{i,t}^{D}
=
R_{i,t}^{F}
+
R_{FX,t}
+
R_{i,t}^{F}R_{FX,t}.
}
\end{equation}

The exchange-rate quotation convention must be specified consistently.


% ------------------------------------------------------------
\subsection{Industry and Country Classifications}
\label{subsec:industry_country_classifications}
% ------------------------------------------------------------

Industry, sector, and country classifications provide categorical information
about securities and firms.

They may later be used for:

\begin{itemize}
    \item universe construction;
    \item neutralization;
    \item portfolio constraints;
    \item risk modeling;
    \item performance attribution.
\end{itemize}

Let

\begin{equation}
g(i,t)
\end{equation}

denote the industry classification of security $i$ at time $t$.


\subsubsection{Industry Representation}

For $K$ mutually exclusive industries, define

\begin{equation}
D_{i,k,t}
=
\begin{cases}
1,
&
\text{if security }i
\text{ belongs to industry }k,
\\
0,
&
\text{otherwise}.
\end{cases}
\end{equation}

The classification matrix is therefore

\begin{equation}
\boxed{
D_t
\in
\mathbb{R}^{N_t\times K}.
}
\end{equation}

Industry classifications are often hierarchical:

\begin{equation}
\boxed{
\text{Sector}
\rightarrow
\text{Industry Group}
\rightarrow
\text{Industry}
\rightarrow
\text{Sub-Industry}.
}
\end{equation}

The appropriate level of granularity depends on the intended application.


\subsubsection{Time Variation}

Industry membership may vary through time:

\begin{equation}
g(i,t_1)
\neq
g(i,t_2).
\end{equation}

Changes may result from:

\begin{itemize}
    \item changes in business activity;
    \item restructurings;
    \item mergers and acquisitions;
    \item classification-provider revisions.
\end{itemize}

Historical classification changes must therefore be retained rather than
replaced by the latest classification.


\subsubsection{Country Classification}

Country assignment may also be defined in several ways, including:

\begin{itemize}
    \item country of incorporation;
    \item headquarters location;
    \item country of primary listing;
    \item provider classification;
    \item dominant economic exposure.
\end{itemize}

The appropriate convention should be defined according to the economic purpose
of the strategy.


% ------------------------------------------------------------
\subsection{Corporate Actions}
\label{subsec:corporate_actions}
% ------------------------------------------------------------

Corporate actions modify the economic value, number, or identity of securities
and must therefore be represented explicitly in the data architecture.

Relevant events include:

\begin{itemize}
    \item cash dividends;
    \item special dividends;
    \item stock splits;
    \item reverse stock splits;
    \item stock dividends;
    \item rights issues;
    \item spin-offs;
    \item mergers;
    \item acquisitions;
    \item tender offers;
    \item share-class conversions;
    \item delistings.
\end{itemize}


\subsubsection{Stock Splits}

Suppose a company performs an $m$-for-$n$ stock split.

Ignoring other market movements,

\begin{equation}
Q_{i,t}^{\mathrm{after}}
=
Q_{i,t^-}
\frac{m}{n},
\end{equation}

while

\begin{equation}
P_{i,t}^{\mathrm{after}}
=
P_{i,t^-}
\frac{n}{m}.
\end{equation}

Hence,

\begin{equation}
Q_{i,t}^{\mathrm{after}}
P_{i,t}^{\mathrm{after}}
=
Q_{i,t^-}
P_{i,t^-}.
\end{equation}

The stock split changes the number of shares and quoted price but does not
itself create shareholder wealth.


\subsubsection{Cash Dividends}

A cash dividend represents an economic distribution to shareholders.

If $D_{i,t}$ is paid during the period, the total return is

\begin{equation}
R_{i,t}^{\mathrm{total}}
=
\frac{
P_{i,t}+D_{i,t}
}{
P_{i,t-1}
}
-1.
\end{equation}


\subsubsection{Mergers, Acquisitions, and Delistings}

A corporate event may replace an existing security with:

\begin{itemize}
    \item cash;
    \item another security;
    \item a combination of cash and securities.
\end{itemize}

The economic conversion must be represented before the original security is
removed from portfolio accounting.

For delistings, the required sequence is

\begin{equation}
\boxed{
\text{Existing Position}
\rightarrow
\text{Terminal Economic Payoff}
\rightarrow
\text{Position Removal}.
}
\end{equation}

as discussed in Section~\ref{subsubsec:delisting_returns}.


\subsubsection{Corporate-Action Event Table}

A generic corporate-action record may contain

\begin{equation}
\boxed{
\left(
\mathrm{SecurityID},
\mathrm{EventType},
t^{\mathrm{announcement}},
t^{\mathrm{ex}},
t^{\mathrm{effective}},
t^{\mathrm{payment}},
\mathrm{Terms}
\right).
}
\end{equation}

These timestamps have distinct economic meanings and must therefore be stored
separately.


% ------------------------------------------------------------
\subsection{Liquidity, Borrowing, and Transaction Cost Data}
\label{subsec:liquidity_borrowing_transaction_cost_data}
% ------------------------------------------------------------

A predictive strategy is economically meaningful only if the resulting
portfolio can be implemented under realistic trading conditions.

The raw-data layer should therefore include information describing:

\begin{itemize}
    \item trading liquidity;
    \item execution conditions;
    \item securities lending;
    \item explicit market fees.
\end{itemize}


\subsubsection{Liquidity Measures}

Relevant variables include:

\begin{itemize}
    \item traded volume;
    \item traded notional;
    \item Average Daily Volume;
    \item bid--ask spread;
    \item quoted depth;
    \item turnover ratio;
    \item number of trades.
\end{itemize}

Average Daily Volume in shares over lookback window $L$ may be defined as

\begin{equation}
\mathrm{ADV}_{i,t}^{\mathrm{shares}}
=
\frac{1}{L}
\sum_{\ell=0}^{L-1}
V_{i,t-\ell}.
\end{equation}

Currency-denominated ADV is

\begin{equation}
\boxed{
\mathrm{ADV}_{i,t}^{\$}
=
\frac{1}{L}
\sum_{\ell=0}^{L-1}
P_{i,t-\ell}
V_{i,t-\ell}.
}
\end{equation}

If the strategy trades notional $Q_{i,t}^{\$}$, its participation rate is

\begin{equation}
\boxed{
PR_{i,t}
=
\frac{
|Q_{i,t}^{\$}|
}{
\mathrm{ADV}_{i,t}^{\$}
}.
}
\end{equation}


\subsubsection{Bid--Ask Data}

Let

\begin{equation}
P_{i,t}^{\mathrm{bid}}
\qquad\text{and}\qquad
P_{i,t}^{\mathrm{ask}}
\end{equation}

denote the best bid and ask.

The mid-price is

\begin{equation}
P_{i,t}^{\mathrm{mid}}
=
\frac{
P_{i,t}^{\mathrm{ask}}
+
P_{i,t}^{\mathrm{bid}}
}{2}.
\end{equation}

The quoted spread is

\begin{equation}
\mathrm{Spread}_{i,t}
=
P_{i,t}^{\mathrm{ask}}
-
P_{i,t}^{\mathrm{bid}},
\end{equation}

and the relative quoted spread is

\begin{equation}
\boxed{
\mathrm{RelativeSpread}_{i,t}
=
\frac{
P_{i,t}^{\mathrm{ask}}
-
P_{i,t}^{\mathrm{bid}}
}{
P_{i,t}^{\mathrm{mid}}
}.
}
\end{equation}

Under the simplified assumption of execution at the best bid or ask, the
one-way spread cost relative to the mid-price is approximately

\begin{equation}
c_{i,t}^{\mathrm{spread}}
\approx
\frac{1}{2}
\mathrm{RelativeSpread}_{i,t}.
\end{equation}


\subsubsection{Securities-Lending Data}

A short position generally requires borrowing the security before it can be
sold.

Relevant variables include:

\begin{itemize}
    \item borrow availability;
    \item available quantity;
    \item annualized borrow fee;
    \item utilization;
    \item hard-to-borrow status;
    \item recalls.
\end{itemize}

Let $b_{i,t}$ denote the annualized borrow rate.

For a short position with absolute market value
$V_{i,t}^{\mathrm{short}}$ held for $\Delta t$ years, a simplified borrowing
cost is

\begin{equation}
\boxed{
BC_{i,t}
\approx
b_{i,t}
V_{i,t}^{\mathrm{short}}
\Delta t.
}
\end{equation}


\subsubsection{Transaction-Cost Inputs}

Transaction costs themselves are outputs of an implementation model.

The raw data used to estimate them may include

\begin{equation}
\mathcal{D}_{i,t}^{TC}
=
\left\{
\mathrm{Spread}_{i,t},
\mathrm{ADV}_{i,t},
\sigma_{i,t},
\mathrm{Depth}_{i,t},
\mathrm{Fees}_{i,t},
\ldots
\right\}.
\end{equation}

The conceptual sequence is therefore

\begin{equation}
\boxed{
\underbrace{
\mathcal{D}_{i,t}^{TC}
}_{\text{Raw Implementation Data}}
\longrightarrow
\underbrace{
C_{i,t}(\cdot)
}_{\text{Transaction-Cost Model}}
\longrightarrow
\underbrace{
TC_{i,t}
}_{\text{Estimated Cost}}.
}
\end{equation}

The transaction-cost model itself will be developed in the portfolio
implementation section.


\subsubsection{Liquidity and Capacity}

If portfolio NAV is $A_t$, a weight change $\Delta w_{i,t}$ corresponds to
traded notional

\begin{equation}
Q_{i,t}^{\$}
=
A_t
|\Delta w_{i,t}|.
\end{equation}

Hence,

\begin{equation}
PR_{i,t}
=
\frac{
A_t|\Delta w_{i,t}|
}{
\mathrm{ADV}_{i,t}^{\$}
}.
\end{equation}

Thus, holding portfolio construction fixed,

\begin{equation}
\boxed{
A_t \uparrow
\quad\Rightarrow\quad
PR_{i,t} \uparrow
\quad\Rightarrow\quad
\text{Implementation Difficulty Typically Increases}.
}
\end{equation}

This establishes the direct connection between raw liquidity data and strategy
capacity.


% ------------------------------------------------------------
\subsection{Raw Data Architecture}
\label{subsec:raw_data_architecture}
% ------------------------------------------------------------

The raw-data layer combines observations measured at different entity levels:

\begin{equation}
\boxed{
\begin{array}{ccc}
\text{Market Data}
&
\rightarrow
&
\text{Security / Listing Level}
\\[0.15cm]
\text{Fundamentals}
&
\rightarrow
&
\text{Company / Accounting-Period Level}
\\[0.15cm]
\text{Analyst Estimates}
&
\rightarrow
&
\text{Company / Forecast-Horizon Level}
\\[0.15cm]
\text{Alternative Data}
&
\rightarrow
&
\text{Entity / Event Level}
\\[0.15cm]
\text{Risk Factors}
&
\rightarrow
&
\text{Market / Factor Level}
\\[0.15cm]
\text{Classifications}
&
\rightarrow
&
\text{Company / Security Level}
\\[0.15cm]
\text{Corporate Actions}
&
\rightarrow
&
\text{Security / Event Level}
\\[0.15cm]
\text{Implementation Data}
&
\rightarrow
&
\text{Security / Market Level}
\end{array}
}
\end{equation}

These sources cannot generally be merged through a simple equality join because
they differ both in entity structure and in time representation.

A robust backtesting architecture must therefore preserve two dimensions:

\begin{equation}
\boxed{
\text{Entity Identity}
\qquad\text{and}\qquad
\text{Information Timing}.
}
\end{equation}

The raw-data stage answers the question

\begin{equation}
\boxed{
\text{What historical observations do we possess?}
}
\end{equation}

The next stage, Point-in-Time Data Construction, answers the distinct and more
important backtesting question

\begin{equation}
\boxed{
\text{Which of these observations were actually available to the investor at time }t?
}
\end{equation}


% ============================================================
\section{Point-in-Time Data Construction}
\label{sec:point_in_time_data}
% ============================================================

Raw historical databases generally describe what is known \emph{today} about
past observations. A backtest, however, requires a different object: the
information that would actually have been available to an investor at each
historical date.

The objective of point-in-time construction is therefore to transform the
complete historical database

\begin{equation}
\mathcal{D}
\end{equation}

into a sequence of historical information sets

\begin{equation}
\boxed{
\mathcal{D}
\longrightarrow
\left\{
\mathcal{I}_t
\right\}_{t=1}^{T},
}
\end{equation}

where $\mathcal{I}_t$ contains only observations that were genuinely observable
by the investment process at time $t$.

For any data observation $j$, define its availability time as

\begin{equation}
t_j^{\mathrm{avail}}.
\end{equation}

Then observation $j$ may enter the information set at time $t$ only if

\begin{equation}
\boxed{
t_j^{\mathrm{avail}}
\leq
t.
}
\end{equation}

Thus,

\begin{equation}
\boxed{
\mathcal{I}_t
=
\left\{
D_j\in\mathcal{D}
:
t_j^{\mathrm{avail}}
\leq t
\right\}.
}
\end{equation}

This condition is the fundamental rule underlying point-in-time data
engineering.

Its practical implementation requires careful treatment of dates, reporting
lags, revisions, entity membership, and historical joins.


% ------------------------------------------------------------
\subsection{Observation Dates, Publication Dates, and Effective Dates}
\label{subsec:pit_dates}
% ------------------------------------------------------------

A single financial observation may be associated with several different dates.
These dates describe different economic concepts and should not be treated as
interchangeable.

For observation $j$, define:

\begin{equation}
t_j^{\mathrm{obs}},
\qquad
t_j^{\mathrm{pub}},
\qquad
t_j^{\mathrm{eff}},
\qquad
t_j^{\mathrm{avail}}.
\end{equation}

These correspond respectively to:

\begin{itemize}
    \item $t_j^{\mathrm{obs}}$: the date or period to which the observation
    economically refers;

    \item $t_j^{\mathrm{pub}}$: the date on which the information is officially
    published or announced;

    \item $t_j^{\mathrm{eff}}$: the date on which an event or classification
    becomes economically or operationally effective;

    \item $t_j^{\mathrm{avail}}$: the earliest date and time at which the
    investment process could actually have used the information.
\end{itemize}

These dates may coincide for some datasets but differ substantially for others.


\subsubsection{Market Data Example}

For a daily closing price,

\begin{equation}
P_{i,t}^{\mathrm{close}},
\end{equation}

the observation date and publication date may both correspond approximately to
trading day $t$.

However, the information becomes usable only after the closing price has
actually been observed.

Hence,

\begin{equation}
t^{\mathrm{obs}}
=
t,
\end{equation}

but the exact

\begin{equation}
t^{\mathrm{avail}}
\end{equation}

depends on the timestamp of the signal and execution process.


\subsubsection{Fundamental Data Example}

Suppose a company reports financial results for the fiscal year ending on
December 31.

Then

\begin{equation}
t^{\mathrm{obs}}
=
\text{December 31}.
\end{equation}

Suppose the results are released on March 15 of the following year.

Then

\begin{equation}
t^{\mathrm{pub}}
=
\text{March 15}.
\end{equation}

The value cannot be used by the strategy between December 31 and March 14,
even though the accounting observation refers economically to December 31.

Thus,

\begin{equation}
\boxed{
t^{\mathrm{obs}}
<
t^{\mathrm{pub}}
\leq
t^{\mathrm{avail}}.
}
\end{equation}


\subsubsection{Corporate Action Example}

Suppose an index provider announces on June 1 that stock $i$ will enter the
index on June 10.

Then

\begin{equation}
t^{\mathrm{pub}}
=
\text{June 1},
\end{equation}

while

\begin{equation}
t^{\mathrm{eff}}
=
\text{June 10}.
\end{equation}

The investor may know on June 1 that the future membership change will occur,
but the security does not become an index constituent until June 10.

Publication and effective dates therefore answer different questions.


\subsubsection{General Principle}

For every dataset, the backtester should identify the date corresponding to the
question

\begin{equation}
\boxed{
\text{When could the strategy first have used this observation?}
}
\end{equation}

This date is $t^{\mathrm{avail}}$.

The point-in-time dataset should then be indexed primarily by availability
rather than purely by economic reference date.


% ------------------------------------------------------------
\subsection{Reporting Lags}
\label{subsec:reporting_lags}
% ------------------------------------------------------------

Many variables become observable only after a delay relative to the period they
describe.

Define the reporting lag of observation $j$ as

\begin{equation}
\boxed{
L_j
=
t_j^{\mathrm{avail}}
-
t_j^{\mathrm{obs}}.
}
\end{equation}

For fundamental information, this lag can be substantial.

For example, a quarterly accounting observation may correspond to

\begin{equation}
t_j^{\mathrm{obs}}
=
\text{March 31},
\end{equation}

but only become available on

\begin{equation}
t_j^{\mathrm{avail}}
=
\text{May 5}.
\end{equation}

The reporting lag is therefore approximately

\begin{equation}
L_j
=
35\text{ days}.
\end{equation}


\subsubsection{Actual vs.\ Assumed Reporting Lags}

The preferred approach is to use the true historical publication or
availability timestamp for each observation.

If these timestamps are unavailable, a conservative lag may be imposed
artificially.

For example,

\begin{equation}
X_{i,\tau}
\text{ becomes usable at }
\tau+L,
\end{equation}

where $L$ is a predefined reporting lag.

The backtesting rule becomes

\begin{equation}
\boxed{
X_{i,\tau}
\in
\mathcal{I}_t
\quad
\Longleftrightarrow
\quad
\tau+L
\leq
t.
}
\end{equation}

This approach is less precise than using historical release dates but is
preferable to assuming that accounting data were known at the fiscal
period-end date.


\subsubsection{Cross-Sectional Variation in Reporting Lags}

Reporting lags need not be identical across companies.

Firm $i$ may report after

\begin{equation}
L_i=25\text{ days},
\end{equation}

while firm $j$ may report after

\begin{equation}
L_j=55\text{ days}.
\end{equation}

Applying a common lag to both companies can therefore distort the historical
information set.

Whenever available, observation-specific timestamps are preferable:

\begin{equation}
\boxed{
L_{i,\tau}
=
t_{i,\tau}^{\mathrm{avail}}
-
t_{i,\tau}^{\mathrm{obs}}.
}
\end{equation}


\subsubsection{Intraday Reporting Lags}

For higher-frequency strategies, dates alone may be insufficient.

Suppose an earnings announcement occurs on day $t$ at 17:30, after the market
close.

A daily database may record only

\begin{equation}
t^{\mathrm{pub}}=t,
\end{equation}

but a strategy trading at the close of day $t$ could not have used the
announcement.

In such applications, availability should be represented by a timestamp:

\begin{equation}
\boxed{
t^{\mathrm{avail}}
=
(\text{date},\text{time},\text{timezone}).
}
\end{equation}

Ignoring intraday timestamps can create look-ahead bias even when calendar
dates appear correct.


% ------------------------------------------------------------
\subsection{Restatements}
\label{subsec:restatements}
% ------------------------------------------------------------

Historical observations may be revised after their initial publication.

This is common for:

\begin{itemize}
    \item accounting statements;
    \item macroeconomic releases;
    \item index and classification data;
    \item analyst databases;
    \item corporate-action records.
\end{itemize}

Suppose a company initially reports accounting value

\begin{equation}
X_{i,\tau}^{(0)}
\end{equation}

at time

\begin{equation}
t_{i,\tau}^{(0)}.
\end{equation}

The same observation is later restated as

\begin{equation}
X_{i,\tau}^{(1)}
\end{equation}

at time

\begin{equation}
t_{i,\tau}^{(1)}
>
t_{i,\tau}^{(0)}.
\end{equation}

Then the historically correct value depends on the backtest date:

\begin{equation}
X_{i,\tau}^{\mathrm{PIT}}(t)
=
\begin{cases}
\text{unavailable},
&
t<t_{i,\tau}^{(0)},
\\[0.2cm]
X_{i,\tau}^{(0)},
&
t_{i,\tau}^{(0)}
\leq
t
<
t_{i,\tau}^{(1)},
\\[0.2cm]
X_{i,\tau}^{(1)},
&
t
\geq
t_{i,\tau}^{(1)}.
\end{cases}
\end{equation}

More generally, if an observation has revisions

\begin{equation}
X_j^{(0)},
X_j^{(1)},
\ldots,
X_j^{(M)}
\end{equation}

with availability times

\begin{equation}
t_j^{(0)}
<
t_j^{(1)}
<
\cdots
<
t_j^{(M)},
\end{equation}

the valid value at time $t$ is

\begin{equation}
\boxed{
X_j^{\mathrm{PIT}}(t)
=
X_j^{(m^*)},
\qquad
m^*
=
\max
\left\{
m:
t_j^{(m)}
\leq
t
\right\}.
}
\end{equation}


\subsubsection{Vintage Data}

A database containing all historical versions of an observation is commonly
referred to as a vintage or point-in-time database.

Conceptually,

\begin{equation}
\boxed{
\text{Observation}
+
\text{Revision Version}
+
\text{Availability Timestamp}.
}
\end{equation}

Using the latest value stored today,

\begin{equation}
X_j^{\mathrm{latest}},
\end{equation}

for all historical dates incorrectly gives the backtest information from future
revisions.


\subsubsection{Restatement Policy}

The backtester should therefore establish an explicit restatement policy.

The preferred rule is:

\begin{equation}
\boxed{
\text{At time }t,
\text{ use the latest version that was available at }t.
}
\end{equation}

If historical vintages are unavailable, the limitation should be documented
explicitly, since the resulting backtest may contain restatement bias.


% ------------------------------------------------------------
\subsection{As-Of Joins}
\label{subsec:asof_joins}
% ------------------------------------------------------------

Point-in-time datasets frequently combine observations recorded at different
frequencies.

For example:

\begin{itemize}
    \item prices may be daily;
    \item fundamentals may be quarterly;
    \item analyst estimates may change irregularly;
    \item classifications may change occasionally.
\end{itemize}

A standard equality join,

\begin{equation}
t_X=t_Y,
\end{equation}

is therefore generally inappropriate.

Instead, the backtester often requires an \emph{as-of join}.

For security $i$ and portfolio date $t$, define the latest admissible
observation time as

\begin{equation}
\boxed{
\tau_i^*(t)
=
\max
\left\{
\tau:
t_{i,\tau}^{\mathrm{avail}}
\leq
t
\right\}.
}
\end{equation}

The point-in-time value used by the strategy is then

\begin{equation}
\boxed{
X_{i,t}^{\mathrm{PIT}}
=
X_{i,\tau_i^*(t)}.
}
\end{equation}

Thus, the as-of join answers the question:

\begin{equation}
\boxed{
\text{What was the most recent observation available as of time }t?
}
\end{equation}


\subsubsection{Backward As-Of Join}

For predictive backtesting, the join direction should normally be backward:

\begin{equation}
\boxed{
t^{\mathrm{source}}
\leq
t^{\mathrm{decision}}.
}
\end{equation}

A forward join,

\begin{equation}
t^{\mathrm{source}}
>
t^{\mathrm{decision}},
\end{equation}

would introduce future information.


\subsubsection{Example}

Suppose a company reports quarterly earnings on

\begin{equation}
t_1=\text{May 5}
\end{equation}

and subsequently on

\begin{equation}
t_2=\text{August 4}.
\end{equation}

For every backtest date satisfying

\begin{equation}
\text{May 5}
\leq
t
<
\text{August 4},
\end{equation}

the valid fundamental observation is the May 5 release.

Thus,

\begin{equation}
X_{i,t}^{\mathrm{PIT}}
=
X_{i,t_1}.
\end{equation}

Only from August 4 onward may the newer observation replace it.


\subsubsection{Maximum Staleness}

The fact that an observation is the latest historically available value does
not necessarily mean that it remains economically usable indefinitely.

A maximum age may therefore be imposed.

Let

\begin{equation}
A_{i,t}
=
t
-
t_{i,\tau_i^*(t)}^{\mathrm{avail}}
\end{equation}

denote data age.

Then one may require

\begin{equation}
\boxed{
A_{i,t}
\leq
A_{\max}.
}
\end{equation}

Otherwise the observation is treated as stale or missing.

This prevents an extremely old fundamental, estimate, or classification record
from being carried forward indefinitely.


\subsubsection{Entity Keys}

An as-of join must generally be performed within the appropriate entity.

For example,

\begin{equation}
\boxed{
\text{Join Key}
=
(\mathrm{SecurityID},t)
}
\end{equation}

for security-level data, while company-level fundamental data may require

\begin{equation}
\boxed{
\text{CompanyID}
\longrightarrow
\text{SecurityID}
}
\end{equation}

before being joined to market observations.

Entity resolution and temporal alignment are therefore jointly necessary.


% ------------------------------------------------------------
\subsection{Historical Universe Membership}
\label{subsec:historical_universe_membership}
% ------------------------------------------------------------

The investment universe itself must be reconstructed point-in-time.

Let

\begin{equation}
I_{i,t}^{\mathrm{member}}
=
\begin{cases}
1,
&
\text{if security }i
\text{ belongs to the relevant universe at }t,
\\
0,
&
\text{otherwise}.
\end{cases}
\end{equation}

Then

\begin{equation}
\boxed{
\mathcal{U}_t
=
\left\{
i:
I_{i,t}^{\mathrm{member}}=1
\right\}.
}
\end{equation}

Historical membership can vary because of:

\begin{itemize}
    \item IPOs;
    \item delistings;
    \item mergers and acquisitions;
    \item index additions and deletions;
    \item changes in exchange eligibility;
    \item changes in liquidity or size requirements;
    \item changes in security type;
    \item shortability constraints.
\end{itemize}


\subsubsection{Entry and Exit Intervals}

A convenient representation stores a membership interval for each security:

\begin{equation}
\left[
t_i^{\mathrm{entry}},
t_i^{\mathrm{exit}}
\right).
\end{equation}

Then

\begin{equation}
I_{i,t}^{\mathrm{member}}
=
\begin{cases}
1,
&
t_i^{\mathrm{entry}}
\leq
t
<
t_i^{\mathrm{exit}},
\\
0,
&
\text{otherwise}.
\end{cases}
\end{equation}


\subsubsection{Index Membership}

If a strategy trades an index universe, current constituents must not be
projected backward through history.

Instead,

\begin{equation}
\boxed{
\mathcal{U}_t
=
\text{constituents valid at time }t.
}
\end{equation}

A distinction may also be required between announcement and effective dates.

Suppose an addition is announced at

\begin{equation}
t^{\mathrm{announce}}
\end{equation}

and becomes effective at

\begin{equation}
t^{\mathrm{effective}}.
\end{equation}

Then whether the security can enter the strategy before the effective date
depends on whether the strategy is explicitly designed to trade announced
future index changes.

The rule should be defined ex ante and applied consistently.


\subsubsection{Rule-Based Universes}

For a rule-based universe, membership may itself be computed from historical
variables.

Suppose securities must satisfy

\begin{equation}
\mathrm{MCap}_{i,t}
>
M_{\min}
\end{equation}

and

\begin{equation}
\mathrm{ADV}_{i,t}
>
ADV_{\min}.
\end{equation}

Then

\begin{equation}
I_{i,t}^{\mathrm{eligible}}
=
\mathbbm{1}
\left(
\mathrm{MCap}_{i,t}>M_{\min}
\right)
\mathbbm{1}
\left(
\mathrm{ADV}_{i,t}>ADV_{\min}
\right).
\end{equation}

Every variable entering the eligibility rule must itself be point-in-time.

Thus, universe construction satisfies the same information constraint as signal
construction:

\begin{equation}
\boxed{
\mathcal{U}_t
=
g(\mathcal{I}_t).
}
\end{equation}


% ------------------------------------------------------------
\subsection{Delisted Securities}
\label{subsec:pit_delisted_securities}
% ------------------------------------------------------------

Delisted securities require special treatment because their disappearance
affects both historical universe membership and portfolio accounting.

A delisted security should not simply disappear retrospectively from the
historical database.

Before delisting,

\begin{equation}
i\in\mathcal{U}_t.
\end{equation}

After the delisting becomes effective,

\begin{equation}
i\notin\mathcal{U}_t.
\end{equation}

The transition must occur at the historically correct date.


\subsubsection{Universe Treatment}

Suppose security $i$ delists at time $T_i$.

Its historical membership may be represented as

\begin{equation}
I_{i,t}^{\mathrm{member}}
=
\begin{cases}
1,
&
t<T_i,
\\
0,
&
t\geq T_i.
\end{cases}
\end{equation}

However, removing the security from future universes does not by itself solve
the portfolio-accounting problem.


\subsubsection{Economic Payoff Before Removal}

If the strategy holds the security immediately before delisting, the terminal
economic payoff must first be incorporated.

As established in Section~\ref{subsubsec:delisting_returns},

\begin{equation}
\boxed{
\text{Existing Position}
\rightarrow
R_{i,T_i}^{\mathrm{delist}}
\rightarrow
\text{Terminal Proceeds}
\rightarrow
\text{Position Removal}.
}
\end{equation}

The historical data pipeline should therefore retain:

\begin{itemize}
    \item the final tradable price;
    \item the delisting or terminal return, when available;
    \item the delisting date;
    \item the reason or event type, when relevant.
\end{itemize}

Replacing the delisting event with a missing value can materially bias
historical returns.


\subsubsection{Missing Delisting Information}

In some datasets, the final economic payoff is unavailable.

The backtester must then adopt an explicit policy rather than silently dropping
the observation.

Possible treatments depend on the database and research objective, but any
assumption should be documented because the resulting performance may be
sensitive to the treatment of missing terminal returns.


% ------------------------------------------------------------
\subsection{Preventing Look-Ahead Bias}
\label{subsec:pit_preventing_lookahead}
% ------------------------------------------------------------

The complete purpose of point-in-time construction is to guarantee that no
future information enters a historical investment decision.

Let $X_{i,t}^{(k)}$ denote feature input $k$ used for security $i$ at time $t$.

The fundamental validity condition is

\begin{equation}
\boxed{
t_{i}^{(k),\mathrm{avail}}
\leq
t
\qquad
\forall i,k.
}
\end{equation}

Equivalently,

\begin{equation}
X_t
\in
\sigma(\mathcal{I}_t),
\end{equation}

where $\sigma(\mathcal{I}_t)$ denotes the information generated by observations
available up to time $t$.


\subsubsection{Point-in-Time Construction Rule}

For each required variable, the backtester should conceptually execute:

\begin{equation}
\boxed{
X_{i,t}^{\mathrm{PIT}}
=
\text{latest valid version of }
X_i
\text{ whose availability time satisfies }
t^{\mathrm{avail}}\leq t.
}
\end{equation}

This single rule simultaneously governs:

\begin{itemize}
    \item fundamental data;
    \item analyst estimates;
    \item alternative data;
    \item classifications;
    \item macroeconomic variables;
    \item corporate-action information;
    \item universe membership.
\end{itemize}


\subsubsection{Decision Time vs.\ Calendar Date}

A calendar date alone may not define the information set sufficiently.

Let

\begin{equation}
t^{\mathrm{decision}}
\end{equation}

denote the exact decision timestamp.

Then the admissibility rule becomes

\begin{equation}
\boxed{
t_j^{\mathrm{avail}}
\leq
t^{\mathrm{decision}}.
}
\end{equation}

For example, information released after the market close on day $t$ may belong
to the information set for day $t+1$, but not to a strategy executed at the
close of day $t$.


\subsubsection{Safe Pipeline Ordering}

A robust backtesting process should preserve the chronological order

\begin{equation}
\boxed{
\begin{aligned}
\text{Raw Observation}
&\rightarrow
\text{Availability Timestamp}
\rightarrow
\text{Point-in-Time Join}
\\[0.15cm]
&\rightarrow
\text{Feature Construction}
\rightarrow
\text{Signal}
\rightarrow
\text{Portfolio Decision}
\end{aligned}
}
\end{equation}

Crucially, point-in-time filtering should occur \emph{before} any operation
that could combine past and future observations.


\subsubsection{Common Point-in-Time Failures}

Common implementation errors include:

\begin{itemize}
    \item joining fundamentals by fiscal-period date rather than publication
    date;

    \item forward-filling observations before imposing availability dates;

    \item using the latest restated accounting value historically;

    \item using today's industry classification throughout history;

    \item using current index constituents for past dates;

    \item including a security before its historical listing or eligibility
    date;

    \item removing securities retrospectively because they later delisted;

    \item using macroeconomic values after revisions rather than real-time
    vintages;

    \item performing nearest-date joins that select future observations;

    \item ignoring intraday release timestamps when the strategy trades at
    intraday or close-to-close frequency.
\end{itemize}


\subsubsection{Point-in-Time Validation Tests}

Because point-in-time errors can be difficult to detect from performance alone,
the data pipeline should include explicit validation tests.

For every joined observation $j$ used at decision time $t$,

\begin{equation}
\boxed{
t_j^{\mathrm{avail}}
\leq
t
}
\end{equation}

should hold programmatically.

One may define the violation indicator

\begin{equation}
V_j
=
\begin{cases}
1,
&
t_j^{\mathrm{avail}}
>
t_j^{\mathrm{decision}},
\\
0,
&
\text{otherwise}.
\end{cases}
\end{equation}

A valid dataset should satisfy

\begin{equation}
\boxed{
\sum_j V_j
=
0.
}
\end{equation}

Additional validation checks should verify that:

\begin{itemize}
    \item source observations are sorted chronologically;
    \item as-of joins always operate backward;
    \item revisions are selected according to historical availability;
    \item entity mappings are valid at the historical date;
    \item no security enters the universe before its eligibility date;
    \item delisted securities remain present until their historical exit date;
    \item stale observations are identified according to predefined rules.
\end{itemize}


\subsubsection{The Point-in-Time Information Set}

After all temporal filters, revisions, memberships, and joins have been
resolved, the final point-in-time dataset at decision date $t$ can be written
abstractly as

\begin{equation}
\boxed{
\mathcal{D}_t^{\mathrm{PIT}}
=
\left\{
D_j:
t_j^{\mathrm{avail}}
\leq
t,
\;
D_j
\text{ is the historically valid version at }t
\right\}.
}
\end{equation}

The investment process can then safely construct features as

\begin{equation}
X_t
=
g
\left(
\mathcal{D}_t^{\mathrm{PIT}}
\right),
\end{equation}

signals as

\begin{equation}
s_t
=
f(X_t),
\end{equation}

and portfolio weights as

\begin{equation}
w_t
=
\phi(s_t).
\end{equation}

Thus,

\begin{equation}
\boxed{
\mathcal{D}
\rightarrow
\mathcal{D}_t^{\mathrm{PIT}}
\rightarrow
X_t
\rightarrow
s_t
\rightarrow
w_t
}
\end{equation}

preserves the fundamental information constraint

\begin{equation}
\boxed{
w_t=f(\mathcal{I}_t).
}
\end{equation}

Point-in-time construction therefore represents the temporal firewall between
the historical database known today and the historical information set that
could genuinely have been used by the strategy.

% ============================================================
\section{Investment Universe Construction}
\label{sec:investment_universe}
% ============================================================

Before constructing investment signals, the backtester must determine which
securities are eligible to participate in the strategy at each decision date.

Let

\begin{equation}
\mathcal{U}_t^{\mathrm{raw}}
\end{equation}

denote the complete set of securities represented in the point-in-time
database at time $t$.

The investment universe is obtained by applying a collection of eligibility
rules:

\begin{equation}
\boxed{
\mathcal{U}_t
=
\left\{
i\in\mathcal{U}_t^{\mathrm{raw}}
:
E_{i,t}=1
\right\},
}
\end{equation}

where

\begin{equation}
E_{i,t}
\in
\{0,1\}
\end{equation}

is the overall eligibility indicator for security $i$.

If the universe is defined through $K$ individual filters, let

\begin{equation}
E_{i,t}^{(k)}
\in
\{0,1\},
\qquad
k=1,\ldots,K,
\end{equation}

denote whether security $i$ satisfies eligibility condition $k$.

If all conditions are required simultaneously,

\begin{equation}
\boxed{
E_{i,t}
=
\prod_{k=1}^{K}
E_{i,t}^{(k)}.
}
\end{equation}

Thus,

\begin{equation}
E_{i,t}=1
\end{equation}

if and only if security $i$ satisfies every required eligibility condition.

Typical filters concern:

\begin{itemize}
    \item geography and exchange;
    \item security type;
    \item market capitalization;
    \item liquidity;
    \item price;
    \item listing history;
    \item data availability;
    \item shortability.
\end{itemize}

The investment universe must be reconstructed dynamically using only
point-in-time information. Universe construction is therefore itself part of
the investment process rather than a static preprocessing operation.

A crucial distinction is

\begin{equation}
\boxed{
\text{Eligible Security}
\neq
\text{Selected Security}.
}
\end{equation}

Universe construction determines which securities \emph{may} enter the
strategy. Signal construction and portfolio construction subsequently determine
which eligible securities actually receive non-zero portfolio weights.


% ------------------------------------------------------------
\subsection{Geographical and Exchange Filters}
\label{subsec:geographical_exchange_filters}
% ------------------------------------------------------------

A strategy may restrict its investment universe according to geographical
exposure, listing location, or trading venue.

Let

\begin{equation}
C_{i,t}
\end{equation}

denote the relevant country classification of security $i$ at time $t$ and let

\begin{equation}
\mathcal{C}^{\mathrm{eligible}}
\end{equation}

denote the set of eligible countries.

The geographical eligibility indicator is

\begin{equation}
E_{i,t}^{\mathrm{geo}}
=
\begin{cases}
1,
&
C_{i,t}\in\mathcal{C}^{\mathrm{eligible}},
\\
0,
&
\text{otherwise}.
\end{cases}
\end{equation}

For example, a European equity strategy might define

\begin{equation}
\mathcal{C}^{\mathrm{eligible}}
=
\{
\text{France},
\text{Germany},
\text{Italy},
\text{Spain},
\ldots
\}.
\end{equation}

However, the definition of country is not unique.

Possible conventions include:

\begin{itemize}
    \item country of incorporation;
    \item headquarters location;
    \item country of primary listing;
    \item provider-defined country;
    \item country of dominant economic exposure.
\end{itemize}

The convention should be specified explicitly and applied consistently.


\subsubsection{Exchange Filters}

The strategy may additionally restrict securities to approved exchanges.

Let

\begin{equation}
X_{i,t}
\end{equation}

denote the exchange on which security $i$ is listed and let

\begin{equation}
\mathcal{X}^{\mathrm{eligible}}
\end{equation}

denote the set of eligible exchanges.

Then

\begin{equation}
E_{i,t}^{\mathrm{exchange}}
=
\begin{cases}
1,
&
X_{i,t}\in\mathcal{X}^{\mathrm{eligible}},
\\
0,
&
\text{otherwise}.
\end{cases}
\end{equation}

Exchange filters may be useful for excluding securities traded on venues with
insufficient liquidity, unreliable data, limited accessibility, or different
market structures.

For securities with multiple listings, the backtester should define whether the
strategy operates on:

\begin{itemize}
    \item the primary listing;
    \item the most liquid listing;
    \item a predefined preferred listing;
    \item all eligible listings.
\end{itemize}

Care must be taken to avoid representing the same underlying economic exposure
multiple times unintentionally.


% ------------------------------------------------------------
\subsection{Security Type Filters}
\label{subsec:security_type_filters}
% ------------------------------------------------------------

Historical security databases may contain instruments that are not part of the
intended investment universe.

Depending on the strategy, these may include:

\begin{itemize}
    \item ordinary common shares;
    \item preferred shares;
    \item depositary receipts;
    \item ETFs;
    \item closed-end funds;
    \item REITs;
    \item units;
    \item warrants;
    \item rights;
    \item convertible securities;
    \item other structured instruments.
\end{itemize}

Let

\begin{equation}
S_{i,t}^{\mathrm{type}}
\end{equation}

denote the security type and let

\begin{equation}
\mathcal{S}^{\mathrm{eligible}}
\end{equation}

denote the admissible set.

Then

\begin{equation}
E_{i,t}^{\mathrm{type}}
=
\begin{cases}
1,
&
S_{i,t}^{\mathrm{type}}
\in
\mathcal{S}^{\mathrm{eligible}},
\\
0,
&
\text{otherwise}.
\end{cases}
\end{equation}

For a conventional common-equity strategy, one may for example restrict the
universe to ordinary common shares.

The precise rule depends on the investment mandate and should be defined before
examining strategy performance.


% ------------------------------------------------------------
\subsection{Market Capitalization Filters}
\label{subsec:market_capitalization_filters}
% ------------------------------------------------------------

Market capitalization filters are commonly used to exclude securities that are
too small to trade realistically or that fall outside the intended investment
segment.

The equity market capitalization of company $i$ can be written as

\begin{equation}
\boxed{
\mathrm{MCap}_{i,t}
=
P_{i,t}
N_{i,t}^{\mathrm{shares}},
}
\end{equation}

where $N_{i,t}^{\mathrm{shares}}$ denotes the relevant number of shares
outstanding.

A minimum capitalization rule is

\begin{equation}
E_{i,t}^{\mathrm{MCap}}
=
\begin{cases}
1,
&
\mathrm{MCap}_{i,t}
\geq
M_{\min},
\\
0,
&
\text{otherwise}.
\end{cases}
\end{equation}


\subsubsection{Full vs.\ Free-Float Market Capitalization}

For investability purposes, free-float capitalization may be more informative
than total capitalization.

Let

\begin{equation}
f_{i,t}^{\mathrm{float}}
\in[0,1]
\end{equation}

denote the fraction of shares considered freely tradable.

Then

\begin{equation}
\boxed{
\mathrm{FFMCap}_{i,t}
=
f_{i,t}^{\mathrm{float}}
\mathrm{MCap}_{i,t}.
}
\end{equation}

A strategy may therefore impose

\begin{equation}
\mathrm{FFMCap}_{i,t}
\geq
M_{\min}^{\mathrm{float}}.
\end{equation}


\subsubsection{Absolute and Relative Thresholds}

A capitalization filter need not use a fixed monetary threshold.

A relative rule may retain, for example, securities above a given
cross-sectional percentile:

\begin{equation}
\mathrm{MCap}_{i,t}
\geq
Q_{q,t}^{\mathrm{MCap}},
\end{equation}

where

\begin{equation}
Q_{q,t}^{\mathrm{MCap}}
\end{equation}

is the $q$-th percentile of the point-in-time market-capitalization
distribution.

Alternatively, the strategy may retain the $N$ largest eligible securities.

All quantities used in the ranking must be computed using the historical
cross-section available at time $t$.


% ------------------------------------------------------------
\subsection{Liquidity Filters}
\label{subsec:liquidity_filters}
% ------------------------------------------------------------

Market capitalization alone does not guarantee that a security can be traded
efficiently.

Liquidity filters attempt to remove securities for which realistic execution
would be difficult or excessively costly.

A common measure is Average Daily Volume in currency units:

\begin{equation}
\mathrm{ADV}_{i,t}^{\$}
=
\frac{1}{L}
\sum_{\ell=1}^{L}
P_{i,t-\ell}
V_{i,t-\ell}.
\end{equation}

A minimum-liquidity rule is

\begin{equation}
E_{i,t}^{\mathrm{ADV}}
=
\begin{cases}
1,
&
\mathrm{ADV}_{i,t}^{\$}
\geq
ADV_{\min},
\\
0,
&
\text{otherwise}.
\end{cases}
\end{equation}

Notice that the window above ends at $t-1$. This is appropriate when the
universe is constructed before trading on day $t$ and the complete trading
volume of day $t$ is not yet observable.


\subsubsection{Alternative Liquidity Measures}

Other liquidity filters may use:

\begin{itemize}
    \item median daily traded notional;
    \item turnover ratio;
    \item bid--ask spread;
    \item number of trading days;
    \item percentage of zero-return days;
    \item quoted depth;
    \item Amihud-type illiquidity measures.
\end{itemize}

For example, the Amihud illiquidity measure over $L$ observations can be written
as

\begin{equation}
\boxed{
\mathrm{ILLIQ}_{i,t}
=
\frac{1}{L}
\sum_{\ell=1}^{L}
\frac{
|R_{i,t-\ell}|
}{
\mathrm{DollarVolume}_{i,t-\ell}
}.
}
\end{equation}

Higher values indicate greater price movement per unit of traded notional and
therefore lower liquidity.


\subsubsection{Universe Filter vs.\ Portfolio Constraint}

Liquidity can enter the investment process at two distinct stages.

First, a universe filter may impose

\begin{equation}
\mathrm{ADV}_{i,t}^{\$}
\geq
ADV_{\min}.
\end{equation}

Second, portfolio construction may impose a position- or trade-size constraint
such as

\begin{equation}
\frac{
|\Delta w_{i,t}|A_t
}{
\mathrm{ADV}_{i,t}^{\$}
}
\leq
PR_{\max}.
\end{equation}

These operations should not be confused.

The first determines whether the security is eligible at all. The second
determines how much of an eligible security may realistically be traded.


% ------------------------------------------------------------
\subsection{Price Filters}
\label{subsec:price_filters}
% ------------------------------------------------------------

Strategies may exclude securities whose prices fall below a minimum threshold.

Let

\begin{equation}
P_{i,t}^{\mathrm{filter}}
\end{equation}

denote the price observable at the universe-construction time.

A simple rule is

\begin{equation}
E_{i,t}^{\mathrm{price}}
=
\begin{cases}
1,
&
P_{i,t}^{\mathrm{filter}}
\geq
P_{\min},
\\
0,
&
\text{otherwise}.
\end{cases}
\end{equation}

Low-priced securities may exhibit:

\begin{itemize}
    \item larger proportional bid--ask spreads;
    \item stronger price discreteness;
    \item lower institutional capacity;
    \item greater sensitivity to microstructure effects.
\end{itemize}

The threshold should be expressed in the appropriate local or common currency,
depending on the design of the universe.

Care must also be taken when using historical prices around stock splits. A
price filter should represent the actual economically relevant price at the
historical date rather than an improperly adjusted value that incorporates
future corporate actions.


% ------------------------------------------------------------
\subsection{Listing History Requirements}
\label{subsec:listing_history_requirements}
% ------------------------------------------------------------

Newly listed securities may not possess sufficient historical information for
signal estimation, risk estimation, or portfolio construction.

Let

\begin{equation}
t_i^{\mathrm{list}}
\end{equation}

denote the historical listing date of security $i$.

Its listing age at decision date $t$ is

\begin{equation}
\boxed{
A_{i,t}^{\mathrm{listing}}
=
t-t_i^{\mathrm{list}}.
}
\end{equation}

A minimum-history rule may require

\begin{equation}
A_{i,t}^{\mathrm{listing}}
\geq
L_{\min}.
\end{equation}

Equivalently,

\begin{equation}
E_{i,t}^{\mathrm{history}}
=
\begin{cases}
1,
&
A_{i,t}^{\mathrm{listing}}
\geq
L_{\min},
\\
0,
&
\text{otherwise}.
\end{cases}
\end{equation}


\subsubsection{Calendar History vs.\ Valid Observations}

Listing age alone may be insufficient.

A stock listed for one year may have experienced trading suspensions or missing
observations.

It may therefore be preferable to require a minimum number of valid historical
observations:

\begin{equation}
\boxed{
N_{i,t}^{\mathrm{valid}}
\geq
N_{\min}.
}
\end{equation}

For example, estimating a rolling beta over 252 trading days may require a
minimum number of valid return observations within that window rather than
merely requiring that the stock was listed 252 calendar days ago.


\subsubsection{IPO Treatment}

A minimum listing-history rule implicitly excludes recent IPOs.

This may be necessary because:

\begin{itemize}
    \item historical returns are insufficient;
    \item accounting histories may be incomplete;
    \item liquidity may be unstable immediately after listing;
    \item certain rolling features cannot yet be computed.
\end{itemize}

However, such exclusions also modify the economic universe of the strategy and
should therefore be documented explicitly.


% ------------------------------------------------------------
\subsection{Data Availability Requirements}
\label{subsec:data_availability_requirements}
% ------------------------------------------------------------

A security may satisfy all economic and liquidity criteria while lacking the
data required by the investment model.

Let the signal require $K$ inputs

\begin{equation}
X_{i,t}^{(1)},
\ldots,
X_{i,t}^{(K)}.
\end{equation}

A strict complete-case eligibility rule is

\begin{equation}
E_{i,t}^{\mathrm{data}}
=
\begin{cases}
1,
&
X_{i,t}^{(k)}
\text{ is available for all }
k=1,\ldots,K,
\\
0,
&
\text{otherwise}.
\end{cases}
\end{equation}

Equivalently,

\begin{equation}
\boxed{
E_{i,t}^{\mathrm{data}}
=
\prod_{k=1}^{K}
\mathbbm{1}
\left(
X_{i,t}^{(k)}
\text{ available}
\right).
}
\end{equation}


\subsubsection{Core vs.\ Optional Variables}

Not every missing variable necessarily requires removing a security.

The feature set can be divided into:

\begin{equation}
\mathcal{K}
=
\mathcal{K}^{\mathrm{required}}
\cup
\mathcal{K}^{\mathrm{optional}}.
\end{equation}

Eligibility then requires only

\begin{equation}
X_{i,t}^{(k)}
\text{ available}
\qquad
\forall
k\in
\mathcal{K}^{\mathrm{required}}.
\end{equation}

Optional variables may subsequently be treated through an explicitly defined
missing-data procedure.


\subsubsection{Missingness as a Selection Mechanism}

Data-availability filters can themselves create selection effects.

For example, small or recently listed firms may have:

\begin{itemize}
    \item fewer analyst estimates;
    \item shorter accounting histories;
    \item poorer alternative-data coverage.
\end{itemize}

Therefore,

\begin{equation}
\boxed{
\text{Missing Data}
\not\Rightarrow
\text{Random Missingness}.
}
\end{equation}

Requiring complete observations can systematically alter the economic
characteristics of the universe.

The number and characteristics of securities excluded because of missing data
should therefore be monitored.


% ------------------------------------------------------------
\subsection{Shortability Constraints}
\label{subsec:shortability_constraints}
% ------------------------------------------------------------

For long--short strategies, eligibility on the long side does not necessarily
imply eligibility on the short side.

Define

\begin{equation}
\mathcal{U}_t^{L}
\end{equation}

and

\begin{equation}
\mathcal{U}_t^{S}
\end{equation}

as the long-eligible and short-eligible universes respectively.

Typically,

\begin{equation}
\boxed{
\mathcal{U}_t^{S}
\subseteq
\mathcal{U}_t^{L}
}
\end{equation}

or, more generally, both are subsets of a common investment universe.


\subsubsection{Borrow Availability}

Let

\begin{equation}
B_{i,t}^{\mathrm{avail}}
\end{equation}

denote the number or notional amount of shares available to borrow.

A minimum borrow-availability condition may be written as

\begin{equation}
B_{i,t}^{\mathrm{avail}}
\geq
B_{\min}.
\end{equation}


\subsubsection{Borrow Cost}

Let

\begin{equation}
b_{i,t}
\end{equation}

denote the annualized borrow fee.

A strategy may exclude securities satisfying

\begin{equation}
b_{i,t}
>
b_{\max}.
\end{equation}

The short-eligibility indicator can therefore combine both conditions:

\begin{equation}
E_{i,t}^{\mathrm{short}}
=
\begin{cases}
1,
&
B_{i,t}^{\mathrm{avail}}
\geq
B_{\min}
\quad\text{and}\quad
b_{i,t}
\leq
b_{\max},
\\
0,
&
\text{otherwise}.
\end{cases}
\end{equation}


\subsubsection{Filter vs.\ Cost Treatment}

A high borrow fee does not necessarily imply that a security should be
excluded.

Two different approaches are possible:

\begin{equation}
\boxed{
\begin{aligned}
\text{Hard Constraint:}
&\quad
b_{i,t}>b_{\max}
\Rightarrow
i\notin\mathcal{U}_t^{S},
\\[0.15cm]
\text{Soft Treatment:}
&\quad
i\in\mathcal{U}_t^{S},
\text{ but borrow cost enters expected net return.}
\end{aligned}
}
\end{equation}

The appropriate choice depends on the portfolio-construction framework and
availability of reliable securities-lending data.


\subsubsection{Asymmetric Long and Short Universes}

The distinction between long and short eligibility becomes important when
interpreting a long--short strategy.

For example, the raw signal may rank a difficult-to-borrow security as one of
the strongest short candidates.

If that security cannot actually be borrowed, the theoretical short portfolio
and implementable short portfolio differ.

Thus,

\begin{equation}
\boxed{
\text{Predictive Opportunity Set}
\neq
\text{Implementable Short Opportunity Set}.
}
\end{equation}


% ------------------------------------------------------------
\subsection{Dynamic Universe Construction}
\label{subsec:dynamic_universe_construction}
% ------------------------------------------------------------

The investment universe should generally be reconstructed at every relevant
decision date rather than defined once over the complete historical sample.

Let

\begin{equation}
\mathcal{U}_t^{(0)}
\end{equation}

denote the point-in-time raw universe.

A sequential representation of the filtering process is

\begin{equation}
\begin{aligned}
\mathcal{U}_t^{(1)}
&=
\mathcal{F}_{\mathrm{geo}}
\left(
\mathcal{U}_t^{(0)}
\right),
\\
\mathcal{U}_t^{(2)}
&=
\mathcal{F}_{\mathrm{type}}
\left(
\mathcal{U}_t^{(1)}
\right),
\\
\mathcal{U}_t^{(3)}
&=
\mathcal{F}_{\mathrm{MCap}}
\left(
\mathcal{U}_t^{(2)}
\right),
\\
\mathcal{U}_t^{(4)}
&=
\mathcal{F}_{\mathrm{liq}}
\left(
\mathcal{U}_t^{(3)}
\right),
\\
&\vdots
\\
\mathcal{U}_t
&=
\mathcal{F}_{K}
\left(
\mathcal{U}_t^{(K-1)}
\right).
\end{aligned}
\end{equation}

Equivalently, the final universe can be expressed directly as

\begin{equation}
\boxed{
\mathcal{U}_t
=
\left\{
i\in\mathcal{U}_t^{\mathrm{raw}}
:
E_{i,t}^{\mathrm{geo}}
E_{i,t}^{\mathrm{exchange}}
E_{i,t}^{\mathrm{type}}
E_{i,t}^{\mathrm{MCap}}
E_{i,t}^{\mathrm{liq}}
E_{i,t}^{\mathrm{price}}
E_{i,t}^{\mathrm{history}}
E_{i,t}^{\mathrm{data}}
=1
\right\}.
}
\end{equation}

Short eligibility may then be imposed separately.


\subsubsection{Universe Evolution}

Because eligibility characteristics change through time,

\begin{equation}
\boxed{
\mathcal{U}_t
\neq
\mathcal{U}_{t+1}
}
\end{equation}

in general.

A security may enter the universe because:

\begin{itemize}
    \item it completes an IPO and satisfies the minimum-history requirement;
    \item its capitalization exceeds the required threshold;
    \item its liquidity improves;
    \item sufficient fundamental or alternative data become available.
\end{itemize}

A security may leave because:

\begin{itemize}
    \item it delists;
    \item its capitalization falls below the threshold;
    \item liquidity deteriorates;
    \item its price falls below the minimum;
    \item required data become unavailable;
    \item its security or geographical classification changes.
\end{itemize}


\subsubsection{Universe Reconstitution Frequency}

The universe does not necessarily need to be reconstructed every day.

Let

\begin{equation}
\mathcal{T}^{U}
=
\left\{
t_1^{U},
t_2^{U},
\ldots
\right\}
\end{equation}

denote universe-reconstitution dates.

Possible frequencies include:

\begin{itemize}
    \item daily;
    \item weekly;
    \item monthly;
    \item quarterly;
    \item index-reconstitution dates.
\end{itemize}

Between two reconstitution dates, the universe may either remain fixed or be
updated for exceptional events such as delistings and trading suspensions.

This convention should be specified explicitly.


\subsubsection{Buffer Rules and Universe Turnover}

Hard thresholds can generate excessive universe turnover.

Suppose the minimum market capitalization is

\begin{equation}
M_{\min}.
\end{equation}

A security whose capitalization fluctuates repeatedly around this threshold may
enter and leave the universe frequently.

One solution is to introduce separate entry and exit thresholds:

\begin{equation}
M_{\mathrm{entry}}
>
M_{\mathrm{exit}}.
\end{equation}

A security outside the universe enters only if

\begin{equation}
\mathrm{MCap}_{i,t}
\geq
M_{\mathrm{entry}},
\end{equation}

whereas an existing constituent remains eligible provided

\begin{equation}
\mathrm{MCap}_{i,t}
\geq
M_{\mathrm{exit}}.
\end{equation}

Hence,

\begin{equation}
\boxed{
M_{\mathrm{entry}}
>
M_{\mathrm{exit}}
}
\end{equation}

creates a buffer region that reduces mechanical universe turnover.

Similar buffer rules can be applied to liquidity, rankings, and other
eligibility criteria.


\subsubsection{Universe Diagnostics}

Universe construction should be treated as an observable component of the
backtest rather than as an invisible preprocessing step.

At each date $t$, useful diagnostics include:

\begin{itemize}
    \item number of securities in the raw universe;
    \item number removed by each filter;
    \item final universe size;
    \item median and distribution of market capitalization;
    \item median and distribution of liquidity;
    \item sector and country composition;
    \item number of shortable securities;
    \item universe entries and exits.
\end{itemize}

Let

\begin{equation}
N_t
=
|\mathcal{U}_t|
\end{equation}

denote the number of eligible securities.

The time series

\begin{equation}
\left\{
N_t
\right\}_{t=1}^{T}
\end{equation}

should be monitored for unexpected discontinuities that may reveal data or
filtering errors.


\subsubsection{Final Point-in-Time Universe}

The final investment universe should depend exclusively on information
available at the universe-construction time.

Formally,

\begin{equation}
\boxed{
\mathcal{U}_t
=
G
\left(
\mathcal{I}_t;
\theta^{U}
\right),
}
\end{equation}

where:

\begin{itemize}
    \item $\mathcal{I}_t$ is the point-in-time information set;
    \item $G(\cdot)$ is the universe-construction procedure;
    \item $\theta^{U}$ contains the predefined eligibility parameters.
\end{itemize}

The resulting universe is then passed to the subsequent stages of the
investment process:

\begin{equation}
\boxed{
\begin{aligned}
\mathcal{I}_t
&\rightarrow
\mathcal{U}_t
\rightarrow
\text{Data Cleaning}
\rightarrow
\text{Feature Construction}
\\[0.15cm]
&\rightarrow
\text{Signal}
\rightarrow
\text{Portfolio Construction}.
\end{aligned}
}
\end{equation}

Universe construction therefore acts as the bridge between the complete
point-in-time database and the cross-section on which the quantitative
investment strategy is actually defined.


% ============================================================
\section{Data Cleaning}
\label{sec:data_cleaning}
% ============================================================

Once the point-in-time investment universe has been constructed, the raw
observations associated with each eligible security must be transformed into a
consistent and economically meaningful dataset.

Let

\begin{equation}
\mathcal{D}_t^{\mathrm{raw}}
\end{equation}

denote the point-in-time data available for the investment universe at time
$t$. The cleaning operator may be represented abstractly as

\begin{equation}
\boxed{
\mathcal{D}_t^{\mathrm{clean}}
=
\mathcal{C}
\left(
\mathcal{D}_t^{\mathrm{raw}}
\right),
}
\end{equation}

where $\mathcal{C}$ denotes the collection of deterministic data-validation and
cleaning rules.

The objective of this stage is not to improve the apparent statistical
properties of the data. Rather, it is to ensure that observations entering
subsequent calculations have a valid economic interpretation.

A robust cleaning pipeline should satisfy several principles:

\begin{enumerate}
    \item \textbf{Determinism:} identical inputs and cleaning rules must produce
    identical outputs.

    \item \textbf{Causality:} cleaning at time $t$ must not use information
    unavailable at time $t$.

    \item \textbf{Traceability:} modifications and exclusions should be
    identifiable and auditable.

    \item \textbf{Minimal intervention:} observations should not be modified
    unless there is a clearly defined reason.

    \item \textbf{Separation of concerns:} invalid-data treatment, missing-data
    treatment, and outlier treatment should remain conceptually distinct.
\end{enumerate}

In particular,

\begin{equation}
\boxed{
\text{Missing}
\neq
\text{Invalid}
\neq
\text{Stale}
\neq
\text{Extreme}.
}
\end{equation}

An extreme but economically valid observation should not automatically be
removed during data cleaning. Statistical treatment of extreme observations is
developed separately in the outlier-handling section.


% ------------------------------------------------------------
\subsection{Duplicate Observations}
\label{subsec:duplicate_observations}
% ------------------------------------------------------------

The first requirement of a panel dataset is that its observational key be
well-defined.

Suppose the intended security-level key is

\begin{equation}
K
=
(i,t),
\end{equation}

where $i$ identifies the security and $t$ the observation timestamp.

The dataset should satisfy

\begin{equation}
\boxed{
\#\left\{
j:
K_j=(i,t)
\right\}
\leq 1
}
\end{equation}

for every $(i,t)$ pair, unless multiple observations at the same timestamp are
explicitly part of the data model.


\subsubsection{Exact Duplicates}

Two rows may be identical across all fields:

\begin{equation}
D_j=D_k.
\end{equation}

If the duplication is known to result from ingestion or concatenation, retaining
a single observation is generally sufficient.


\subsubsection{Conflicting Duplicates}

The more difficult case occurs when

\begin{equation}
(i_j,t_j)
=
(i_k,t_k)
\end{equation}

but

\begin{equation}
X_j
\neq
X_k.
\end{equation}

For example, the same security and timestamp may appear with two different
prices.

Such observations should not be resolved through an arbitrary operation such
as

\begin{equation}
\texttt{drop\_duplicates(keep="last")}
\end{equation}

unless the ordering itself has a documented economic meaning.

A deterministic priority rule should instead be defined. Possible criteria
include:

\begin{itemize}
    \item vendor revision timestamp;
    \item source priority;
    \item exchange or listing priority;
    \item data-quality flag;
    \item observation status;
    \item ingestion version.
\end{itemize}

If no defensible rule exists, the observation should be flagged for
investigation rather than silently selected.


\subsubsection{Uniqueness Invariant}

After duplicate resolution, the pipeline should verify

\begin{equation}
\boxed{
\operatorname{Unique}
\left(
\mathrm{SecurityID},
\mathrm{Timestamp}
\right)
=
\mathrm{True}.
}
\end{equation}

The same principle applies to other datasets, although their keys differ. For
example,

\begin{equation}
(\mathrm{CompanyID},
\mathrm{FiscalPeriod},
\mathrm{Vintage})
\end{equation}

may be appropriate for fundamental data, while analyst-level data may require

\begin{equation}
(\mathrm{CompanyID},
\mathrm{AnalystID},
\mathrm{ForecastPeriod},
\mathrm{Timestamp}).
\end{equation}


% ------------------------------------------------------------
\subsection{Missing Values}
\label{subsec:missing_values}
% ------------------------------------------------------------

Missing observations are unavoidable in financial datasets.

Let

\begin{equation}
M_{i,t}^{(k)}
=
\begin{cases}
1,
&
X_{i,t}^{(k)}
\text{ is missing},
\\
0,
&
\text{otherwise},
\end{cases}
\end{equation}

denote the missingness indicator for variable $k$.

A crucial principle is that

\begin{equation}
\boxed{
X_{i,t}=\mathrm{NA}
}
\end{equation}

is not itself an instruction describing how the observation should be treated.


\subsubsection{Sources of Missingness}

A value may be missing because:

\begin{itemize}
    \item the security did not yet exist;
    \item the market was closed;
    \item trading was suspended;
    \item the variable had not yet been published;
    \item the variable is economically not applicable;
    \item the vendor failed to collect the observation;
    \item security coverage is incomplete;
    \item the observation failed a previous validation rule.
\end{itemize}

These situations are economically different and should not automatically
receive identical treatment.


\subsubsection{Missingness State}

Whenever possible, the pipeline should preserve not only the missing value but
also its reason.

Conceptually,

\begin{equation}
\boxed{
\mathrm{MissingState}_{i,t}
\in
\left\{
\begin{array}{l}
\mathrm{NotListed},\\
\mathrm{MarketClosed},\\
\mathrm{Suspended},\\
\mathrm{NotReported},\\
\mathrm{NotApplicable},\\
\mathrm{VendorMissing},\\
\mathrm{Invalidated},\\
\mathrm{Unknown}
\end{array}
\right\}.
}
\end{equation}

This distinction becomes particularly important when missing values are later
used in feature construction or machine-learning models.


\subsubsection{Dropping Observations}

A complete-case procedure retains security $i$ only when all required
variables are observed:

\begin{equation}
\boxed{
\sum_{k=1}^{K}
M_{i,t}^{(k)}
=
0.
}
\end{equation}

Although simple, this approach can substantially modify the investment
universe if missingness is correlated with size, liquidity, geography, analyst
coverage, or firm age.


\subsubsection{Forward Filling}

For slowly changing variables, a previous observation may sometimes be carried
forward.

Define

\begin{equation}
\tau_i^*(t)
=
\max
\left\{
\tau\leq t:
X_{i,\tau}
\text{ observed}
\right\}.
\end{equation}

Then

\begin{equation}
X_{i,t}^{\mathrm{ffill}}
=
X_{i,\tau_i^*(t)}.
\end{equation}

However, forward filling is valid only when the previous observation remains
economically meaningful.

A maximum staleness condition should generally be imposed:

\begin{equation}
\boxed{
t-\tau_i^*(t)
\leq
L_{\max}.
}
\end{equation}

Otherwise,

\begin{equation}
X_{i,t}^{\mathrm{ffill}}
=
\mathrm{NA}.
\end{equation}


\subsubsection{Interpolation}

Time-series interpolation such as

\begin{equation}
\widetilde{X}_{i,t}
=
\lambda X_{i,t_1}
+
(1-\lambda)X_{i,t_2},
\qquad
t_1<t<t_2,
\end{equation}

is dangerous in predictive backtesting whenever $X_{i,t_2}$ was not known at
time $t$.

In that case,

\begin{equation}
\boxed{
\text{Interpolation using future observations}
\Rightarrow
\text{Look-Ahead Bias}.
}
\end{equation}

Interpolation should therefore not be treated as an innocuous data-cleaning
operation.


\subsubsection{Imputation}

Statistical imputation belongs conceptually closer to feature preprocessing
than to raw data cleaning.

For example, replacing a missing feature with a cross-sectional median,

\begin{equation}
X_{i,t}^{\mathrm{imp}}
=
\operatorname{Median}_{j\in\mathcal{U}_t}
X_{j,t},
\end{equation}

changes the information supplied to the predictive model.

Such transformations should therefore be explicit, fitted only on admissible
information, and documented as part of feature preprocessing rather than hidden
inside the raw cleaning layer.


% ------------------------------------------------------------
\subsection{Stale Prices}
\label{subsec:stale_prices}
% ------------------------------------------------------------

A recorded price may be numerically valid while failing to represent a recent
market transaction.

This situation is referred to as a stale price.

Let

\begin{equation}
t_{i}^{\mathrm{lasttrade}}(t)
\end{equation}

denote the timestamp of the most recent valid trade observed for security $i$
as of time $t$.

Define price age as

\begin{equation}
\boxed{
A_{i,t}^{P}
=
t
-
t_{i}^{\mathrm{lasttrade}}(t).
}
\end{equation}

A price may be classified as stale if

\begin{equation}
A_{i,t}^{P}
>
A_{\max}^{P}.
\end{equation}


\subsubsection{Repeated Prices}

A simple diagnostic is the number of consecutive observations with an
unchanged price.

Define

\begin{equation}
Z_{i,t}
=
\begin{cases}
1,
&
P_{i,t}=P_{i,t-1},
\\
0,
&
\text{otherwise}.
\end{cases}
\end{equation}

A long sequence

\begin{equation}
Z_{i,t}
=
Z_{i,t-1}
=
\cdots
=
1
\end{equation}

may indicate stale data.

However,

\begin{equation}
\boxed{
P_{i,t}=P_{i,t-1}
\not\Rightarrow
\text{Stale Price}.
}
\end{equation}

A security can genuinely close at the same price on consecutive days.
Additional information such as trading volume, number of trades, quotes, and
market status should therefore be considered.


\subsubsection{Zero Return and Zero Volume}

A stronger warning signal is

\begin{equation}
R_{i,t}=0
\qquad\text{and}\qquad
V_{i,t}=0.
\end{equation}

Even this condition does not uniquely identify an error: the market may have
been closed or the security may have been suspended.

Staleness detection should therefore distinguish

\begin{equation}
\boxed{
\text{No Price Movement}
\neq
\text{No Trading}
\neq
\text{Missing Observation}.
}
\end{equation}


\subsubsection{Effect on Backtests}

Stale prices can distort:

\begin{itemize}
    \item return calculations;
    \item volatility estimates;
    \item covariance matrices;
    \item beta estimates;
    \item momentum signals;
    \item liquidity measures;
    \item portfolio valuations.
\end{itemize}

For example, repeated stale prices artificially generate

\begin{equation}
R_{i,t}=0,
\end{equation}

which can mechanically reduce estimated volatility.

The correct treatment depends on the cause of staleness and should be
determined before risk and signal calculations are performed.


% ------------------------------------------------------------
\subsection{Invalid Observations}
\label{subsec:invalid_observations}
% ------------------------------------------------------------

An invalid observation violates a known logical, economic, or structural
constraint.

Unlike an outlier, an invalid observation is not merely unusual: it is
inconsistent with the definition of the variable or the data-generating
process.


\subsubsection{Hard Validation Rules}

Examples include

\begin{equation}
P_{i,t}\leq0,
\end{equation}

\begin{equation}
V_{i,t}<0,
\end{equation}

or, for OHLC data,

\begin{equation}
P_{i,t}^{\mathrm{high}}
<
P_{i,t}^{\mathrm{low}}.
\end{equation}

Additional consistency conditions include

\begin{equation}
P_{i,t}^{\mathrm{high}}
\geq
\max
\left(
P_{i,t}^{\mathrm{open}},
P_{i,t}^{\mathrm{close}}
\right),
\end{equation}

and

\begin{equation}
P_{i,t}^{\mathrm{low}}
\leq
\min
\left(
P_{i,t}^{\mathrm{open}},
P_{i,t}^{\mathrm{close}}
\right).
\end{equation}

If bid and ask quotes are used, one would normally expect

\begin{equation}
P_{i,t}^{\mathrm{bid}}
\leq
P_{i,t}^{\mathrm{ask}}.
\end{equation}


\subsubsection{Cross-Field Consistency}

Validation should not be restricted to individual variables.

Suppose

\begin{equation}
\mathrm{DollarVolume}_{i,t}
\end{equation}

is provided directly by the vendor.

It can be compared with an approximate reconstructed value

\begin{equation}
\widehat{\mathrm{DollarVolume}}_{i,t}
=
P_{i,t}V_{i,t}.
\end{equation}

Large discrepancies may indicate inconsistent units, adjustment conventions, or
data errors.


\subsubsection{Hard Errors vs.\ Extreme Observations}

Consider a one-day return of

\begin{equation}
R_{i,t}=+150\%.
\end{equation}

This observation is extreme but not necessarily invalid.

It may result from:

\begin{itemize}
    \item genuine market movement;
    \item an acquisition;
    \item a corporate restructuring;
    \item an unadjusted stock split;
    \item a data error.
\end{itemize}

Therefore,

\begin{equation}
\boxed{
\text{Large Magnitude}
\not\Rightarrow
\text{Invalid Observation}.
}
\end{equation}

The pipeline should first investigate structural explanations, particularly
corporate actions, before treating the observation statistically.


\subsubsection{Validation Flags}

Rather than immediately deleting suspicious observations, a robust system may
attach validation flags:

\begin{equation}
F_{i,t}^{\mathrm{valid}}
\in
\{
\mathrm{Valid},
\mathrm{Warning},
\mathrm{Invalid}
\}.
\end{equation}

This preserves the original observation and allows the downstream policy to be
defined separately from detection.


% ------------------------------------------------------------
\subsection{Corporate Action Adjustments}
\label{subsec:corporate_action_adjustments}
% ------------------------------------------------------------

Corporate actions create discontinuities in prices, shares outstanding, and
other security-level variables that must be distinguished from genuine economic
returns.

Let

\begin{equation}
A_{i,t}
\end{equation}

denote the cumulative adjustment factor applied to the raw price.

An adjusted price may be represented as

\begin{equation}
\boxed{
P_{i,t}^{\mathrm{adj}}
=
A_{i,t}
P_{i,t}^{\mathrm{raw}}.
}
\end{equation}


\subsubsection{Stock Splits}

Consider an $m$-for-$n$ split.

The split ratio is

\begin{equation}
s
=
\frac{m}{n}.
\end{equation}

Ignoring market movements,

\begin{equation}
P_{i,t}^{\mathrm{after}}
=
\frac{1}{s}
P_{i,t^-}^{\mathrm{before}},
\end{equation}

while

\begin{equation}
Q_{i,t}^{\mathrm{after}}
=
s
Q_{i,t^-}^{\mathrm{before}}.
\end{equation}

Thus,

\begin{equation}
P_{i,t}^{\mathrm{after}}
Q_{i,t}^{\mathrm{after}}
=
P_{i,t^-}^{\mathrm{before}}
Q_{i,t^-}^{\mathrm{before}}.
\end{equation}

If the price series is not adjusted, a split may incorrectly appear as a large
investment return.


\subsubsection{Dividends}

For a cash dividend $D_{i,t}$, total return is

\begin{equation}
\boxed{
R_{i,t}^{\mathrm{total}}
=
\frac{
P_{i,t}
+
D_{i,t}
}{
P_{i,t-1}
}
-1.
}
\end{equation}

Using only

\begin{equation}
\frac{P_{i,t}}{P_{i,t-1}}-1
\end{equation}

would incorrectly interpret the mechanical ex-dividend price decline as an
economic loss.


\subsubsection{Adjustment Consistency}

A particularly important engineering rule is that adjustment conventions must
be internally consistent.

For example, combining

\begin{equation}
P_{i,t}^{\mathrm{adj}}
\end{equation}

with unadjusted shares outstanding or unadjusted volume can produce inconsistent
quantities.

For a split factor $s_{i,t}$, if historical prices are multiplied by an
adjustment factor, historical share quantities generally require the inverse
economic treatment.

Thus, variables such as

\begin{equation}
P,
\qquad
V,
\qquad
N^{\mathrm{shares}},
\qquad
\mathrm{MCap}
\end{equation}

must be checked for consistent adjustment conventions.


\subsubsection{Avoiding Double Adjustment}

If a vendor already provides a total-return or fully adjusted series, applying
corporate-action adjustments again will double count the event.

Therefore,

\begin{equation}
\boxed{
\text{Vendor-Adjusted Data}
+
\text{Manual Adjustment}
\not\equiv
\text{Correctly Adjusted Data}.
}
\end{equation}

The pipeline should explicitly record which variables are raw and which are
already adjusted.


% ------------------------------------------------------------
\subsection{Currency Conversion}
\label{subsec:currency_conversion_cleaning}
% ------------------------------------------------------------

For multi-country strategies, monetary quantities must be expressed under a
consistent currency convention.

Let

\begin{equation}
C_0
\end{equation}

denote the portfolio base currency.

Suppose

\begin{equation}
X_{i,t}^{(C_i)}
\end{equation}

is a monetary quantity expressed in local currency $C_i$.

Let

\begin{equation}
FX_t^{C_i\rightarrow C_0}
\end{equation}

denote the number of units of base currency obtained for one unit of local
currency.

Then

\begin{equation}
\boxed{
X_{i,t}^{(C_0)}
=
X_{i,t}^{(C_i)}
FX_t^{C_i\rightarrow C_0}.
}
\end{equation}


\subsubsection{Quotation Convention}

FX quotation conventions must be explicit.

If instead the database stores

\begin{equation}
FX_t^{C_0/C_i}
\end{equation}

as units of local currency per unit of base currency, conversion requires

\begin{equation}
X_{i,t}^{(C_0)}
=
\frac{
X_{i,t}^{(C_i)}
}{
FX_t^{C_0/C_i}
}.
\end{equation}

Confusing these two conventions creates multiplicative errors that may not be
obvious from the resulting dataset.


\subsubsection{Prices vs.\ Returns}

Currency conversion of price levels and currency conversion of returns are
different operations.

For returns,

\begin{equation}
1+R_{i,t}^{C_0}
=
\left(
1+R_{i,t}^{C_i}
\right)
\left(
1+R_{FX,t}^{C_i\rightarrow C_0}
\right).
\end{equation}

Hence,

\begin{equation}
\boxed{
R_{i,t}^{C_0}
=
R_{i,t}^{C_i}
+
R_{FX,t}^{C_i\rightarrow C_0}
+
R_{i,t}^{C_i}
R_{FX,t}^{C_i\rightarrow C_0}.
}
\end{equation}

The cross-product term should not be discarded when exact returns are required.


\subsubsection{Currency Metadata}

Every monetary variable should ideally carry an explicit currency identifier:

\begin{equation}
\boxed{
\left(
X_{i,t},
C_{i,t}
\right).
}
\end{equation}

This is especially important for fundamentals because reporting currencies can
change through time.

Currency should therefore be treated as data rather than inferred permanently
from a company's country.


% ------------------------------------------------------------
\subsection{Frequency Alignment}
\label{subsec:frequency_alignment}
% ------------------------------------------------------------

Quantitative strategies frequently combine variables observed at different
frequencies.

For example,

\begin{equation}
\begin{aligned}
\text{Market Data}
&:\quad
\text{Daily or Intraday},
\\
\text{Fundamentals}
&:\quad
\text{Quarterly or Semi-Annual},
\\
\text{Analyst Estimates}
&:\quad
\text{Irregular},
\\
\text{Macroeconomic Data}
&:\quad
\text{Monthly or Quarterly}.
\end{aligned}
\end{equation}

These observations must ultimately be mapped to a common set of strategy
decision times

\begin{equation}
\mathcal{T}^{D}
=
\left\{
t_1,t_2,\ldots,t_T
\right\}.
\end{equation}


\subsubsection{Downsampling}

Suppose daily returns are converted into monthly returns.

If month $m$ contains trading days

\begin{equation}
\mathcal{T}_m,
\end{equation}

the correct simple monthly return is obtained geometrically:

\begin{equation}
\boxed{
R_{i,m}
=
\prod_{t\in\mathcal{T}_m}
(1+R_{i,t})
-1.
}
\end{equation}

In general,

\begin{equation}
R_{i,m}
\neq
\sum_{t\in\mathcal{T}_m}
R_{i,t}.
\end{equation}

For log returns, however,

\begin{equation}
r_{i,m}
=
\sum_{t\in\mathcal{T}_m}
r_{i,t}.
\end{equation}


\subsubsection{Aggregation Depends on Variable Type}

Different variables require different aggregation operators.

For example:

\begin{equation}
\boxed{
\begin{array}{lll}
\text{Returns}
&
\rightarrow
&
\text{Geometric Compounding},
\\
\text{Volume}
&
\rightarrow
&
\text{Summation or Average},
\\
\text{Closing Price}
&
\rightarrow
&
\text{Last Valid Observation},
\\
\text{Volatility}
&
\rightarrow
&
\text{Model-Dependent Aggregation},
\\
\text{Fundamental State}
&
\rightarrow
&
\text{Latest Available Observation}.
\end{array}
}
\end{equation}

A generic resampling operation should therefore not apply the same aggregation
function to every column.


\subsubsection{Upsampling}

Lower-frequency variables are often mapped onto higher-frequency decision
grids.

For example, a quarterly accounting variable may be carried forward to daily
dates:

\begin{equation}
X_{i,t}
=
X_{i,\tau^*(t)},
\end{equation}

where

\begin{equation}
\tau^*(t)
=
\max
\left\{
\tau:
t_{\tau}^{\mathrm{avail}}
\leq t
\right\}.
\end{equation}

This is fundamentally an as-of operation, not interpolation.

No value from the next reporting period should enter before its availability
date.


\subsubsection{Information Frequency vs.\ Trading Frequency}

The frequency of the input data does not determine the portfolio rebalancing
frequency.

A strategy may use daily information but rebalance monthly:

\begin{equation}
\boxed{
f_{\mathrm{data}}
\neq
f_{\mathrm{signal}}
\neq
f_{\mathrm{rebalance}}.
}
\end{equation}

These frequencies should be represented separately in the backtesting
architecture.


% ------------------------------------------------------------
\subsection{Trading Calendar Alignment}
\label{subsec:trading_calendar_alignment}
% ------------------------------------------------------------

Financial datasets do not necessarily share the same calendar.

Different exchanges may have different:

\begin{itemize}
    \item weekends;
    \item public holidays;
    \item exceptional closures;
    \item trading suspensions;
    \item market opening and closing times;
    \item daylight-saving conventions.
\end{itemize}

Let

\begin{equation}
\mathcal{T}_i
\end{equation}

denote the valid trading calendar associated with security $i$.


\subsubsection{Local Trading Calendars}

For a single-market strategy, a master calendar may be defined as

\begin{equation}
\mathcal{T}
=
\left\{
t:
\text{the relevant exchange is open at }t
\right\}.
\end{equation}

For a multi-market strategy, however,

\begin{equation}
\mathcal{T}_i
\neq
\mathcal{T}_j
\end{equation}

may hold for securities listed in different countries.


\subsubsection{Union and Intersection Calendars}

A global backtester may construct either the union

\begin{equation}
\mathcal{T}^{\cup}
=
\bigcup_i
\mathcal{T}_i
\end{equation}

or the intersection

\begin{equation}
\mathcal{T}^{\cap}
=
\bigcap_i
\mathcal{T}_i.
\end{equation}

The union retains every date on which at least one market is open, whereas the
intersection retains only dates on which all relevant markets are open.

Neither choice is universally correct.

For global portfolios, the union calendar is often necessary for accurate NAV
accounting, but individual securities must retain their own tradability status.


\subsubsection{Market Closure vs.\ Missing Data}

Suppose market $M$ is closed on day $t$.

Then the absence of a new price is not equivalent to a vendor data error.

Define

\begin{equation}
O_{i,t}
=
\begin{cases}
1,
&
\text{security }i\text{'s market is open at }t,
\\
0,
&
\text{otherwise}.
\end{cases}
\end{equation}

Then

\begin{equation}
O_{i,t}=0
\end{equation}

and

\begin{equation}
P_{i,t}=\mathrm{NA}
\end{equation}

may be completely valid.

The cleaning pipeline should therefore distinguish

\begin{equation}
\boxed{
\text{Scheduled Market Closure}
\neq
\text{Missing Market Observation}.
}
\end{equation}


\subsubsection{Month-End Alignment}

Calendar month-end and trading month-end are not necessarily identical.

Let

\begin{equation}
T_m^{\mathrm{cal}}
\end{equation}

denote the final calendar day of month $m$, and

\begin{equation}
T_m^{\mathrm{trade}}
=
\max
\left\{
t\in\mathcal{T}:
t\leq T_m^{\mathrm{cal}}
\right\}
\end{equation}

the final valid trading day.

Then a month-end strategy should normally operate on

\begin{equation}
\boxed{
T_m^{\mathrm{trade}}
}
\end{equation}

rather than mechanically requiring an observation on the final calendar day.


\subsubsection{Time Zones}

For global or intraday strategies, timestamps should be stored in an
unambiguous time-zone convention.

A robust architecture may store timestamps internally in UTC,

\begin{equation}
t^{UTC},
\end{equation}

while retaining the corresponding local exchange timestamp

\begin{equation}
t^{\mathrm{local}}.
\end{equation}

The mapping depends on the exchange time zone and daylight-saving regime.

This becomes essential when determining whether information from one market
was available before a trading decision in another.


\subsubsection{Tradability Matrix}

For multi-market strategies, it is useful to define a tradability indicator

\begin{equation}
H_{i,t}
=
\begin{cases}
1,
&
\text{security }i\text{ can be traded at }t,
\\
0,
&
\text{otherwise}.
\end{cases}
\end{equation}

Collecting these indicators across securities produces the tradability matrix

\begin{equation}
\boxed{
H
=
\left[
H_{i,t}
\right]_{i,t}.
}
\end{equation}

Portfolio implementation must respect

\begin{equation}
H_{i,t}=0
\quad\Rightarrow\quad
\Delta q_{i,t}=0,
\end{equation}

where $\Delta q_{i,t}$ denotes the requested trade in security $i$.

A security may therefore remain part of the portfolio while being temporarily
untradable.


% ------------------------------------------------------------
\subsection{Cleaning Pipeline and Validation Invariants}
\label{subsec:cleaning_pipeline}
% ------------------------------------------------------------

The individual cleaning operations should be organized into a deterministic
pipeline.

A possible ordering is

\begin{equation}
\boxed{
\begin{aligned}
\text{Raw PIT Data}
&\rightarrow
\text{Key Validation}
\rightarrow
\text{Duplicate Resolution}
\\[0.1cm]
&\rightarrow
\text{Structural Validation}
\rightarrow
\text{Corporate-Action Adjustment}
\\[0.1cm]
&\rightarrow
\text{Currency Normalization}
\rightarrow
\text{Calendar Alignment}
\\[0.1cm]
&\rightarrow
\text{Frequency Alignment}
\rightarrow
\text{Missing/Stale Flags}.
\end{aligned}
}
\end{equation}

The exact ordering may depend on the dataset, but dependencies between
operations must be explicit.

For example, an apparent $-50\%$ price return should not be classified as a
statistical anomaly before verifying whether a two-for-one stock split occurred.


\subsubsection{Validation Invariants}

At the output of the cleaning layer, a set of invariants should hold.

For example:

\begin{equation}
\boxed{
\begin{aligned}
&\text{Unique observational keys},\\
&\text{No structurally impossible prices or volumes},\\
&\text{Consistent corporate-action adjustments},\\
&\text{Explicit currency units},\\
&\text{Valid trading-calendar mapping},\\
&\text{Explicit missingness and staleness states},\\
&\text{No use of information after the decision timestamp}.
\end{aligned}
}
\end{equation}


\subsubsection{Audit Trail}

Cleaning should ideally be non-destructive.

Let

\begin{equation}
X_{i,t}^{\mathrm{raw}}
\end{equation}

denote the original observation and

\begin{equation}
X_{i,t}^{\mathrm{clean}}
\end{equation}

the value exposed to the downstream research pipeline.

The system should retain, whenever feasible,

\begin{equation}
\boxed{
\left(
X_{i,t}^{\mathrm{raw}},
X_{i,t}^{\mathrm{clean}},
F_{i,t}^{\mathrm{clean}},
R_{i,t}^{\mathrm{clean}}
\right),
}
\end{equation}

where:

\begin{itemize}
    \item $F_{i,t}^{\mathrm{clean}}$ identifies the cleaning action;
    \item $R_{i,t}^{\mathrm{clean}}$ records the reason for that action.
\end{itemize}

This makes it possible to distinguish an original vendor observation from a
value that has been adjusted, invalidated, carried forward, or otherwise
modified.


\subsubsection{Output of the Cleaning Layer}

The output of this stage should be a validated but not yet statistically
transformed dataset:

\begin{equation}
\boxed{
\mathcal{D}_t^{\mathrm{clean}}
=
\mathcal{C}
\left(
\mathcal{D}_t^{\mathrm{PIT}},
\mathcal{U}_t
\right).
}
\end{equation}

The distinction between cleaning and subsequent statistical preprocessing is
fundamental:

\begin{equation}
\boxed{
\begin{aligned}
\text{Cleaning}
&:
\text{ Is the observation economically and structurally valid?}
\\[0.15cm]
\text{Outlier Treatment}
&:
\text{ How should valid extreme observations be handled?}
\\[0.15cm]
\text{Feature Engineering}
&:
\text{ How should valid data be transformed into predictors?}
\end{aligned}
}
\end{equation}

Thus, the core data path becomes

\begin{equation}
\boxed{
\begin{aligned}
\mathcal{D}^{\mathrm{raw}}
&\rightarrow
\mathcal{D}^{\mathrm{PIT}}
\rightarrow
\mathcal{U}_t
\rightarrow
\mathcal{D}_t^{\mathrm{clean}}
\\[0.1cm]
&\rightarrow
\text{Outlier Treatment}
\rightarrow
\text{Feature Engineering}
\rightarrow
\text{Signal Construction}.
\end{aligned}
}
\end{equation}


% ============================================================
\section{Outlier Detection and Treatment}
\label{sec:outlier_detection}
% ============================================================

After structural data cleaning, the remaining dataset may still contain
observations whose magnitude is unusually large relative to the rest of the
sample.

Such observations are commonly referred to as \emph{outliers}.

However, an outlier is not necessarily an erroneous observation. In financial
data, extreme observations may reflect genuine economic events such as:

\begin{itemize}
    \item earnings surprises;
    \item bankruptcies;
    \item acquisitions;
    \item market crashes;
    \item liquidity shocks;
    \item large analyst revisions;
    \item extreme valuation changes.
\end{itemize}

Consequently,

\begin{equation}
\boxed{
\text{Extreme Observation}
\not\Rightarrow
\text{Invalid Observation}.
}
\end{equation}

The purpose of outlier treatment is therefore generally not to remove all
extreme observations, but rather to prevent a small number of observations
from exerting disproportionate influence on statistical estimates, signals, or
portfolio weights.

Let

\begin{equation}
X_{i,t}
\end{equation}

denote a cleaned variable for security $i$ at time $t$.

An outlier-treatment operator can be represented as

\begin{equation}
\boxed{
\widetilde{X}_{i,t}
=
\mathcal{O}
\left(
X_{i,t};
\theta_t
\right),
}
\end{equation}

where $\theta_t$ contains the parameters required by the chosen treatment,
such as quantile thresholds, medians, or robust scale estimates.

The parameters themselves must respect the information constraint.

For a cross-sectional transformation at time $t$,

\begin{equation}
\theta_t
=
g
\left(
\left\{
X_{j,t}
:
j\in\mathcal{U}_t
\right\}
\right),
\end{equation}

while for a time-series transformation,

\begin{equation}
\theta_{i,t}
=
g
\left(
X_{i,t-1},
X_{i,t-2},
\ldots
\right).
\end{equation}

Future observations must never enter the estimation of the treatment
parameters.


% ------------------------------------------------------------
\subsection{Definition of an Outlier}
\label{subsec:definition_outlier}
% ------------------------------------------------------------

There is no universal definition of an outlier.

An observation is generally considered an outlier if it is unusually distant
from a reference distribution or inconsistent with the typical behavior of the
variable under consideration.

Let

\begin{equation}
\mathcal{X}
=
\left\{
X_1,\ldots,X_N
\right\}
\end{equation}

denote a sample.

An outlier-detection rule can be written abstractly as

\begin{equation}
O_i
=
\begin{cases}
1,
&
X_i
\text{ is classified as an outlier},
\\
0,
&
\text{otherwise}.
\end{cases}
\end{equation}

More generally,

\begin{equation}
\boxed{
O_i
=
\mathbb{I}
\left(
D(X_i;\mathcal{X})
>
c
\right),
}
\end{equation}

where:

\begin{itemize}
    \item $D(\cdot)$ is a measure of distance from the reference distribution;
    \item $c$ is a predefined threshold.
\end{itemize}


\subsubsection{Univariate vs.\ Multivariate Outliers}

An observation may be extreme with respect to a single variable:

\begin{equation}
|X_i|
\gg
|X_j|,
\end{equation}

or extreme only through a combination of variables.

Suppose

\begin{equation}
X_i
\in
\mathbb{R}^{K}.
\end{equation}

A multivariate observation may appear normal along each dimension separately
while being unusual jointly.

A classical multivariate distance is the Mahalanobis distance:

\begin{equation}
\boxed{
D_i^2
=
(X_i-\mu)^\top
\Sigma^{-1}
(X_i-\mu),
}
\end{equation}

where $\mu$ and $\Sigma$ denote the location vector and covariance matrix.

However, both $\mu$ and $\Sigma$ can themselves be highly sensitive to
outliers. Robust location and covariance estimators may therefore be preferable
when multivariate outlier detection is required.


\subsubsection{Distribution Dependence}

Whether an observation should be considered unusual depends on the underlying
distribution.

For example,

\begin{equation}
|Z_i|>3
\end{equation}

may be rare under a Gaussian distribution, but financial variables often
exhibit:

\begin{itemize}
    \item heavy tails;
    \item skewness;
    \item multimodality;
    \item heteroskedasticity.
\end{itemize}

Therefore, fixed Gaussian thresholds should not automatically be interpreted as
universal definitions of financial outliers.


% ------------------------------------------------------------
\subsection{Cross-Sectional vs.\ Time-Series Outliers}
\label{subsec:cross_sectional_time_series_outliers}
% ------------------------------------------------------------

Outlier detection in quantitative equity research generally occurs along one
of two dimensions:

\begin{equation}
\boxed{
\text{Cross-Sectional}
\qquad\text{or}\qquad
\text{Time-Series}.
}
\end{equation}

The distinction is fundamental because the corresponding reference
distributions are different.


\subsubsection{Cross-Sectional Outliers}

At a fixed date $t$, define the cross-section

\begin{equation}
\mathcal{X}_t
=
\left\{
X_{i,t}
:
i\in\mathcal{U}_t
\right\}.
\end{equation}

A cross-sectional outlier is unusual relative to other securities observed at
the same date.

For example, a valuation ratio may satisfy

\begin{equation}
X_{i,t}
\gg
\operatorname{Median}_{j\in\mathcal{U}_t}
X_{j,t}.
\end{equation}

Cross-sectional treatment is particularly relevant for variables used to rank
securities against one another, such as:

\begin{itemize}
    \item valuation ratios;
    \item profitability measures;
    \item analyst revisions;
    \item earnings surprises;
    \item momentum signals;
    \item quality metrics.
\end{itemize}

The treatment parameters should normally be estimated separately at each date:

\begin{equation}
\boxed{
\theta_t^{CS}
=
g
\left(
\mathcal{X}_t
\right).
}
\end{equation}


\subsubsection{Time-Series Outliers}

For a fixed security $i$, define the historical series

\begin{equation}
\mathcal{X}_{i,t}^{TS}
=
\left\{
X_{i,t-L},
\ldots,
X_{i,t}
\right\}.
\end{equation}

A time-series outlier is unusual relative to the historical behavior of the
same security or variable.

A generic rule may compare

\begin{equation}
X_{i,t}
\end{equation}

with a historical location estimate

\begin{equation}
\widehat{\mu}_{i,t}
\end{equation}

and scale estimate

\begin{equation}
\widehat{\sigma}_{i,t}.
\end{equation}

For example,

\begin{equation}
\left|
\frac{
X_{i,t}-\widehat{\mu}_{i,t}
}{
\widehat{\sigma}_{i,t}
}
\right|
>
c
\end{equation}

may identify an unusually large deviation.


\subsubsection{Different Economic Meaning}

Consider an earnings-growth rate of $50\%$.

It may be completely normal for company $i$ relative to its own history but
extreme relative to the current cross-section.

Alternatively, it may be ordinary relative to other firms but highly unusual
for company $i$ itself.

Thus,

\begin{equation}
\boxed{
\text{Cross-Sectional Extremeness}
\neq
\text{Time-Series Extremeness}.
}
\end{equation}

The appropriate dimension should therefore be chosen according to the
investment hypothesis.


\subsubsection{Panel Data}

For panel data, both treatments can in principle be applied:

\begin{equation}
X_{i,t}
\rightarrow
\widetilde{X}_{i,t}^{TS}
\rightarrow
\widetilde{X}_{i,t}^{CS}.
\end{equation}

However, the ordering matters.

A time-series treatment modifies the historical behavior of each security,
whereas a cross-sectional treatment modifies relative dispersion across
securities.

The sequence should therefore be economically motivated and explicitly
documented.


% ------------------------------------------------------------
\subsection{Quantile-Based Winsorization}
\label{subsec:quantile_winsorization}
% ------------------------------------------------------------

Quantile-based winsorization is one of the most common outlier treatments in
cross-sectional financial research.

At time $t$, let

\begin{equation}
Q_{\alpha,t}
\end{equation}

and

\begin{equation}
Q_{1-\alpha,t}
\end{equation}

denote the lower and upper empirical quantiles of the cross-sectional
distribution.

For example, with

\begin{equation}
\alpha=0.01,
\end{equation}

the thresholds correspond to the $1$st and $99$th percentiles.


\subsubsection{Winsorization Operator}

The winsorized observation is

\begin{equation}
\boxed{
\widetilde{X}_{i,t}
=
\min
\left[
\max
\left(
X_{i,t},
Q_{\alpha,t}
\right),
Q_{1-\alpha,t}
\right].
}
\end{equation}

Equivalently,

\begin{equation}
\widetilde{X}_{i,t}
=
\begin{cases}
Q_{\alpha,t},
&
X_{i,t}
<
Q_{\alpha,t},
\\[0.15cm]
X_{i,t},
&
Q_{\alpha,t}
\leq
X_{i,t}
\leq
Q_{1-\alpha,t},
\\[0.15cm]
Q_{1-\alpha,t},
&
X_{i,t}
>
Q_{1-\alpha,t}.
\end{cases}
\end{equation}

The observation remains in the dataset, but its magnitude is capped.


\subsubsection{Example}

Suppose the cross-sectional thresholds are

\begin{equation}
Q_{0.01,t}
=
-3
\end{equation}

and

\begin{equation}
Q_{0.99,t}
=
4.
\end{equation}

Then

\begin{equation}
X_{i,t}=10
\end{equation}

becomes

\begin{equation}
\widetilde{X}_{i,t}=4,
\end{equation}

while

\begin{equation}
X_{j,t}=-8
\end{equation}

becomes

\begin{equation}
\widetilde{X}_{j,t}=-3.
\end{equation}


\subsubsection{Point-in-Time Computation}

For a cross-sectional signal, the thresholds should be estimated independently
at each time $t$:

\begin{equation}
Q_{\alpha,t}
=
Q_{\alpha}
\left(
\left\{
X_{j,t}
:
j\in\mathcal{U}_t
\right\}
\right).
\end{equation}

One should not compute

\begin{equation}
Q_{\alpha}
\end{equation}

once over the full panel

\begin{equation}
\left\{
X_{i,t}
:
i=1,\ldots,N,
\;
t=1,\ldots,T
\right\}
\end{equation}

and then apply it historically.

Doing so allows future observations to influence historical thresholds and may
also ignore changes in the cross-sectional distribution through time.


\subsubsection{Advantages and Limitations}

Quantile winsorization has several advantages:

\begin{itemize}
    \item it makes few distributional assumptions;
    \item it is easy to implement;
    \item it is robust to extreme magnitudes;
    \item it adapts to the cross-sectional scale of the variable.
\end{itemize}

However, it also has limitations:

\begin{itemize}
    \item the choice of $\alpha$ is arbitrary;
    \item it treats all observations beyond the threshold identically;
    \item it may suppress economically meaningful extreme information;
    \item small cross-sections can produce unstable quantiles.
\end{itemize}


% ------------------------------------------------------------
\subsection{MAD-Based Methods}
\label{subsec:mad_methods}
% ------------------------------------------------------------

The Median Absolute Deviation provides a robust alternative to the standard
deviation for measuring dispersion.

For a sample

\begin{equation}
X_1,\ldots,X_N,
\end{equation}

define the median

\begin{equation}
m
=
\operatorname{Median}(X_i).
\end{equation}

The Median Absolute Deviation is

\begin{equation}
\boxed{
MAD
=
\operatorname{Median}
\left(
|X_i-m|
\right).
}
\end{equation}

Unlike the sample mean and standard deviation, the median and MAD are much less
sensitive to extreme observations.


\subsubsection{Gaussian-Consistency Scaling}

For normally distributed data,

\begin{equation}
MAD
\approx
0.67449\,\sigma.
\end{equation}

Therefore, a consistent estimator of Gaussian scale is

\begin{equation}
\boxed{
\widehat{\sigma}_{MAD}
=
1.4826
\,MAD.
}
\end{equation}

since approximately

\begin{equation}
1.4826
=
\frac{1}{0.67449}.
\end{equation}


\subsubsection{MAD-Based Thresholds}

A robust outlier rule may be defined as

\begin{equation}
\boxed{
|X_i-m|
>
c
\widehat{\sigma}_{MAD}.
}
\end{equation}

Equivalently,

\begin{equation}
X_i
<
m-c\widehat{\sigma}_{MAD}
\end{equation}

or

\begin{equation}
X_i
>
m+c\widehat{\sigma}_{MAD}.
\end{equation}

The threshold $c$ determines the aggressiveness of the rule.


\subsubsection{Degenerate MAD}

A practical issue arises when

\begin{equation}
MAD=0.
\end{equation}

This may occur when a large fraction of observations have identical values.

In that case,

\begin{equation}
\frac{X_i-m}{MAD}
\end{equation}

is undefined.

The backtester should therefore specify an explicit fallback policy, such as:

\begin{itemize}
    \item no MAD-based treatment;
    \item alternative robust scale estimator;
    \item quantile-based treatment;
    \item explicit handling of discrete variables.
\end{itemize}

A silent division by zero should never determine the resulting signal.


% ------------------------------------------------------------
\subsection{Robust Z-Scores}
\label{subsec:robust_z_scores}
% ------------------------------------------------------------

The conventional z-score is

\begin{equation}
Z_i
=
\frac{
X_i-\bar{X}
}{
s_X
},
\end{equation}

where

\begin{equation}
\bar{X}
=
\frac{1}{N}
\sum_{i=1}^{N}
X_i
\end{equation}

and

\begin{equation}
s_X
=
\sqrt{
\frac{1}{N-1}
\sum_{i=1}^{N}
(X_i-\bar{X})^2
}.
\end{equation}

Both $\bar{X}$ and $s_X$ can be strongly influenced by outliers.


\subsubsection{Robust Z-Score}

Replacing the mean and standard deviation with the median and MAD gives

\begin{equation}
\boxed{
Z_i^{\mathrm{rob}}
=
\frac{
X_i-m
}{
1.4826\,MAD
}.
}
\end{equation}

An observation may then be flagged whenever

\begin{equation}
\boxed{
|Z_i^{\mathrm{rob}}|
>
c.
}
\end{equation}

Alternatively, the robust z-score itself may be used as the transformed
feature.


\subsubsection{Cross-Sectional Robust Standardization}

At time $t$,

\begin{equation}
m_t
=
\operatorname{Median}_{i\in\mathcal{U}_t}
X_{i,t},
\end{equation}

and

\begin{equation}
MAD_t
=
\operatorname{Median}_{i\in\mathcal{U}_t}
|X_{i,t}-m_t|.
\end{equation}

Then

\begin{equation}
\boxed{
Z_{i,t}^{\mathrm{rob}}
=
\frac{
X_{i,t}-m_t
}{
1.4826\,MAD_t
}.
}
\end{equation}

This transformation provides both centering and robust cross-sectional scaling.


\subsubsection{Detection vs.\ Transformation}

It is important to distinguish using the robust z-score to identify outliers

\begin{equation}
|Z_{i,t}^{\mathrm{rob}}|>c,
\end{equation}

from using it as the actual transformed signal:

\begin{equation}
\widetilde{X}_{i,t}
=
Z_{i,t}^{\mathrm{rob}}.
\end{equation}

These are different operations.

The first produces a diagnostic flag. The second changes the scale and
cross-sectional interpretation of the variable.


% ------------------------------------------------------------
\subsection{Truncation, Winsorization, and Observation Removal}
\label{subsec:truncation_winsorization_removal}
% ------------------------------------------------------------

Once an observation has been identified as extreme, several treatments are
possible.

The three most common approaches are:

\begin{equation}
\boxed{
\text{Winsorization},
\qquad
\text{Truncation},
\qquad
\text{Removal}.
}
\end{equation}


\subsubsection{Winsorization}

As defined previously, winsorization caps values at predefined boundaries.

For lower and upper bounds $L_t$ and $U_t$,

\begin{equation}
\boxed{
\widetilde{X}_{i,t}
=
\begin{cases}
L_t,
&
X_{i,t}<L_t,
\\
X_{i,t},
&
L_t\leq X_{i,t}\leq U_t,
\\
U_t,
&
X_{i,t}>U_t.
\end{cases}
}
\end{equation}

The observation remains in the sample.


\subsubsection{Truncation}

Terminology varies across fields. In this document, \emph{truncation} refers to
treating observations outside the admissible interval as missing:

\begin{equation}
\boxed{
\widetilde{X}_{i,t}
=
\begin{cases}
X_{i,t},
&
L_t\leq X_{i,t}\leq U_t,
\\
\mathrm{NA},
&
\text{otherwise}.
\end{cases}
}
\end{equation}

The security itself need not necessarily be removed; only the affected variable
is invalidated for that observation.


\subsubsection{Observation Removal}

A more aggressive procedure removes the entire observation from the
cross-section:

\begin{equation}
i
\notin
\mathcal{S}_t
\end{equation}

if

\begin{equation}
X_{i,t}
\notin
[L_t,U_t].
\end{equation}

This means that all variables associated with the corresponding security-date
observation are excluded from that particular analysis.


\subsubsection{Comparison}

The three approaches have different consequences:

\begin{equation}
\boxed{
\begin{array}{lll}
\text{Winsorization}
&
\rightarrow
&
\text{retain observation, reduce magnitude},
\\[0.15cm]
\text{Truncation}
&
\rightarrow
&
\text{retain row, invalidate variable},
\\[0.15cm]
\text{Removal}
&
\rightarrow
&
\text{remove observation entirely}.
\end{array}
}
\end{equation}

Winsorization generally preserves cross-sectional coverage better, whereas
removal may introduce sample-selection effects if extreme observations are
economically non-random.


\subsubsection{Clipping}

A deterministic bound may also be imposed using economically defined limits.

For example,

\begin{equation}
\widetilde{X}_{i,t}
=
\min
\left[
\max
\left(
X_{i,t},
L
\right),
U
\right].
\end{equation}

Unlike quantile winsorization, the thresholds $L$ and $U$ are fixed rather than
estimated from the data.

This may be appropriate when the variable has known economic or structural
bounds.


\subsubsection{Effect on Ranking}

For rank-based strategies, winsorization may have less impact than for
magnitude-based strategies.

Suppose

\begin{equation}
X_1<X_2<\cdots<X_N.
\end{equation}

If the largest values are winsorized to a common upper threshold, several
securities may become tied:

\begin{equation}
\widetilde{X}_{N-k}
=
\cdots
=
\widetilde{X}_{N}.
\end{equation}

The ranking method must therefore specify how ties are handled.

By contrast, if only the order matters and winsorization does not change the
ordering, portfolio membership may remain unchanged.


\subsubsection{Effect on Linear Models}

Magnitude-based models can be much more sensitive.

Consider

\begin{equation}
Y_i
=
\alpha
+
\beta X_i
+
\varepsilon_i.
\end{equation}

An observation with unusually large $|X_i|$ may have high leverage and exert
substantial influence on

\begin{equation}
\widehat{\beta}.
\end{equation}

Winsorization, robust regression, or alternative estimators may therefore
substantially change estimated model parameters.


% ------------------------------------------------------------
\subsection{Economic Sanity Checks}
\label{subsec:economic_sanity_checks}
% ------------------------------------------------------------

Statistical outlier rules should not replace economic reasoning.

Before modifying an extreme observation, the backtester should determine
whether its magnitude is economically plausible and whether it can be explained
by the underlying data-generating process.


\subsubsection{Mechanical Ratio Explosions}

Many financial ratios can become extremely large because their denominator is
close to zero.

Consider

\begin{equation}
X_i
=
\frac{A_i}{B_i}.
\end{equation}

If

\begin{equation}
B_i
\rightarrow
0,
\end{equation}

then

\begin{equation}
|X_i|
\rightarrow
\infty
\end{equation}

even when $A_i$ is economically ordinary.

Examples include:

\begin{itemize}
    \item P/E ratios when earnings approach zero;
    \item growth rates when the lagged value is close to zero;
    \item leverage ratios when equity is very small;
    \item percentage analyst revisions relative to near-zero forecasts.
\end{itemize}

In such cases, winsorization treats the symptom rather than the structural
source of instability.


\subsubsection{Ratio Domain Checks}

Suppose

\begin{equation}
PE_{i,t}
=
\frac{
P_{i,t}
}{
EPS_{i,t}
}.
\end{equation}

If

\begin{equation}
EPS_{i,t}<0,
\end{equation}

the economic interpretation of the P/E ratio changes substantially.

The appropriate treatment may therefore be to classify the ratio as
non-meaningful rather than treating a large negative P/E as a conventional
statistical outlier.


\subsubsection{Corporate Events}

An extreme return may result from:

\begin{itemize}
    \item acquisition announcements;
    \item bankruptcy;
    \item spin-offs;
    \item special dividends;
    \item rights issues;
    \item incorrectly adjusted stock splits.
\end{itemize}

Before applying statistical treatment, the data should therefore be checked
against relevant corporate-action information.


\subsubsection{Units and Scaling}

Extreme observations may indicate inconsistent units.

For example, one provider may store a percentage as

\begin{equation}
5
\end{equation}

while another stores the same quantity as

\begin{equation}
0.05.
\end{equation}

Similarly, monetary quantities may be expressed in:

\begin{itemize}
    \item units;
    \item thousands;
    \item millions;
    \item different currencies.
\end{itemize}

A unit mismatch can easily appear statistically as an extreme observation.

Therefore,

\begin{equation}
\boxed{
\text{Outlier Detection}
\text{ should occur only after }
\text{unit and structural validation}.
}
\end{equation}


\subsubsection{Group-Conditional Distributions}

A variable may have systematically different distributions across industries,
countries, or other groups.

Suppose

\begin{equation}
g(i,t)
\end{equation}

denotes the industry of security $i$.

Instead of computing global thresholds

\begin{equation}
Q_{\alpha,t},
\end{equation}

one may compute group-specific thresholds

\begin{equation}
Q_{\alpha,g,t}.
\end{equation}

Then winsorization becomes

\begin{equation}
\widetilde{X}_{i,t}
=
\min
\left[
\max
\left(
X_{i,t},
Q_{\alpha,g(i,t),t}
\right),
Q_{1-\alpha,g(i,t),t}
\right].
\end{equation}

This can be useful when economic distributions differ structurally across
groups.

However, small groups may produce unstable threshold estimates.


\subsubsection{Treatment Hierarchy}

A robust outlier process can therefore be summarized as

\begin{equation}
\boxed{
\begin{aligned}
\text{Extreme Observation}
&\rightarrow
\text{Structural Validation}
\rightarrow
\text{Economic Interpretation}
\\[0.1cm]
&\rightarrow
\text{Outlier Detection}
\rightarrow
\text{Treatment Decision}.
\end{aligned}
}
\end{equation}

The ordering is important.

Statistical clipping should not be used to conceal errors that should instead
be corrected during data cleaning.


% ------------------------------------------------------------
\subsection{Point-in-Time Outlier Treatment}
\label{subsec:pit_outlier_treatment}
% ------------------------------------------------------------

Outlier treatment itself forms part of the historical information-processing
pipeline and must therefore respect chronology.

For a cross-sectional transformation,

\begin{equation}
\widetilde{X}_{i,t}
=
\mathcal{O}
\left(
X_{i,t};
\theta_t
\right),
\end{equation}

the parameters $\theta_t$ may use the cross-section available at time $t$:

\begin{equation}
\theta_t
=
g
\left(
\left\{
X_{j,t}
:
j\in\mathcal{U}_t
\right\}
\right).
\end{equation}

For a time-series transformation,

\begin{equation}
\theta_{i,t}
=
g
\left(
X_{i,t-1},
X_{i,t-2},
\ldots,
X_{i,t-L}
\right),
\end{equation}

if the signal is constructed before $X_{i,t}$ becomes fully observable.

More generally,

\begin{equation}
\boxed{
\theta_{i,t}
\in
\mathcal{I}_t.
}
\end{equation}

Global full-sample preprocessing violates this principle whenever future
observations influence historical thresholds.


% ------------------------------------------------------------
\subsection{Outlier-Treatment Diagnostics}
\label{subsec:outlier_diagnostics}
% ------------------------------------------------------------

Outlier treatment should be monitored quantitatively.

For each date $t$, define the fraction of observations affected by the
treatment:

\begin{equation}
\boxed{
p_t^{\mathrm{out}}
=
\frac{
1
}{
N_t
}
\sum_{i=1}^{N_t}
\mathbb{I}
\left(
\widetilde{X}_{i,t}
\neq
X_{i,t}
\right).
}
\end{equation}

For quantile winsorization at symmetric level $\alpha$, one would expect
approximately

\begin{equation}
p_t^{\mathrm{out}}
\approx
2\alpha,
\end{equation}

subject to ties and finite-sample effects.

Useful diagnostics include:

\begin{itemize}
    \item number and proportion of treated observations;
    \item lower and upper thresholds through time;
    \item securities repeatedly classified as outliers;
    \item industries or countries disproportionately affected;
    \item pre- and post-treatment distributions;
    \item effect on signal ranks;
    \item effect on model coefficients;
    \item effect on portfolio turnover and performance.
\end{itemize}

A sudden change in

\begin{equation}
p_t^{\mathrm{out}}
\end{equation}

may indicate a regime change, a data-quality problem, or a structural change in
the variable rather than an ordinary increase in outliers.


% ------------------------------------------------------------
\subsection{Output of the Outlier Layer}
\label{subsec:outlier_layer_output}
% ------------------------------------------------------------

The output of the outlier layer should preserve both the original and treated
observations.

For each variable, one may retain

\begin{equation}
\boxed{
\left(
X_{i,t}^{\mathrm{clean}},
X_{i,t}^{\mathrm{treated}},
O_{i,t},
M_{i,t}^{\mathrm{treatment}}
\right),
}
\end{equation}

where:

\begin{itemize}
    \item $X_{i,t}^{\mathrm{clean}}$ is the economically validated observation;
    \item $X_{i,t}^{\mathrm{treated}}$ is the value passed downstream;
    \item $O_{i,t}$ is an outlier indicator;
    \item $M_{i,t}^{\mathrm{treatment}}$ records the applied treatment.
\end{itemize}

The data pipeline therefore becomes

\begin{equation}
\boxed{
\begin{aligned}
\mathcal{D}^{\mathrm{raw}}
&\rightarrow
\mathcal{D}^{\mathrm{PIT}}
\rightarrow
\mathcal{D}^{\mathrm{clean}}
\\[0.1cm]
&\rightarrow
\mathcal{D}^{\mathrm{treated}}
\rightarrow
\text{Feature Engineering}
\rightarrow
\text{Signal Construction}.
\end{aligned}
}
\end{equation}

The distinction between these stages should remain explicit:

\begin{equation}
\boxed{
\begin{aligned}
\text{Data Cleaning}
&:
\text{ Is the observation valid?}
\\
\text{Outlier Detection}
&:
\text{ Is the valid observation unusually extreme?}
\\
\text{Outlier Treatment}
&:
\text{ How much influence should that extreme observation retain?}
\end{aligned}
}
\end{equation}


% ============================================================
\part{Signal Construction}
% ============================================================

% ============================================================
\section{From Firm-Level Characteristics to Raw Scores}
\label{sec:characteristics_to_raw_scores}
% ============================================================

After data cleaning and outlier treatment, the next stage of the investment
process transforms economically meaningful firm-level characteristics into
cross-sectionally comparable investment scores.

Let

\begin{equation}
x_{i,t}^{(k)}
\end{equation}

denote characteristic $k$ for security $i$ at decision time $t$, after all
required point-in-time construction, cleaning, and outlier treatment.

Examples of characteristics include valuation ratios, profitability measures,
earnings revisions, momentum measures, or any other firm-level quantity used by
the investment strategy.

The purpose of the scoring layer is to construct

\begin{equation}
s_{i,t}^{(k),\mathrm{raw}},
\end{equation}

where the score represents the relative attractiveness of security $i$
according to characteristic $k$.

Abstractly,

\begin{equation}
\boxed{
s_{i,t}^{(k),\mathrm{raw}}
=
\mathcal{S}_t^{(k)}
\left(
x_{i,t}^{(k)}
\right),
}
\end{equation}

where $\mathcal{S}_t^{(k)}$ denotes the chosen scoring transformation.

For a cross-sectional strategy, the transformation generally depends on the
complete eligible cross-section:

\begin{equation}
\boxed{
s_{i,t}^{(k),\mathrm{raw}}
=
\mathcal{S}^{(k)}
\left(
x_{i,t}^{(k)};
\left\{
x_{j,t}^{(k)}
:
j\in\mathcal{U}_t
\right\}
\right).
}
\end{equation}

The scoring transformation serves several possible purposes:

\begin{itemize}
    \item placing characteristics measured in different units on comparable
    scales;
    \item expressing the relative position of each security in the
    cross-section;
    \item controlling the influence of characteristic magnitude;
    \item establishing a common sign convention;
    \item facilitating the combination of multiple characteristics.
\end{itemize}

Importantly, the different transformations developed below are generally
\emph{alternative scoring conventions} rather than mandatory sequential steps.

For example,

\begin{equation}
x_{i,t}
\rightarrow
z_{i,t}
\end{equation}

and

\begin{equation}
x_{i,t}
\rightarrow
\operatorname{Rank}_{i,t}
\end{equation}

represent two different ways of constructing a score from the same
characteristic.


% ------------------------------------------------------------
\subsection{Definition of a Firm-Level Characteristic}
\label{subsec:firm_level_characteristic}
% ------------------------------------------------------------

A firm-level characteristic is a measurable quantity associated with a
security or its underlying company at a given point in time.

For characteristic $k$,

\begin{equation}
\boxed{
x_{i,t}^{(k)}
\in
\mathbb{R}
}
\end{equation}

denotes the characteristic value for security $i$ at time $t$.

For $N_t$ eligible securities, the corresponding cross-sectional characteristic
vector is

\begin{equation}
\boxed{
\mathbf{x}_t^{(k)}
=
\begin{bmatrix}
x_{1,t}^{(k)}
\\
x_{2,t}^{(k)}
\\
\vdots
\\
x_{N_t,t}^{(k)}
\end{bmatrix}
\in
\mathbb{R}^{N_t}.
}
\end{equation}

With $K$ characteristics, the complete characteristic matrix is

\begin{equation}
\boxed{
\mathbf{X}_t
=
\begin{bmatrix}
x_{1,t}^{(1)} & \cdots & x_{1,t}^{(K)}
\\
\vdots & \ddots & \vdots
\\
x_{N_t,t}^{(1)} & \cdots & x_{N_t,t}^{(K)}
\end{bmatrix}
\in
\mathbb{R}^{N_t\times K}.
}
\end{equation}


\subsubsection{Characteristic vs.\ Score}

A characteristic should be distinguished from an investment score.

The characteristic

\begin{equation}
x_{i,t}^{(k)}
\end{equation}

has an economic definition and may retain its natural units.

For example,

\begin{equation}
x_{i,t}^{(k)}
=
\frac{
\mathrm{BookEquity}_{i,t}
}{
\mathrm{MarketCapitalization}_{i,t}
}
\end{equation}

is a valuation characteristic.

The score

\begin{equation}
s_{i,t}^{(k)}
\end{equation}

instead represents how that observation is interpreted by the investment
process relative to other securities.

Thus,

\begin{equation}
\boxed{
\text{Economic Characteristic}
\rightarrow
\text{Scoring Transformation}
\rightarrow
\text{Investment Score}.
}
\end{equation}


\subsubsection{Characteristic Definition}

The exact mathematical definition of each characteristic should be specified
before its predictive performance is evaluated.

A characteristic definition includes, where relevant:

\begin{itemize}
    \item the underlying raw variables;
    \item the transformation applied to those variables;
    \item the historical lookback window;
    \item the required lag;
    \item the observation frequency;
    \item the treatment of missing observations;
    \item the sign of the expected relationship with future returns.
\end{itemize}

For example, a generic momentum characteristic may be defined as

\begin{equation}
x_{i,t}^{\mathrm{mom}}
=
\prod_{\tau=t-L_2}^{t-L_1}
\left(
1+R_{i,\tau}
\right)
-1,
\end{equation}

where $L_1$ and $L_2$ determine the exclusion and lookback periods.

Once the characteristic definition is fixed, the scoring layer should not
change its underlying economic meaning.


% ------------------------------------------------------------
\subsection{Cross-Sectional vs.\ Time-Series Transformations}
\label{subsec:cs_vs_ts_transformations}
% ------------------------------------------------------------

A characteristic can be transformed relative to two fundamentally different
reference sets.


\subsubsection{Cross-Sectional Transformation}

A cross-sectional transformation compares security $i$ with other securities
at the same time $t$.

Define

\begin{equation}
\mathcal{X}_t^{(k)}
=
\left\{
x_{j,t}^{(k)}
:
j\in\mathcal{U}_t
\right\}.
\end{equation}

Then

\begin{equation}
\boxed{
s_{i,t}^{(k)}
=
\mathcal{S}^{CS}
\left(
x_{i,t}^{(k)};
\mathcal{X}_t^{(k)}
\right).
}
\end{equation}

The resulting score answers a question of the form:

\begin{equation}
\boxed{
\text{How attractive is security }i
\text{ relative to other securities at time }t?
}
\end{equation}

This is the natural framework for cross-sectional stock-selection strategies.


\subsubsection{Time-Series Transformation}

A time-series transformation instead compares the current characteristic with
its own historical distribution.

Let

\begin{equation}
\mathcal{H}_{i,t}^{(k)}
=
\left\{
x_{i,\tau}^{(k)}
:
\tau\leq t
\right\}.
\end{equation}

Then

\begin{equation}
\boxed{
s_{i,t}^{(k),TS}
=
\mathcal{S}^{TS}
\left(
x_{i,t}^{(k)};
\mathcal{H}_{i,t}^{(k)}
\right).
}
\end{equation}

For example,

\begin{equation}
s_{i,t}^{(k),TS}
=
\frac{
x_{i,t}^{(k)}
-
\widehat{\mu}_{i,t}^{(k)}
}{
\widehat{\sigma}_{i,t}^{(k)}
}.
\end{equation}

This answers a different question:

\begin{equation}
\boxed{
\text{How unusual is the current value for security }i
\text{ relative to its own history?}
}
\end{equation}


\subsubsection{Distinction}

In general,

\begin{equation}
\boxed{
s_{i,t}^{CS}
\neq
s_{i,t}^{TS}.
}
\end{equation}

A firm may appear cheap relative to other firms while simultaneously appearing
expensive relative to its own historical valuation.

The reference distribution must therefore be chosen according to the economic
hypothesis being tested.


% ------------------------------------------------------------
\subsection{Z-Score Transformation}
\label{subsec:z_score_transformation}
% ------------------------------------------------------------

A standard cross-sectional transformation is the z-score.

For characteristic $k$, define the cross-sectional mean

\begin{equation}
\mu_t^{(k)}
=
\frac{1}{N_t}
\sum_{i=1}^{N_t}
x_{i,t}^{(k)},
\end{equation}

and standard deviation

\begin{equation}
\sigma_t^{(k)}
=
\sqrt{
\frac{1}{N_t-1}
\sum_{i=1}^{N_t}
\left(
x_{i,t}^{(k)}
-
\mu_t^{(k)}
\right)^2
}.
\end{equation}

The standardized characteristic is

\begin{equation}
\boxed{
z_{i,t}^{(k)}
=
\frac{
x_{i,t}^{(k)}
-
\mu_t^{(k)}
}{
\sigma_t^{(k)}
}.
}
\end{equation}

By construction,

\begin{equation}
\frac{1}{N_t}
\sum_{i=1}^{N_t}
z_{i,t}^{(k)}
=
0,
\end{equation}

while its cross-sectional variance is normalized according to the chosen
variance convention.


\subsubsection{Interpretation}

A positive z-score indicates that the characteristic lies above the
cross-sectional mean:

\begin{equation}
z_{i,t}^{(k)}>0
\quad\Longleftrightarrow\quad
x_{i,t}^{(k)}
>
\mu_t^{(k)}.
\end{equation}

For example,

\begin{equation}
z_{i,t}^{(k)}=2
\end{equation}

indicates that the observation lies two cross-sectional standard deviations
above the mean.


\subsubsection{Invariance to Affine Rescaling}

Suppose

\begin{equation}
y_{i,t}
=
a
+
b x_{i,t},
\qquad
b>0.
\end{equation}

Then

\begin{equation}
z_{i,t}^{y}
=
z_{i,t}^{x}.
\end{equation}

Thus, positive affine changes in units do not affect the standardized score.

If

\begin{equation}
b<0,
\end{equation}

the sign is reversed.


\subsubsection{Sensitivity to Extreme Observations}

The mean and standard deviation are sensitive to extreme values.

Consequently, even after outlier treatment,

\begin{equation}
\mu_t^{(k)}
\end{equation}

and

\begin{equation}
\sigma_t^{(k)}
\end{equation}

may be influenced by heavy-tailed cross-sectional distributions.

This motivates robust alternatives.


% ------------------------------------------------------------
\subsection{Robust Z-Score Transformation}
\label{subsec:raw_score_robust_z}
% ------------------------------------------------------------

Let

\begin{equation}
m_t^{(k)}
=
\operatorname{Median}_{i\in\mathcal{U}_t}
x_{i,t}^{(k)}
\end{equation}

and

\begin{equation}
MAD_t^{(k)}
=
\operatorname{Median}_{i\in\mathcal{U}_t}
\left|
x_{i,t}^{(k)}
-
m_t^{(k)}
\right|.
\end{equation}

A robust z-score is

\begin{equation}
\boxed{
z_{i,t}^{(k),\mathrm{rob}}
=
\frac{
x_{i,t}^{(k)}
-
m_t^{(k)}
}{
1.4826\,MAD_t^{(k)}
}.
}
\end{equation}

The factor $1.4826$ makes the MAD-based scale approximately consistent with
the standard deviation under Gaussianity.


\subsubsection{Relationship with Outlier Treatment}

The same robust statistic may appear both in outlier detection and in scoring,
but the objectives differ.

In the outlier layer,

\begin{equation}
z_{i,t}^{\mathrm{rob}}
\end{equation}

may be used to decide whether an observation is unusually extreme.

In the scoring layer,

\begin{equation}
\boxed{
s_{i,t}^{\mathrm{raw}}
=
z_{i,t}^{\mathrm{rob}}
}
\end{equation}

may itself be the investment score.

Thus,

\begin{equation}
\boxed{
\text{Same Mathematical Transformation}
\neq
\text{Same Role in the Pipeline}.
}
\end{equation}


% ------------------------------------------------------------
\subsection{Rank Transformation}
\label{subsec:rank_transformation}
% ------------------------------------------------------------

Rank transformations discard the magnitude of cross-sectional differences and
retain only their ordering.

Let

\begin{equation}
r_{i,t}^{(k)}
=
\operatorname{Rank}_t
\left(
x_{i,t}^{(k)}
\right),
\end{equation}

with the convention

\begin{equation}
r_{i,t}^{(k)}
\in
\{1,\ldots,N_t\},
\end{equation}

where rank $1$ corresponds to the smallest characteristic value and rank $N_t$
to the largest.


\subsubsection{Rank Invariance}

For any strictly increasing transformation $g$,

\begin{equation}
\boxed{
\operatorname{Rank}
\left(
g(x_{i,t})
\right)
=
\operatorname{Rank}
\left(
x_{i,t}
\right).
}
\end{equation}

Thus, ranking is invariant to monotonic transformations.

This property makes rank-based signals robust to extreme magnitudes and changes
in the scale of the characteristic.


\subsubsection{Loss of Magnitude Information}

The robustness comes at a cost.

Suppose

\begin{equation}
x_{1,t}=1,
\qquad
x_{2,t}=2,
\qquad
x_{3,t}=100.
\end{equation}

The ranks are simply

\begin{equation}
1,\quad2,\quad3.
\end{equation}

The fact that the third observation is dramatically farther from the second
than the second is from the first is discarded.

Thus,

\begin{equation}
\boxed{
\text{Rank Score}
=
\text{Ordinal Information},
}
\end{equation}

whereas a z-score retains information about relative magnitude.


\subsubsection{Ties}

If multiple observations have identical values, a tie-breaking convention is
required.

Common approaches include:

\begin{itemize}
    \item average ranks;
    \item minimum ranks;
    \item maximum ranks;
    \item deterministic secondary ordering.
\end{itemize}

Random tie breaking should generally be avoided unless the random seed is
explicitly controlled and the procedure is economically justified.


% ------------------------------------------------------------
\subsection{Percentile Transformation}
\label{subsec:percentile_transformation}
% ------------------------------------------------------------

Raw ranks depend mechanically on the number of securities in the universe.

A percentile transformation maps ranks onto a common interval.

One convention is

\begin{equation}
\boxed{
p_{i,t}^{(k)}
=
\frac{
r_{i,t}^{(k)}-1
}{
N_t-1
},
}
\end{equation}

for

\begin{equation}
N_t>1.
\end{equation}

Then

\begin{equation}
p_{i,t}^{(k)}
\in
[0,1].
\end{equation}

The lowest-ranked security receives

\begin{equation}
p_{i,t}^{(k)}=0,
\end{equation}

and the highest-ranked security receives

\begin{equation}
p_{i,t}^{(k)}=1.
\end{equation}


\subsubsection{Centered Percentile Score}

For long--short applications, it is often convenient to center the percentile
score around zero:

\begin{equation}
\boxed{
s_{i,t}^{(k)}
=
2p_{i,t}^{(k)}-1.
}
\end{equation}

Then

\begin{equation}
s_{i,t}^{(k)}
\in
[-1,1].
\end{equation}

The bottom of the distribution receives scores close to $-1$, the median
receives scores close to zero, and the top receives scores close to $+1$.


\subsubsection{Alternative Percentile Conventions}

Other conventions are possible, such as

\begin{equation}
p_{i,t}
=
\frac{r_{i,t}}{N_t+1}
\end{equation}

or

\begin{equation}
p_{i,t}
=
\frac{r_{i,t}-0.5}{N_t}.
\end{equation}

These differ slightly in finite samples.

The precise convention should therefore be specified explicitly, particularly
if percentile scores are subsequently transformed through an inverse
distribution function.


% ------------------------------------------------------------
\subsection{Quantile and Decile Scores}
\label{subsec:quantile_decile_scores}
% ------------------------------------------------------------

A continuous cross-sectional score may be discretized into $Q$ groups.

Let

\begin{equation}
q_{i,t}
\in
\{1,\ldots,Q\}
\end{equation}

denote the quantile group assigned to security $i$.

For deciles,

\begin{equation}
Q=10.
\end{equation}

A simple approximation is

\begin{equation}
\boxed{
q_{i,t}
=
1
+
\left\lfloor
Q
\frac{
r_{i,t}-1
}{
N_t
}
\right\rfloor,
}
\end{equation}

with the final value restricted to

\begin{equation}
q_{i,t}
\in
\{1,\ldots,Q\}.
\end{equation}


\subsubsection{Discrete Scores}

The group label may itself be transformed into a centered score:

\begin{equation}
\boxed{
s_{i,t}^{Q}
=
\frac{
2q_{i,t}
-
(Q+1)
}{
Q-1
}.
}
\end{equation}

This maps

\begin{equation}
q=1
\end{equation}

to

\begin{equation}
s^Q=-1,
\end{equation}

and

\begin{equation}
q=Q
\end{equation}

to

\begin{equation}
s^Q=+1.
\end{equation}


\subsubsection{Score vs.\ Portfolio Formation}

Quantile assignment at this stage should be distinguished from constructing a
quantile portfolio.

Here,

\begin{equation}
q_{i,t}
\end{equation}

is simply a discrete representation of the signal.

Later, portfolio construction may decide to hold only

\begin{equation}
q_{i,t}=1
\end{equation}

and

\begin{equation}
q_{i,t}=Q,
\end{equation}

or to assign weights across all groups.

Thus,

\begin{equation}
\boxed{
\text{Quantile Score}
\neq
\text{Quantile Portfolio}.
}
\end{equation}


% ------------------------------------------------------------
\subsection{Gaussian Rank Transformation}
\label{subsec:gaussian_rank_transformation}
% ------------------------------------------------------------

A Gaussian rank transformation combines the robustness of ranking with a
continuous score having an approximately Gaussian cross-sectional shape.

First construct a percentile strictly inside $(0,1)$, for example

\begin{equation}
u_{i,t}
=
\frac{
r_{i,t}-0.5
}{
N_t
}.
\end{equation}

Then apply the inverse standard-normal cumulative distribution function:

\begin{equation}
\boxed{
g_{i,t}
=
\Phi^{-1}
\left(
u_{i,t}
\right),
}
\end{equation}

where

\begin{equation}
\Phi(\cdot)
\end{equation}

is the standard-normal cumulative distribution function.


\subsubsection{Why Interior Percentiles Are Required}

If percentile ranks equal exactly zero or one, then

\begin{equation}
\Phi^{-1}(0)
=
-\infty
\end{equation}

and

\begin{equation}
\Phi^{-1}(1)
=
+\infty.
\end{equation}

This is why a convention such as

\begin{equation}
u_{i,t}
=
\frac{r_{i,t}-0.5}{N_t}
\end{equation}

is preferable when applying an inverse-normal transformation.


\subsubsection{Interpretation}

The transformation preserves the ordering:

\begin{equation}
x_{i,t}<x_{j,t}
\quad\Rightarrow\quad
g_{i,t}<g_{j,t},
\end{equation}

but imposes a predetermined relationship between rank and score magnitude.

Securities near the center receive scores close to zero, while securities in
the tails receive increasingly large absolute scores.

Thus,

\begin{equation}
\boxed{
\text{Rank Information}
+
\text{Gaussian Score Geometry}.
}
\end{equation}


% ------------------------------------------------------------
\subsection{Direction and Sign Convention}
\label{subsec:direction_sign_convention}
% ------------------------------------------------------------

The raw mathematical transformation does not necessarily have the desired
investment direction.

For each characteristic $k$, define a direction parameter

\begin{equation}
d_k
\in
\{-1,+1\}.
\end{equation}

The direction-adjusted score is

\begin{equation}
\boxed{
s_{i,t}^{(k)}
=
d_k
\widetilde{s}_{i,t}^{(k)},
}
\end{equation}

where $\widetilde{s}_{i,t}^{(k)}$ is the score before sign adjustment.

A useful global convention is

\begin{equation}
\boxed{
s_{i,t}^{(k)}>0
\quad\Longleftrightarrow\quad
\text{higher expected attractiveness}.
}
\end{equation}

Consequently,

\begin{equation}
s_{i,t}^{(k)}<0
\end{equation}

indicates lower expected attractiveness.


\subsubsection{Example}

Suppose higher profitability is expected to predict higher returns.

Then

\begin{equation}
d_{\mathrm{profitability}}=+1.
\end{equation}

Suppose lower valuation multiples are considered more attractive.

If the raw characteristic is P/E, then

\begin{equation}
d_{\mathrm{PE}}=-1.
\end{equation}

The scoring convention therefore ensures that all characteristics point in the
same investment direction before they are combined.


\subsubsection{Direction Must Be Defined Ex Ante}

The sign should follow from the economic hypothesis or predefined research
design.

Selecting

\begin{equation}
d_k
\end{equation}

after observing which direction produced the best historical performance is
itself a form of data snooping.


% ------------------------------------------------------------
\subsection{Score Clipping and Bounding}
\label{subsec:score_clipping}
% ------------------------------------------------------------

Even after outlier treatment at the characteristic level, the scoring
transformation may produce scores with undesirable magnitudes.

For example, a z-score or Gaussian rank score is theoretically unbounded.

A symmetric clipping rule is

\begin{equation}
\boxed{
s_{i,t}^{\mathrm{clip}}
=
\min
\left[
\max
\left(
s_{i,t},
-c
\right),
c
\right].
}
\end{equation}

For example,

\begin{equation}
c=3
\end{equation}

implies

\begin{equation}
s_{i,t}^{\mathrm{clip}}
\in
[-3,3].
\end{equation}


\subsubsection{Characteristic Winsorization vs.\ Score Clipping}

These two operations occur at different levels.

Characteristic-level treatment modifies

\begin{equation}
x_{i,t},
\end{equation}

whereas score clipping modifies

\begin{equation}
s_{i,t}.
\end{equation}

Thus,

\begin{equation}
\boxed{
x_{i,t}
\rightarrow
x_{i,t}^{\mathrm{treated}}
\rightarrow
s_{i,t}
\rightarrow
s_{i,t}^{\mathrm{clip}}.
}
\end{equation}

They should not be treated as interchangeable operations.


\subsubsection{Bounded Nonlinear Transformations}

Instead of hard clipping, an unbounded score may be mapped smoothly into a
bounded interval.

For example,

\begin{equation}
\boxed{
s_{i,t}^{\mathrm{bounded}}
=
\tanh
\left(
\frac{s_{i,t}}{\kappa}
\right),
}
\end{equation}

where

\begin{equation}
\kappa>0
\end{equation}

controls the degree of compression.

Then

\begin{equation}
s_{i,t}^{\mathrm{bounded}}
\in
(-1,1).
\end{equation}

Unlike hard clipping, this transformation compresses large values smoothly.


% ------------------------------------------------------------
\subsection{Combining Multiple Characteristics}
\label{subsec:combining_characteristics}
% ------------------------------------------------------------

A strategy may use several characteristics simultaneously.

Suppose $K$ characteristics have been transformed into directionally consistent
scores:

\begin{equation}
s_{i,t}^{(1)},
\ldots,
s_{i,t}^{(K)}.
\end{equation}

Collect them into

\begin{equation}
\mathbf{s}_{i,t}
=
\begin{bmatrix}
s_{i,t}^{(1)}
\\
\vdots
\\
s_{i,t}^{(K)}
\end{bmatrix}
\in
\mathbb{R}^{K}.
\end{equation}


\subsubsection{Equal-Weighted Combination}

The simplest composite score is

\begin{equation}
\boxed{
s_{i,t}^{\mathrm{comp}}
=
\frac{1}{K}
\sum_{k=1}^{K}
s_{i,t}^{(k)}.
}
\end{equation}

This requires the individual scores to be on sufficiently comparable scales.


\subsubsection{Weighted Combination}

More generally,

\begin{equation}
\boxed{
s_{i,t}^{\mathrm{comp}}
=
\sum_{k=1}^{K}
\omega_{k,t}
s_{i,t}^{(k)},
}
\end{equation}

where

\begin{equation}
\boldsymbol{\omega}_t
=
\begin{bmatrix}
\omega_{1,t}
&
\cdots
&
\omega_{K,t}
\end{bmatrix}^{\top}.
\end{equation}

In vector form,

\begin{equation}
\boxed{
s_{i,t}^{\mathrm{comp}}
=
\boldsymbol{\omega}_t^{\top}
\mathbf{s}_{i,t}.
}
\end{equation}

A common normalization is

\begin{equation}
\sum_{k=1}^{K}
\omega_{k,t}
=
1,
\end{equation}

although this is not mathematically required.


\subsubsection{Static vs.\ Dynamic Combination Weights}

Combination weights may be static:

\begin{equation}
\omega_{k,t}
=
\omega_k,
\end{equation}

or dynamic:

\begin{equation}
\omega_{k,t}
=
g_k
\left(
\mathcal{I}_t
\right).
\end{equation}

Dynamic weights may depend on historically estimated predictive performance,
risk, regime variables, or other information available at time $t$.

Any estimated weighting rule must itself respect the point-in-time constraint.


\subsubsection{Missing Component Scores}

If some component scores are unavailable, the aggregation rule must specify
how the composite score is constructed.

Let

\begin{equation}
\mathcal{K}_{i,t}^{\mathrm{avail}}
=
\left\{
k:
s_{i,t}^{(k)}
\text{ is available}
\right\}.
\end{equation}

One possible renormalized combination is

\begin{equation}
\boxed{
s_{i,t}^{\mathrm{comp}}
=
\frac{
\sum_{k\in\mathcal{K}_{i,t}^{\mathrm{avail}}}
\omega_{k,t}s_{i,t}^{(k)}
}{
\sum_{k\in\mathcal{K}_{i,t}^{\mathrm{avail}}}
\omega_{k,t}
},
}
\end{equation}

provided the denominator is non-zero.

However, this changes the effective characteristic mix across securities.

The resulting scores may therefore no longer be directly comparable if
missingness is substantial.


\subsubsection{Correlation Between Characteristics}

A simple average implicitly treats each component as a separate contribution.

If two characteristics are highly correlated,

\begin{equation}
\operatorname{Corr}
\left(
s_t^{(k)},
s_t^{(\ell)}
\right)
\approx1,
\end{equation}

including both may effectively overweight the same underlying information.

Thus, signal combination should consider not only individual predictive
strength but also redundancy across characteristics.

More sophisticated combination methods may therefore use covariance,
regression, optimization, or predictive models.

Such approaches connect naturally to the model-based signal construction
developed in the following section.


% ------------------------------------------------------------
\subsection{Final Raw Score}
\label{subsec:final_raw_score}
% ------------------------------------------------------------

The output of the characteristic-based signal-construction layer is a raw score
for every eligible security.

Define

\begin{equation}
\boxed{
s_{i,t}^{\mathrm{raw}}
}
\end{equation}

as the final characteristic-based raw score before explicit risk or factor
neutralization.

Collecting all securities gives

\begin{equation}
\boxed{
\mathbf{s}_t^{\mathrm{raw}}
=
\begin{bmatrix}
s_{1,t}^{\mathrm{raw}}
\\
s_{2,t}^{\mathrm{raw}}
\\
\vdots
\\
s_{N_t,t}^{\mathrm{raw}}
\end{bmatrix}
\in
\mathbb{R}^{N_t}.
}
\end{equation}


\subsubsection{Interpretation}

Under the sign convention adopted above,

\begin{equation}
s_{i,t}^{\mathrm{raw}}>0
\end{equation}

represents relatively attractive securities, whereas

\begin{equation}
s_{i,t}^{\mathrm{raw}}<0
\end{equation}

represents relatively unattractive securities.

However,

\begin{equation}
\boxed{
s_{i,t}^{\mathrm{raw}}
\neq
w_{i,t}.
}
\end{equation}

A score is an investment signal, not a portfolio weight.

Its magnitude expresses relative conviction under the chosen scoring
convention, but does not yet specify:

\begin{itemize}
    \item portfolio gross exposure;
    \item net exposure;
    \item position size;
    \item beta exposure;
    \item industry exposure;
    \item factor exposure;
    \item leverage;
    \item risk contribution;
    \item trading constraints.
\end{itemize}

These are handled by subsequent stages of the investment process.


\subsubsection{Characteristic-Based Signal Pipeline}

For a single characteristic, the complete transformation can be represented as

\begin{equation}
\boxed{
\begin{aligned}
x_{i,t}^{\mathrm{raw}}
&\rightarrow
x_{i,t}^{\mathrm{clean}}
\rightarrow
x_{i,t}^{\mathrm{treated}}
\\[0.1cm]
&\rightarrow
\mathcal{S}_t
\left(
x_{i,t}^{\mathrm{treated}}
\right)
\rightarrow
s_{i,t}^{\mathrm{raw}}.
\end{aligned}
}
\end{equation}

For multiple characteristics,

\begin{equation}
\boxed{
\begin{aligned}
\mathbf{x}_{i,t}
&\rightarrow
\mathbf{x}_{i,t}^{\mathrm{clean}}
\rightarrow
\mathbf{x}_{i,t}^{\mathrm{treated}}
\\[0.1cm]
&\rightarrow
\mathbf{s}_{i,t}
\rightarrow
\mathcal{A}
\left(
\mathbf{s}_{i,t}
\right)
\rightarrow
s_{i,t}^{\mathrm{raw}},
\end{aligned}
}
\end{equation}

where $\mathcal{A}$ denotes the chosen score-aggregation operator.


\subsubsection{Interface with the Downstream Engine}

The key architectural principle is that the subsequent backtesting engine
receives

\begin{equation}
\boxed{
\mathbf{s}_t^{\mathrm{raw}}
\in
\mathbb{R}^{N_t}
}
\end{equation}

regardless of how the score was produced.

A characteristic-based strategy may generate

\begin{equation}
\mathbf{s}_t^{\mathrm{raw}}
=
\mathcal{S}
\left(
\mathbf{X}_t
\right),
\end{equation}

while a predictive model developed in the next section may generate

\begin{equation}
\mathbf{s}_t^{\mathrm{raw}}
=
\mathcal{P}
\left(
\widehat{\mathbf{y}}_{t+h|t}
\right).
\end{equation}

Both objects then enter the same downstream pipeline:

\begin{equation}
\boxed{
\begin{aligned}
\mathbf{s}_t^{\mathrm{raw}}
&\rightarrow
\text{Signal Neutralization}
\rightarrow
\mathbf{s}_t^{\mathrm{neutral}}
\\[0.1cm]
&\rightarrow
\text{Final Signal}
\rightarrow
\mathbf{s}_t^{\mathrm{final}}
\rightarrow
\text{Portfolio Construction}.
\end{aligned}
}
\end{equation}

This separation allows signal generation and portfolio construction to remain
modular: the downstream portfolio engine need not know whether the raw score
originated from a single firm characteristic, a composite score, an
econometric model, or a machine-learning model.
```latex
% ============================================================
\section{Model-Based Signal Construction}
\label{sec:model_based_signal_construction}
% ============================================================

A characteristic-based strategy directly transforms one or several firm-level
characteristics into investment scores. A model-based strategy instead estimates
a predictive relationship between a set of characteristics and a future target.

The general architecture is

\begin{equation}
\boxed{
\mathbf{X}_t
\rightarrow
\widehat{\mathbf{y}}_{t+h|t}
\rightarrow
\mathbf{s}_t^{\mathrm{raw}},
}
\end{equation}

where:

\begin{itemize}
    \item $\mathbf{X}_t$ denotes the information available at time $t$;
    \item $\widehat{\mathbf{y}}_{t+h|t}$ denotes model predictions for horizon
    $t+h$ formed using information available at $t$;
    \item $\mathbf{s}_t^{\mathrm{raw}}$ denotes the final model-based raw
    investment score.
\end{itemize}

The fundamental distinction from direct characteristic scoring is therefore

\begin{equation}
\boxed{
\begin{aligned}
\text{Characteristic-Based:}
&\quad
x_{i,t}
\rightarrow
s_{i,t}^{\mathrm{raw}},
\\[0.15cm]
\text{Model-Based:}
&\quad
\mathbf{x}_{i,t}
\rightarrow
\widehat{y}_{i,t+h|t}
\rightarrow
s_{i,t}^{\mathrm{raw}}.
\end{aligned}
}
\end{equation}

Once $\mathbf{s}_t^{\mathrm{raw}}$ has been produced, the downstream
neutralization and portfolio-construction engine is identical in both cases.


% ------------------------------------------------------------
\subsection{General Predictive Framework}
\label{subsec:general_predictive_framework}
% ------------------------------------------------------------

Let

\begin{equation}
\mathbf{x}_{i,t}
=
\begin{bmatrix}
x_{i,t}^{(1)}
&
\cdots
&
x_{i,t}^{(K)}
\end{bmatrix}^{\top}
\in
\mathbb{R}^{K}
\end{equation}

denote the vector of $K$ characteristics available for security $i$ at time
$t$.

Let

\begin{equation}
y_{i,t+h}
\end{equation}

denote the future target variable over prediction horizon $h$.

The general predictive problem is

\begin{equation}
\boxed{
y_{i,t+h}
=
f
\left(
\mathbf{x}_{i,t};
\theta
\right)
+
\varepsilon_{i,t+h},
}
\end{equation}

where:

\begin{itemize}
    \item $f(\cdot)$ is the predictive model;
    \item $\theta$ denotes model parameters;
    \item $\varepsilon_{i,t+h}$ denotes the prediction error.
\end{itemize}

At decision time $t$, model parameters must be estimated using only data
available before or at that time:

\begin{equation}
\widehat{\theta}_t
=
\mathcal{A}
\left(
\mathcal{D}_t^{\mathrm{train}}
\right),
\end{equation}

where $\mathcal{A}$ denotes the estimation algorithm.

The prediction is then

\begin{equation}
\boxed{
\widehat{y}_{i,t+h|t}
=
f
\left(
\mathbf{x}_{i,t};
\widehat{\theta}_t
\right).
}
\end{equation}

The notation

\begin{equation}
\widehat{y}_{i,t+h|t}
\end{equation}

explicitly emphasizes that the prediction concerns time $t+h$ but is formed
using information available at time $t$.


% ------------------------------------------------------------
\subsection{Target Variable Definition}
\label{subsec:model_target_variable}
% ------------------------------------------------------------

The target variable determines what the model is attempting to predict.

For return-prediction strategies, a natural target is the future simple return

\begin{equation}
\boxed{
y_{i,t+h}
=
R_{i,t\rightarrow t+h}.
}
\end{equation}

However, alternative targets may include:

\begin{itemize}
    \item future excess return;
    \item benchmark-relative return;
    \item market- or factor-adjusted return;
    \item future return rank;
    \item future return sign;
    \item probability of outperforming;
    \item future volatility;
    \item future drawdown;
    \item future risk-adjusted return.
\end{itemize}


\subsubsection{Absolute vs.\ Relative Targets}

An absolute return target is

\begin{equation}
y_{i,t+h}
=
R_{i,t\rightarrow t+h}.
\end{equation}

A benchmark-relative target is

\begin{equation}
\boxed{
y_{i,t+h}
=
R_{i,t\rightarrow t+h}
-
R_{B,t\rightarrow t+h}.
}
\end{equation}

A cross-sectionally demeaned target may be defined as

\begin{equation}
y_{i,t+h}
=
R_{i,t\rightarrow t+h}
-
\frac{1}{N_t}
\sum_{j=1}^{N_t}
R_{j,t\rightarrow t+h}.
\end{equation}

Relative targets may be more closely aligned with cross-sectional stock
selection because they remove part of the common market component.


\subsubsection{Continuous vs.\ Discrete Targets}

For regression,

\begin{equation}
y_{i,t+h}
\in
\mathbb{R}.
\end{equation}

For binary classification,

\begin{equation}
y_{i,t+h}
=
\mathbb{I}
\left(
R_{i,t\rightarrow t+h}>0
\right).
\end{equation}

For relative classification,

\begin{equation}
y_{i,t+h}
=
\mathbb{I}
\left(
R_{i,t\rightarrow t+h}
>
R_{B,t\rightarrow t+h}
\right).
\end{equation}

For ranking problems, the target may instead be

\begin{equation}
y_{i,t+h}
=
\operatorname{Rank}_t
\left(
R_{i,t\rightarrow t+h}
\right).
\end{equation}


\subsubsection{Prediction Horizon}

The target horizon $h$ should correspond to the economic horizon of the
strategy.

A mismatch between target horizon and portfolio implementation can materially
affect both the statistical and economic interpretation of the model.

For example, predicting

\begin{equation}
R_{i,t\rightarrow t+20}
\end{equation}

while rebuilding the portfolio every day generates overlapping target periods.

This creates dependence across training observations and must be considered in
validation and inference.


% ------------------------------------------------------------
\subsection{Input Characteristics and Design Matrix}
\label{subsec:input_design_matrix}
% ------------------------------------------------------------

Suppose $K$ firm-level characteristics are available.

For security $i$ at time $t$,

\begin{equation}
\mathbf{x}_{i,t}
=
\begin{bmatrix}
x_{i,t}^{(1)}
\\
\vdots
\\
x_{i,t}^{(K)}
\end{bmatrix}.
\end{equation}

Across the $N_t$ securities observed at time $t$, the design matrix is

\begin{equation}
\boxed{
\mathbf{X}_t
=
\begin{bmatrix}
(\mathbf{x}_{1,t})^\top
\\
\vdots
\\
(\mathbf{x}_{N_t,t})^\top
\end{bmatrix}
\in
\mathbb{R}^{N_t\times K}.
}
\end{equation}

The corresponding target vector is

\begin{equation}
\mathbf{y}_{t+h}
=
\begin{bmatrix}
y_{1,t+h}
\\
\vdots
\\
y_{N_t,t+h}
\end{bmatrix}.
\end{equation}


\subsubsection{Panel Representation}

Across multiple dates, observations may be stacked into a panel:

\begin{equation}
\mathcal{P}
=
\left\{
(\mathbf{x}_{i,t},y_{i,t+h})
:
i\in\mathcal{U}_t,
\;
t\in\mathcal{T}
\right\}.
\end{equation}

The corresponding stacked design matrix can be represented as

\begin{equation}
\mathbf{X}
\in
\mathbb{R}^{M\times K},
\end{equation}

where

\begin{equation}
M
=
\sum_t N_t
\end{equation}

is the total number of security-date observations.


\subsubsection{Preprocessing Parameters}

Some models require feature transformations such as scaling or centering.

If feature $k$ is standardized using training-sample parameters,

\begin{equation}
\widetilde{x}_{i,t}^{(k)}
=
\frac{
x_{i,t}^{(k)}
-
\widehat{\mu}_{k,\mathrm{train}}
}{
\widehat{\sigma}_{k,\mathrm{train}}
}.
\end{equation}

Crucially,

\begin{equation}
\boxed{
\widehat{\mu}_{k,\mathrm{train}},
\widehat{\sigma}_{k,\mathrm{train}}
}
\end{equation}

must be estimated exclusively from the corresponding training sample.

Preprocessing is therefore part of the fitted model pipeline rather than an
operation that may be estimated once using the entire dataset.


% ------------------------------------------------------------
\subsection{Econometric Models}
\label{subsec:econometric_models}
% ------------------------------------------------------------

Econometric models specify an explicit parametric relationship between the
target and explanatory characteristics.


\subsubsection{Linear Regression}

A cross-sectional or pooled linear model takes the form

\begin{equation}
\boxed{
y_{i,t+h}
=
\alpha
+
\mathbf{x}_{i,t}^{\top}\beta
+
\varepsilon_{i,t+h}.
}
\end{equation}

Stacking observations gives

\begin{equation}
\mathbf{y}
=
\mathbf{X}\beta
+
\varepsilon,
\end{equation}

after incorporating an intercept into the design matrix if required.

The Ordinary Least Squares estimator is

\begin{equation}
\boxed{
\widehat{\beta}^{OLS}
=
\left(
\mathbf{X}^{\top}\mathbf{X}
\right)^{-1}
\mathbf{X}^{\top}\mathbf{y},
}
\end{equation}

when $\mathbf{X}^{\top}\mathbf{X}$ is invertible.


\subsubsection{Weighted Least Squares}

Observations may receive different estimation weights.

Let

\begin{equation}
\mathbf{W}
=
\operatorname{diag}
\left(
\omega_1,\ldots,\omega_M
\right).
\end{equation}

Then

\begin{equation}
\boxed{
\widehat{\beta}^{WLS}
=
\left(
\mathbf{X}^{\top}
\mathbf{W}
\mathbf{X}
\right)^{-1}
\mathbf{X}^{\top}
\mathbf{W}
\mathbf{y}.
}
\end{equation}

Weights may reflect:

\begin{itemize}
    \item recency;
    \item market capitalization;
    \item observation reliability;
    \item volatility;
    \item sampling design.
\end{itemize}


\subsubsection{Regularized Linear Models}

Ridge regression solves

\begin{equation}
\boxed{
\widehat{\beta}^{\mathrm{Ridge}}
=
\underset{\beta}{\arg\min}
\left[
\|\mathbf{y}-\mathbf{X}\beta\|_2^2
+
\lambda
\|\beta\|_2^2
\right].
}
\end{equation}

Lasso regression solves

\begin{equation}
\boxed{
\widehat{\beta}^{\mathrm{Lasso}}
=
\underset{\beta}{\arg\min}
\left[
\|\mathbf{y}-\mathbf{X}\beta\|_2^2
+
\lambda
\|\beta\|_1
\right].
}
\end{equation}

Elastic Net combines both penalties:

\begin{equation}
\boxed{
\widehat{\beta}^{EN}
=
\underset{\beta}{\arg\min}
\left[
\|\mathbf{y}-\mathbf{X}\beta\|_2^2
+
\lambda
\left(
\alpha\|\beta\|_1
+
(1-\alpha)\|\beta\|_2^2
\right)
\right].
}
\end{equation}

Regularization can reduce estimation variance and improve stability when
characteristics are numerous or highly correlated.


\subsubsection{Panel Models}

Panel regressions may incorporate entity or time effects:

\begin{equation}
y_{i,t+h}
=
\alpha_i
+
\gamma_t
+
\mathbf{x}_{i,t}^{\top}\beta
+
\varepsilon_{i,t+h},
\end{equation}

where:

\begin{itemize}
    \item $\alpha_i$ denotes security-specific effects;
    \item $\gamma_t$ denotes time-specific effects.
\end{itemize}

The appropriate specification depends on whether the objective is prediction,
economic interpretation, or both.


% ------------------------------------------------------------
\subsection{Machine Learning Models}
\label{subsec:machine_learning_models}
% ------------------------------------------------------------

Machine-learning models relax the assumption of a fixed linear functional form.

The general prediction remains

\begin{equation}
\widehat{y}_{i,t+h|t}
=
f
\left(
\mathbf{x}_{i,t};
\widehat{\theta}_t
\right),
\end{equation}

but $f(\cdot)$ may be nonlinear and contain interactions that are not specified
manually.


\subsubsection{Decision Trees}

A regression tree partitions the feature space into regions

\begin{equation}
\mathcal{R}_1,
\ldots,
\mathcal{R}_M
\end{equation}

and predicts

\begin{equation}
\widehat{y}(x)
=
\sum_{m=1}^{M}
c_m
\mathbb{I}
\left(
x\in\mathcal{R}_m
\right).
\end{equation}

Trees naturally capture nonlinearities and feature interactions.


\subsubsection{Random Forests}

A Random Forest combines $B$ trees:

\begin{equation}
\boxed{
\widehat{y}(x)
=
\frac{1}{B}
\sum_{b=1}^{B}
\widehat{f}_b(x).
}
\end{equation}

Bootstrap sampling and random feature selection reduce correlation between
individual trees and can improve generalization.


\subsubsection{Gradient Boosting}

Boosting constructs an additive model:

\begin{equation}
\boxed{
\widehat{f}_M(x)
=
\sum_{m=1}^{M}
\eta
f_m(x),
}
\end{equation}

where:

\begin{itemize}
    \item $f_m$ is typically a weak learner;
    \item $\eta$ is the learning rate.
\end{itemize}

Each new learner is chosen to reduce the residual error of the existing
ensemble.

Gradient-boosted trees are particularly common in tabular financial prediction
because they can capture nonlinearities while remaining relatively efficient.


\subsubsection{Model Complexity}

Unlike linear models, machine-learning models may contain a large number of
effective degrees of freedom.

Their economic usefulness therefore depends critically on:

\begin{itemize}
    \item temporal validation;
    \item regularization;
    \item hyperparameter control;
    \item sufficient sample size;
    \item robustness across periods.
\end{itemize}

Higher in-sample predictive power does not necessarily imply higher
out-of-sample investment performance.


% ------------------------------------------------------------
\subsection{Deep Learning Models}
\label{subsec:deep_learning_models}
% ------------------------------------------------------------

Deep-learning models represent a particularly flexible class of nonlinear
predictive functions.

A feed-forward neural network can be written recursively as

\begin{equation}
\mathbf{h}^{(1)}
=
\phi
\left(
\mathbf{W}^{(1)}\mathbf{x}
+
\mathbf{b}^{(1)}
\right),
\end{equation}

and

\begin{equation}
\mathbf{h}^{(\ell)}
=
\phi
\left(
\mathbf{W}^{(\ell)}
\mathbf{h}^{(\ell-1)}
+
\mathbf{b}^{(\ell)}
\right),
\end{equation}

for hidden layers $\ell=2,\ldots,L$.

The final prediction may be

\begin{equation}
\boxed{
\widehat{y}
=
\mathbf{W}^{(L+1)}
\mathbf{h}^{(L)}
+
\mathbf{b}^{(L+1)}.
}
\end{equation}


\subsubsection{Sequential Models}

When the complete historical path is an input, the model may receive

\begin{equation}
\mathbf{X}_{i,t}^{\mathrm{seq}}
=
\left[
\mathbf{x}_{i,t-L+1},
\ldots,
\mathbf{x}_{i,t}
\right].
\end{equation}

Architectures may then include recurrent networks, temporal convolutional
models, or attention-based models.


\subsubsection{Cross-Sectional Models}

Deep-learning models may also process the entire cross-section jointly:

\begin{equation}
\mathbf{X}_t
\rightarrow
\widehat{\mathbf{y}}_{t+h|t}.
\end{equation}

Such models can potentially learn interactions across securities, but require
particular care because the cross-sectional universe changes through time.


\subsubsection{Higher Capacity and Overfitting}

Deep models generally have substantially greater representational capacity than
simple econometric models.

Consequently,

\begin{equation}
\boxed{
\text{Model Capacity}
\uparrow
\quad\Rightarrow\quad
\text{Need for Regularization and Validation}
\uparrow.
}
\end{equation}

Appropriate controls may include:

\begin{itemize}
    \item weight decay;
    \item dropout;
    \item early stopping;
    \item architecture constraints;
    \item temporal cross-validation;
    \item ensembling.
\end{itemize}


% ------------------------------------------------------------
\subsection{Regression, Classification, and Ranking Objectives}
\label{subsec:prediction_objectives}
% ------------------------------------------------------------

The model objective should be aligned with the economic decision the strategy
ultimately makes.


\subsubsection{Regression}

For continuous return prediction, a common loss is Mean Squared Error:

\begin{equation}
\boxed{
\mathcal{L}_{MSE}
=
\frac{1}{M}
\sum_{j=1}^{M}
\left(
y_j-\widehat{y}_j
\right)^2.
}
\end{equation}

Alternatively,

\begin{equation}
\mathcal{L}_{MAE}
=
\frac{1}{M}
\sum_{j=1}^{M}
|y_j-\widehat{y}_j|.
\end{equation}


\subsubsection{Classification}

For binary classification,

\begin{equation}
y_j\in\{0,1\},
\end{equation}

and the model predicts

\begin{equation}
p_j
=
P(y_j=1|\mathbf{x}_j).
\end{equation}

Binary cross-entropy is

\begin{equation}
\boxed{
\mathcal{L}_{BCE}
=
-
\frac{1}{M}
\sum_{j=1}^{M}
\left[
y_j\log(p_j)
+
(1-y_j)\log(1-p_j)
\right].
}
\end{equation}


\subsubsection{Ranking}

For cross-sectional stock selection, the exact numerical return prediction may
be less important than the ordering of securities.

The objective may therefore emphasize

\begin{equation}
\operatorname{Rank}
\left(
\widehat{y}_{i,t+h|t}
\right)
\end{equation}

relative to

\begin{equation}
\operatorname{Rank}
\left(
y_{i,t+h}
\right).
\end{equation}

A rank-based evaluation metric is the cross-sectional Spearman correlation:

\begin{equation}
\boxed{
IC_t^{\mathrm{rank}}
=
\operatorname{Corr}_{\mathrm{Spearman}}
\left(
\widehat{\mathbf{y}}_{t+h|t},
\mathbf{y}_{t+h}
\right).
}
\end{equation}

Ranking losses may be pairwise or listwise and attempt to optimize ordering
directly.


\subsubsection{Statistical Loss vs.\ Investment Objective}

The statistical training objective need not coincide with the final portfolio
objective.

For example,

\begin{equation}
\mathcal{L}_{MSE}
\end{equation}

optimizes squared prediction errors, whereas the investment process may care
primarily about:

\begin{itemize}
    \item cross-sectional ranking;
    \item long--short spread;
    \item turnover;
    \item transaction costs;
    \item net Sharpe ratio.
\end{itemize}

Thus,

\begin{equation}
\boxed{
\text{Statistical Accuracy}
\neq
\text{Economic Performance}.
}
\end{equation}

The connection between predictions and portfolio performance is therefore
evaluated only after the complete downstream portfolio-construction process has
been applied.


% ------------------------------------------------------------
\subsection{Training, Validation, and Test Samples}
\label{subsec:train_validation_test}
% ------------------------------------------------------------

Model estimation and model evaluation must be separated temporally.

Let the full historical sample be divided into:

\begin{equation}
\boxed{
\mathcal{D}^{\mathrm{train}},
\qquad
\mathcal{D}^{\mathrm{val}},
\qquad
\mathcal{D}^{\mathrm{test}}.
}
\end{equation}

These datasets serve different roles.


\subsubsection{Training Sample}

The training sample is used to estimate model parameters:

\begin{equation}
\widehat{\theta}
=
\mathcal{A}
\left(
\mathcal{D}^{\mathrm{train}}
\right).
\end{equation}


\subsubsection{Validation Sample}

The validation sample is used for:

\begin{itemize}
    \item hyperparameter selection;
    \item model selection;
    \item early stopping;
    \item architecture selection;
    \item feature-selection decisions.
\end{itemize}


\subsubsection{Test Sample}

The test sample is reserved for final evaluation.

It should not influence:

\begin{equation}
\widehat{\theta},
\qquad
\lambda,
\qquad
\text{feature selection},
\qquad
\text{model architecture}.
\end{equation}

Repeatedly using test performance to modify the model converts the test sample
into another validation sample.


\subsubsection{Chronological Ordering}

For financial prediction,

\begin{equation}
\boxed{
\mathcal{T}^{\mathrm{train}}
<
\mathcal{T}^{\mathrm{val}}
<
\mathcal{T}^{\mathrm{test}}.
}
\end{equation}

Random train-test splitting is generally inappropriate when it allows future
observations to influence models evaluated on earlier periods.


% ------------------------------------------------------------
\subsection{Expanding and Rolling Estimation Windows}
\label{subsec:expanding_rolling_windows}
% ------------------------------------------------------------

Model parameters may be estimated using either expanding or rolling historical
windows.


\subsubsection{Expanding Window}

Let the first training observation occur at $t_0$.

At decision time $t$, an expanding training set is

\begin{equation}
\boxed{
\mathcal{D}_t^{\mathrm{train,exp}}
=
\left\{
(\mathbf{x}_{i,\tau},y_{i,\tau+h})
:
t_0
\leq
\tau
\leq
t-h
\right\}.
}
\end{equation}

As $t$ increases, the training sample grows.

The corresponding estimator is

\begin{equation}
\widehat{\theta}_t^{\mathrm{exp}}
=
\mathcal{A}
\left(
\mathcal{D}_t^{\mathrm{train,exp}}
\right).
\end{equation}


\subsubsection{Rolling Window}

For a fixed window length $L$,

\begin{equation}
\boxed{
\mathcal{D}_t^{\mathrm{train,roll}}
=
\left\{
(\mathbf{x}_{i,\tau},y_{i,\tau+h})
:
t-L
\leq
\tau
\leq
t-h
\right\}.
}
\end{equation}

Older observations are discarded as new ones enter.


\subsubsection{Trade-Off}

Expanding windows use more observations and may reduce estimation variance.

Rolling windows may adapt more quickly to structural change.

Thus,

\begin{equation}
\boxed{
\begin{aligned}
\text{Expanding Window}
&\rightarrow
\text{More Data, Greater Historical Memory},
\\
\text{Rolling Window}
&\rightarrow
\text{Less Data, Greater Adaptability}.
\end{aligned}
}
\end{equation}


% ------------------------------------------------------------
\subsection{Walk-Forward Validation}
\label{subsec:walk_forward_validation}
% ------------------------------------------------------------

Walk-forward validation reproduces the chronological process through which a
model would have been estimated and deployed historically.

Let

\begin{equation}
t_1,
t_2,
\ldots,
t_M
\end{equation}

denote successive prediction dates.

For each date $t_m$:

\begin{enumerate}
    \item construct the admissible training sample;
    \item estimate preprocessing parameters;
    \item fit the model;
    \item optionally select hyperparameters using historical validation data;
    \item generate predictions for $t_m$;
    \item move forward to $t_{m+1}$.
\end{enumerate}

The resulting predictions are

\begin{equation}
\boxed{
\left\{
\widehat{\mathbf{y}}_{t_m+h|t_m}
\right\}_{m=1}^{M}.
}
\end{equation}

Each prediction is therefore genuinely out of sample relative to its own
estimation window.


\subsubsection{Pseudo-Out-of-Sample Predictions}

The sequence

\begin{equation}
\widehat{\mathbf{y}}_{t+h|t}
\end{equation}

is often referred to as a pseudo-out-of-sample prediction because the model is
evaluated historically under the same information restrictions that would have
applied in real time.


% ------------------------------------------------------------
\subsection{Embargo and Leakage Prevention}
\label{subsec:embargo_leakage}
% ------------------------------------------------------------

When targets overlap in time, simply ending the training sample immediately
before the prediction date may still create leakage.

Suppose the target is

\begin{equation}
y_{i,\tau+h}
=
R_{i,\tau\rightarrow\tau+h}.
\end{equation}

An observation whose feature date is $\tau<t$ may nevertheless have a target
period extending beyond the decision date $t$.

To prevent this, the training sample should contain only observations whose
targets are fully realized and admissible before the model is estimated.


\subsubsection{Embargo}

Let $E$ denote an embargo interval.

The last admissible training observation may be restricted to

\begin{equation}
\boxed{
\tau+h
\leq
t-E.
}
\end{equation}

Equivalently,

\begin{equation}
\tau
\leq
t-E-h.
\end{equation}

The embargo introduces a temporal buffer between training outcomes and the
prediction date.


\subsubsection{Purging Overlapping Labels}

If training and validation observations contain overlapping target intervals,
observations whose labels overlap the validation period may be removed.

This is particularly relevant when

\begin{equation}
h
\end{equation}

is long relative to the rebalancing frequency.


\subsubsection{Preprocessing Leakage}

All data-dependent preprocessing must be estimated inside the training window.

This includes:

\begin{itemize}
    \item feature means;
    \item feature standard deviations;
    \item imputation parameters;
    \item dimensionality reduction;
    \item feature selection;
    \item learned embeddings.
\end{itemize}

Thus,

\begin{equation}
\boxed{
\text{Fit preprocessing on train}
\rightarrow
\text{Apply unchanged to validation/test}.
}
\end{equation}


% ------------------------------------------------------------
\subsection{Hyperparameter Selection}
\label{subsec:hyperparameter_selection}
% ------------------------------------------------------------

Many models depend on hyperparameters that are not estimated directly through
the primary training objective.

Let

\begin{equation}
\lambda
\in
\Lambda
\end{equation}

denote a hyperparameter configuration.

Examples include:

\begin{itemize}
    \item Ridge penalty;
    \item Lasso penalty;
    \item tree depth;
    \item number of trees;
    \item learning rate;
    \item neural-network architecture;
    \item dropout rate;
    \item regularization strength.
\end{itemize}

Hyperparameter selection can be expressed as

\begin{equation}
\boxed{
\lambda_t^*
=
\underset{\lambda\in\Lambda}{\arg\min}
\;
\mathcal{L}_{\mathrm{val},t}(\lambda).
}
\end{equation}

Alternatively, if a higher evaluation metric is better,

\begin{equation}
\lambda_t^*
=
\underset{\lambda\in\Lambda}{\arg\max}
\;
M_{\mathrm{val},t}(\lambda).
\end{equation}


\subsubsection{Statistical vs.\ Economic Selection Metrics}

Possible validation criteria include:

\begin{itemize}
    \item RMSE;
    \item MAE;
    \item classification accuracy;
    \item cross-entropy;
    \item Pearson IC;
    \item rank IC;
    \item validation long--short spread;
    \item validation net Sharpe ratio.
\end{itemize}

If an economic performance metric is used, the complete validation portfolio
construction and cost model must itself remain strictly inside the validation
procedure.


\subsubsection{Nested Selection}

When hyperparameter tuning is intensive, an inner validation loop may be used
inside each outer walk-forward period.

Conceptually,

\begin{equation}
\boxed{
\text{Outer Training Window}
\supset
\left[
\text{Inner Train}
+
\text{Inner Validation}
\right].
}
\end{equation}

The outer prediction date remains untouched until the final selected model is
applied.


% ------------------------------------------------------------
\subsection{Model Refitting}
\label{subsec:model_refitting}
% ------------------------------------------------------------

The model does not necessarily need to be re-estimated at every portfolio
rebalance.

Let

\begin{equation}
\mathcal{T}^{\mathrm{refit}}
=
\left\{
t_1^{R},
t_2^{R},
\ldots
\right\}
\end{equation}

denote model-refitting dates.

Between two refitting dates,

\begin{equation}
t_m^{R}
\leq
t
<
t_{m+1}^{R},
\end{equation}

the parameters may remain fixed:

\begin{equation}
\boxed{
\widehat{\theta}_t
=
\widehat{\theta}_{t_m^{R}}.
}
\end{equation}


\subsubsection{Refit Frequency vs.\ Prediction Frequency}

The following frequencies are distinct:

\begin{equation}
\boxed{
f_{\mathrm{refit}}
\neq
f_{\mathrm{prediction}}
\neq
f_{\mathrm{rebalance}}.
}
\end{equation}

For example, a model may:

\begin{itemize}
    \item refit monthly;
    \item generate predictions daily;
    \item rebalance weekly.
\end{itemize}


\subsubsection{Full Refitting vs.\ Parameter Updates}

A refit may involve:

\begin{itemize}
    \item re-estimating model parameters only;
    \item re-estimating preprocessing parameters;
    \item reselecting hyperparameters;
    \item reselecting features;
    \item retraining the complete model architecture.
\end{itemize}

These choices should be distinguished because they imply different levels of
adaptation and computational cost.


% ------------------------------------------------------------
\subsection{Prediction-to-Score Transformation}
\label{subsec:prediction_to_score}
% ------------------------------------------------------------

The model prediction

\begin{equation}
\widehat{y}_{i,t+h|t}
\end{equation}

is not necessarily used directly as the investment score.

The prediction vector is

\begin{equation}
\widehat{\mathbf{y}}_{t+h|t}
=
\begin{bmatrix}
\widehat{y}_{1,t+h|t}
\\
\vdots
\\
\widehat{y}_{N_t,t+h|t}
\end{bmatrix}.
\end{equation}

A scoring transformation produces

\begin{equation}
\boxed{
s_{i,t}^{\mathrm{raw}}
=
\mathcal{P}_t
\left(
\widehat{y}_{i,t+h|t}
\right).
}
\end{equation}


\subsubsection{Direct Prediction Score}

The simplest choice is

\begin{equation}
\boxed{
s_{i,t}^{\mathrm{raw}}
=
\widehat{y}_{i,t+h|t}.
}
\end{equation}

This preserves the model's predicted magnitude.


\subsubsection{Cross-Sectional Standardization}

Predictions may instead be standardized:

\begin{equation}
s_{i,t}^{\mathrm{raw}}
=
\frac{
\widehat{y}_{i,t+h|t}
-
\widehat{\mu}_t
}{
\widehat{\sigma}_t
},
\end{equation}

where

\begin{equation}
\widehat{\mu}_t
=
\frac{1}{N_t}
\sum_i
\widehat{y}_{i,t+h|t}.
\end{equation}


\subsubsection{Rank-Based Prediction Score}

Alternatively,

\begin{equation}
r_{i,t}^{\mathrm{pred}}
=
\operatorname{Rank}_t
\left(
\widehat{y}_{i,t+h|t}
\right),
\end{equation}

followed by

\begin{equation}
\boxed{
s_{i,t}^{\mathrm{raw}}
=
2
\frac{
r_{i,t}^{\mathrm{pred}}-1
}{
N_t-1
}
-1.
}
\end{equation}

This discards predicted magnitude and retains only ordering.


\subsubsection{Classification Probabilities}

For a classifier predicting

\begin{equation}
p_{i,t}
=
P
\left(
y_{i,t+h}=1
\mid
\mathbf{x}_{i,t}
\right),
\end{equation}

a centered score may be

\begin{equation}
\boxed{
s_{i,t}^{\mathrm{raw}}
=
2p_{i,t}-1.
}
\end{equation}

Then

\begin{equation}
s_{i,t}^{\mathrm{raw}}
\in
[-1,1].
\end{equation}


\subsubsection{Prediction Calibration}

The numerical magnitude of model predictions should not automatically be
interpreted as economically calibrated expected returns.

For example,

\begin{equation}
\widehat{y}_{i,t+h|t}=0.02
\end{equation}

does not necessarily imply that the true expected return is exactly $2\%$.

Models optimized for ranking may produce useful cross-sectional ordering even
when prediction magnitudes are poorly calibrated.

Thus,

\begin{equation}
\boxed{
\text{Predictive Ranking Quality}
\neq
\text{Prediction Calibration}.
}
\end{equation}


% ------------------------------------------------------------
\subsection{Final Model-Based Raw Score}
\label{subsec:final_model_based_raw_score}
% ------------------------------------------------------------

The final output of the model-based signal layer is

\begin{equation}
\boxed{
\mathbf{s}_t^{\mathrm{raw}}
=
\begin{bmatrix}
s_{1,t}^{\mathrm{raw}}
\\
\vdots
\\
s_{N_t,t}^{\mathrm{raw}}
\end{bmatrix}
\in
\mathbb{R}^{N_t}.
}
\end{equation}

The complete model-based pipeline can therefore be represented as

\begin{equation}
\boxed{
\begin{aligned}
\mathcal{D}_t^{\mathrm{PIT}}
&\rightarrow
\mathbf{X}_t
\rightarrow
\widehat{\theta}_t
\\[0.1cm]
&\rightarrow
\widehat{\mathbf{y}}_{t+h|t}
\rightarrow
\mathcal{P}_t
\left(
\widehat{\mathbf{y}}_{t+h|t}
\right)
\rightarrow
\mathbf{s}_t^{\mathrm{raw}}.
\end{aligned}
}
\end{equation}

The estimation operator itself depends only on historical admissible data:

\begin{equation}
\widehat{\theta}_t
=
\mathcal{A}
\left(
\mathcal{D}_t^{\mathrm{train}}
\right),
\end{equation}

so that

\begin{equation}
\boxed{
\widehat{\theta}_t
\in
\mathcal{I}_t.
}
\end{equation}


\subsubsection{Unified Signal Interface}

The characteristic-based route produces

\begin{equation}
\mathbf{s}_t^{\mathrm{raw}}
=
\mathcal{S}
\left(
\mathbf{X}_t
\right),
\end{equation}

while the model-based route produces

\begin{equation}
\mathbf{s}_t^{\mathrm{raw}}
=
\mathcal{P}
\left[
f
\left(
\mathbf{X}_t;
\widehat{\theta}_t
\right)
\right].
\end{equation}

In both cases, the downstream engine receives exactly the same mathematical
object:

\begin{equation}
\boxed{
\mathbf{s}_t^{\mathrm{raw}}
\in
\mathbb{R}^{N_t}.
}
\end{equation}

The subsequent processing is therefore model agnostic:

\begin{equation}
\boxed{
\begin{aligned}
\mathbf{s}_t^{\mathrm{raw}}
&\rightarrow
\text{Signal Neutralization}
\rightarrow
\mathbf{s}_t^{\mathrm{neutral}}
\\[0.1cm]
&\rightarrow
\text{Final Signal Construction}
\rightarrow
\mathbf{s}_t^{\mathrm{final}}
\\[0.1cm]
&\rightarrow
\text{Portfolio Construction}
\rightarrow
\mathbf{w}_t.
\end{aligned}
}
\end{equation}


\subsubsection{Backtest Invariant}

At every historical prediction date $t$, the model-based engine should satisfy

\begin{equation}
\boxed{
\mathbf{s}_t^{\mathrm{raw}}
=
g
\left(
\mathcal{I}_t
\right).
}
\end{equation}

No characteristic transformation, model parameter, hyperparameter,
preprocessing statistic, or model-selection decision used to generate the raw
score may depend on information outside the historical information set
$\mathcal{I}_t$.

This condition is the fundamental invariant of model-based signal generation.
```

```latex
% ============================================================
\part{Signal Neutralization and Final Signal}
% ============================================================

% ============================================================
\section{Signal Neutralization}
\label{sec:signal_neutralization}
% ============================================================

The raw investment score

\begin{equation}
\mathbf{s}_t^{\mathrm{raw}}
=
\begin{bmatrix}
s_{1,t}^{\mathrm{raw}}
\\
\vdots
\\
s_{N_t,t}^{\mathrm{raw}}
\end{bmatrix}
\in
\mathbb{R}^{N_t}
\end{equation}

may contain unintended exposures to systematic characteristics such as:

\begin{itemize}
    \item market beta;
    \item industry or sector membership;
    \item country;
    \item size;
    \item volatility;
    \item momentum;
    \item other style factors.
\end{itemize}

For example, a valuation score may unintentionally assign systematically higher
scores to small-cap stocks, or an analyst-revision signal may be concentrated
in particular industries.

Signal neutralization attempts to remove these systematic components before the
signal is mapped into portfolio weights.

The generic objective is

\begin{equation}
\boxed{
\mathbf{s}_t^{\mathrm{raw}}
\rightarrow
\mathbf{s}_t^{\mathrm{neutral}},
}
\end{equation}

where the neutralized score has zero exposure, according to the chosen
neutralization metric, to a predefined set of variables.

The distinction between signal neutralization and portfolio neutralization is
fundamental:

\begin{equation}
\boxed{
\text{Signal Neutralization}
\neq
\text{Portfolio Neutralization}.
}
\end{equation}

Signal neutralization modifies the alpha signal before portfolio construction.
Portfolio neutralization instead imposes constraints directly on final
portfolio weights.

Both may be used in the same strategy, but they operate at different stages of
the investment process.


% ------------------------------------------------------------
\subsection{Motivation for Neutralization}
\label{subsec:motivation_neutralization}
% ------------------------------------------------------------

Suppose the raw signal is predictive because

\begin{equation}
\mathbb{E}
\left[
R_{i,t+1}
\mid
s_{i,t}^{\mathrm{raw}}
\right]
\end{equation}

varies across securities.

However, assume that the signal is correlated with a systematic factor
exposure

\begin{equation}
b_{i,t}.
\end{equation}

Then part of the apparent predictive power of the signal may simply reflect
exposure to that factor.

For example,

\begin{equation}
\operatorname{Corr}_t
\left(
s_{i,t}^{\mathrm{raw}},
\beta_{i,t}
\right)
\neq
0.
\end{equation}

A strategy that buys high-score stocks and shorts low-score stocks may therefore
generate a portfolio with systematic market-beta exposure even if the original
objective was to capture stock-specific alpha.


\subsubsection{Signal Decomposition}

A simple decomposition is

\begin{equation}
\boxed{
\mathbf{s}_t^{\mathrm{raw}}
=
\mathbf{B}_t
\boldsymbol{\gamma}_t
+
\boldsymbol{\varepsilon}_t,
}
\end{equation}

where:

\begin{itemize}
    \item $\mathbf{B}_t$ contains systematic exposures;
    \item $\boldsymbol{\gamma}_t$ measures the component of the raw score
    associated with those exposures;
    \item $\boldsymbol{\varepsilon}_t$ is the residual component.
\end{itemize}

The neutralized score is then defined as

\begin{equation}
\boxed{
\mathbf{s}_t^{\mathrm{neutral}}
=
\boldsymbol{\varepsilon}_t.
}
\end{equation}

The purpose of neutralization is therefore to preserve the component of the
signal that cannot be explained by the selected systematic exposures.


\subsubsection{Economic Interpretation}

Neutralization answers the question:

\begin{equation}
\boxed{
\begin{aligned}
&\text{How attractive is security }i
\\
&\text{after controlling for its systematic characteristics?}
\end{aligned}
}
\end{equation}

For example, industry-neutral value does not compare a bank directly with a
technology company according to raw valuation levels. Instead, it attempts to
identify securities that are attractive relative to what would be expected
given their industry characteristics.


% ------------------------------------------------------------
\subsection{General Exposure Matrix}
\label{subsec:general_exposure_matrix}
% ------------------------------------------------------------

Let

\begin{equation}
K_t
\end{equation}

denote the number of exposures that should be removed at time $t$.

For security $i$, define

\begin{equation}
\mathbf{b}_{i,t}
=
\begin{bmatrix}
b_{i,t}^{(1)}
\\
\vdots
\\
b_{i,t}^{(K_t)}
\end{bmatrix}
\in
\mathbb{R}^{K_t}.
\end{equation}

Stacking the exposures across securities gives the exposure matrix

\begin{equation}
\boxed{
\mathbf{B}_t
=
\begin{bmatrix}
(\mathbf{b}_{1,t})^\top
\\
\vdots
\\
(\mathbf{b}_{N_t,t})^\top
\end{bmatrix}
\in
\mathbb{R}^{N_t\times K_t}.
}
\end{equation}

Possible columns of $\mathbf{B}_t$ include:

\begin{equation}
\mathbf{B}_t
=
\begin{bmatrix}
\mathbf{1}
&
\boldsymbol{\beta}_t
&
\mathbf{D}_t^{\mathrm{industry}}
&
\log(\mathbf{MCap}_t)
&
\boldsymbol{\sigma}_t
&
\mathbf{F}_t^{\mathrm{style}}
\end{bmatrix}.
\end{equation}

The exact matrix depends on the desired neutralization.


\subsubsection{Intercept}

An intercept column

\begin{equation}
\mathbf{1}
=
\begin{bmatrix}
1\\
\vdots\\
1
\end{bmatrix}
\end{equation}

is frequently included.

Then the residualized signal satisfies

\begin{equation}
\mathbf{1}^{\top}
\mathbf{s}_t^{\mathrm{neutral}}
=
0
\end{equation}

under ordinary least squares.

Thus, the cross-sectional mean of the residual signal is zero.


\subsubsection{Continuous and Categorical Exposures}

Continuous exposures may include

\begin{equation}
\beta_{i,t},
\qquad
\log(\mathrm{MCap}_{i,t}),
\qquad
\sigma_{i,t}.
\end{equation}

Categorical variables such as industry membership are represented through dummy
variables.

If there are $G$ industries,

\begin{equation}
D_{i,g,t}
=
\begin{cases}
1,
&
\text{if security }i\text{ belongs to industry }g,
\\
0,
&
\text{otherwise}.
\end{cases}
\end{equation}

The corresponding matrix is

\begin{equation}
\mathbf{D}_t
\in
\mathbb{R}^{N_t\times G}.
\end{equation}


\subsubsection{Dummy Variable Collinearity}

If an intercept is included together with all mutually exclusive industry
dummy variables,

\begin{equation}
\mathbf{1}
=
\sum_{g=1}^{G}
\mathbf{D}_{g,t}.
\end{equation}

The design matrix is therefore perfectly collinear.

One category must be omitted, or another identification convention must be
used.

For example,

\begin{equation}
\mathbf{B}_t
=
\begin{bmatrix}
\mathbf{1}
&
\mathbf{D}_{1,t}
&
\cdots
&
\mathbf{D}_{G-1,t}
\end{bmatrix}.
\end{equation}


% ------------------------------------------------------------
\subsection{Cross-Sectional Regression Framework}
\label{subsec:cross_sectional_neutralization}
% ------------------------------------------------------------

The standard neutralization method is a cross-sectional regression performed
independently at each decision date $t$.

The model is

\begin{equation}
\boxed{
\mathbf{s}_t^{\mathrm{raw}}
=
\mathbf{B}_t
\boldsymbol{\gamma}_t
+
\boldsymbol{\varepsilon}_t.
}
\end{equation}

Under Ordinary Least Squares,

\begin{equation}
\widehat{\boldsymbol{\gamma}}_t
=
\left(
\mathbf{B}_t^\top
\mathbf{B}_t
\right)^{-1}
\mathbf{B}_t^\top
\mathbf{s}_t^{\mathrm{raw}},
\end{equation}

assuming the matrix is invertible.

The fitted systematic component is

\begin{equation}
\widehat{\mathbf{s}}_t^{\mathrm{systematic}}
=
\mathbf{B}_t
\widehat{\boldsymbol{\gamma}}_t.
\end{equation}

The neutralized score is

\begin{equation}
\boxed{
\mathbf{s}_t^{\mathrm{neutral}}
=
\mathbf{s}_t^{\mathrm{raw}}
-
\mathbf{B}_t
\widehat{\boldsymbol{\gamma}}_t.
}
\end{equation}


\subsubsection{Projection-Matrix Representation}

Define the projection matrix onto the column space of $\mathbf{B}_t$:

\begin{equation}
\boxed{
\mathbf{P}_{B,t}
=
\mathbf{B}_t
\left(
\mathbf{B}_t^\top
\mathbf{B}_t
\right)^{-1}
\mathbf{B}_t^\top.
}
\end{equation}

Then

\begin{equation}
\widehat{\mathbf{s}}_t^{\mathrm{systematic}}
=
\mathbf{P}_{B,t}
\mathbf{s}_t^{\mathrm{raw}}.
\end{equation}

Define the residual-maker matrix

\begin{equation}
\boxed{
\mathbf{M}_{B,t}
=
\mathbf{I}
-
\mathbf{P}_{B,t}.
}
\end{equation}

The neutralized score is therefore

\begin{equation}
\boxed{
\mathbf{s}_t^{\mathrm{neutral}}
=
\mathbf{M}_{B,t}
\mathbf{s}_t^{\mathrm{raw}}.
}
\end{equation}


\subsubsection{Orthogonality Property}

OLS residuals satisfy

\begin{equation}
\boxed{
\mathbf{B}_t^\top
\mathbf{s}_t^{\mathrm{neutral}}
=
\mathbf{0}.
}
\end{equation}

Thus, the neutralized score is orthogonal to every column of the exposure
matrix under the Euclidean inner product.

This is the mathematical basis of regression-based neutralization.


\subsubsection{Geometric Interpretation}

The raw signal can be decomposed as

\begin{equation}
\boxed{
\mathbf{s}_t^{\mathrm{raw}}
=
\underbrace{
\mathbf{P}_{B,t}
\mathbf{s}_t^{\mathrm{raw}}
}_{\text{component explained by exposures}}
+
\underbrace{
\mathbf{M}_{B,t}
\mathbf{s}_t^{\mathrm{raw}}
}_{\text{orthogonal component}}.
}
\end{equation}

Neutralization removes the projection of the signal onto the subspace spanned
by the unwanted exposures.


% ------------------------------------------------------------
\subsection{Beta Neutralization}
\label{subsec:signal_beta_neutralization}
% ------------------------------------------------------------

Suppose

\begin{equation}
\boldsymbol{\beta}_t
=
\begin{bmatrix}
\beta_{1,t}
\\
\vdots
\\
\beta_{N_t,t}
\end{bmatrix}
\end{equation}

denotes the vector of stock-level market betas.

A beta-neutralized signal can be obtained from

\begin{equation}
\boxed{
s_{i,t}^{\mathrm{raw}}
=
\alpha_t
+
\gamma_{\beta,t}
\beta_{i,t}
+
\varepsilon_{i,t}.
}
\end{equation}

In matrix form,

\begin{equation}
\mathbf{B}_t
=
\begin{bmatrix}
\mathbf{1}
&
\boldsymbol{\beta}_t
\end{bmatrix}.
\end{equation}

The neutralized score is

\begin{equation}
\boxed{
\mathbf{s}_t^{\beta\text{-neutral}}
=
\widehat{\boldsymbol{\varepsilon}}_t.
}
\end{equation}

By construction,

\begin{equation}
\boldsymbol{\beta}_t^\top
\mathbf{s}_t^{\beta\text{-neutral}}
=
0
\end{equation}

under OLS.


\subsubsection{Interpretation}

Suppose high-score securities systematically have high market beta.

Then

\begin{equation}
\gamma_{\beta,t}>0.
\end{equation}

Beta neutralization removes the component of the raw score that can be
explained linearly by beta.

The residual score therefore represents relative attractiveness conditional on
market beta.


\subsubsection{Signal Beta Neutrality vs.\ Portfolio Beta Neutrality}

The condition

\begin{equation}
\boldsymbol{\beta}_t^\top
\mathbf{s}_t^{\mathrm{neutral}}
=
0
\end{equation}

does not automatically imply

\begin{equation}
\boldsymbol{\beta}_t^\top
\mathbf{w}_t
=
0.
\end{equation}

This is because the subsequent signal-to-weight mapping may be nonlinear or may
introduce constraints and scaling.

Therefore,

\begin{equation}
\boxed{
\text{Beta-Neutral Signal}
\not\Rightarrow
\text{Beta-Neutral Portfolio}.
}
\end{equation}

Portfolio beta neutrality must be checked or imposed separately during
portfolio construction.


% ------------------------------------------------------------
\subsection{Industry and Sector Neutralization}
\label{subsec:industry_neutralization}
% ------------------------------------------------------------

Industry neutralization removes systematic differences in signal levels across
industries or sectors.

Suppose there are $G$ industries.

The raw score may be modeled as

\begin{equation}
s_{i,t}^{\mathrm{raw}}
=
\alpha_t
+
\sum_{g=1}^{G-1}
\gamma_{g,t}
D_{i,g,t}
+
\varepsilon_{i,t}.
\end{equation}

The residual

\begin{equation}
\boxed{
s_{i,t}^{\mathrm{industry\text{-}neutral}}
=
\widehat{\varepsilon}_{i,t}
}
\end{equation}

measures the security's score relative to the systematic level associated with
its industry.


\subsubsection{Group-Demeaning Interpretation}

If the model contains only industry dummies and no additional continuous
variables, industry neutralization is closely related to subtracting the
industry mean.

For industry $g$,

\begin{equation}
\bar{s}_{g,t}
=
\frac{
1
}{
N_{g,t}
}
\sum_{i:g(i,t)=g}
s_{i,t}^{\mathrm{raw}}.
\end{equation}

Then a simple industry-demeaned score is

\begin{equation}
\boxed{
s_{i,t}^{\mathrm{neutral}}
=
s_{i,t}^{\mathrm{raw}}
-
\bar{s}_{g(i,t),t}.
}
\end{equation}

Thus, a security is evaluated relative to peers in the same industry rather
than relative to the entire market.


\subsubsection{Sector vs.\ Industry Granularity}

Neutralization may be performed at different hierarchical levels:

\begin{equation}
\text{Sector}
\rightarrow
\text{Industry Group}
\rightarrow
\text{Industry}
\rightarrow
\text{Sub-Industry}.
\end{equation}

A finer classification removes more granular systematic variation but leaves
fewer securities within each group.

Thus,

\begin{equation}
\boxed{
\text{Finer Neutralization}
\rightarrow
\text{Less Cross-Sectional Information Per Group}.
}
\end{equation}


% ------------------------------------------------------------
\subsection{Country Neutralization}
\label{subsec:country_neutralization}
% ------------------------------------------------------------

In international strategies, raw scores may differ systematically across
countries because of:

\begin{itemize}
    \item accounting conventions;
    \item sector composition;
    \item macroeconomic regimes;
    \item valuation levels;
    \item local market structure.
\end{itemize}

Let

\begin{equation}
C_{i,c,t}
\end{equation}

denote country dummy $c$.

The regression becomes

\begin{equation}
s_{i,t}^{\mathrm{raw}}
=
\alpha_t
+
\sum_{c=1}^{C-1}
\gamma_{c,t}^{C}
C_{i,c,t}
+
\varepsilon_{i,t}.
\end{equation}

The country-neutral score is

\begin{equation}
\boxed{
s_{i,t}^{\mathrm{country\text{-}neutral}}
=
\widehat{\varepsilon}_{i,t}.
}
\end{equation}

This prevents systematic differences in country-level signal distributions from
dominating the stock-selection process.


% ------------------------------------------------------------
\subsection{Size Neutralization}
\label{subsec:size_neutralization}
% ------------------------------------------------------------

Many firm characteristics are correlated with company size.

A standard size variable is

\begin{equation}
\boxed{
z_{i,t}^{\mathrm{size}}
=
\log
\left(
\mathrm{MCap}_{i,t}
\right).
}
\end{equation}

The logarithm is typically preferred because market capitalization is highly
right-skewed.

A size-neutralization regression is

\begin{equation}
s_{i,t}^{\mathrm{raw}}
=
\alpha_t
+
\gamma_{\mathrm{size},t}
\log
\left(
\mathrm{MCap}_{i,t}
\right)
+
\varepsilon_{i,t}.
\end{equation}

The residual is

\begin{equation}
\boxed{
s_{i,t}^{\mathrm{size\text{-}neutral}}
=
\widehat{\varepsilon}_{i,t}.
}
\end{equation}


\subsubsection{Nonlinear Size Effects}

The relationship between a raw signal and size need not be linear.

A more flexible specification may include

\begin{equation}
\left[
\log(\mathrm{MCap}_{i,t})
\right]^2,
\end{equation}

or other basis functions.

For example,

\begin{equation}
s_{i,t}^{\mathrm{raw}}
=
\alpha_t
+
\gamma_{1,t}
\log(\mathrm{MCap}_{i,t})
+
\gamma_{2,t}
\left[
\log(\mathrm{MCap}_{i,t})
\right]^2
+
\varepsilon_{i,t}.
\end{equation}

Neutralization removes only the functional relationship explicitly included in
the regression.


% ------------------------------------------------------------
\subsection{Volatility Neutralization}
\label{subsec:volatility_neutralization}
% ------------------------------------------------------------

Signals may also be systematically related to stock-level risk.

Let

\begin{equation}
\sigma_{i,t}
\end{equation}

denote an ex-ante volatility estimate.

A volatility-neutralization regression is

\begin{equation}
\boxed{
s_{i,t}^{\mathrm{raw}}
=
\alpha_t
+
\gamma_{\sigma,t}
\sigma_{i,t}
+
\varepsilon_{i,t}.
}
\end{equation}

The residual signal is orthogonal to the included volatility measure.

This can be useful when the strategy intends to separate alpha information from
a low-volatility or high-volatility tilt.


\subsubsection{Choice of Volatility Measure}

Possible measures include:

\begin{itemize}
    \item realized historical volatility;
    \item exponentially weighted volatility;
    \item idiosyncratic volatility;
    \item risk-model forecast volatility.
\end{itemize}

The choice determines the exposure being removed.

For example,

\begin{equation}
\sigma_{i,t}^{\mathrm{total}}
\end{equation}

and

\begin{equation}
\sigma_{i,t}^{\mathrm{idiosyncratic}}
\end{equation}

represent economically different risk concepts.


% ------------------------------------------------------------
\subsection{Momentum and Style Factor Neutralization}
\label{subsec:style_neutralization}
% ------------------------------------------------------------

The same framework extends naturally to additional style characteristics.

Suppose

\begin{equation}
\mathbf{f}_{i,t}^{\mathrm{style}}
=
\begin{bmatrix}
f_{i,t}^{\mathrm{momentum}}
\\
f_{i,t}^{\mathrm{value}}
\\
f_{i,t}^{\mathrm{quality}}
\\
f_{i,t}^{\mathrm{lowvol}}
\\
\vdots
\end{bmatrix}.
\end{equation}

Then

\begin{equation}
s_{i,t}^{\mathrm{raw}}
=
\alpha_t
+
\left(
\mathbf{f}_{i,t}^{\mathrm{style}}
\right)^\top
\boldsymbol{\gamma}_t
+
\varepsilon_{i,t}.
\end{equation}

The residual

\begin{equation}
\boxed{
s_{i,t}^{\mathrm{style\text{-}neutral}}
=
\widehat{\varepsilon}_{i,t}
}
\end{equation}

represents the component of the signal orthogonal to the included style
exposures.


\subsubsection{Avoiding Mechanical Removal of the Alpha}

Neutralization should be applied only to exposures that are genuinely unwanted.

If the alpha hypothesis itself is a momentum strategy, neutralizing the signal
against momentum may remove the intended source of return.

Thus,

\begin{equation}
\boxed{
\text{Neutralization Variable}
\neq
\text{Automatically Desirable Control}.
}
\end{equation}

Every exposure included in the neutralization matrix should therefore have a
clear economic justification.

```latex
% ------------------------------------------------------------
\subsection{Joint vs.\ Sequential Neutralization}
\label{subsec:joint_sequential_neutralization}
% ------------------------------------------------------------

When several groups of exposures must be removed from a raw signal, two broad
approaches may be considered:

\begin{equation}
\boxed{
\text{Joint Neutralization}
\qquad\text{or}\qquad
\text{Sequential Neutralization}.
}
\end{equation}

The distinction must be handled carefully. A naive sequential residualization
is generally not equivalent to a joint regression. However, a properly
constructed sequential procedure based on the
\emph{Frisch--Waugh--Lovell theorem} is exactly equivalent to joint
neutralization.

Let

\begin{equation}
\mathbf{s}_t^{\mathrm{raw}}
\in
\mathbb{R}^{N_t}
\end{equation}

denote the raw cross-sectional signal at time $t$, where $N_t$ is the number of
securities in the investment universe.

Suppose the exposures to be removed are divided into two blocks:

\begin{equation}
\mathbf{B}_{1,t}
\in
\mathbb{R}^{N_t\times K_1},
\end{equation}

and

\begin{equation}
\mathbf{B}_{2,t}
\in
\mathbb{R}^{N_t\times K_2},
\end{equation}

where $K_1$ and $K_2$ denote the number of exposures in the first and second
blocks respectively.

Examples could be:

\begin{equation}
\mathbf{B}_{1,t}
=
\begin{bmatrix}
\mathbf{1}
&
\boldsymbol{\beta}_t
\end{bmatrix},
\end{equation}

and

\begin{equation}
\mathbf{B}_{2,t}
=
\begin{bmatrix}
\mathbf{D}_t^{\mathrm{industry}}
&
\log(\mathbf{MCap}_t)
\end{bmatrix}.
\end{equation}


% ------------------------------------------------------------
\subsubsection{Joint Neutralization}
% ------------------------------------------------------------

Define the complete exposure matrix

\begin{equation}
\boxed{
\mathbf{B}_t
=
\begin{bmatrix}
\mathbf{B}_{1,t}
&
\mathbf{B}_{2,t}
\end{bmatrix}
\in
\mathbb{R}^{N_t\times(K_1+K_2)}.
}
\end{equation}

Joint neutralization estimates the cross-sectional regression

\begin{equation}
\boxed{
\mathbf{s}_t^{\mathrm{raw}}
=
\mathbf{B}_{1,t}
\boldsymbol{\gamma}_{1,t}
+
\mathbf{B}_{2,t}
\boldsymbol{\gamma}_{2,t}
+
\boldsymbol{\varepsilon}_t,
}
\end{equation}

where

\begin{equation}
\boldsymbol{\gamma}_{1,t}
\in
\mathbb{R}^{K_1},
\qquad
\boldsymbol{\gamma}_{2,t}
\in
\mathbb{R}^{K_2},
\end{equation}

and

\begin{equation}
\boldsymbol{\varepsilon}_t
\in
\mathbb{R}^{N_t}.
\end{equation}

Equivalently,

\begin{equation}
\mathbf{s}_t^{\mathrm{raw}}
=
\mathbf{B}_t
\boldsymbol{\gamma}_t
+
\boldsymbol{\varepsilon}_t,
\end{equation}

with

\begin{equation}
\boldsymbol{\gamma}_t
=
\begin{bmatrix}
\boldsymbol{\gamma}_{1,t}
\\
\boldsymbol{\gamma}_{2,t}
\end{bmatrix}
\in
\mathbb{R}^{K_1+K_2}.
\end{equation}

Under Ordinary Least Squares,

\begin{equation}
\widehat{\boldsymbol{\gamma}}_t
=
\left(
\mathbf{B}_t^\top
\mathbf{B}_t
\right)^{-1}
\mathbf{B}_t^\top
\mathbf{s}_t^{\mathrm{raw}},
\end{equation}

assuming full column rank.

The joint-neutralized signal is

\begin{equation}
\boxed{
\mathbf{s}_t^{\mathrm{joint}}
=
\mathbf{s}_t^{\mathrm{raw}}
-
\mathbf{B}_t
\widehat{\boldsymbol{\gamma}}_t.
}
\end{equation}

Define the projection matrix onto the column space of $\mathbf{B}_t$ as

\begin{equation}
\mathbf{P}_{B,t}
=
\mathbf{B}_t
\left(
\mathbf{B}_t^\top
\mathbf{B}_t
\right)^{-1}
\mathbf{B}_t^\top
\in
\mathbb{R}^{N_t\times N_t},
\end{equation}

and the corresponding residual-maker matrix as

\begin{equation}
\boxed{
\mathbf{M}_{B,t}
=
\mathbf{I}_{N_t}
-
\mathbf{P}_{B,t}.
}
\end{equation}

Then

\begin{equation}
\boxed{
\mathbf{s}_t^{\mathrm{joint}}
=
\mathbf{M}_{B,t}
\mathbf{s}_t^{\mathrm{raw}}.
}
\end{equation}

The OLS normal equations imply

\begin{equation}
\boxed{
\mathbf{B}_{1,t}^\top
\mathbf{s}_t^{\mathrm{joint}}
=
\mathbf{0}_{K_1},
}
\end{equation}

and

\begin{equation}
\boxed{
\mathbf{B}_{2,t}^\top
\mathbf{s}_t^{\mathrm{joint}}
=
\mathbf{0}_{K_2}.
}
\end{equation}

Hence, the resulting signal is simultaneously orthogonal to every exposure
included in both blocks.


% ------------------------------------------------------------
\subsubsection{Naive Sequential Neutralization}
% ------------------------------------------------------------

A seemingly natural alternative is to neutralize the signal against the first
block and then neutralize the resulting residual against the second block.

Define

\begin{equation}
\mathbf{P}_{1,t}
=
\mathbf{B}_{1,t}
\left(
\mathbf{B}_{1,t}^\top
\mathbf{B}_{1,t}
\right)^{-1}
\mathbf{B}_{1,t}^\top,
\end{equation}

and

\begin{equation}
\boxed{
\mathbf{M}_{1,t}
=
\mathbf{I}_{N_t}
-
\mathbf{P}_{1,t}.
}
\end{equation}

The first-stage residual is

\begin{equation}
\boxed{
\mathbf{s}_t^{(1)}
=
\mathbf{M}_{1,t}
\mathbf{s}_t^{\mathrm{raw}}.
}
\end{equation}

By construction,

\begin{equation}
\mathbf{B}_{1,t}^\top
\mathbf{s}_t^{(1)}
=
\mathbf{0}_{K_1}.
\end{equation}

Now define similarly

\begin{equation}
\mathbf{P}_{2,t}
=
\mathbf{B}_{2,t}
\left(
\mathbf{B}_{2,t}^\top
\mathbf{B}_{2,t}
\right)^{-1}
\mathbf{B}_{2,t}^\top,
\end{equation}

and

\begin{equation}
\mathbf{M}_{2,t}
=
\mathbf{I}_{N_t}
-
\mathbf{P}_{2,t}.
\end{equation}

A naive second step gives

\begin{equation}
\boxed{
\mathbf{s}_t^{\mathrm{naive}}
=
\mathbf{M}_{2,t}
\mathbf{M}_{1,t}
\mathbf{s}_t^{\mathrm{raw}}.
}
\end{equation}

This procedure is generally \emph{not} equivalent to joint neutralization.


% ------------------------------------------------------------
\subsubsection{Why Naive Sequential Neutralization Is Order Dependent}
% ------------------------------------------------------------

In general,

\begin{equation}
\boxed{
\mathbf{M}_{2,t}
\mathbf{M}_{1,t}
\neq
\mathbf{M}_{1,t}
\mathbf{M}_{2,t}.
}
\end{equation}

Therefore, reversing the order of the two neutralizations generally produces a
different final signal.

Moreover, although

\begin{equation}
\mathbf{B}_{1,t}^\top
\mathbf{s}_t^{(1)}
=
0,
\end{equation}

after the second projection one generally has

\begin{equation}
\boxed{
\mathbf{B}_{1,t}^\top
\mathbf{s}_t^{\mathrm{naive}}
\neq
0.
}
\end{equation}

Thus, the second projection can reintroduce exposure to the first block.

The reason is that $\mathbf{B}_{2,t}$ may contain components lying in the
column space of $\mathbf{B}_{1,t}$.


% ------------------------------------------------------------
\subsubsection{Frisch--Waugh--Lovell Theorem}
% ------------------------------------------------------------

The correct sequential procedure is given by the
\emph{Frisch--Waugh--Lovell theorem}.

The theorem states that, in the joint regression

\begin{equation}
\mathbf{s}_t^{\mathrm{raw}}
=
\mathbf{B}_{1,t}
\boldsymbol{\gamma}_{1,t}
+
\mathbf{B}_{2,t}
\boldsymbol{\gamma}_{2,t}
+
\boldsymbol{\varepsilon}_t,
\end{equation}

the coefficient on $\mathbf{B}_{2,t}$ can be obtained by:

\begin{enumerate}
    \item residualizing the dependent variable with respect to
    $\mathbf{B}_{1,t}$;

    \item residualizing every column of $\mathbf{B}_{2,t}$ with respect to
    $\mathbf{B}_{1,t}$;

    \item regressing the residualized dependent variable on the residualized
    second exposure block.
\end{enumerate}

The first residualized object is

\begin{equation}
\boxed{
\widetilde{\mathbf{s}}_t
=
\mathbf{M}_{1,t}
\mathbf{s}_t^{\mathrm{raw}}
\in
\mathbb{R}^{N_t}.
}
\end{equation}

The second exposure block must also be residualized:

\begin{equation}
\boxed{
\widetilde{\mathbf{B}}_{2,t}
=
\mathbf{M}_{1,t}
\mathbf{B}_{2,t}
\in
\mathbb{R}^{N_t\times K_2}.
}
\end{equation}

Thus, the relevant second-stage regression is not

\begin{equation}
\widetilde{\mathbf{s}}_t
\text{ on }
\mathbf{B}_{2,t},
\end{equation}

but rather

\begin{equation}
\boxed{
\widetilde{\mathbf{s}}_t
=
\widetilde{\mathbf{B}}_{2,t}
\boldsymbol{\gamma}_{2,t}
+
\boldsymbol{\varepsilon}_t.
}
\end{equation}

The Frisch--Waugh--Lovell theorem guarantees that

\begin{equation}
\boxed{
\widehat{\boldsymbol{\gamma}}_{2,t}^{FWL}
=
\widehat{\boldsymbol{\gamma}}_{2,t}^{\mathrm{joint}}.
}
\end{equation}


% ------------------------------------------------------------
\subsubsection{FWL-Consistent Sequential Neutralization}
% ------------------------------------------------------------

Define the projection matrix onto the residualized exposure block

\begin{equation}
\widetilde{\mathbf{B}}_{2,t}
=
\mathbf{M}_{1,t}
\mathbf{B}_{2,t}
\end{equation}

as

\begin{equation}
\mathbf{P}_{\widetilde{B}_2,t}
=
\widetilde{\mathbf{B}}_{2,t}
\left(
\widetilde{\mathbf{B}}_{2,t}^{\top}
\widetilde{\mathbf{B}}_{2,t}
\right)^{-1}
\widetilde{\mathbf{B}}_{2,t}^{\top}.
\end{equation}

The associated residual-maker matrix is

\begin{equation}
\boxed{
\mathbf{M}_{\widetilde{B}_2,t}
=
\mathbf{I}_{N_t}
-
\mathbf{P}_{\widetilde{B}_2,t}.
}
\end{equation}

The correctly sequentially neutralized signal is then

\begin{equation}
\boxed{
\mathbf{s}_t^{FWL}
=
\mathbf{M}_{\widetilde{B}_2,t}
\mathbf{M}_{1,t}
\mathbf{s}_t^{\mathrm{raw}}.
}
\end{equation}

The Frisch--Waugh--Lovell theorem implies

\begin{equation}
\boxed{
\mathbf{s}_t^{FWL}
=
\mathbf{s}_t^{\mathrm{joint}}.
}
\end{equation}

Equivalently,

\begin{equation}
\boxed{
\mathbf{M}_{\widetilde{B}_2,t}
\mathbf{M}_{1,t}
\mathbf{s}_t^{\mathrm{raw}}
=
\mathbf{M}_{B,t}
\mathbf{s}_t^{\mathrm{raw}}.
}
\end{equation}

Hence, properly constructed sequential neutralization is exactly equivalent to
joint neutralization.


% ------------------------------------------------------------
\subsubsection{Geometric Interpretation}
% ------------------------------------------------------------

The exposure matrix

\begin{equation}
\mathbf{B}_{2,t}
\end{equation}

may contain two components:

\begin{equation}
\boxed{
\mathbf{B}_{2,t}
=
\mathbf{P}_{1,t}
\mathbf{B}_{2,t}
+
\mathbf{M}_{1,t}
\mathbf{B}_{2,t}.
}
\end{equation}

The first component,

\begin{equation}
\mathbf{P}_{1,t}
\mathbf{B}_{2,t},
\end{equation}

already lies in the space spanned by $\mathbf{B}_{1,t}$.

The second component,

\begin{equation}
\boxed{
\mathbf{M}_{1,t}
\mathbf{B}_{2,t},
}
\end{equation}

is orthogonal to $\mathbf{B}_{1,t}$ and represents the genuinely new directions
added by the second exposure block.

Therefore,

\begin{equation}
\boxed{
\operatorname{Col}
\left(
\begin{bmatrix}
\mathbf{B}_{1,t}
&
\mathbf{B}_{2,t}
\end{bmatrix}
\right)
=
\operatorname{Col}
\left(
\begin{bmatrix}
\mathbf{B}_{1,t}
&
\mathbf{M}_{1,t}\mathbf{B}_{2,t}
\end{bmatrix}
\right),
}
\end{equation}

where $\operatorname{Col}(\cdot)$ denotes the column space.

This is the geometric intuition behind the Frisch--Waugh--Lovell result.


% ------------------------------------------------------------
\subsubsection{Connection with Gram--Schmidt Orthogonalization}
% ------------------------------------------------------------

The same procedure can be interpreted as a block version of Gram--Schmidt
orthogonalization.

The first exposure block is retained as

\begin{equation}
\widetilde{\mathbf{B}}_{1,t}
=
\mathbf{B}_{1,t}.
\end{equation}

The second block is replaced by its component orthogonal to the first:

\begin{equation}
\widetilde{\mathbf{B}}_{2,t}
=
\mathbf{M}_{1,t}
\mathbf{B}_{2,t}.
\end{equation}

The joint exposure space can then be written as

\begin{equation}
\operatorname{Col}(\mathbf{B}_t)
=
\operatorname{Col}
\left(
\widetilde{\mathbf{B}}_{1,t},
\widetilde{\mathbf{B}}_{2,t}
\right).
\end{equation}

Sequentially removing these orthogonalized exposure blocks therefore removes
exactly the same subspace as the original joint regression.


% ------------------------------------------------------------
\subsubsection{Special Case: Orthogonal Exposure Blocks}
% ------------------------------------------------------------

Suppose the two exposure blocks are already orthogonal:

\begin{equation}
\boxed{
\mathbf{B}_{1,t}^{\top}
\mathbf{B}_{2,t}
=
\mathbf{0}_{K_1\times K_2}.
}
\end{equation}

Then

\begin{equation}
\mathbf{P}_{1,t}
\mathbf{B}_{2,t}
=
0,
\end{equation}

and therefore

\begin{equation}
\boxed{
\mathbf{M}_{1,t}
\mathbf{B}_{2,t}
=
\mathbf{B}_{2,t}.
}
\end{equation}

In this special case, the raw second exposure block is already equal to its
FWL-residualized version.

Therefore,

\begin{equation}
\mathbf{M}_{2,t}
\mathbf{M}_{1,t}
=
\mathbf{M}_{1,t}
\mathbf{M}_{2,t},
\end{equation}

and naive sequential neutralization becomes equivalent to joint
neutralization.


% ------------------------------------------------------------
\subsubsection{Extension to More Than Two Exposure Blocks}
% ------------------------------------------------------------

The same logic extends to $J$ exposure blocks

\begin{equation}
\mathbf{B}_{1,t},
\mathbf{B}_{2,t},
\ldots,
\mathbf{B}_{J,t}.
\end{equation}

The second block is orthogonalized with respect to the first:

\begin{equation}
\widetilde{\mathbf{B}}_{2,t}
=
\mathbf{M}_{B_1,t}
\mathbf{B}_{2,t}.
\end{equation}

The third block must be orthogonalized with respect to the entire space already
removed:

\begin{equation}
\widetilde{\mathbf{B}}_{3,t}
=
\mathbf{M}_{[B_1,B_2],t}
\mathbf{B}_{3,t}.
\end{equation}

More generally,

\begin{equation}
\boxed{
\widetilde{\mathbf{B}}_{j,t}
=
\mathbf{M}_{[B_1,\ldots,B_{j-1}],t}
\mathbf{B}_{j,t}.
}
\end{equation}

Sequential projection onto these orthogonalized exposure blocks reproduces the
residual from the joint regression containing all exposures.


% ------------------------------------------------------------
\subsubsection{Practical Implication for Signal Neutralization}
% ------------------------------------------------------------

It is therefore important to distinguish three procedures:

\begin{equation}
\boxed{
\begin{aligned}
\text{Joint Regression}
&:
\quad
[\mathbf{B}_{1,t},\mathbf{B}_{2,t}]
\text{ included simultaneously},
\\[0.15cm]
\text{Naive Sequential}
&:
\quad
\mathbf{M}_{2,t}\mathbf{M}_{1,t}
\mathbf{s}_t^{\mathrm{raw}},
\\[0.15cm]
\text{FWL Sequential}
&:
\quad
\mathbf{M}_{\widetilde{B}_2,t}
\mathbf{M}_{1,t}
\mathbf{s}_t^{\mathrm{raw}}.
\end{aligned}
}
\end{equation}

In general,

\begin{equation}
\boxed{
\text{Naive Sequential Neutralization}
\neq
\text{Joint Neutralization},
}
\end{equation}

whereas

\begin{equation}
\boxed{
\text{FWL-Consistent Sequential Neutralization}
=
\text{Joint Neutralization}.
}
\end{equation}

For implementation purposes, a single joint regression is usually the simplest
and least error-prone approach when all exposures are known simultaneously.

Sequential neutralization is mathematically equivalent only when each new
exposure block is first residualized with respect to all exposure blocks that
have already been removed.

% ------------------------------------------------------------
\subsection{Weighted Neutralization}
\label{subsec:weighted_neutralization}
% ------------------------------------------------------------

Ordinary Least Squares gives every security equal importance in the
cross-sectional regression.

In some applications, a weighted neutralization may be preferable.

Let

\begin{equation}
\mathbf{W}_t
=
\operatorname{diag}
\left(
\omega_{1,t},
\ldots,
\omega_{N_t,t}
\right),
\end{equation}

with

\begin{equation}
\omega_{i,t}>0.
\end{equation}

The Weighted Least Squares estimator is

\begin{equation}
\boxed{
\widehat{\boldsymbol{\gamma}}_t^{WLS}
=
\left(
\mathbf{B}_t^\top
\mathbf{W}_t
\mathbf{B}_t
\right)^{-1}
\mathbf{B}_t^\top
\mathbf{W}_t
\mathbf{s}_t^{\mathrm{raw}}.
}
\end{equation}

The residualized signal is

\begin{equation}
\boxed{
\mathbf{s}_t^{\mathrm{neutral}}
=
\mathbf{s}_t^{\mathrm{raw}}
-
\mathbf{B}_t
\widehat{\boldsymbol{\gamma}}_t^{WLS}.
}
\end{equation}


\subsubsection{Weighted Orthogonality}

WLS residuals satisfy

\begin{equation}
\boxed{
\mathbf{B}_t^\top
\mathbf{W}_t
\mathbf{s}_t^{\mathrm{neutral}}
=
\mathbf{0}.
}
\end{equation}

Notice that this is weighted orthogonality.

In general,

\begin{equation}
\mathbf{B}_t^\top
\mathbf{s}_t^{\mathrm{neutral}}
\neq
0.
\end{equation}

Thus, WLS neutralization changes the definition of neutrality.


\subsubsection{Possible Weighting Schemes}

Possible weights include:

\begin{itemize}
    \item market capitalization;
    \item square-root market capitalization;
    \item liquidity;
    \item inverse variance;
    \item portfolio relevance;
    \item data-quality confidence.
\end{itemize}

For example,

\begin{equation}
\omega_{i,t}
=
\sqrt{
\mathrm{MCap}_{i,t}
}
\end{equation}

gives greater influence to larger securities without allowing the largest firms
to dominate as strongly as full market-cap weighting.


\subsubsection{Economic Meaning of the Metric}

The choice of $\mathbf{W}_t$ defines the inner product

\begin{equation}
\langle
\mathbf{x},
\mathbf{y}
\rangle_W
=
\mathbf{x}^{\top}
\mathbf{W}_t
\mathbf{y}.
\end{equation}

Neutralization therefore depends not only on the exposure matrix but also on
the metric under which orthogonality is defined.

This distinction becomes important when comparing equal-weight and
capitalization-weighted neutralization.


% ------------------------------------------------------------
\subsection{Residualized Signal Interpretation}
\label{subsec:residualized_signal_interpretation}
% ------------------------------------------------------------

After neutralization,

\begin{equation}
\boxed{
\mathbf{s}_t^{\mathrm{neutral}}
=
\mathbf{M}_{B,t}
\mathbf{s}_t^{\mathrm{raw}}
}
\end{equation}

under OLS, or its weighted analogue under WLS.

The residual signal should be interpreted as the component of the original
score that is unexplained by the chosen exposure matrix.


\subsubsection{Not Necessarily Pure Alpha}

It is tempting to interpret the residual as ``pure alpha.''

This interpretation should be used cautiously.

The residual is orthogonal only to the variables explicitly included in

\begin{equation}
\mathbf{B}_t.
\end{equation}

It may remain exposed to omitted factors.

Thus,

\begin{equation}
\boxed{
\text{Residualized Signal}
\neq
\text{Guaranteed Pure Alpha}.
}
\end{equation}

More precisely,

\begin{equation}
\boxed{
\mathbf{s}_t^{\mathrm{neutral}}
=
\text{component unexplained by the specified controls}.
}
\end{equation}


\subsubsection{Loss of Signal Magnitude}

Because part of the raw signal is removed,

\begin{equation}
\left\|
\mathbf{s}_t^{\mathrm{neutral}}
\right\|_2
\leq
\left\|
\mathbf{s}_t^{\mathrm{raw}}
\right\|_2
\end{equation}

under an orthogonal OLS projection.

Indeed,

\begin{equation}
\left\|
\mathbf{s}_t^{\mathrm{raw}}
\right\|_2^2
=
\left\|
\mathbf{P}_{B,t}
\mathbf{s}_t^{\mathrm{raw}}
\right\|_2^2
+
\left\|
\mathbf{M}_{B,t}
\mathbf{s}_t^{\mathrm{raw}}
\right\|_2^2.
\end{equation}

Therefore,

\begin{equation}
\boxed{
\left\|
\mathbf{s}_t^{\mathrm{neutral}}
\right\|_2^2
=
\left\|
\mathbf{s}_t^{\mathrm{raw}}
\right\|_2^2
-
\left\|
\mathbf{s}_t^{\mathrm{systematic}}
\right\|_2^2.
}
\end{equation}

Neutralization mechanically reduces signal variation whenever the raw score is
correlated with the controlled exposures.


\subsubsection{Explained Fraction of the Signal}

The cross-sectional regression $R^2$ provides a useful diagnostic:

\begin{equation}
\boxed{
R_t^2
=
1
-
\frac{
\sum_i
\left(
s_{i,t}^{\mathrm{neutral}}
\right)^2
}{
\sum_i
\left(
s_{i,t}^{\mathrm{raw}}
-
\bar{s}_t^{\mathrm{raw}}
\right)^2
}.
}
\end{equation}

A high

\begin{equation}
R_t^2
\end{equation}

indicates that a large fraction of the raw signal is explained by the
neutralization variables.

If

\begin{equation}
R_t^2
\approx
1,
\end{equation}

the chosen neutralization removes almost all cross-sectional variation in the
signal.

This may indicate that the supposed alpha is largely a disguised systematic
exposure.


\subsubsection{Preserving the Raw and Neutralized Signals}

The backtesting architecture should preserve both

\begin{equation}
\mathbf{s}_t^{\mathrm{raw}}
\end{equation}

and

\begin{equation}
\mathbf{s}_t^{\mathrm{neutral}}.
\end{equation}

This allows direct comparison of:

\begin{itemize}
    \item predictive performance;
    \item factor exposures;
    \item cross-sectional dispersion;
    \item turnover;
    \item portfolio performance;
    \item performance attribution.
\end{itemize}

The complete signal path is therefore

\begin{equation}
\boxed{
\begin{aligned}
\mathbf{s}_t^{\mathrm{raw}}
&=
\underbrace{
\mathbf{P}_{B,t}
\mathbf{s}_t^{\mathrm{raw}}
}_{\text{controlled component}}
+
\underbrace{
\mathbf{M}_{B,t}
\mathbf{s}_t^{\mathrm{raw}}
}_{\text{residual component}}
\\[0.15cm]
&\rightarrow
\mathbf{s}_t^{\mathrm{neutral}}
=
\mathbf{M}_{B,t}
\mathbf{s}_t^{\mathrm{raw}}.
\end{aligned}
}
\end{equation}


\subsubsection{Interface with Final Signal Construction}

The neutralized score is not necessarily the final signal used for portfolio
construction.

Further transformations may include:

\begin{itemize}
    \item re-standardization;
    \item combination with other neutralized signals;
    \item temporal smoothing;
    \item turnover-aware adjustments;
    \item confidence weighting.
\end{itemize}

Thus,

\begin{equation}
\boxed{
\mathbf{s}_t^{\mathrm{raw}}
\rightarrow
\mathbf{s}_t^{\mathrm{neutral}}
\rightarrow
\mathbf{s}_t^{\mathrm{final}}
\rightarrow
\mathbf{w}_t.
}
\end{equation}

The neutralization layer therefore acts as the interface between alpha
generation and the construction of the final investable signal.


% ============================================================
\section{Final Signal Construction}
\label{sec:final_signal_construction}
% ============================================================

After neutralization, the signal is not necessarily ready to be converted
directly into portfolio weights.

Additional transformations may be required to ensure that the signal has a
stable cross-sectional scale, combines multiple sources of information
consistently, avoids unnecessary temporal instability, incorporates trading
considerations, and reflects differences in signal confidence.

The objective of this section is therefore to transform the neutralized signal

\begin{equation}
\mathbf{s}_t^{\mathrm{neutral}}
\in
\mathbb{R}^{N_t}
\end{equation}

into a final investable signal

\begin{equation}
\boxed{
\mathbf{s}_t^{\mathrm{final}}
\in
\mathbb{R}^{N_t}.
}
\end{equation}

The general architecture may be represented schematically as

\begin{equation}
\boxed{
\begin{aligned}
\mathbf{s}_t^{\mathrm{neutral}}
&\rightarrow
\text{Standardization}
\rightarrow
\text{Signal Combination}
\\
&\rightarrow
\text{Smoothing}
\rightarrow
\text{Turnover Adjustment}
\\
&\rightarrow
\text{Confidence Adjustment}
\rightarrow
\mathbf{s}_t^{\mathrm{final}}.
\end{aligned}
}
\end{equation}

Not every strategy requires every transformation. The precise sequence should
be determined by the economic interpretation of the signal and the intended
portfolio-construction methodology.

Moreover, the order of transformations may matter because several of these
operators are nonlinear. Consequently, the signal-construction pipeline should
be explicitly defined and reproduced identically in research and production.


% ------------------------------------------------------------
\subsection{Post-Neutralization Standardization}
\label{subsec:post_neutralization_standardization}
% ------------------------------------------------------------

Neutralization generally changes the cross-sectional location and dispersion of
the original signal.

Let

\begin{equation}
\mathbf{s}_t^{\mathrm{neutral}}
=
\begin{bmatrix}
s_{1,t}^{\mathrm{neutral}}
\\
\vdots
\\
s_{N_t,t}^{\mathrm{neutral}}
\end{bmatrix}.
\end{equation}

Even if the raw signal was standardized before neutralization, there is no
reason for the residualized signal to retain the same variance.

It is therefore common to standardize the signal again after neutralization.


\subsubsection{Cross-Sectional Z-Score}

Define

\begin{equation}
\mu_t^{s}
=
\frac{1}{N_t}
\sum_{i=1}^{N_t}
s_{i,t}^{\mathrm{neutral}},
\end{equation}

and

\begin{equation}
\sigma_t^{s}
=
\sqrt{
\frac{1}{N_t}
\sum_{i=1}^{N_t}
\left(
s_{i,t}^{\mathrm{neutral}}
-
\mu_t^{s}
\right)^2
}.
\end{equation}

The standardized signal is

\begin{equation}
\boxed{
z_{i,t}^{s}
=
\frac{
s_{i,t}^{\mathrm{neutral}}
-
\mu_t^{s}
}{
\sigma_t^{s}
}.
}
\end{equation}

In vector notation,

\begin{equation}
\mathbf{z}_t^{s}
\in
\mathbb{R}^{N_t}.
\end{equation}

If an intercept was included in an OLS neutralization regression, then

\begin{equation}
\mu_t^{s}=0
\end{equation}

up to numerical precision, and the transformation primarily rescales the
cross-sectional dispersion.


\subsubsection{Why Re-Standardize?}

Suppose two dates have

\begin{equation}
\operatorname{Std}_i
\left(
s_{i,t}^{\mathrm{neutral}}
\right)
=
0.25
\end{equation}

and

\begin{equation}
\operatorname{Std}_i
\left(
s_{i,t+1}^{\mathrm{neutral}}
\right)
=
1.50.
\end{equation}

Without re-standardization, the numerical magnitude of the signal changes
substantially through time.

If the downstream signal-to-weight mapping depends on signal magnitude, this
may mechanically alter portfolio exposure even though the relative
cross-sectional information content is similar.

Re-standardization separates

\begin{equation}
\boxed{
\text{Cross-Sectional Ordering}
}
\end{equation}

from

\begin{equation}
\boxed{
\text{Cross-Sectional Signal Scale}.
}
\end{equation}


\subsubsection{Alternative Standardizations}

Post-neutralization normalization need not use a classical z-score.

Possible alternatives include:

\begin{itemize}
    \item robust z-scores;
    \item ranks;
    \item percentile scores;
    \item Gaussian rank transformations;
    \item bounded transformations.
\end{itemize}

These transformations were introduced previously in the raw-score construction
section and need not be re-derived here.

The important point is that post-neutralization normalization acts on

\begin{equation}
\mathbf{s}_t^{\mathrm{neutral}},
\end{equation}

rather than on the original firm-level characteristic.


% ------------------------------------------------------------
\subsection{Multiple-Signal Combination}
\label{subsec:multiple_signal_combination}
% ------------------------------------------------------------

Many investment strategies combine several alpha signals rather than relying on
a single predictor.

Suppose there are $K$ signals.

For security $i$, define

\begin{equation}
\mathbf{s}_{i,t}
=
\begin{bmatrix}
s_{i,t}^{(1)}
\\
\vdots
\\
s_{i,t}^{(K)}
\end{bmatrix}
\in
\mathbb{R}^{K}.
\end{equation}

Across all securities, define the signal matrix

\begin{equation}
\boxed{
\mathbf{S}_t
=
\begin{bmatrix}
s_{1,t}^{(1)} & \cdots & s_{1,t}^{(K)}
\\
\vdots & \ddots & \vdots
\\
s_{N_t,t}^{(1)} & \cdots & s_{N_t,t}^{(K)}
\end{bmatrix}
\in
\mathbb{R}^{N_t\times K}.
}
\end{equation}

Let

\begin{equation}
\boldsymbol{\alpha}_t
=
\begin{bmatrix}
\alpha_{1,t}
\\
\vdots
\\
\alpha_{K,t}
\end{bmatrix}
\in
\mathbb{R}^{K}
\end{equation}

denote the signal-combination weights.

The combined signal is

\begin{equation}
\boxed{
\mathbf{s}_t^{\mathrm{comb}}
=
\mathbf{S}_t
\boldsymbol{\alpha}_t
\in
\mathbb{R}^{N_t}.
}
\end{equation}

At the individual-security level,

\begin{equation}
\boxed{
s_{i,t}^{\mathrm{comb}}
=
\sum_{k=1}^{K}
\alpha_{k,t}
s_{i,t}^{(k)}.
}
\end{equation}


\subsubsection{Equal-Weighted Combination}

The simplest specification is

\begin{equation}
\alpha_{k,t}
=
\frac{1}{K}.
\end{equation}

Then

\begin{equation}
\boxed{
s_{i,t}^{\mathrm{comb}}
=
\frac{1}{K}
\sum_{k=1}^{K}
s_{i,t}^{(k)}.
}
\end{equation}

This approach is simple, transparent, and generally robust to estimation error.

However, the individual signals should usually be placed on comparable scales
before averaging them.


\subsubsection{Fixed Unequal Weights}

If economic reasoning suggests that some signals should receive greater
importance, fixed weights may be used:

\begin{equation}
\sum_{k=1}^{K}
\alpha_k
=
1.
\end{equation}

For example,

\begin{equation}
\boldsymbol{\alpha}
=
\begin{bmatrix}
0.50\\
0.30\\
0.20
\end{bmatrix}.
\end{equation}

Such weights should ideally be determined ex ante rather than selected to
maximize full-sample backtest performance.


\subsubsection{Performance-Based Weights}

Weights may also depend on historically estimated signal quality.

Let

\begin{equation}
q_{k,t}
\end{equation}

denote a point-in-time estimate of the quality of signal $k$.

A simple normalized weighting scheme is

\begin{equation}
\boxed{
\alpha_{k,t}
=
\frac{
q_{k,t}
}{
\sum_{j=1}^{K} q_{j,t}
}.
}
\end{equation}

Possible quality measures include historical:

\begin{itemize}
    \item information coefficients;
    \item rank information coefficients;
    \item long--short spreads;
    \item information ratios;
    \item net portfolio Sharpe ratios.
\end{itemize}

Any dynamically estimated combination weight must itself respect the
point-in-time information constraint:

\begin{equation}
\boxed{
\boldsymbol{\alpha}_t
\in
\mathcal{I}_t.
}
\end{equation}


\subsubsection{Correlation Between Signals}

Signal quality should not be considered independently of signal redundancy.

Suppose two signals have high individual predictive power but are almost
perfectly correlated:

\begin{equation}
\operatorname{Corr}
\left(
s^{(1)},
s^{(2)}
\right)
\approx 1.
\end{equation}

Combining them with equal weights does not provide the same diversification
benefit as combining two signals with similar quality but low correlation.

Define the signal covariance matrix

\begin{equation}
\boldsymbol{\Sigma}_{s,t}
\in
\mathbb{R}^{K\times K}.
\end{equation}

More sophisticated signal-combination schemes may depend jointly on expected
signal quality and the covariance structure between signals.

Let

\begin{equation}
\widehat{\boldsymbol{\mu}}_{s,t}
=
\begin{bmatrix}
\widehat{\mu}_{1,t}
\\
\vdots
\\
\widehat{\mu}_{K,t}
\end{bmatrix}
\in
\mathbb{R}^{K}
\end{equation}

denote estimates of the expected performance of the $K$ signals, and let

\begin{equation}
\widehat{\boldsymbol{\Sigma}}_{s,t}
\in
\mathbb{R}^{K\times K}
\end{equation}

denote their estimated covariance matrix.

The combination weights may then be written abstractly as

\begin{equation}
\boxed{
\boldsymbol{\alpha}_t
=
g
\left(
\widehat{\boldsymbol{\mu}}_{s,t},
\widehat{\boldsymbol{\Sigma}}_{s,t}
\right).
}
\end{equation}


\subsubsection{Mean--Variance Signal Combination}

A direct analogue of mean--variance portfolio construction can be applied at
the signal level.

Suppose each signal $k$ generates a historical return series

\begin{equation}
R_{k,\tau}^{S},
\end{equation}

for example from a standardized long--short portfolio constructed from that
signal.

Define

\begin{equation}
\widehat{\mu}_{k,t}^{S}
=
\mathbb{E}_t
\left[
R_{k}^{S}
\right],
\end{equation}

and let

\begin{equation}
\widehat{\boldsymbol{\Sigma}}_{s,t}
=
\operatorname{Cov}_t
\left(
\mathbf{R}^{S}
\right)
\end{equation}

denote the estimated covariance matrix of the signal return streams.

The signal-combination problem can then be formulated as

\begin{equation}
\boxed{
\boldsymbol{\alpha}_t^{*}
=
\underset{\boldsymbol{\alpha}}{\arg\max}
\left[
\boldsymbol{\alpha}^{\top}
\widehat{\boldsymbol{\mu}}_{s,t}
-
\frac{\lambda}{2}
\boldsymbol{\alpha}^{\top}
\widehat{\boldsymbol{\Sigma}}_{s,t}
\boldsymbol{\alpha}
\right],
}
\end{equation}

where $\lambda>0$ controls the penalty applied to variation in the combined
signal return.

Without additional constraints, the solution is proportional to

\begin{equation}
\boxed{
\boldsymbol{\alpha}_t^{*}
\propto
\widehat{\boldsymbol{\Sigma}}_{s,t}^{-1}
\widehat{\boldsymbol{\mu}}_{s,t}.
}
\end{equation}

Thus, a signal receives a larger weight when:

\begin{itemize}
    \item its expected historical performance is stronger;
    \item its variance is lower;
    \item its return stream provides diversification relative to the other
    signals.
\end{itemize}

For two signals, the covariance matrix is

\begin{equation}
\widehat{\boldsymbol{\Sigma}}_{s,t}
=
\begin{bmatrix}
\widehat{\sigma}_{1,t}^{2}
&
\widehat{\rho}_{12,t}
\widehat{\sigma}_{1,t}
\widehat{\sigma}_{2,t}
\\
\widehat{\rho}_{12,t}
\widehat{\sigma}_{1,t}
\widehat{\sigma}_{2,t}
&
\widehat{\sigma}_{2,t}^{2}
\end{bmatrix}.
\end{equation}

Consequently, two individually strong but highly correlated signals receive
less combined weight than two equally strong signals that provide more
independent sources of alpha.


\subsubsection{Maximum Sharpe Signal Combination}

A closely related objective is to maximize the estimated Sharpe ratio of the
combined signal return:

\begin{equation}
\boxed{
SR
\left(
\boldsymbol{\alpha}
\right)
=
\frac{
\boldsymbol{\alpha}^{\top}
\widehat{\boldsymbol{\mu}}_{s,t}
}{
\sqrt{
\boldsymbol{\alpha}^{\top}
\widehat{\boldsymbol{\Sigma}}_{s,t}
\boldsymbol{\alpha}
}
}.
}
\end{equation}

The optimization problem is

\begin{equation}
\boxed{
\boldsymbol{\alpha}_t^{*}
=
\underset{\boldsymbol{\alpha}}{\arg\max}
\;
\frac{
\boldsymbol{\alpha}^{\top}
\widehat{\boldsymbol{\mu}}_{s,t}
}{
\sqrt{
\boldsymbol{\alpha}^{\top}
\widehat{\boldsymbol{\Sigma}}_{s,t}
\boldsymbol{\alpha}
}
}.
}
\end{equation}

Under the usual unconstrained assumptions, the direction of the optimal
solution is again

\begin{equation}
\boxed{
\boldsymbol{\alpha}_t^{*}
\propto
\widehat{\boldsymbol{\Sigma}}_{s,t}^{-1}
\widehat{\boldsymbol{\mu}}_{s,t}.
}
\end{equation}

The weights may subsequently be normalized, for example by imposing

\begin{equation}
\sum_{k=1}^{K}
\alpha_{k,t}=1.
\end{equation}


\subsubsection{Performance-Adjusted Diversification Weighting}

A more robust alternative can separate signal quality from diversification.

Let

\begin{equation}
q_{k,t}>0
\end{equation}

denote a historical quality measure for signal $k$, such as its rolling
Information Ratio or Sharpe ratio.

First define preliminary performance weights

\begin{equation}
\widetilde{\alpha}_{k,t}
\propto
q_{k,t}.
\end{equation}

These weights can then be adjusted for redundancy with other signals.

For example, define the average absolute correlation of signal $k$ with the
other signals as

\begin{equation}
\boxed{
\bar{\rho}_{k,t}
=
\frac{1}{K-1}
\sum_{\substack{j=1\\j\neq k}}^{K}
\left|
\rho_{kj,t}
\right|.
}
\end{equation}

A simple diversification-adjusted quality score is

\begin{equation}
\boxed{
d_{k,t}
=
\frac{
q_{k,t}
}{
1+\bar{\rho}_{k,t}
}.
}
\end{equation}

The final combination weights can then be normalized as

\begin{equation}
\boxed{
\alpha_{k,t}
=
\frac{
d_{k,t}
}{
\sum_{j=1}^{K}d_{j,t}
}.
}
\end{equation}

Thus, a signal receives a larger allocation when it has strong historical
performance but a smaller allocation when it is highly redundant with the
other signals.

This procedure is heuristic rather than the solution to a formal
mean--variance optimization problem, but it can be more robust when the
estimated covariance matrix is noisy.


\subsubsection{Combination Before or After Neutralization}

Signals may be combined before neutralization:

\begin{equation}
\mathbf{s}_t^{\mathrm{comb,raw}}
=
\sum_{k=1}^{K}
\alpha_{k,t}
\mathbf{s}_t^{(k),\mathrm{raw}},
\end{equation}

followed by

\begin{equation}
\mathbf{s}_t^{\mathrm{final}}
=
\mathbf{M}_{B,t}
\mathbf{s}_t^{\mathrm{comb,raw}}.
\end{equation}

Alternatively, each signal may first be neutralized:

\begin{equation}
\mathbf{s}_t^{(k),\mathrm{neutral}}
=
\mathbf{M}_{B,t}
\mathbf{s}_t^{(k),\mathrm{raw}},
\end{equation}

and then combined:

\begin{equation}
\mathbf{s}_t^{\mathrm{final}}
=
\sum_{k=1}^{K}
\alpha_{k,t}
\mathbf{s}_t^{(k),\mathrm{neutral}}.
\end{equation}

If the same linear neutralization operator $\mathbf{M}_{B,t}$ and fixed
cross-sectional combination weights are applied to every signal, linearity
implies

\begin{equation}
\boxed{
\mathbf{M}_{B,t}
\left(
\sum_{k=1}^{K}
\alpha_{k,t}
\mathbf{s}_t^{(k),\mathrm{raw}}
\right)
=
\sum_{k=1}^{K}
\alpha_{k,t}
\mathbf{M}_{B,t}
\mathbf{s}_t^{(k),\mathrm{raw}}.
}
\end{equation}

Thus, under these conditions, combination and neutralization commute.

This equivalence generally disappears if different neutralization operators,
nonlinear transformations, signal-specific universes, or security-specific
combination weights are used.


% ------------------------------------------------------------
\subsection{Signal Smoothing}
\label{subsec:signal_smoothing}
% ------------------------------------------------------------

Cross-sectional signals may fluctuate substantially between consecutive
decision dates.

Rapid changes in the signal may lead to:

\begin{itemize}
    \item unstable rankings;
    \item frequent position reversals;
    \item high turnover;
    \item larger transaction costs;
    \item sensitivity to measurement noise.
\end{itemize}

Signal smoothing attempts to retain persistent information while reducing
high-frequency variation.


\subsubsection{Simple Moving Average}

For a smoothing window of length $L$,

\begin{equation}
\boxed{
\widetilde{s}_{i,t}
=
\frac{1}{L}
\sum_{\ell=0}^{L-1}
s_{i,t-\ell}.
}
\end{equation}

The smoothed signal averages the most recent $L$ observations.


\subsubsection{Exponentially Weighted Smoothing}

A common alternative is

\begin{equation}
\boxed{
\widetilde{s}_{i,t}
=
\lambda
s_{i,t}
+
(1-\lambda)
\widetilde{s}_{i,t-1},
}
\end{equation}

where

\begin{equation}
0<\lambda\leq1.
\end{equation}

A large $\lambda$ gives greater importance to the current signal and therefore
produces faster adaptation.

A small $\lambda$ produces stronger smoothing.


\subsubsection{Half-Life Parameterization}

Exponential smoothing may also be expressed using a half-life $H$.

The decay factor is

\begin{equation}
\rho
=
2^{-1/H},
\end{equation}

so that

\begin{equation}
\boxed{
\widetilde{s}_{i,t}
=
(1-\rho)s_{i,t}
+
\rho\widetilde{s}_{i,t-1}.
}
\end{equation}

The weight assigned to an observation $H$ periods old is one half of the
corresponding current-period weight.


\subsubsection{Economic Trade-Off}

Smoothing introduces a fundamental trade-off:

\begin{equation}
\boxed{
\text{Lower Noise and Turnover}
\quad\longleftrightarrow\quad
\text{Slower Reaction to New Information}.
}
\end{equation}

Therefore, smoothing should not be interpreted as a universally beneficial
operation.

Its usefulness depends on the persistence and decay rate of the underlying
alpha signal.


\subsubsection{Point-in-Time Requirement}

Any smoothing procedure must be backward looking.

At time $t$,

\begin{equation}
\widetilde{s}_{i,t}
=
g
\left(
s_{i,t},
s_{i,t-1},
\ldots
\right),
\end{equation}

and must not depend on

\begin{equation}
s_{i,t+1},
s_{i,t+2},
\ldots.
\end{equation}


% ------------------------------------------------------------
\subsection{Turnover-Aware Signal Adjustments}
\label{subsec:turnover_aware_signal_adjustments}
% ------------------------------------------------------------

Signal smoothing indirectly reduces turnover by reducing temporal variation in
the signal.

A more explicit approach is to account for the current portfolio or previous
signal when determining whether a new signal change is economically large
enough to justify trading.

The key idea is that a small change in expected attractiveness may not justify
the transaction costs required to modify an existing position.


\subsubsection{Signal Change}

Define

\begin{equation}
\Delta s_{i,t}
=
s_{i,t}
-
s_{i,t-1}.
\end{equation}

A large value of

\begin{equation}
|\Delta s_{i,t}|
\end{equation}

indicates a substantial change in the investment view, whereas a small change
may primarily reflect noise.


\subsubsection{No-Trade Threshold}

A simple turnover-aware rule introduces a threshold $\tau_i\geq0$:

\begin{equation}
\boxed{
s_{i,t}^{\mathrm{adj}}
=
\begin{cases}
s_{i,t-1}^{\mathrm{adj}},
&
|s_{i,t}-s_{i,t-1}^{\mathrm{adj}}|
\leq
\tau_i,
\\[0.15cm]
s_{i,t},
&
|s_{i,t}-s_{i,t-1}^{\mathrm{adj}}|
>
\tau_i.
\end{cases}
}
\end{equation}

Small signal changes therefore do not modify the investable signal.


\subsubsection{Security-Specific Thresholds}

The threshold may depend on implementation costs:

\begin{equation}
\boxed{
\tau_{i,t}
=
g
\left(
\mathrm{Spread}_{i,t},
\mathrm{ADV}_{i,t},
\mathrm{Volatility}_{i,t},
\mathrm{MarketImpact}_{i,t},
\ldots
\right).
}
\end{equation}

Securities that are expensive to trade may therefore require stronger signal
changes before the desired position is modified.


\subsubsection{Signal Inertia}

Another approach blends the current desired signal with the previous
implemented signal:

\begin{equation}
\boxed{
s_{i,t}^{\mathrm{adj}}
=
\lambda_{i,t}s_{i,t}
+
(1-\lambda_{i,t})
s_{i,t-1}^{\mathrm{adj}},
}
\end{equation}

where

\begin{equation}
0\leq\lambda_{i,t}\leq1.
\end{equation}

For costly securities, $\lambda_{i,t}$ may be smaller, producing greater
inertia.


\subsubsection{Signal-Layer vs.\ Portfolio-Layer Turnover Control}

A distinction should be made between

\begin{equation}
\boxed{
\text{Turnover-Aware Signal Construction}
}
\end{equation}

and

\begin{equation}
\boxed{
\text{Turnover-Constrained Portfolio Optimization}.
}
\end{equation}

The former modifies the alpha input before portfolio construction.

The latter directly penalizes or constrains changes in portfolio weights, for
example through an objective of the form

\begin{equation}
\max_{\mathbf{w}_t}
\left[
\text{Expected Alpha}
-
\text{Risk Penalty}
-
\text{Trading-Cost Penalty}
\right].
\end{equation}

Detailed turnover constraints and transaction-cost-aware optimization belong to
the portfolio-construction layer.

In sophisticated implementations, explicit portfolio-level treatment is often
preferable because actual transaction costs depend on trades in portfolio
weights rather than merely changes in signal values.


% ------------------------------------------------------------
\subsection{Confidence Weighting}
\label{subsec:confidence_weighting}
% ------------------------------------------------------------

Not all signal observations necessarily have the same degree of reliability.

Let

\begin{equation}
c_{i,t}
\end{equation}

denote a confidence measure associated with the signal for security $i$ at time
$t$.

A simple confidence-adjusted signal is

\begin{equation}
\boxed{
s_{i,t}^{\mathrm{conf}}
=
c_{i,t}
s_{i,t},
}
\end{equation}

where, for example,

\begin{equation}
c_{i,t}
\in[0,1].
\end{equation}

A value close to one leaves the signal largely unchanged, whereas a value close
to zero shrinks the signal toward zero.


\subsubsection{Sources of Confidence}

Confidence may depend on:

\begin{itemize}
    \item data quality;
    \item number of observations used to estimate a characteristic;
    \item number of analyst estimates;
    \item dispersion of analyst estimates;
    \item model prediction uncertainty;
    \item historical signal reliability;
    \item model ensemble agreement;
    \item distance from the training-data distribution.
\end{itemize}


\subsubsection{Prediction Uncertainty}

Suppose a model produces both a prediction

\begin{equation}
\widehat{y}_{i,t+h|t}
\end{equation}

and an estimated prediction uncertainty

\begin{equation}
\widehat{\sigma}_{i,t}^{\mathrm{pred}}.
\end{equation}

A confidence-adjusted score could depend on a signal-to-uncertainty ratio:

\begin{equation}
\boxed{
s_{i,t}^{\mathrm{conf}}
=
\frac{
\widehat{y}_{i,t+h|t}
}{
\widehat{\sigma}_{i,t}^{\mathrm{pred}}
}.
}
\end{equation}

This resembles a forecast-level risk adjustment: large predictions supported
by high uncertainty receive less importance than equally large predictions
estimated more precisely.


\subsubsection{Ensemble Agreement}

Suppose $M$ models generate predictions

\begin{equation}
\widehat{y}_{i,t}^{(1)},
\ldots,
\widehat{y}_{i,t}^{(M)}.
\end{equation}

The ensemble mean is

\begin{equation}
\overline{y}_{i,t}
=
\frac{1}{M}
\sum_{m=1}^{M}
\widehat{y}_{i,t}^{(m)},
\end{equation}

while disagreement may be measured by

\begin{equation}
d_{i,t}
=
\sqrt{
\frac{1}{M}
\sum_{m=1}^{M}
\left(
\widehat{y}_{i,t}^{(m)}
-
\overline{y}_{i,t}
\right)^2
}.
\end{equation}

Confidence may then be specified as a decreasing function of disagreement:

\begin{equation}
\boxed{
c_{i,t}
=
g(d_{i,t}),
\qquad
g'(\cdot)<0.
}
\end{equation}

Thus, securities for which the underlying models strongly disagree receive a
smaller final signal.


\subsubsection{Avoiding Double Counting}

Confidence weighting should not mechanically duplicate adjustments that will
later be performed during portfolio construction.

For example, if portfolio weights will already be scaled inversely with
security volatility, defining

\begin{equation}
c_{i,t}
\propto
\frac{1}{\sigma_{i,t}}
\end{equation}

at the signal level may introduce the same low-volatility preference twice.

Every confidence adjustment should therefore have a distinct economic
interpretation.


% ------------------------------------------------------------
\subsection{Final Investable Signal}
\label{subsec:final_investable_signal}
% ------------------------------------------------------------

The final output of the complete signal-generation engine is

\begin{equation}
\boxed{
\mathbf{s}_t^{\mathrm{final}}
=
\begin{bmatrix}
s_{1,t}^{\mathrm{final}}
\\
\vdots
\\
s_{N_t,t}^{\mathrm{final}}
\end{bmatrix}
\in
\mathbb{R}^{N_t}.
}
\end{equation}

This vector summarizes the relative investment attractiveness of all securities
in the investable universe at decision time $t$ after all desired signal-level
transformations have been applied.


\subsubsection{Generic Operator Representation}

The entire final-signal layer can be represented by an operator

\begin{equation}
\boxed{
\mathbf{s}_t^{\mathrm{final}}
=
\mathcal{F}_t
\left(
\mathbf{s}_t^{\mathrm{neutral}},
\mathbf{s}_{t-1}^{\mathrm{final}},
\mathcal{I}_t
\right),
}
\end{equation}

where $\mathcal{F}_t$ may incorporate:

\begin{itemize}
    \item post-neutralization normalization;
    \item multiple-signal combination;
    \item temporal smoothing;
    \item turnover-aware adjustments;
    \item confidence weighting.
\end{itemize}

The information-set requirement remains

\begin{equation}
\boxed{
\mathbf{s}_t^{\mathrm{final}}
\in
\mathcal{I}_t.
}
\end{equation}

No component of the final signal may depend on information unavailable at the
historical decision time.


\subsubsection{Signal Scale and Economic Meaning}

The interpretation of the numerical magnitude of

\begin{equation}
s_{i,t}^{\mathrm{final}}
\end{equation}

must be explicitly defined.

Depending on the strategy, the final signal may represent:

\begin{itemize}
    \item a standardized relative attractiveness score;
    \item a percentile or rank;
    \item a calibrated expected return;
    \item a predicted probability;
    \item a confidence-adjusted expected return;
    \item an abstract alpha score.
\end{itemize}

This distinction is essential because the portfolio-construction layer must know
whether signal magnitudes have cardinal economic meaning or merely ordinal
meaning.


\subsubsection{Ordinal vs.\ Cardinal Signals}

An ordinal signal provides information primarily about ordering:

\begin{equation}
s_{i,t}^{\mathrm{final}}
>
s_{j,t}^{\mathrm{final}}
\end{equation}

means that security $i$ is preferred to security $j$.

However, the difference

\begin{equation}
s_{i,t}^{\mathrm{final}}
-
s_{j,t}^{\mathrm{final}}
\end{equation}

need not have a direct economic interpretation.

A cardinal signal has meaningful magnitude. For example, if

\begin{equation}
s_{i,t}^{\mathrm{final}}
=
\widehat{\mu}_{i,t},
\end{equation}

where $\widehat{\mu}_{i,t}$ is a calibrated expected return, then differences in
signal magnitude can enter portfolio optimization directly.

Thus,

\begin{equation}
\boxed{
\text{Ordinal Signal}
\neq
\text{Cardinal Expected Return}.
}
\end{equation}

This distinction will determine which signal-to-weight mappings are
economically appropriate.


\subsubsection{Final Signal Diagnostics}

Before passing the signal to portfolio construction, the backtester should
store and monitor diagnostics such as:

\begin{itemize}
    \item cross-sectional mean;
    \item cross-sectional standard deviation;
    \item minimum and maximum;
    \item quantiles;
    \item fraction of positive and negative scores;
    \item correlation with the previous-period signal;
    \item correlation with each underlying signal;
    \item exposure to neutralized factors;
    \item signal turnover;
    \item number of valid securities.
\end{itemize}

For a factor exposure vector $\mathbf{b}_t$, for example, one may verify

\begin{equation}
\mathbf{b}_t^\top
\mathbf{s}_t^{\mathrm{final}}
\approx
0
\end{equation}

if the subsequent transformations are intended to preserve the corresponding
neutrality condition.

This check is important because nonlinear post-neutralization transformations
can reintroduce exposures that were previously removed.


\subsubsection{Neutrality After Nonlinear Transformations}

Suppose

\begin{equation}
\mathbf{s}_t^{\mathrm{neutral}}
\end{equation}

satisfies

\begin{equation}
\mathbf{B}_t^\top
\mathbf{s}_t^{\mathrm{neutral}}
=
\mathbf{0}.
\end{equation}

If the subsequent transformation is a common linear scaling,

\begin{equation}
\mathbf{s}_t^{\mathrm{final}}
=
a_t
\mathbf{s}_t^{\mathrm{neutral}},
\end{equation}

then

\begin{equation}
\mathbf{B}_t^\top
\mathbf{s}_t^{\mathrm{final}}
=
a_t
\mathbf{B}_t^\top
\mathbf{s}_t^{\mathrm{neutral}}
=
\mathbf{0}.
\end{equation}

Neutrality is preserved.

However, for a nonlinear element-wise transformation

\begin{equation}
s_{i,t}^{\mathrm{final}}
=
g
\left(
s_{i,t}^{\mathrm{neutral}}
\right),
\end{equation}

one generally has

\begin{equation}
\boxed{
\mathbf{B}_t^\top
\mathbf{s}_t^{\mathrm{neutral}}
=
\mathbf{0}
\;\not\Rightarrow\;
\mathbf{B}_t^\top
\mathbf{s}_t^{\mathrm{final}}
=
\mathbf{0}.
}
\end{equation}

Similarly, security-specific confidence weights

\begin{equation}
s_{i,t}^{\mathrm{final}}
=
c_{i,t}
s_{i,t}^{\mathrm{neutral}}
\end{equation}

may reintroduce systematic exposures.

Consequently, factor exposures should be checked again after the complete
final-signal transformation.


\subsubsection{Interface with Portfolio Construction}

The signal-construction layer terminates with

\begin{equation}
\boxed{
\mathbf{s}_t^{\mathrm{final}}.
}
\end{equation}

The portfolio-construction layer then applies a mapping

\begin{equation}
\boxed{
\mathbf{w}_t^{\mathrm{target}}
=
\mathcal{G}_t
\left(
\mathbf{s}_t^{\mathrm{final}},
\mathbf{\Sigma}_t,
\mathbf{B}_t,
\mathbf{w}_{t^-},
\mathbf{c}_t,
\ldots
\right),
}
\end{equation}

where:

\begin{itemize}
    \item $\mathbf{w}_t^{\mathrm{target}}\in\mathbb{R}^{N_t}$ denotes target
    portfolio weights;
    \item $\mathbf{\Sigma}_t\in\mathbb{R}^{N_t\times N_t}$ denotes a risk or
    covariance matrix when required;
    \item $\mathbf{B}_t$ contains relevant portfolio exposure constraints;
    \item $\mathbf{w}_{t^-}$ denotes pre-rebalance portfolio weights;
    \item $\mathbf{c}_t$ summarizes transaction-cost or liquidity information.
\end{itemize}

This separation provides a clean architectural boundary:

\begin{equation}
\boxed{
\begin{aligned}
\text{Signal Engine}
&:
\quad
\text{Information}
\rightarrow
\mathbf{s}_t^{\mathrm{final}},
\\[0.15cm]
\text{Portfolio Engine}
&:
\quad
\mathbf{s}_t^{\mathrm{final}}
\rightarrow
\mathbf{w}_t^{\mathrm{target}}.
\end{aligned}
}
\end{equation}

The signal engine determines \emph{which securities are attractive} and by how
much according to the chosen signal convention.

The portfolio engine determines \emph{how much capital should actually be
allocated} to those securities subject to risk, exposure, leverage, liquidity,
turnover, and implementation constraints.

% ------------------------------------------------------------
\subsubsection{Position in the Complete Backtesting Pipeline}
% ------------------------------------------------------------

At this stage, the backtester has transformed heterogeneous raw information into
a single cross-sectional investment signal.

The complete path developed so far is

\begin{equation}
\boxed{
\begin{aligned}
\text{Raw Data}
\;(\mathcal{D}^{\mathrm{raw}})
&\rightarrow
\text{Point-in-Time Data}
\;(\mathcal{D}_t^{\mathrm{PIT}})
\\
&\rightarrow
\text{Investment Universe}
\;(\mathcal{U}_t)
\\
&\rightarrow
\text{Cleaned Data}
\;(\mathcal{D}_t^{\mathrm{clean}})
\\
&\rightarrow
\text{Outlier-Treated Data}
\;(\mathcal{D}_t^{\mathrm{treated}})
\\
&\rightarrow
\text{Firm Characteristics / Model Predictions}
\;(\mathbf{x}_t \text{ or } \widehat{\mathbf{y}}_t)
\\
&\rightarrow
\text{Raw Score}
\;(\mathbf{s}_t^{\mathrm{raw}})
\\
&\rightarrow
\text{Neutralized Signal}
\;(\mathbf{s}_t^{\mathrm{neutral}})
\\
&\rightarrow
\text{Final Investable Signal}
\;(\mathbf{s}_t^{\mathrm{final}}).
\end{aligned}
}
\end{equation}

Thus, the process started from raw historical information

\begin{equation}
\mathcal{D}^{\mathrm{raw}},
\end{equation}

and now ends with

\begin{equation}
\boxed{
\mathbf{s}_t^{\mathrm{final}}
\in
\mathbb{R}^{N_t},
}
\end{equation}

a point-in-time, cleaned, transformed, neutralized, and economically
interpretable signal defined for every eligible security in the investment
universe $\mathcal{U}_t$.

The next stage of the backtester converts this signal into actual target
portfolio weights:

\begin{equation}
\boxed{
\mathbf{s}_t^{\mathrm{final}}
\in
\mathbb{R}^{N_t}
\quad
\longrightarrow
\quad
\mathbf{w}_t^{\mathrm{target}}
\in
\mathbb{R}^{N_t}.
}
\end{equation}


% ============================================================
\section{Signal Diagnostics}
\label{sec:signal_diagnostics}
% ============================================================

Before converting the final investable signal into portfolio weights, it is
useful to introduce a short diagnostic step.

At this stage of the backtesting pipeline, we have constructed

\begin{equation}
\boxed{
\mathbf{s}_t^{\mathrm{final}}
=
\begin{bmatrix}
s_{1,t}^{\mathrm{final}}
\\
\vdots
\\
s_{N_t,t}^{\mathrm{final}}
\end{bmatrix}
\in
\mathbb{R}^{N_t},
}
\end{equation}

but we have not yet determined how much capital should be allocated to each
security.

The reader should therefore view this section as a temporary diagnostic layer
between signal construction and portfolio construction:

\begin{equation}
\boxed{
\begin{aligned}
\text{Signal Construction}
&\rightarrow
\mathbf{s}_t^{\mathrm{final}}
\\
&\rightarrow
\underbrace{
\text{Signal Diagnostics}
}_{\text{current section}}
\\
&\rightarrow
\text{Portfolio Construction}
\rightarrow
\mathbf{w}_t.
\end{aligned}
}
\end{equation}

The purpose is to answer a simple question before introducing position sizing,
risk models, leverage, constraints, and transaction costs:

\begin{equation}
\boxed{
\text{Does the signal itself contain useful information about future returns?}
}
\end{equation}

This distinction is important because portfolio performance depends jointly on

\begin{equation}
\boxed{
\text{Signal Quality}
+
\text{Portfolio Construction}
+
\text{Implementation}.
}
\end{equation}

A weak signal may occasionally produce an attractive backtest because of
unintended factor exposures or a favorable portfolio-construction choice.
Conversely, a genuinely predictive signal may produce disappointing portfolio
performance if it is mapped inefficiently into positions or generates excessive
turnover.

Signal diagnostics therefore attempt to study the statistical and economic
properties of the signal before these additional layers are introduced.


% ------------------------------------------------------------
\subsection{Cross-Sectional Signal Distribution}
\label{subsec:cross_sectional_signal_distribution}
% ------------------------------------------------------------

The first diagnostic is simply to examine the cross-sectional distribution of
the final signal.

At each decision date $t$, define the cross-sectional mean

\begin{equation}
\boxed{
\bar{s}_t
=
\frac{1}{N_t}
\sum_{i=1}^{N_t}
s_{i,t}^{\mathrm{final}}.
}
\end{equation}

The cross-sectional standard deviation is

\begin{equation}
\boxed{
\sigma_{s,t}
=
\sqrt{
\frac{1}{N_t}
\sum_{i=1}^{N_t}
\left(
s_{i,t}^{\mathrm{final}}
-
\bar{s}_t
\right)^2
}.
}
\end{equation}

Other useful diagnostics include:

\begin{itemize}
    \item minimum and maximum;
    \item median;
    \item selected quantiles;
    \item skewness;
    \item kurtosis;
    \item fraction of positive and negative observations;
    \item fraction of missing observations.
\end{itemize}

These statistics help detect abnormal changes in the signal-generation
pipeline.

For example, a sudden collapse in

\begin{equation}
\sigma_{s,t}
\end{equation}

may indicate that the signal has lost cross-sectional discriminatory power or
that an upstream data-processing step has malfunctioned.

Conversely, an unusually large dispersion may indicate extreme observations,
a change in data coverage, or instability in a predictive model.


\subsubsection{Cross-Sectional Dispersion Through Time}

The time series

\begin{equation}
\left\{
\sigma_{s,t}
\right\}_{t=1}^{T}
\end{equation}

is itself an important diagnostic.

If the final signal has been standardized at every date, one may expect

\begin{equation}
\sigma_{s,t}
\approx 1.
\end{equation}

If it has not been standardized, changes in signal dispersion may contain
economic information but may also mechanically change downstream portfolio
exposure when the signal-to-weight mapping depends on signal magnitude.


% ------------------------------------------------------------
\subsection{Information Coefficient}
\label{subsec:information_coefficient}
% ------------------------------------------------------------

One of the most common measures of cross-sectional predictive power is the
Information Coefficient, or IC.

Let

\begin{equation}
R_{i,t\rightarrow t+h}
\end{equation}

denote the realized return of security $i$ over the future horizon
$[t,t+h]$.

At decision date $t$, the Information Coefficient is defined as the
cross-sectional correlation between the signal and subsequent returns:

\begin{equation}
\boxed{
IC_t
=
\operatorname{Corr}_i
\left(
s_{i,t}^{\mathrm{final}},
R_{i,t\rightarrow t+h}
\right).
}
\end{equation}

Using the Pearson correlation coefficient,

\begin{equation}
IC_t
=
\frac{
\sum_{i=1}^{N_t}
\left(
s_{i,t}^{\mathrm{final}}
-
\bar{s}_t
\right)
\left(
R_{i,t\rightarrow t+h}
-
\bar{R}_{t\rightarrow t+h}
\right)
}{
\sqrt{
\sum_{i=1}^{N_t}
\left(
s_{i,t}^{\mathrm{final}}
-
\bar{s}_t
\right)^2
}
\sqrt{
\sum_{i=1}^{N_t}
\left(
R_{i,t\rightarrow t+h}
-
\bar{R}_{t\rightarrow t+h}
\right)^2
}
}.
\end{equation}

A positive IC means that securities with higher signals tended, on average, to
subsequently outperform securities with lower signals.

A negative IC indicates the opposite relationship.


\subsubsection{Average Information Coefficient}

Across $T$ decision dates,

\begin{equation}
\boxed{
\overline{IC}
=
\frac{1}{T}
\sum_{t=1}^{T}
IC_t.
}
\end{equation}

A positive average IC indicates persistent positive cross-sectional predictive
power.

However, the magnitude of $\overline{IC}$ should not be interpreted in
isolation. Its stability through time is equally important.


\subsubsection{Prediction Horizon}

The IC depends explicitly on the future-return horizon $h$:

\begin{equation}
IC_t^{(h)}
=
\operatorname{Corr}_i
\left(
s_{i,t}^{\mathrm{final}},
R_{i,t\rightarrow t+h}
\right).
\end{equation}

Computing ICs for several horizons,

\begin{equation}
h
\in
\{1,5,20,60,\ldots\},
\end{equation}

can therefore help characterize the decay of the signal's predictive content.

A short-lived signal may have

\begin{equation}
IC^{(1)}>0
\end{equation}

but

\begin{equation}
IC^{(60)}
\approx 0.
\end{equation}

This information will later be relevant when choosing the portfolio
rebalancing frequency.


% ------------------------------------------------------------
\subsection{Rank Information Coefficient}
\label{subsec:rank_information_coefficient}
% ------------------------------------------------------------

The Rank Information Coefficient replaces the raw values of the signal and
future returns with their cross-sectional ranks.

It is therefore the cross-sectional Spearman correlation:

\begin{equation}
\boxed{
RIC_t
=
\operatorname{Corr}_i^{\mathrm{Spearman}}
\left(
s_{i,t}^{\mathrm{final}},
R_{i,t\rightarrow t+h}
\right).
}
\end{equation}

Equivalently,

\begin{equation}
RIC_t
=
\operatorname{Corr}_i
\left(
\operatorname{Rank}_i
\left(
s_{i,t}^{\mathrm{final}}
\right),
\operatorname{Rank}_i
\left(
R_{i,t\rightarrow t+h}
\right)
\right).
\end{equation}


\subsubsection{Interpretation}

The Pearson IC measures linear association between signal magnitude and future
returns.

The Rank IC instead measures whether higher-ranked securities tend to have
higher-ranked subsequent returns.

Therefore,

\begin{equation}
\boxed{
\text{IC}
\rightarrow
\text{Linear Predictive Relationship},
}
\end{equation}

whereas

\begin{equation}
\boxed{
\text{Rank IC}
\rightarrow
\text{Monotonic Predictive Relationship}.
}
\end{equation}

Rank IC is particularly useful when the signal is ordinal, such as a percentile
score, rank transformation, or nonlinear model output whose absolute magnitude
has no direct economic interpretation.

It is also less sensitive than Pearson correlation to extreme realized returns.


% ------------------------------------------------------------
\subsection{Information Coefficient Information Ratio}
\label{subsec:information_coefficient_information_ratio}
% ------------------------------------------------------------

The average IC measures the average strength of the relationship, but does not
capture its stability.

Define the time-series standard deviation of the IC as

\begin{equation}
\sigma_{IC}
=
\sqrt{
\frac{1}{T-1}
\sum_{t=1}^{T}
\left(
IC_t-\overline{IC}
\right)^2
}.
\end{equation}

The Information Coefficient Information Ratio is

\begin{equation}
\boxed{
ICIR
=
\frac{
\overline{IC}
}{
\sigma_{IC}
}.
}
\end{equation}

If the IC is measured at a frequency with $F$ observations per year, an
annualized version may be reported as

\begin{equation}
\boxed{
ICIR_{\mathrm{ann}}
=
\sqrt{F}
\frac{
\overline{IC}
}{
\sigma_{IC}
}.
}
\end{equation}

The same construction can be applied to Rank IC:

\begin{equation}
\boxed{
RICIR
=
\frac{
\overline{RIC}
}{
\sigma_{RIC}
}.
}
\end{equation}

A signal with a moderate but stable IC may be economically more useful than a
signal with a larger average IC driven by a small number of exceptional
periods.


\subsubsection{Serial Dependence}

The usual annualization

\begin{equation}
\sqrt{F}
\end{equation}

implicitly treats successive IC observations as sufficiently independent.

This assumption may be inappropriate when future-return horizons overlap or
when the signal is highly persistent.

For example, if a daily IC is computed against future 20-day returns, adjacent
IC observations use strongly overlapping return windows.

Consequently, annualized ICIR values should be interpreted carefully in the
presence of serial dependence.


% ------------------------------------------------------------
\subsection{Quantile and Decile Returns}
\label{subsec:quantile_decile_returns}
% ------------------------------------------------------------

Correlation-based diagnostics summarize the entire cross-section with a single
number.

A complementary approach is to divide securities into groups according to
their signal values and examine their subsequent returns.

Suppose the cross-section is divided into $Q$ quantiles.

Define

\begin{equation}
\mathcal{Q}_{q,t}
\subseteq
\mathcal{U}_t,
\qquad
q=1,\ldots,Q,
\end{equation}

where $\mathcal{Q}_{1,t}$ contains the lowest-signal securities and
$\mathcal{Q}_{Q,t}$ contains the highest-signal securities.

For deciles,

\begin{equation}
Q=10.
\end{equation}

The future equal-weighted return of quantile $q$ is

\begin{equation}
\boxed{
R_{q,t\rightarrow t+h}
=
\frac{1}{
|\mathcal{Q}_{q,t}|
}
\sum_{i\in\mathcal{Q}_{q,t}}
R_{i,t\rightarrow t+h}.
}
\end{equation}

More generally, securities within each quantile may be weighted according to a
specified weighting rule.


\subsubsection{Average Quantile Returns}

Over the backtest,

\begin{equation}
\boxed{
\overline{R}_q
=
\frac{1}{T}
\sum_{t=1}^{T}
R_{q,t\rightarrow t+h}.
}
\end{equation}

The sequence

\begin{equation}
\overline{R}_1,
\overline{R}_2,
\ldots,
\overline{R}_Q
\end{equation}

provides a direct visualization of how future returns vary with the signal.


% ------------------------------------------------------------
\subsection{Long--Short Spread}
\label{subsec:signal_long_short_spread}
% ------------------------------------------------------------

The extreme quantiles can be used to construct a simple diagnostic long--short
spread.

For $Q$ signal portfolios,

\begin{equation}
\boxed{
R_{t}^{LS}
=
R_{Q,t}
-
R_{1,t}.
}
\end{equation}

For deciles,

\begin{equation}
R_t^{D10-D1}
=
R_{10,t}
-
R_{1,t}.
\end{equation}

The average spread is

\begin{equation}
\boxed{
\overline{R}^{LS}
=
\frac{1}{T}
\sum_{t=1}^{T}
R_t^{LS}.
}
\end{equation}

This provides an intuitive measure of the economic separation between
high-signal and low-signal securities.


\subsubsection{Diagnostic Portfolio vs.\ Final Portfolio}

The long--short spread constructed here should not be confused with the final
investment portfolio.

It is deliberately simple and is primarily intended to diagnose the signal.

It may ignore:

\begin{itemize}
    \item optimal position sizing;
    \item covariance information;
    \item leverage targets;
    \item portfolio-level factor constraints;
    \item liquidity constraints;
    \item turnover optimization;
    \item transaction costs.
\end{itemize}

These elements will be introduced in the portfolio-construction and
implementation sections.


% ------------------------------------------------------------
\subsection{Monotonicity}
\label{subsec:signal_monotonicity}
% ------------------------------------------------------------

A useful signal should often produce an economically coherent relationship
between signal strength and subsequent returns.

For a positively oriented signal, one may expect

\begin{equation}
\boxed{
\overline{R}_1
\leq
\overline{R}_2
\leq
\cdots
\leq
\overline{R}_Q.
}
\end{equation}

This property is referred to as monotonicity.


\subsubsection{Why Monotonicity Matters}

Suppose the top decile strongly outperforms the bottom decile, but the
intermediate deciles exhibit no systematic ordering.

The signal may still be economically useful, but its predictive relationship
may be concentrated only in the tails.

Conversely, a smooth increase in returns across quantiles provides stronger
evidence that the signal captures a broad cross-sectional relationship.


\subsubsection{Simple Monotonicity Measures}

One possible diagnostic is the correlation between quantile number and average
quantile return:

\begin{equation}
\boxed{
M
=
\operatorname{Corr}
\left(
\begin{bmatrix}
1\\
2\\
\vdots\\
Q
\end{bmatrix},
\begin{bmatrix}
\overline{R}_1\\
\overline{R}_2\\
\vdots\\
\overline{R}_Q
\end{bmatrix}
\right).
}
\end{equation}

A value close to one indicates a strongly increasing relationship.

A rank-based version may also be used when only the ordering of quantile
returns is of interest.


% ------------------------------------------------------------
\subsection{Hit Rate}
\label{subsec:signal_hit_rate}
% ------------------------------------------------------------

The hit rate measures how frequently the signal produces the expected
directional relationship.


\subsubsection{IC Hit Rate}

A simple definition is the fraction of periods with positive IC:

\begin{equation}
\boxed{
HR_{IC}
=
\frac{1}{T}
\sum_{t=1}^{T}
\mathbb{I}
\left(
IC_t>0
\right),
}
\end{equation}

where $\mathbb{I}(\cdot)$ denotes the indicator function.

Similarly,

\begin{equation}
HR_{RIC}
=
\frac{1}{T}
\sum_{t=1}^{T}
\mathbb{I}
\left(
RIC_t>0
\right).
\end{equation}


\subsubsection{Long--Short Hit Rate}

For the diagnostic long--short spread,

\begin{equation}
\boxed{
HR_{LS}
=
\frac{1}{T}
\sum_{t=1}^{T}
\mathbb{I}
\left(
R_t^{LS}>0
\right).
}
\end{equation}

A high hit rate indicates that the signal works frequently, whereas a lower hit
rate combined with strong average performance may indicate that returns are
generated by relatively infrequent but large successful periods.

Hit rate should therefore always be interpreted jointly with the magnitude and
distribution of returns.


% ------------------------------------------------------------
\subsection{Signal Persistence and Autocorrelation}
\label{subsec:signal_persistence_autocorrelation}
% ------------------------------------------------------------

Signal persistence measures how rapidly the cross-sectional investment view
changes through time.

A natural measure is the cross-sectional correlation between signals at
different dates.

For lag $\ell$,

\begin{equation}
\boxed{
\rho_s(\ell)
=
\operatorname{Corr}_{i,t}
\left(
s_{i,t}^{\mathrm{final}},
s_{i,t-\ell}^{\mathrm{final}}
\right).
}
\end{equation}

More commonly, the correlation is first computed cross-sectionally at each
date:

\begin{equation}
\rho_{s,t}^{(\ell)}
=
\operatorname{Corr}_i
\left(
s_{i,t}^{\mathrm{final}},
s_{i,t-\ell}^{\mathrm{final}}
\right),
\end{equation}

and then averaged through time:

\begin{equation}
\boxed{
\overline{\rho}_s^{(\ell)}
=
\frac{1}{T_\ell}
\sum_t
\rho_{s,t}^{(\ell)}.
}
\end{equation}

A rank-based persistence measure can similarly be defined using Spearman
correlations.


\subsubsection{Interpretation}

A highly persistent signal satisfies approximately

\begin{equation}
\rho_s(1)
\approx 1.
\end{equation}

Such a signal changes slowly and will generally require less portfolio
rebalancing.

A rapidly changing signal may have

\begin{equation}
\rho_s(1)
\approx 0,
\end{equation}

which may imply substantially greater turnover.


\subsubsection{Persistence Curve}

Evaluating

\begin{equation}
\overline{\rho}_s^{(1)},
\overline{\rho}_s^{(2)},
\ldots,
\overline{\rho}_s^{(L)}
\end{equation}

provides a persistence curve.

This helps characterize the effective memory of the signal and can inform the
choice of rebalance frequency and smoothing parameters.


% ------------------------------------------------------------
\subsection{Signal Turnover}
\label{subsec:signal_turnover}
% ------------------------------------------------------------

Signal turnover measures how much the signal changes between consecutive
decision dates.

It should not be confused with portfolio turnover, which is defined from actual
changes in portfolio weights.

A simple measure is

\begin{equation}
\boxed{
TO_t^{s}
=
\frac{1}{N_t}
\sum_{i=1}^{N_t}
\left|
s_{i,t}^{\mathrm{final}}
-
s_{i,t-1}^{\mathrm{final}}
\right|.
}
\end{equation}

If signal scales vary through time, this measure should be interpreted
carefully.

For rank-based signals, one may instead examine rank changes:

\begin{equation}
\boxed{
TO_t^{\mathrm{rank}}
=
\frac{1}{N_t}
\sum_{i=1}^{N_t}
\left|
\operatorname{Rank}_{i,t}
-
\operatorname{Rank}_{i,t-1}
\right|.
}
\end{equation}


\subsubsection{Signal Turnover vs.\ Portfolio Turnover}

In general,

\begin{equation}
\boxed{
TO_t^{s}
\neq
TO_t^{p}.
}
\end{equation}

Portfolio turnover depends on the signal-to-weight mapping, constraints,
portfolio drift, and previous holdings.

Nevertheless, high signal turnover is an early warning that implementation
costs may become important.


% ------------------------------------------------------------
\subsection{Factor Exposure Diagnostics}
\label{subsec:factor_exposure_diagnostics}
% ------------------------------------------------------------

If the signal has been neutralized against a set of systematic exposures, these
exposures should be verified after the complete final-signal construction
pipeline.

Let

\begin{equation}
\mathbf{B}_t
\in
\mathbb{R}^{N_t\times K}
\end{equation}

denote the relevant exposure matrix.

Under OLS neutralization, one may examine

\begin{equation}
\boxed{
\mathbf{e}_t^{s}
=
\mathbf{B}_t^\top
\mathbf{s}_t^{\mathrm{final}}
\in
\mathbb{R}^{K}.
}
\end{equation}

If the final signal is intended to remain neutral,

\begin{equation}
\mathbf{e}_t^{s}
\approx
\mathbf{0}.
\end{equation}

Under WLS neutralization with weighting matrix

\begin{equation}
\mathbf{W}_t
\in
\mathbb{R}^{N_t\times N_t},
\end{equation}

the relevant condition becomes

\begin{equation}
\boxed{
\mathbf{B}_t^\top
\mathbf{W}_t
\mathbf{s}_t^{\mathrm{final}}
\approx
\mathbf{0}.
}
\end{equation}


\subsubsection{Why Recheck Exposures?}

Neutrality may have been imposed earlier on

\begin{equation}
\mathbf{s}_t^{\mathrm{neutral}},
\end{equation}

but subsequent nonlinear transformations can reintroduce systematic exposures.

For example,

\begin{equation}
s_{i,t}^{\mathrm{final}}
=
g
\left(
s_{i,t}^{\mathrm{neutral}}
\right)
\end{equation}

does not generally preserve

\begin{equation}
\mathbf{B}_t^\top
\mathbf{s}_t^{\mathrm{neutral}}
=
0.
\end{equation}

Factor exposures should therefore be monitored on the actual signal passed to
the portfolio-construction engine.


% ------------------------------------------------------------
\subsection{Cross-Sectional and Temporal Stability}
\label{subsec:cross_sectional_temporal_stability}
% ------------------------------------------------------------

A final diagnostic concerns whether signal properties remain stable across
different parts of the investment universe and through different periods.


\subsubsection{Cross-Sectional Stability}

Signal performance may be evaluated separately across groups such as:

\begin{itemize}
    \item large-, mid-, and small-cap securities;
    \item liquidity buckets;
    \item industries and sectors;
    \item countries and regions;
    \item high- and low-volatility securities.
\end{itemize}

For subgroup $g$, one may compute

\begin{equation}
IC_{g,t},
\end{equation}

and compare

\begin{equation}
\overline{IC}_g
\end{equation}

across groups.

A signal whose apparent predictive power is entirely concentrated in a small,
illiquid subset of the universe may be difficult to implement even if its
aggregate IC is attractive.


\subsubsection{Temporal Stability}

The same diagnostics should be examined through time.

For example, define a rolling average IC over a window of length $L$:

\begin{equation}
\boxed{
\overline{IC}_{t}^{(L)}
=
\frac{1}{L}
\sum_{\ell=0}^{L-1}
IC_{t-\ell}.
}
\end{equation}

Similarly, one may examine rolling:

\begin{itemize}
    \item Rank IC;
    \item ICIR;
    \item long--short returns;
    \item hit rates;
    \item signal dispersion;
    \item persistence;
    \item factor exposures.
\end{itemize}

This can reveal regime dependence or structural deterioration that is hidden by
full-sample averages.


\subsubsection{Subperiod Analysis}

The sample may also be divided into economically meaningful subperiods:

\begin{equation}
\mathcal{T}
=
\mathcal{T}_1
\cup
\mathcal{T}_2
\cup
\cdots
\cup
\mathcal{T}_J.
\end{equation}

Diagnostics can then be computed independently within each subperiod.

The objective is not to require identical performance in every regime, but to
understand whether the signal's historical performance is broad and persistent
or concentrated in a particular environment.


\subsubsection{Final Diagnostic Summary}

At the end of the signal-diagnostic stage, the backtester should have answered
four distinct questions:

\begin{equation}
\boxed{
\begin{aligned}
\text{Predictiveness:}
&\quad
\text{Does the signal predict future cross-sectional returns?}
\\
\text{Economic Shape:}
&\quad
\text{Are returns monotonic across signal levels?}
\\
\text{Persistence:}
&\quad
\text{Is the signal sufficiently stable through time?}
\\
\text{Robustness:}
&\quad
\text{Does the relationship survive across periods and subuniverses?}
\end{aligned}
}
\end{equation}

These diagnostics characterize the signal itself, but they do not determine
whether the resulting investment strategy is attractive after portfolio
construction and implementation.

We now return to the main backtesting pipeline:

\begin{equation}
\boxed{
\mathbf{s}_t^{\mathrm{final}}
\quad
\longrightarrow
\quad
\text{Portfolio Construction}
\quad
\longrightarrow
\quad
\mathbf{w}_t^{\mathrm{target}}.
}
\end{equation}

The next step determines how the information contained in the signal should be
translated into actual capital allocations.

% ============================================================
\part{Portfolio Construction}
% ============================================================

% ============================================================
\section{From Signals to Portfolio Weights}
\label{sec:signals_to_weights}
% ============================================================

The previous parts of the backtester transformed raw information into a final
investable signal

\begin{equation}
\mathbf{s}_t^{\mathrm{final}}
=
\begin{bmatrix}
s_{1,t}^{\mathrm{final}}
\\
\vdots
\\
s_{N_t,t}^{\mathrm{final}}
\end{bmatrix}
\in
\mathbb{R}^{N_t}.
\end{equation}

The signal describes the relative attractiveness of securities, but it does not
yet determine how much capital should be allocated to each security.

Portfolio construction begins by defining a mapping from signals to portfolio
weights:

\begin{equation}
\boxed{
\mathbf{s}_t^{\mathrm{final}}
\quad
\longrightarrow
\quad
\mathbf{w}_t^{\mathrm{target}}.
}
\end{equation}

Let

\begin{equation}
\mathbf{w}_t^{\mathrm{target}}
=
\begin{bmatrix}
w_{1,t}^{\mathrm{target}}
\\
\vdots
\\
w_{N_t,t}^{\mathrm{target}}
\end{bmatrix}
\in
\mathbb{R}^{N_t}
\end{equation}

denote the target portfolio weights immediately after the rebalance at time
$t$.

A positive weight represents a long position,

\begin{equation}
w_{i,t}^{\mathrm{target}}>0,
\end{equation}

a negative weight represents a short position,

\begin{equation}
w_{i,t}^{\mathrm{target}}<0,
\end{equation}

and

\begin{equation}
w_{i,t}^{\mathrm{target}}=0
\end{equation}

means that no position is desired in security $i$.

At this stage, the objective is to study direct signal-to-weight mappings.
Portfolio-level risk optimization, factor constraints, leverage constraints,
turnover penalties, and implementation costs will be introduced subsequently.

Thus, the current section focuses on

\begin{equation}
\boxed{
\text{Signal Strength}
\rightarrow
\text{Desired Relative Position Size}.
}
\end{equation}


% ------------------------------------------------------------
\subsection{Signal-to-Weight Mapping Functions}
\label{subsec:signal_weight_mapping}
% ------------------------------------------------------------

Let

\begin{equation}
\mathcal{G}_t
\end{equation}

denote a signal-to-weight mapping.

In its most general form,

\begin{equation}
\boxed{
\mathbf{w}_t^{\mathrm{raw}}
=
\mathcal{G}_t
\left(
\mathbf{s}_t^{\mathrm{final}}
\right),
}
\end{equation}

where

\begin{equation}
\mathbf{w}_t^{\mathrm{raw}}
\in
\mathbb{R}^{N_t}
\end{equation}

denotes the portfolio weights implied directly by the signal before subsequent
portfolio-level constraints or risk adjustments.

The mapping may depend only on the ordering of the signal, or it may use its
numerical magnitude.


\subsubsection{Ordinal vs.\ Cardinal Mapping}

If the signal is ordinal, only relative ordering is economically meaningful.

For example,

\begin{equation}
s_{i,t}^{\mathrm{final}}
>
s_{j,t}^{\mathrm{final}}
\end{equation}

means that security $i$ is preferred to security $j$, but the difference

\begin{equation}
s_{i,t}^{\mathrm{final}}
-
s_{j,t}^{\mathrm{final}}
\end{equation}

need not have a meaningful economic interpretation.

Suitable mappings therefore include:

\begin{itemize}
    \item quantile portfolios;
    \item top/bottom $N$ portfolios;
    \item rank-based portfolios.
\end{itemize}

If the signal is cardinal, its magnitude contains information that may be used
directly for position sizing.

Examples include:

\begin{itemize}
    \item standardized alpha scores;
    \item calibrated expected returns;
    \item confidence-adjusted expected returns.
\end{itemize}

In this case, linear or nonlinear mappings may be appropriate.

Hence,

\begin{equation}
\boxed{
\text{Nature of Signal}
\rightarrow
\text{Choice of Mapping Function}.
}
\end{equation}


\subsubsection{Raw and Normalized Weights}

It is useful to distinguish the raw output of the mapping function from the
subsequently normalized portfolio.

Let

\begin{equation}
\widetilde{\mathbf{w}}_t
=
\mathcal{G}_t
\left(
\mathbf{s}_t^{\mathrm{final}}
\right)
\end{equation}

denote unnormalized desired positions.

A normalization operator $\mathcal{N}$ can then produce

\begin{equation}
\boxed{
\mathbf{w}_t^{\mathrm{raw}}
=
\mathcal{N}
\left(
\widetilde{\mathbf{w}}_t
\right).
}
\end{equation}

For example, a unit-gross normalization is

\begin{equation}
\boxed{
w_{i,t}^{\mathrm{raw}}
=
\frac{
\widetilde{w}_{i,t}
}{
\sum_{j=1}^{N_t}
\left|
\widetilde{w}_{j,t}
\right|
}.
}
\end{equation}

Then

\begin{equation}
\sum_{i=1}^{N_t}
\left|
w_{i,t}^{\mathrm{raw}}
\right|
=
1.
\end{equation}

The desired gross and net exposure conventions will be treated explicitly in a
later section.

% ------------------------------------------------------------
\subsection{Direct Pure Alpha Construction and Risk Scaling}
\label{subsec:direct_pure_alpha_construction}
% ------------------------------------------------------------

The signal-to-weight mappings introduced above can be used directly to construct
a long--short \emph{Pure Alpha} portfolio without requiring a numerical
portfolio optimizer.

This provides an important alternative to optimization-based portfolio
construction.

Let

\begin{equation}
\mathbf{s}_t^{\mathrm{final}}
\in
\mathbb{R}^{N_t}
\end{equation}

denote the final cross-sectional signal after the transformations and
neutralizations developed previously.

A mapping function

\begin{equation}
\mathcal{G}_t:
\mathbb{R}^{N_t}
\rightarrow
\mathbb{R}^{N_t}
\end{equation}

can transform this signal directly into Pure Alpha portfolio weights:

\begin{equation}
\boxed{
\mathbf{w}_t^{PA}
=
\mathcal{G}_t
\left(
\mathbf{s}_t^{\mathrm{final}}
\right),
}
\end{equation}

where

\begin{equation}
\mathbf{w}_t^{PA}
=
\begin{bmatrix}
w_{1,t}^{PA}\\
\vdots\\
w_{N_t,t}^{PA}
\end{bmatrix}
\in
\mathbb{R}^{N_t}.
\end{equation}

The superscript $PA$ denotes \emph{Pure Alpha}.

The mapping $\mathcal{G}_t(\cdot)$ may correspond, for example, to a linear
signal mapping, rank weighting, quantile portfolios, or another deterministic
mapping rule. These alternatives are developed in the following subsections.

If the signal and mapping procedure have already been designed to satisfy the
desired exposure properties, the resulting portfolio can be constructed
directly without solving an optimization problem.

Conceptually,

\begin{equation}
\boxed{
\mathbf{s}_t^{\mathrm{final}}
\rightarrow
\mathcal{G}_t
\rightarrow
\mathbf{w}_t^{PA}.
}
\end{equation}


\subsubsection{Direct Construction Without Numerical Optimization}

A particularly simple example is a linear mapping.

Suppose that the final signal satisfies

\begin{equation}
\mathbf{1}^{\top}
\mathbf{s}_t^{\mathrm{final}}
=
0.
\end{equation}

A gross-normalized Pure Alpha portfolio can be defined as

\begin{equation}
\boxed{
\mathbf{w}_t^{PA}
=
G^{PA}
\frac{
\mathbf{s}_t^{\mathrm{final}}
}{
\left\|
\mathbf{s}_t^{\mathrm{final}}
\right\|_1
},
}
\end{equation}

where

\begin{equation}
G^{PA}>0
\end{equation}

is the desired gross exposure.

By construction,

\begin{equation}
\left\|
\mathbf{w}_t^{PA}
\right\|_1
=
G^{PA}.
\end{equation}

Moreover, because the signal sums to zero,

\begin{equation}
\mathbf{1}^{\top}
\mathbf{w}_t^{PA}
=
0.
\end{equation}

The portfolio is therefore dollar neutral.

More generally, deterministic transformations can be used to impose the
desired portfolio structure analytically whenever closed-form mappings are
available.

This can be attractive in large-scale backtests because it avoids solving a
numerical optimization problem at every rebalance date.

Thus, portfolio construction does not necessarily imply portfolio
optimization:

\begin{equation}
\boxed{
\text{Signal}
\rightarrow
\text{Closed-Form Mapping}
\rightarrow
\text{Pure Alpha Portfolio}
}
\end{equation}

may itself constitute a complete portfolio-construction procedure.


\subsubsection{Alpha-Level Volatility Scaling}

The risk of the resulting Pure Alpha portfolio may vary substantially through
time even if its gross exposure is kept constant.

Let

\begin{equation}
\widehat{\boldsymbol{\Sigma}}_t
\in
\mathbb{R}^{N_t\times N_t}
\end{equation}

denote the point-in-time covariance estimate.

The ex-ante volatility of the Pure Alpha portfolio is

\begin{equation}
\boxed{
\widehat{\sigma}_{PA,t}
=
\sqrt{
\left(
\mathbf{w}_t^{PA}
\right)^{\top}
\widehat{\boldsymbol{\Sigma}}_t
\mathbf{w}_t^{PA}
}.
}
\end{equation}

Let

\begin{equation}
\sigma_{PA}^{*}>0
\end{equation}

denote the desired Pure Alpha volatility.

An optional isovol scaling coefficient is

\begin{equation}
\boxed{
\lambda_t^{PA}
=
\frac{
\sigma_{PA}^{*}
}{
\widehat{\sigma}_{PA,t}
}.
}
\end{equation}

The risk-scaled Pure Alpha portfolio is then

\begin{equation}
\boxed{
\mathbf{w}_t^{PA,\mathrm{iso}}
=
\lambda_t^{PA}
\mathbf{w}_t^{PA}.
}
\end{equation}

Because $\lambda_t^{PA}$ is a common scalar applied to all securities, this
operation changes the overall scale of the Pure Alpha portfolio without
changing its relative composition:

\begin{equation}
\frac{
w_{i,t}^{PA,\mathrm{iso}}
}{
w_{j,t}^{PA,\mathrm{iso}}
}
=
\frac{
w_{i,t}^{PA}
}{
w_{j,t}^{PA}
}.
\end{equation}

The direct Pure Alpha construction may therefore be summarized as

\begin{equation}
\boxed{
\mathbf{s}_t^{\mathrm{final}}
\rightarrow
\mathbf{w}_t^{PA}
\rightarrow
\mathbf{w}_t^{PA,\mathrm{iso}},
}
\end{equation}

where the final isovol step is optional.


\subsubsection{Preservation of Neutrality}

Common volatility scaling preserves homogeneous zero-exposure constraints.

Suppose the Pure Alpha portfolio satisfies

\begin{equation}
\mathbf{B}_t^\top
\mathbf{w}_t^{PA}
=
\mathbf{0},
\end{equation}

where

\begin{equation}
\mathbf{B}_t
\in
\mathbb{R}^{N_t\times K}
\end{equation}

contains the exposures against which neutrality is required.

Then

\begin{equation}
\begin{aligned}
\mathbf{B}_t^\top
\mathbf{w}_t^{PA,\mathrm{iso}}
&=
\mathbf{B}_t^\top
\left(
\lambda_t^{PA}
\mathbf{w}_t^{PA}
\right)
\\
&=
\lambda_t^{PA}
\mathbf{B}_t^\top
\mathbf{w}_t^{PA}
\\
&=
\mathbf{0}.
\end{aligned}
\end{equation}

Therefore,

\begin{equation}
\boxed{
\text{Isovol Scaling}
\quad\text{preserves exact zero-exposure constraints}.
}
\end{equation}

This includes, when already satisfied by the Pure Alpha weights, dollar
neutrality, beta neutrality, industry neutrality, country neutrality, and
style-factor neutrality.


\subsubsection{Pure Alpha as a Standalone Portfolio}

The risk-scaled Pure Alpha portfolio may itself be the final strategy.

In this case,

\begin{equation}
\boxed{
\mathbf{w}_t^{\mathrm{final}}
=
\mathbf{w}_t^{PA,\mathrm{iso}},
}
\end{equation}

and no optimization engine is required.

This architecture is particularly useful when the desired portfolio can be
constructed through simple and transparent closed-form rules:

\begin{equation}
\boxed{
\text{Final Signal}
\rightarrow
\text{Direct L/S Mapping}
\rightarrow
\text{Optional Isovol}
\rightarrow
\text{Final Pure Alpha Portfolio}.
}
\end{equation}


\subsubsection{Pure Alpha as an Input to a Subsequent Optimizer}

Alternatively, the Pure Alpha construction can represent an intermediate
research object rather than the final investable portfolio.

For example, the objective may be to translate a long--short Pure Alpha signal
into a constrained long-only portfolio.

Information contained in the Pure Alpha strategy can then be used to define an
expected-return vector

\begin{equation}
\widehat{\boldsymbol{\mu}}_t
\in
\mathbb{R}^{N_t}
\end{equation}

that subsequently enters the portfolio optimization engine.

Schematically,

\begin{equation}
\boxed{
\begin{aligned}
\mathbf{s}_t^{\mathrm{final}}
&\rightarrow
\mathbf{w}_t^{PA}
\\
&\rightarrow
\text{Pure Alpha Information}
\\
&\rightarrow
\widehat{\boldsymbol{\mu}}_t
\\
&\rightarrow
\text{Constrained Portfolio Optimization}
\\
&\rightarrow
\mathbf{w}_t^{\mathrm{target}}.
\end{aligned}
}
\end{equation}

For example, a simple calibration may take the form

\begin{equation}
\boxed{
\widehat{\boldsymbol{\mu}}_t
=
\kappa_t
\mathbf{w}_t^{PA},
}
\end{equation}

where $\kappa_t$ converts the relative Pure Alpha direction into expected-return
units.

However, an important distinction must be maintained between the
\emph{cross-sectional alpha information} and the \emph{risk scaling} applied to
the Pure Alpha portfolio.

Indeed,

\begin{equation}
\mathbf{w}_t^{PA,\mathrm{iso}}
=
\lambda_t^{PA}
\mathbf{w}_t^{PA},
\end{equation}

where $\lambda_t^{PA}$ depends on forecast portfolio volatility.

The fact that the Pure Alpha portfolio receives more leverage when forecast
volatility is low does not, by itself, imply that individual expected returns
are proportionally larger.

Therefore, when the Pure Alpha object is used to construct expected returns for
a subsequent optimizer, the alpha direction and its portfolio-level risk scale
should be conceptually distinguished.


\subsubsection{Two Different Roles of Isovol Scaling}

It is useful to distinguish the alpha-level isovol operation introduced here
from the final portfolio-level volatility scaling developed later.

The first operates on the directly constructed Pure Alpha portfolio:

\begin{equation}
\boxed{
\mathbf{w}_t^{PA}
\rightarrow
\mathbf{w}_t^{PA,\mathrm{iso}}.
}
\end{equation}

Its purpose is to normalize the risk of the Pure Alpha strategy itself.

The second operates on the target portfolio produced after the complete
portfolio-construction process:

\begin{equation}
\boxed{
\mathbf{w}_t^{\mathrm{target}}
\rightarrow
\mathbf{w}_t^{\mathrm{final}}.
}
\end{equation}

Its purpose is to control the absolute risk of the final investable portfolio.

These two operations need not both be performed.

If the directly constructed Pure Alpha portfolio is itself the final portfolio,
alpha-level isovol may already provide the desired risk scaling.

If a subsequent optimizer materially changes the portfolio composition,
however, the volatility of the optimized portfolio generally differs from that
of the original Pure Alpha portfolio. A new portfolio-level risk forecast may
therefore be required.

Thus,

\begin{equation}
\boxed{
\begin{aligned}
\text{Alpha-Level Isovol}
&:
&
\mathbf{w}_t^{PA}
&\rightarrow
\mathbf{w}_t^{PA,\mathrm{iso}},
\\[3pt]
\text{Final Portfolio Isovol}
&:
&
\mathbf{w}_t^{\mathrm{target}}
&\rightarrow
\mathbf{w}_t^{\mathrm{final}}.
\end{aligned}
}
\end{equation}

The appropriate architecture depends on whether the Pure Alpha portfolio is
used directly as the investment strategy or as an intermediate input to a
subsequent portfolio-construction problem.


% ------------------------------------------------------------
\subsection{Equal-Weighted Quantile Portfolios}
\label{subsec:equal_weighted_quantile_portfolios}
% ------------------------------------------------------------

One of the simplest mappings is to divide the investment universe into
quantiles according to the final signal.

Let

\begin{equation}
\mathcal{Q}_{q,t}
\subseteq
\mathcal{U}_t,
\qquad
q=1,\ldots,Q,
\end{equation}

denote quantile $q$, where $\mathcal{Q}_{1,t}$ contains the lowest-signal
securities and $\mathcal{Q}_{Q,t}$ contains the highest-signal securities.

For a long-only top-quantile portfolio,

\begin{equation}
\boxed{
w_{i,t}
=
\begin{cases}
\dfrac{1}{|\mathcal{Q}_{Q,t}|},
&
i\in\mathcal{Q}_{Q,t},
\\[0.25cm]
0,
&
\text{otherwise}.
\end{cases}
}
\end{equation}


\subsubsection{Long--Short Quantile Portfolio}

A simple dollar-neutral long--short portfolio can buy the highest quantile and
short the lowest quantile.

If each side is allocated gross exposure $G_t/2$, then

\begin{equation}
\boxed{
w_{i,t}
=
\begin{cases}
\dfrac{G_t/2}{|\mathcal{Q}_{Q,t}|},
&
i\in\mathcal{Q}_{Q,t},
\\[0.30cm]
-\dfrac{G_t/2}{|\mathcal{Q}_{1,t}|},
&
i\in\mathcal{Q}_{1,t},
\\[0.30cm]
0,
&
\text{otherwise}.
\end{cases}
}
\end{equation}

The long-leg gross exposure is

\begin{equation}
\sum_{i:w_{i,t}>0}w_{i,t}
=
\frac{G_t}{2},
\end{equation}

while the absolute short-leg exposure is

\begin{equation}
\sum_{i:w_{i,t}<0}|w_{i,t}|
=
\frac{G_t}{2}.
\end{equation}

Therefore,

\begin{equation}
\sum_i |w_{i,t}|=G_t
\end{equation}

and

\begin{equation}
\sum_i w_{i,t}=0.
\end{equation}

For example, if

\begin{equation}
G_t=2,
\end{equation}

the portfolio is $100\%$ long and $100\%$ short.


\subsubsection{Properties}

Equal-weighted quantile portfolios are:

\begin{itemize}
    \item simple;
    \item transparent;
    \item relatively insensitive to extreme signal magnitudes;
    \item appropriate for ordinal signals.
\end{itemize}

However, all selected securities receive identical weights regardless of the
strength of their signals.


% ------------------------------------------------------------
\subsection{Top/Bottom $N$ Portfolios}
\label{subsec:top_bottom_n}
% ------------------------------------------------------------

Instead of selecting a fixed fraction of the universe, the strategy may select
a fixed number of securities.

Let

\begin{equation}
\mathcal{L}_t^{N}
\end{equation}

denote the set containing the $N_L$ highest-signal securities and

\begin{equation}
\mathcal{S}_t^{N}
\end{equation}

the set containing the $N_S$ lowest-signal securities.

An equal-weighted long--short specification is

\begin{equation}
\boxed{
w_{i,t}
=
\begin{cases}
\dfrac{G_t/2}{N_L},
&
i\in\mathcal{L}_t^{N},
\\[0.30cm]
-\dfrac{G_t/2}{N_S},
&
i\in\mathcal{S}_t^{N},
\\[0.30cm]
0,
&
\text{otherwise}.
\end{cases}
}
\end{equation}

Unlike quantile selection, the number of positions remains constant when the
size of the investment universe changes.


\subsubsection{Concentration Trade-Off}

A smaller $N$ produces a more concentrated portfolio:

\begin{equation}
N\downarrow
\quad\Rightarrow\quad
\text{Higher Concentration}.
\end{equation}

This may increase exposure to the strongest signals but also increases
idiosyncratic risk and sensitivity to individual-security errors.

Conversely,

\begin{equation}
N\uparrow
\quad\Rightarrow\quad
\text{Greater Diversification},
\end{equation}

but potentially weaker average signal strength among selected securities.


% ------------------------------------------------------------
\subsection{Rank-Weighted Portfolios}
\label{subsec:rank_weighted_portfolios}
% ------------------------------------------------------------

Equal-weighted portfolios use the signal to select securities but ignore
differences in relative rank once selection has occurred.

Rank-weighted portfolios preserve more information about the ordering of
securities without relying directly on the magnitude of the original signal.

Let

\begin{equation}
r_{i,t}
\in
\{1,\ldots,N_t\}
\end{equation}

denote the cross-sectional rank of security $i$, with larger ranks corresponding
to more attractive securities.

A centered rank score can be defined as

\begin{equation}
\boxed{
\widetilde{r}_{i,t}
=
r_{i,t}
-
\frac{N_t+1}{2}.
}
\end{equation}

Then

\begin{equation}
\sum_{i=1}^{N_t}
\widetilde{r}_{i,t}
=
0.
\end{equation}

A rank-weighted portfolio can be constructed as

\begin{equation}
\boxed{
w_{i,t}
=
G_t
\frac{
\widetilde{r}_{i,t}
}{
\sum_{j=1}^{N_t}
|\widetilde{r}_{j,t}|
}.
}
\end{equation}

Therefore,

\begin{equation}
\sum_i |w_{i,t}|=G_t,
\end{equation}

while

\begin{equation}
\sum_i w_{i,t}=0.
\end{equation}


\subsubsection{Percentile Representation}

Alternatively, define the percentile rank

\begin{equation}
p_{i,t}
=
\frac{
r_{i,t}-1
}{
N_t-1
}
\in[0,1].
\end{equation}

A centered percentile score is

\begin{equation}
\boxed{
q_{i,t}
=
p_{i,t}
-
\frac{1}{2}.
}
\end{equation}

Weights may then be defined proportionally to $q_{i,t}$.

This representation is convenient because its scale does not depend directly
on the number of securities in the universe.


% ------------------------------------------------------------
\subsection{Linear Signal Mapping}
\label{subsec:linear_signal_mapping}
% ------------------------------------------------------------

If signal magnitude is meaningful, the simplest cardinal mapping assigns
positions directly proportional to the signal:

\begin{equation}
\boxed{
\widetilde{w}_{i,t}
=
\kappa_t
s_{i,t}^{\mathrm{final}},
}
\end{equation}

where $\kappa_t>0$ is a scaling parameter.

If only relative weights matter before normalization, $\kappa_t$ can be omitted:

\begin{equation}
\widetilde{w}_{i,t}
=
s_{i,t}^{\mathrm{final}}.
\end{equation}

Under unit-gross normalization,

\begin{equation}
\boxed{
w_{i,t}
=
\frac{
s_{i,t}^{\mathrm{final}}
}{
\sum_{j=1}^{N_t}
|s_{j,t}^{\mathrm{final}}|
}.
}
\end{equation}

More generally, for target gross exposure $G_t$,

\begin{equation}
\boxed{
w_{i,t}
=
G_t
\frac{
s_{i,t}^{\mathrm{final}}
}{
\sum_{j=1}^{N_t}
|s_{j,t}^{\mathrm{final}}|
}.
}
\end{equation}


\subsubsection{Net Exposure}

Notice that gross normalization alone does not guarantee

\begin{equation}
\sum_i w_{i,t}=0.
\end{equation}

Indeed,

\begin{equation}
\sum_i w_{i,t}
=
G_t
\frac{
\sum_i s_{i,t}^{\mathrm{final}}
}{
\sum_i |s_{i,t}^{\mathrm{final}}|
}.
\end{equation}

Thus, a zero-net portfolio is obtained automatically only if

\begin{equation}
\boxed{
\sum_i s_{i,t}^{\mathrm{final}}=0.
}
\end{equation}

Otherwise, the signal must be centered or the resulting weights must be
adjusted if dollar neutrality is required.


\subsubsection{Interpretation}

Linear mapping preserves relative differences in signal magnitude:

\begin{equation}
\frac{
w_{i,t}
}{
w_{j,t}
}
=
\frac{
s_{i,t}^{\mathrm{final}}
}{
s_{j,t}^{\mathrm{final}}
}
\end{equation}

for two securities with non-zero signals before additional constraints.

This is appropriate only when such magnitude differences have a meaningful
interpretation.


% ------------------------------------------------------------
\subsection{Nonlinear Signal Mapping}
\label{subsec:nonlinear_signal_mapping}
% ------------------------------------------------------------

A nonlinear mapping applies a function

\begin{equation}
g:\mathbb{R}\rightarrow\mathbb{R}
\end{equation}

to each signal:

\begin{equation}
\boxed{
\widetilde{w}_{i,t}
=
g
\left(
s_{i,t}^{\mathrm{final}}
\right).
}
\end{equation}

Normalized weights are then

\begin{equation}
\boxed{
w_{i,t}
=
G_t
\frac{
g(s_{i,t}^{\mathrm{final}})
}{
\sum_{j=1}^{N_t}
|g(s_{j,t}^{\mathrm{final}})|
}.
}
\end{equation}


\subsubsection{Power Mapping}

One example is the signed power transformation

\begin{equation}
\boxed{
g(s)
=
\operatorname{sign}(s)
|s|^{p},
}
\end{equation}

where

\begin{equation}
p>0.
\end{equation}

If

\begin{equation}
p>1,
\end{equation}

the mapping increases the relative importance of extreme signals.

If

\begin{equation}
0<p<1,
\end{equation}

the mapping compresses differences between large and small signals.

The linear mapping is recovered when

\begin{equation}
p=1.
\end{equation}


\subsubsection{Bounded Mapping}

A bounded function may be used to prevent extreme signals from generating
disproportionately large raw positions.

For example,

\begin{equation}
\boxed{
g(s)
=
\tanh(\kappa s),
}
\end{equation}

where $\kappa>0$ controls the steepness of the transformation.

For small values of $s$,

\begin{equation}
\tanh(\kappa s)
\approx
\kappa s,
\end{equation}

so the mapping is approximately linear.

For large positive or negative signals,

\begin{equation}
\tanh(\kappa s)
\rightarrow
\pm1.
\end{equation}

Thus, the marginal effect of increasingly extreme signals decreases.


\subsubsection{Exponential Mapping}

For a long-only strategy, an exponential transformation may be used:

\begin{equation}
\widetilde{w}_{i,t}
=
\exp
\left(
\kappa s_{i,t}^{\mathrm{final}}
\right).
\end{equation}

Normalized weights become

\begin{equation}
\boxed{
w_{i,t}
=
\frac{
\exp(\kappa s_{i,t}^{\mathrm{final}})
}{
\sum_{j=1}^{N_t}
\exp(\kappa s_{j,t}^{\mathrm{final}})
}.
}
\end{equation}

This is equivalent to a softmax mapping.

The parameter $\kappa$ controls concentration:

\begin{equation}
\kappa\rightarrow0
\quad\Rightarrow\quad
w_{i,t}\rightarrow\frac{1}{N_t},
\end{equation}

whereas larger values of $\kappa$ increasingly concentrate the portfolio in
high-signal securities.


% ------------------------------------------------------------
\subsection{Threshold-Based Mapping}
\label{subsec:threshold_based_mapping}
% ------------------------------------------------------------

A strategy may choose to hold positions only when the signal exceeds a minimum
conviction threshold.

Let

\begin{equation}
\tau_t>0
\end{equation}

denote the threshold.

A simple symmetric rule is

\begin{equation}
\boxed{
\widetilde{w}_{i,t}
=
\begin{cases}
s_{i,t}^{\mathrm{final}},
&
|s_{i,t}^{\mathrm{final}}|>\tau_t,
\\[0.20cm]
0,
&
|s_{i,t}^{\mathrm{final}}|\leq\tau_t.
\end{cases}
}
\end{equation}

The surviving positions can subsequently be normalized to the desired gross
exposure.


\subsubsection{Binary Threshold Mapping}

A more aggressive simplification ignores signal magnitude once the threshold is
crossed:

\begin{equation}
\boxed{
\widetilde{w}_{i,t}
=
\begin{cases}
+1,
&
s_{i,t}^{\mathrm{final}}>\tau_t,
\\
0,
&
-\tau_t
\leq
s_{i,t}^{\mathrm{final}}
\leq
\tau_t,
\\
-1,
&
s_{i,t}^{\mathrm{final}}<-\tau_t.
\end{cases}
}
\end{equation}

This creates three regions:

\begin{equation}
\boxed{
\text{Short}
\qquad
\text{No Position}
\qquad
\text{Long}.
}
\end{equation}


\subsubsection{Soft Thresholding}

Instead of introducing a discontinuity at $\tau_t$, one may use

\begin{equation}
\boxed{
g(s)
=
\operatorname{sign}(s)
\max
\left(
|s|-\tau_t,
0
\right).
}
\end{equation}

Thus,

\begin{equation}
\widetilde{w}_{i,t}
=
g
\left(
s_{i,t}^{\mathrm{final}}
\right).
\end{equation}

Small signals are mapped to zero, while larger signals increase continuously
beyond the threshold.


\subsubsection{Economic Motivation}

Thresholding can be useful when small signals are believed to contain little
economic information.

It may also reduce the number of positions and avoid trading on weak changes in
estimated attractiveness.

However, thresholding introduces an additional parameter and can create
instability around the threshold.

A security moving from

\begin{equation}
s_{i,t}
=
\tau_t-\epsilon
\end{equation}

to

\begin{equation}
s_{i,t+1}
=
\tau_t+\epsilon
\end{equation}

may suddenly enter the portfolio despite an economically negligible change in
the underlying signal.

Such discontinuities can contribute to turnover.


% ------------------------------------------------------------
\subsection{Conviction-Weighted Portfolios}
\label{subsec:conviction_weighted_portfolios}
% ------------------------------------------------------------

A conviction-weighted portfolio explicitly allocates more capital to securities
for which the investment signal is stronger.

In its simplest form,

\begin{equation}
\boxed{
|\widetilde{w}_{i,t}|
=
g
\left(
|s_{i,t}^{\mathrm{final}}|
\right),
}
\end{equation}

where $g(\cdot)$ is increasing.

The sign of the position is determined by the sign of the signal:

\begin{equation}
\boxed{
\widetilde{w}_{i,t}
=
\operatorname{sign}
\left(
s_{i,t}^{\mathrm{final}}
\right)
g
\left(
|s_{i,t}^{\mathrm{final}}|
\right).
}
\end{equation}


\subsubsection{Conviction Within Selected Tails}

Conviction weighting can also be combined with quantile or top/bottom
selection.

Suppose

\begin{equation}
\mathcal{L}_t
\end{equation}

and

\begin{equation}
\mathcal{S}_t
\end{equation}

denote the selected long and short sets.

For the long leg,

\begin{equation}
\boxed{
w_{i,t}^{L}
=
\frac{G_t^{L}
g(s_{i,t}^{\mathrm{final}})}
{
\sum_{j\in\mathcal{L}_t}
g(s_{j,t}^{\mathrm{final}})
},
\qquad
i\in\mathcal{L}_t,
}
\end{equation}

where

\begin{equation}
G_t^{L}>0
\end{equation}

is the desired long gross exposure.

For the short leg, using a positive conviction function of signal magnitude,

\begin{equation}
\boxed{
w_{i,t}^{S}
=
-
\frac{
G_t^{S}
g(|s_{i,t}^{\mathrm{final}}|)
}{
\sum_{j\in\mathcal{S}_t}
g(|s_{j,t}^{\mathrm{final}}|)
},
\qquad
i\in\mathcal{S}_t,
}
\end{equation}

where

\begin{equation}
G_t^{S}>0
\end{equation}

denotes the absolute short gross exposure.

Then

\begin{equation}
\sum_{i\in\mathcal{L}_t}
w_{i,t}^{L}
=
G_t^{L},
\end{equation}

and

\begin{equation}
\sum_{i\in\mathcal{S}_t}
|w_{i,t}^{S}|
=
G_t^{S}.
\end{equation}


\subsubsection{Signal Strength vs.\ Forecast Confidence}

Conviction should not necessarily be identified with signal magnitude alone.

A large signal may have low estimation confidence.

Suppose

\begin{equation}
c_{i,t}\in[0,1]
\end{equation}

denotes a confidence measure.

One may define effective conviction as

\begin{equation}
\boxed{
\chi_{i,t}
=
c_{i,t}
\left|
s_{i,t}^{\mathrm{final}}
\right|.
}
\end{equation}

Then

\begin{equation}
\widetilde{w}_{i,t}
=
\operatorname{sign}
\left(
s_{i,t}^{\mathrm{final}}
\right)
g(\chi_{i,t}).
\end{equation}

This distinguishes

\begin{equation}
\boxed{
\text{Strength of Investment View}
}
\end{equation}

from

\begin{equation}
\boxed{
\text{Confidence in the Investment View}.
}
\end{equation}


% ------------------------------------------------------------
\subsubsection{Comparison of Direct Mapping Approaches}
% ------------------------------------------------------------

The main direct signal-to-weight mappings can be summarized as

\begin{equation}
\boxed{
\begin{array}{lll}
\text{Quantile Mapping}
&
\rightarrow
&
\text{Group membership determines position},
\\[0.15cm]
\text{Top/Bottom }N
&
\rightarrow
&
\text{Fixed number of extreme securities},
\\[0.15cm]
\text{Rank Mapping}
&
\rightarrow
&
\text{Relative ordering determines position size},
\\[0.15cm]
\text{Linear Mapping}
&
\rightarrow
&
w_i\propto s_i,
\\[0.15cm]
\text{Nonlinear Mapping}
&
\rightarrow
&
w_i\propto g(s_i),
\\[0.15cm]
\text{Threshold Mapping}
&
\rightarrow
&
\text{Weak signals receive zero weight},
\\[0.15cm]
\text{Conviction Mapping}
&
\rightarrow
&
\text{Position size increases with conviction}.
\end{array}
}
\end{equation}

There is no universally optimal mapping.

The appropriate choice depends on the economic meaning of the signal, its
predictive relationship with future returns, its stability, and the intended
portfolio architecture.


% ------------------------------------------------------------
\subsubsection{From Raw Mapped Weights to Portfolio Construction}
% ------------------------------------------------------------

The mappings developed in this section produce a first portfolio representation

\begin{equation}
\boxed{
\mathbf{s}_t^{\mathrm{final}}
\quad
\overset{\mathcal{G}_t}{\longrightarrow}
\quad
\mathbf{w}_t^{\mathrm{raw}}.
}
\end{equation}

However,

\begin{equation}
\mathbf{w}_t^{\mathrm{raw}}
\end{equation}

should not necessarily be interpreted as the final implemented portfolio.

It may still need to satisfy requirements concerning:

\begin{itemize}
    \item long and short gross exposure;
    \item net exposure;
    \item leverage;
    \item individual position limits;
    \item sector, country, beta, and factor exposures;
    \item portfolio volatility;
    \item diversification;
    \item liquidity;
    \item turnover;
    \item transaction costs.
\end{itemize}

The portfolio-construction problem therefore continues from

\begin{equation}
\boxed{
\mathbf{s}_t^{\mathrm{final}}
\rightarrow
\mathbf{w}_t^{\mathrm{raw}}
\rightarrow
\mathbf{w}_t^{\mathrm{target}}
\rightarrow
\mathbf{w}_t^{\mathrm{implemented}}.
}
\end{equation}

The next sections determine how the raw signal-implied positions should be
scaled, constrained, and optimized before they become actual portfolio
holdings.


% ============================================================
\section{Portfolio Risk Model}
\label{sec:portfolio_risk_model}
% ============================================================

The signal-construction process determines which securities are attractive, but
does not describe the risk created by holding them.

Two securities with identical expected attractiveness may have very different
volatilities, and two individually risky securities may together form a
relatively low-risk portfolio if their returns are weakly or negatively
correlated.

The portfolio-construction engine therefore requires a model of the joint
distribution of asset returns.

Let

\begin{equation}
\mathbf{R}_{t+1}
=
\begin{bmatrix}
R_{1,t+1}\\
\vdots\\
R_{N_t,t+1}
\end{bmatrix}
\in
\mathbb{R}^{N_t}
\end{equation}

denote the vector of future security returns.

The central object of the portfolio risk model is the conditional covariance
matrix

\begin{equation}
\boxed{
\boldsymbol{\Sigma}_t
=
\operatorname{Cov}
\left(
\mathbf{R}_{t+1}
\mid
\mathcal{I}_t
\right)
\in
\mathbb{R}^{N_t\times N_t}.
}
\end{equation}

Its elements are

\begin{equation}
\Sigma_{ij,t}
=
\operatorname{Cov}_t
\left(
R_{i,t+1},
R_{j,t+1}
\right).
\end{equation}

The diagonal elements contain individual asset variances:

\begin{equation}
\Sigma_{ii,t}
=
\sigma_{i,t}^{2},
\end{equation}

while the off-diagonal elements contain pairwise covariances.

The risk model therefore provides the mapping

\begin{equation}
\boxed{
\text{Security Positions}
+
\text{Joint Return Distribution}
\rightarrow
\text{Portfolio Risk}.
}
\end{equation}

At this stage of the backtester, the two principal inputs to portfolio
construction are therefore

\begin{equation}
\boxed{
\begin{aligned}
\mathbf{s}_t^{\mathrm{final}}
&\in
\mathbb{R}^{N_t}
&&
\text{: investment attractiveness},
\\
\boldsymbol{\Sigma}_t
&\in
\mathbb{R}^{N_t\times N_t}
&&
\text{: estimated joint risk}.
\end{aligned}
}
\end{equation}

Additional exposure matrices, liquidity variables, and portfolio constraints
will subsequently complement these two objects.


% ------------------------------------------------------------
\subsection{Portfolio Variance and Covariance}
\label{subsec:portfolio_variance_covariance}
% ------------------------------------------------------------

Let

\begin{equation}
\mathbf{w}_t
=
\begin{bmatrix}
w_{1,t}\\
\vdots\\
w_{N_t,t}
\end{bmatrix}
\in
\mathbb{R}^{N_t}
\end{equation}

denote portfolio weights.

The future portfolio return is

\begin{equation}
\boxed{
R_{p,t+1}
=
\mathbf{w}_t^{\top}
\mathbf{R}_{t+1}.
}
\end{equation}

Conditional on the information available at time $t$, portfolio variance is

\begin{equation}
\operatorname{Var}_t
\left(
R_{p,t+1}
\right)
=
\operatorname{Var}_t
\left(
\mathbf{w}_t^\top
\mathbf{R}_{t+1}
\right).
\end{equation}

Using the covariance matrix,

\begin{equation}
\boxed{
\sigma_{p,t}^{2}
=
\mathbf{w}_t^\top
\boldsymbol{\Sigma}_t
\mathbf{w}_t.
}
\end{equation}

Hence portfolio volatility is

\begin{equation}
\boxed{
\sigma_{p,t}
=
\sqrt{
\mathbf{w}_t^\top
\boldsymbol{\Sigma}_t
\mathbf{w}_t
}.
}
\end{equation}


\subsubsection{Expanded Representation}

Portfolio variance can also be written as

\begin{equation}
\sigma_{p,t}^{2}
=
\sum_{i=1}^{N_t}
w_{i,t}^{2}
\sigma_{i,t}^{2}
+
\sum_{\substack{i,j=1\\i\neq j}}^{N_t}
w_{i,t}w_{j,t}
\sigma_{ij,t}.
\end{equation}

Equivalently,

\begin{equation}
\boxed{
\sigma_{p,t}^{2}
=
\sum_{i=1}^{N_t}
w_{i,t}^{2}
\sigma_{i,t}^{2}
+
2
\sum_{i<j}
w_{i,t}w_{j,t}
\sigma_{ij,t}.
}
\end{equation}

Portfolio risk therefore depends not only on individual-security volatility but
also on the dependence structure between securities.


\subsubsection{Covariance and Correlation}

For securities $i$ and $j$,

\begin{equation}
\boxed{
\sigma_{ij,t}
=
\rho_{ij,t}
\sigma_{i,t}
\sigma_{j,t},
}
\end{equation}

where

\begin{equation}
\rho_{ij,t}
\in[-1,1]
\end{equation}

is their correlation.

Define the diagonal volatility matrix

\begin{equation}
\mathbf{D}_{\sigma,t}
=
\operatorname{diag}
\left(
\sigma_{1,t},
\ldots,
\sigma_{N_t,t}
\right).
\end{equation}

If

\begin{equation}
\mathbf{C}_t
\in
\mathbb{R}^{N_t\times N_t}
\end{equation}

denotes the correlation matrix, then

\begin{equation}
\boxed{
\boldsymbol{\Sigma}_t
=
\mathbf{D}_{\sigma,t}
\mathbf{C}_t
\mathbf{D}_{\sigma,t}.
}
\end{equation}

This decomposition separates the estimation of marginal volatility from the
estimation of cross-security dependence.


\subsubsection{Properties of a Valid Covariance Matrix}

A valid covariance matrix must satisfy

\begin{equation}
\boldsymbol{\Sigma}_t
=
\boldsymbol{\Sigma}_t^\top
\end{equation}

and must be positive semidefinite:

\begin{equation}
\boxed{
\mathbf{x}^\top
\boldsymbol{\Sigma}_t
\mathbf{x}
\geq
0
\qquad
\forall
\mathbf{x}\in\mathbb{R}^{N_t}.
}
\end{equation}

This property guarantees that estimated portfolio variance cannot be negative.

For some optimization problems, positive definiteness,

\begin{equation}
\mathbf{x}^\top
\boldsymbol{\Sigma}_t
\mathbf{x}
>
0
\qquad
\forall\mathbf{x}\neq0,
\end{equation}

is particularly useful because it implies invertibility.


% ------------------------------------------------------------
\subsection{Historical Covariance Estimation}
\label{subsec:historical_covariance}
% ------------------------------------------------------------

The simplest covariance estimator uses historical returns.

Suppose a lookback window contains $L$ observations:

\begin{equation}
\mathbf{R}_{t-L+1},
\ldots,
\mathbf{R}_{t}.
\end{equation}

Define the historical mean-return vector

\begin{equation}
\overline{\mathbf{R}}_t
=
\frac{1}{L}
\sum_{\ell=0}^{L-1}
\mathbf{R}_{t-\ell}.
\end{equation}

The sample covariance matrix is

\begin{equation}
\boxed{
\widehat{\boldsymbol{\Sigma}}_t^{\mathrm{sample}}
=
\frac{1}{L-1}
\sum_{\ell=0}^{L-1}
\left(
\mathbf{R}_{t-\ell}
-
\overline{\mathbf{R}}_t
\right)
\left(
\mathbf{R}_{t-\ell}
-
\overline{\mathbf{R}}_t
\right)^\top.
}
\end{equation}

Its dimension is

\begin{equation}
\widehat{\boldsymbol{\Sigma}}_t^{\mathrm{sample}}
\in
\mathbb{R}^{N_t\times N_t}.
\end{equation}


\subsubsection{Point-in-Time Requirement}

The covariance estimate used for portfolio construction at time $t$ must depend
only on returns already observed at the decision time.

If the portfolio decision occurs before the return for date $t$ is known, the
estimation window must instead terminate at $t-1$.

Thus, schematically,

\begin{equation}
\boxed{
\widehat{\boldsymbol{\Sigma}}_t
=
g
\left(
\mathbf{R}_{t-1},
\mathbf{R}_{t-2},
\ldots
\right).
}
\end{equation}


\subsubsection{Estimation Error}

The covariance matrix contains

\begin{equation}
\frac{
N_t(N_t+1)
}{
2
}
\end{equation}

distinct parameters.

Therefore, estimation becomes difficult when the number of securities is large
relative to the number of historical observations.

For example,

\begin{equation}
N_t=500
\end{equation}

requires estimating

\begin{equation}
\frac{500\times501}{2}
=
125{,}250
\end{equation}

distinct covariance elements.

This is one of the principal reasons why naive historical covariance matrices
can be unstable in large equity universes.


\subsubsection{Singularity}

Let

\begin{equation}
\mathbf{X}
\in
\mathbb{R}^{L\times N_t}
\end{equation}

denote the centered historical return matrix.

The sample covariance is proportional to

\begin{equation}
\mathbf{X}^\top\mathbf{X}.
\end{equation}

Its rank satisfies

\begin{equation}
\operatorname{rank}
\left(
\widehat{\boldsymbol{\Sigma}}
\right)
\leq
\min
\left(
N_t,L-1
\right).
\end{equation}

Consequently, if

\begin{equation}
N_t>L-1,
\end{equation}

then

\begin{equation}
\boxed{
\widehat{\boldsymbol{\Sigma}}
\text{ is necessarily singular}.
}
\end{equation}

This is a major practical limitation for portfolio optimization.


\subsubsection{Rolling vs.\ Expanding Estimation}

A rolling estimator uses only the most recent $L$ observations:

\begin{equation}
\widehat{\boldsymbol{\Sigma}}_t
=
g
\left(
\mathbf{R}_{t-L+1},
\ldots,
\mathbf{R}_{t}
\right).
\end{equation}

An expanding estimator uses the complete history up to $t$.

The same trade-off encountered in predictive modeling applies:

\begin{equation}
\boxed{
\begin{aligned}
\text{Longer History}
&\rightarrow
\text{Lower Sampling Noise, Slower Adaptation},
\\
\text{Shorter History}
&\rightarrow
\text{Higher Sampling Noise, Faster Adaptation}.
\end{aligned}
}
\end{equation}


\subsubsection{Exponentially Weighted Covariance}

Recent observations may be assigned larger weights.

Let

\begin{equation}
0<\lambda<1
\end{equation}

denote the decay parameter.

Normalized historical weights may be defined as

\begin{equation}
\omega_\ell
=
\frac{
(1-\lambda)\lambda^\ell
}{
1-\lambda^L
},
\qquad
\ell=0,\ldots,L-1.
\end{equation}

Then

\begin{equation}
\sum_{\ell=0}^{L-1}
\omega_\ell
=
1.
\end{equation}

The exponentially weighted covariance estimator is

\begin{equation}
\boxed{
\widehat{\boldsymbol{\Sigma}}_t^{EW}
=
\sum_{\ell=0}^{L-1}
\omega_\ell
\left(
\mathbf{R}_{t-\ell}
-
\boldsymbol{\mu}_{t}^{EW}
\right)
\left(
\mathbf{R}_{t-\ell}
-
\boldsymbol{\mu}_{t}^{EW}
\right)^\top,
}
\end{equation}

where

\begin{equation}
\boldsymbol{\mu}_{t}^{EW}
=
\sum_{\ell=0}^{L-1}
\omega_\ell
\mathbf{R}_{t-\ell}.
\end{equation}

This allows the covariance estimate to react more quickly to changing market
conditions.


% ------------------------------------------------------------
\subsection{Shrinkage Estimators}
\label{subsec:shrinkage_covariance}
% ------------------------------------------------------------

Sample covariance matrices can be noisy, particularly when

\begin{equation}
N_t
\end{equation}

is large relative to

\begin{equation}
L.
\end{equation}

Shrinkage reduces estimation error by combining the sample covariance matrix
with a more structured target matrix.

Let

\begin{equation}
\mathbf{S}_t
=
\widehat{\boldsymbol{\Sigma}}_t^{\mathrm{sample}}
\end{equation}

denote the sample covariance matrix and let

\begin{equation}
\mathbf{F}_t
\end{equation}

denote a structured target.

A general shrinkage estimator is

\begin{equation}
\boxed{
\widehat{\boldsymbol{\Sigma}}_t^{\mathrm{shrunk}}
=
(1-\delta_t)
\mathbf{S}_t
+
\delta_t
\mathbf{F}_t,
}
\end{equation}

where

\begin{equation}
0\leq\delta_t\leq1.
\end{equation}


\subsubsection{Interpretation of the Shrinkage Intensity}

If

\begin{equation}
\delta_t=0,
\end{equation}

then

\begin{equation}
\widehat{\boldsymbol{\Sigma}}_t^{\mathrm{shrunk}}
=
\mathbf{S}_t.
\end{equation}

If

\begin{equation}
\delta_t=1,
\end{equation}

then

\begin{equation}
\widehat{\boldsymbol{\Sigma}}_t^{\mathrm{shrunk}}
=
\mathbf{F}_t.
\end{equation}

Intermediate values trade off sample-specific information against structural
stability.


\subsubsection{Possible Shrinkage Targets}

Possible targets include:

\begin{itemize}
    \item diagonal covariance matrices;
    \item constant-correlation models;
    \item single-factor market models;
    \item multi-factor risk models.
\end{itemize}

For example, a diagonal target is

\begin{equation}
\boxed{
\mathbf{F}_t
=
\operatorname{diag}
\left(
\widehat{\sigma}_{1,t}^{2},
\ldots,
\widehat{\sigma}_{N_t,t}^{2}
\right).
}
\end{equation}

This retains estimated individual variances while shrinking all pairwise
covariances toward zero.


\subsubsection{Why Shrinkage Can Improve Portfolio Construction}

Portfolio optimization frequently depends on

\begin{equation}
\widehat{\boldsymbol{\Sigma}}_t^{-1}.
\end{equation}

Small estimation errors in covariance eigenvalues can become greatly amplified
through matrix inversion.

Shrinkage typically reduces extreme estimated eigenvalues and improves
conditioning.

Thus,

\begin{equation}
\boxed{
\text{Slightly Biased Risk Estimate}
\quad\text{may be preferable to}\quad
\text{Highly Noisy Unbiased Estimate}.
}
\end{equation}


% ------------------------------------------------------------
\subsection{Factor Risk Models}
\label{subsec:factor_risk_models}
% ------------------------------------------------------------

A factor risk model reduces the dimensionality of the covariance-estimation
problem by assuming that security returns are driven by a smaller number of
systematic factors plus security-specific residuals.

Let

\begin{equation}
K
\ll
N_t
\end{equation}

denote the number of risk factors.

For security $i$,

\begin{equation}
\boxed{
R_{i,t}
=
\mathbf{b}_{i,t}^{\top}
\mathbf{f}_t
+
\varepsilon_{i,t},
}
\end{equation}

where

\begin{equation}
\mathbf{b}_{i,t}
\in
\mathbb{R}^{K}
\end{equation}

is the vector of factor exposures and

\begin{equation}
\mathbf{f}_t
\in
\mathbb{R}^{K}
\end{equation}

is the vector of factor returns.

Across all securities,

\begin{equation}
\boxed{
\mathbf{R}_t
=
\mathbf{B}_t
\mathbf{f}_t
+
\boldsymbol{\varepsilon}_t,
}
\end{equation}

with

\begin{equation}
\mathbf{B}_t
\in
\mathbb{R}^{N_t\times K}.
\end{equation}


\subsubsection{Factor Covariance Matrix}

Let

\begin{equation}
\boldsymbol{\Omega}_t
=
\operatorname{Cov}_t
\left(
\mathbf{f}_{t+1}
\right)
\in
\mathbb{R}^{K\times K}
\end{equation}

denote the factor covariance matrix.

Let

\begin{equation}
\mathbf{D}_t
=
\operatorname{Cov}_t
\left(
\boldsymbol{\varepsilon}_{t+1}
\right)
\in
\mathbb{R}^{N_t\times N_t}.
\end{equation}

Under the common assumption that specific returns are mutually uncorrelated,

\begin{equation}
\boxed{
\mathbf{D}_t
=
\operatorname{diag}
\left(
\sigma_{\varepsilon,1,t}^{2},
\ldots,
\sigma_{\varepsilon,N_t,t}^{2}
\right).
}
\end{equation}

The asset covariance matrix becomes

\begin{equation}
\boxed{
\boldsymbol{\Sigma}_t
=
\mathbf{B}_t
\boldsymbol{\Omega}_t
\mathbf{B}_t^\top
+
\mathbf{D}_t.
}
\end{equation}


\subsubsection{Dimensionality Reduction}

The factor model replaces direct estimation of approximately

\begin{equation}
\frac{N_t(N_t+1)}{2}
\end{equation}

asset-covariance parameters with estimation of:

\begin{itemize}
    \item $N_tK$ factor exposures;
    \item $K(K+1)/2$ factor covariances;
    \item $N_t$ specific variances.
\end{itemize}

When

\begin{equation}
K\ll N_t,
\end{equation}

this imposes substantial structure on the risk model.


\subsubsection{Types of Risk Factors}

Risk factors may be:

\begin{itemize}
    \item statistical factors, such as principal components;
    \item macroeconomic factors;
    \item market and country factors;
    \item industry factors;
    \item style factors such as size, value, momentum, or volatility.
\end{itemize}

The choice of risk factors need not coincide with the factors used to construct
the alpha signal.


% ------------------------------------------------------------
\subsection{Systematic and Specific Risk}
\label{subsec:systematic_specific_risk}
% ------------------------------------------------------------

The factor model provides a natural decomposition of portfolio risk.

For portfolio weights

\begin{equation}
\mathbf{w}_t
\in
\mathbb{R}^{N_t},
\end{equation}

portfolio factor exposure is

\begin{equation}
\boxed{
\mathbf{b}_{p,t}
=
\mathbf{B}_t^\top
\mathbf{w}_t
\in
\mathbb{R}^{K}.
}
\end{equation}

Using

\begin{equation}
\boldsymbol{\Sigma}_t
=
\mathbf{B}_t
\boldsymbol{\Omega}_t
\mathbf{B}_t^\top
+
\mathbf{D}_t,
\end{equation}

portfolio variance becomes

\begin{equation}
\sigma_{p,t}^{2}
=
\mathbf{w}_t^\top
\mathbf{B}_t
\boldsymbol{\Omega}_t
\mathbf{B}_t^\top
\mathbf{w}_t
+
\mathbf{w}_t^\top
\mathbf{D}_t
\mathbf{w}_t.
\end{equation}

Since

\begin{equation}
\mathbf{b}_{p,t}
=
\mathbf{B}_t^\top\mathbf{w}_t,
\end{equation}

we obtain

\begin{equation}
\boxed{
\sigma_{p,t}^{2}
=
\underbrace{
\mathbf{b}_{p,t}^{\top}
\boldsymbol{\Omega}_t
\mathbf{b}_{p,t}
}_{\text{systematic variance}}
+
\underbrace{
\mathbf{w}_t^\top
\mathbf{D}_t
\mathbf{w}_t
}_{\text{specific variance}}.
}
\end{equation}


\subsubsection{Systematic Risk}

Systematic risk results from portfolio exposure to common factors:

\begin{equation}
\boxed{
\sigma_{p,t,\mathrm{sys}}^{2}
=
\mathbf{b}_{p,t}^{\top}
\boldsymbol{\Omega}_t
\mathbf{b}_{p,t}.
}
\end{equation}

It cannot generally be eliminated merely by increasing the number of securities
while maintaining the same factor exposures.


\subsubsection{Specific Risk}

Specific risk is

\begin{equation}
\boxed{
\sigma_{p,t,\mathrm{spec}}^{2}
=
\mathbf{w}_t^\top
\mathbf{D}_t
\mathbf{w}_t.
}
\end{equation}

If specific returns are independent and positions are sufficiently
diversified, this component can decrease substantially as the number of
holdings increases.


\subsubsection{Portfolio Risk Decomposition}

Thus,

\begin{equation}
\boxed{
\text{Total Portfolio Risk}
=
\text{Systematic Risk}
+
\text{Specific Risk}
}
\end{equation}

when expressed in variance terms under the factor-model assumptions.

The corresponding volatilities are not additive:

\begin{equation}
\sigma_{p,t}
\neq
\sigma_{p,t,\mathrm{sys}}
+
\sigma_{p,t,\mathrm{spec}}.
\end{equation}

Instead,

\begin{equation}
\boxed{
\sigma_{p,t}
=
\sqrt{
\sigma_{p,t,\mathrm{sys}}^{2}
+
\sigma_{p,t,\mathrm{spec}}^{2}
}.
}
\end{equation}


% ------------------------------------------------------------
\subsection{Beta Estimation}
\label{subsec:portfolio_beta_estimation}
% ------------------------------------------------------------

Market beta is one of the most common security-level risk exposures.

For security $i$, consider

\begin{equation}
\boxed{
R_{i,\tau}^{e}
=
\alpha_{i,t}
+
\beta_{i,t}
R_{M,\tau}^{e}
+
\varepsilon_{i,\tau},
}
\end{equation}

estimated over a historical window

\begin{equation}
\tau
\in
\mathcal{W}_t.
\end{equation}

Here:

\begin{itemize}
    \item $R_{i,\tau}^{e}$ is the security excess return;
    \item $R_{M,\tau}^{e}$ is the market excess return;
    \item $\beta_{i,t}$ is the estimated market exposure at decision date $t$.
\end{itemize}


\subsubsection{Covariance Representation}

For a single-factor regression with an intercept, the population beta is

\begin{equation}
\boxed{
\beta_{i,t}
=
\frac{
\operatorname{Cov}_t
\left(
R_i^{e},
R_M^{e}
\right)
}{
\operatorname{Var}_t
\left(
R_M^{e}
\right)
}.
}
\end{equation}

The sample estimator is

\begin{equation}
\widehat{\beta}_{i,t}
=
\frac{
\widehat{\operatorname{Cov}}_t
\left(
R_i^{e},
R_M^{e}
\right)
}{
\widehat{\operatorname{Var}}_t
\left(
R_M^{e}
\right)
}.
\end{equation}


\subsubsection{Portfolio Beta}

For a linear portfolio,

\begin{equation}
R_{p,t}
=
\sum_i
w_{i,t}
R_{i,t},
\end{equation}

portfolio beta is

\begin{equation}
\boxed{
\beta_{p,t}
=
\sum_{i=1}^{N_t}
w_{i,t}
\beta_{i,t}.
}
\end{equation}

In vector notation,

\begin{equation}
\boxed{
\beta_{p,t}
=
\boldsymbol{\beta}_t^\top
\mathbf{w}_t,
}
\end{equation}

where

\begin{equation}
\boldsymbol{\beta}_t
=
\begin{bmatrix}
\beta_{1,t}\\
\vdots\\
\beta_{N_t,t}
\end{bmatrix}
\in
\mathbb{R}^{N_t}.
\end{equation}


\subsubsection{Rolling Beta}

A rolling beta estimate may use the most recent $L$ returns:

\begin{equation}
\widehat{\beta}_{i,t}^{(L)}
=
\frac{
\widehat{\operatorname{Cov}}
\left(
R_{i,t-L+1:t},
R_{M,t-L+1:t}
\right)
}{
\widehat{\operatorname{Var}}
\left(
R_{M,t-L+1:t}
\right)
}.
\end{equation}

A shorter lookback adapts more quickly but produces noisier estimates.

A longer lookback is more stable but may react slowly to structural changes in
the security's market exposure.


\subsubsection{Exponentially Weighted Beta}

Historical observations may receive exponentially decaying weights.

A weighted beta estimator can be represented as

\begin{equation}
\boxed{
\widehat{\beta}_{i,t}^{EW}
=
\frac{
\widehat{\operatorname{Cov}}_{EW,t}
\left(
R_i,R_M
\right)
}{
\widehat{\operatorname{Var}}_{EW,t}
\left(
R_M
\right)
}.
}
\end{equation}

This gives more importance to recent comovements.


\subsubsection{Point-in-Time Beta}

The beta used for portfolio construction must be estimated entirely from
historical observations available at the decision time.

Thus,

\begin{equation}
\boxed{
\widehat{\beta}_{i,t}
\in
\mathcal{I}_t.
}
\end{equation}

Using future returns when estimating historical betas would directly introduce
look-ahead bias.


% ------------------------------------------------------------
\subsection{Factor Exposure Estimation}
\label{subsec:factor_exposure_estimation}
% ------------------------------------------------------------

In a multi-factor risk model, each security has factor exposure vector

\begin{equation}
\boxed{
\mathbf{b}_{i,t}
=
\begin{bmatrix}
b_{i,t}^{(1)}
\\
\vdots\\
b_{i,t}^{(K)}
\end{bmatrix}
\in
\mathbb{R}^{K}.
}
\end{equation}

The exposure matrix is

\begin{equation}
\boxed{
\mathbf{B}_t
=
\begin{bmatrix}
(\mathbf{b}_{1,t})^\top\\
\vdots\\
(\mathbf{b}_{N_t,t})^\top
\end{bmatrix}
\in
\mathbb{R}^{N_t\times K}.
}
\end{equation}

Factor exposures can be estimated or directly observed depending on the type of
factor.


\subsubsection{Time-Series Estimated Exposures}

For traded or return-based factors,

\begin{equation}
R_{i,\tau}
=
\alpha_i
+
\mathbf{b}_{i,t}^\top
\mathbf{f}_\tau
+
\varepsilon_{i,\tau}.
\end{equation}

Given a historical factor-return matrix

\begin{equation}
\mathbf{F}
\in
\mathbb{R}^{L\times K},
\end{equation}

the OLS estimator is

\begin{equation}
\boxed{
\widehat{\mathbf{b}}_{i,t}
=
\left(
\mathbf{F}^\top
\mathbf{F}
\right)^{-1}
\mathbf{F}^\top
\mathbf{r}_i,
}
\end{equation}

after appropriate treatment of the intercept.


\subsubsection{Observed Exposures}

Some risk factors are based directly on observable security characteristics.

Examples include:

\begin{itemize}
    \item industry membership;
    \item country membership;
    \item log market capitalization;
    \item leverage;
    \item liquidity;
    \item volatility.
\end{itemize}

These exposures may be transformed or standardized cross-sectionally before
entering the risk model.


\subsubsection{Portfolio Factor Exposures}

For weights

\begin{equation}
\mathbf{w}_t
\in
\mathbb{R}^{N_t},
\end{equation}

portfolio exposures are

\begin{equation}
\boxed{
\mathbf{b}_{p,t}
=
\mathbf{B}_t^\top
\mathbf{w}_t
\in
\mathbb{R}^{K}.
}
\end{equation}

The $k$-th component is

\begin{equation}
b_{p,t}^{(k)}
=
\sum_{i=1}^{N_t}
w_{i,t}
b_{i,t}^{(k)}.
\end{equation}

This linear aggregation will later allow the portfolio optimizer to impose
factor-exposure constraints directly.


% ------------------------------------------------------------
\subsection{Volatility Forecasting}
\label{subsec:volatility_forecasting}
% ------------------------------------------------------------

Historical volatility is not necessarily the same as expected future
volatility.

Portfolio construction requires an estimate

\begin{equation}
\widehat{\sigma}_{i,t}
\end{equation}

or

\begin{equation}
\widehat{\boldsymbol{\Sigma}}_t
\end{equation}

that is intended to describe risk over the future portfolio holding period.


\subsubsection{Rolling Historical Volatility}

For daily returns and lookback window $L$,

\begin{equation}
\widehat{\sigma}_{i,t}^{\mathrm{daily}}
=
\sqrt{
\frac{1}{L-1}
\sum_{\ell=1}^{L}
\left(
R_{i,t-\ell}
-
\bar{R}_{i,t}
\right)^2
}.
\end{equation}

Assuming $F$ observations per year, annualized volatility is

\begin{equation}
\boxed{
\widehat{\sigma}_{i,t}^{\mathrm{ann}}
=
\sqrt{F}
\,
\widehat{\sigma}_{i,t}^{\mathrm{daily}}.
}
\end{equation}

For daily data, one commonly uses

\begin{equation}
F\approx252.
\end{equation}


\subsubsection{Exponentially Weighted Volatility}

An exponentially weighted variance can be defined recursively as

\begin{equation}
\boxed{
\widehat{\sigma}_{i,t}^{2}
=
\lambda
\widehat{\sigma}_{i,t-1}^{2}
+
(1-\lambda)
R_{i,t-1}^{2},
}
\end{equation}

when returns are treated as approximately mean zero over the relevant horizon.

Here,

\begin{equation}
0<\lambda<1.
\end{equation}

A smaller $\lambda$ gives more weight to recent observations and produces a
more reactive volatility estimate.


\subsubsection{Half-Life}

The decay parameter may alternatively be expressed through a half-life $H$:

\begin{equation}
\boxed{
\lambda
=
2^{-1/H}.
}
\end{equation}

After $H$ periods, the relative weight assigned to an observation falls by
one half.


\subsubsection{Conditional Volatility Models}

More sophisticated models explicitly model time-varying conditional variance.

For example, a GARCH$(1,1)$ model is

\begin{equation}
\boxed{
\sigma_{i,t}^{2}
=
\omega
+
\alpha
\varepsilon_{i,t-1}^{2}
+
\beta
\sigma_{i,t-1}^{2},
}
\end{equation}

where

\begin{equation}
\omega>0,
\qquad
\alpha\geq0,
\qquad
\beta\geq0.
\end{equation}

Such models capture volatility clustering by allowing large recent shocks to
increase the forecast of future risk.


\subsubsection{Forecast Horizon}

Risk estimates should match the relevant portfolio horizon.

If daily covariance is estimated as

\begin{equation}
\boldsymbol{\Sigma}_t^{(1d)},
\end{equation}

then under an independent and identically distributed approximation, an
$h$-day covariance could be approximated by

\begin{equation}
\boxed{
\boldsymbol{\Sigma}_t^{(h)}
\approx
h
\boldsymbol{\Sigma}_t^{(1d)}.
}
\end{equation}

Corresponding volatility scales as

\begin{equation}
\boxed{
\sigma_t^{(h)}
\approx
\sqrt{h}
\sigma_t^{(1d)}.
}
\end{equation}

These square-root-of-time relations rely on strong assumptions and may fail in
the presence of serial correlation or time-varying volatility.


\subsubsection{Forecast Risk vs.\ Realized Risk}

A crucial distinction is

\begin{equation}
\boxed{
\widehat{\sigma}_{p,t}
\neq
\sigma_{p,t}^{\mathrm{realized}}.
}
\end{equation}

The first is an ex-ante forecast available when the portfolio is constructed.

The second is an ex-post quantity observed after returns have occurred.

The accuracy of the risk model can later be evaluated by comparing predicted
and realized risk.


% ------------------------------------------------------------
\subsubsection{Risk Model Output}
% ------------------------------------------------------------

At the end of the risk-model stage, the backtester should provide at least

\begin{equation}
\boxed{
\widehat{\boldsymbol{\Sigma}}_t
\in
\mathbb{R}^{N_t\times N_t}.
}
\end{equation}

If a factor model is used, the underlying components should also be retained:

\begin{equation}
\boxed{
\begin{aligned}
\mathbf{B}_t
&\in
\mathbb{R}^{N_t\times K},
\\
\widehat{\boldsymbol{\Omega}}_t
&\in
\mathbb{R}^{K\times K},
\\
\widehat{\mathbf{D}}_t
&\in
\mathbb{R}^{N_t\times N_t},
\\
\widehat{\boldsymbol{\Sigma}}_t
&=
\mathbf{B}_t
\widehat{\boldsymbol{\Omega}}_t
\mathbf{B}_t^\top
+
\widehat{\mathbf{D}}_t.
\end{aligned}
}
\end{equation}

The portfolio-construction engine now has two central objects:

\begin{equation}
\boxed{
\begin{aligned}
\mathbf{s}_t^{\mathrm{final}}
&\in
\mathbb{R}^{N_t},
&&
\text{expected attractiveness},
\\
\widehat{\boldsymbol{\Sigma}}_t
&\in
\mathbb{R}^{N_t\times N_t},
&&
\text{expected risk}.
\end{aligned}
}
\end{equation}

These objects provide the core inputs for the optimization problems developed in
the following sections:

\begin{equation}
\boxed{
\begin{aligned}
\text{Final Signal}
\;(\mathbf{s}_t^{\mathrm{final}})
&+
\text{Risk Model}
\;(\widehat{\boldsymbol{\Sigma}}_t)
\\
&\rightarrow
\text{Portfolio Optimization}
\\
&\rightarrow
\mathbf{w}_t^{\mathrm{target}}.
\end{aligned}
}
\end{equation}

% ============================================================
\section{Portfolio Optimization}
\label{sec:portfolio_optimization}
% ============================================================

Once the final signal and the portfolio risk model have been constructed, the
backtester can determine the portfolio weights that best trade off expected
return, risk, implementation costs, and portfolio stability.

The central optimization problem can be written abstractly as

\begin{equation}
\boxed{
\mathbf{w}_t^{\mathrm{target}}
=
\underset{\mathbf{w}\in\mathcal{W}_t}{\arg\max}
\;
\mathcal{J}_t
\left(
\mathbf{w}
\right),
}
\end{equation}

where

\begin{equation}
\mathbf{w}
\in
\mathbb{R}^{N_t}
\end{equation}

is a candidate portfolio-weight vector and

\begin{equation}
\mathcal{W}_t
\end{equation}

denotes the feasible set defined by the portfolio constraints.

The objective function may incorporate several components:

\begin{equation}
\boxed{
\begin{aligned}
\mathcal{J}_t(\mathbf{w})
=
&
\underbrace{
\mathcal{A}_t(\mathbf{w})
}_{\text{Expected Alpha}}
-
\underbrace{
\mathcal{R}_t(\mathbf{w})
}_{\text{Risk Penalty}}
\\
&
-
\underbrace{
\mathcal{C}_t(\mathbf{w},\mathbf{w}_{t^-})
}_{\text{Trading Cost}}
-
\underbrace{
\mathcal{P}_t(\mathbf{w})
}_{\text{Regularization}}.
\end{aligned}
}
\end{equation}

Here,

\begin{equation}
\mathbf{w}_{t^-}
\end{equation}

denotes portfolio weights immediately before the rebalance.

The exact formulation depends on the interpretation of the signal.

If

\begin{equation}
\mathbf{s}_t^{\mathrm{final}}
\end{equation}

is a calibrated expected-return vector, it can enter the optimization directly.

If it is only an ordinal or standardized alpha score, an explicit mapping from
score to expected return may be required before using expected-return-based
optimization.

Thus,

\begin{equation}
\boxed{
\text{Signal Score}
\neq
\text{Expected Return}
}
\end{equation}

unless the signal has been explicitly calibrated as such.


% ------------------------------------------------------------
\subsection{General Optimization Framework}
\label{subsec:general_portfolio_optimization}
% ------------------------------------------------------------

Let

\begin{equation}
\widehat{\boldsymbol{\mu}}_t
=
\begin{bmatrix}
\widehat{\mu}_{1,t}
\\
\vdots
\\
\widehat{\mu}_{N_t,t}
\end{bmatrix}
\in
\mathbb{R}^{N_t}
\end{equation}

denote the vector of expected future security returns.

Let

\begin{equation}
\widehat{\boldsymbol{\Sigma}}_t
\in
\mathbb{R}^{N_t\times N_t}
\end{equation}

denote the ex-ante covariance matrix.

A generic optimization problem is

\begin{equation}
\boxed{
\mathbf{w}_t^{*}
=
\underset{\mathbf{w}\in\mathcal{W}_t}{\arg\max}
\left[
\mathbf{w}^{\top}
\widehat{\boldsymbol{\mu}}_t
-
\lambda_R
\mathcal{R}_t(\mathbf{w})
-
\lambda_C
\mathcal{C}_t(\mathbf{w},\mathbf{w}_{t^-})
-
\lambda_P
\mathcal{P}_t(\mathbf{w})
\right].
}
\end{equation}

The coefficients

\begin{equation}
\lambda_R,
\qquad
\lambda_C,
\qquad
\lambda_P
\end{equation}

control the relative importance of risk, trading costs, and regularization.


\subsubsection{Feasible Set}

The feasible set may be written abstractly as

\begin{equation}
\boxed{
\mathcal{W}_t
=
\left\{
\mathbf{w}\in\mathbb{R}^{N_t}
:
\mathbf{w}
\text{ satisfies all portfolio constraints}
\right\}.
}
\end{equation}

Possible restrictions include:

\begin{itemize}
    \item gross exposure limits;
    \item net exposure limits;
    \item individual position bounds;
    \item beta constraints;
    \item sector constraints;
    \item country constraints;
    \item factor-exposure limits;
    \item liquidity constraints;
    \item shortability constraints.
\end{itemize}

These constraints will be developed explicitly in subsequent sections.


\subsubsection{Convexity}

A major practical distinction is whether the optimization problem is convex.

If the objective is concave and the feasible set is convex, the optimizer can
generally identify a global optimum reliably.

For example,

\begin{equation}
\mathbf{w}^{\top}
\widehat{\boldsymbol{\mu}}_t
-
\frac{\lambda}{2}
\mathbf{w}^{\top}
\widehat{\boldsymbol{\Sigma}}_t
\mathbf{w}
\end{equation}

is concave in $\mathbf{w}$ when

\begin{equation}
\widehat{\boldsymbol{\Sigma}}_t
\succeq
0.
\end{equation}

Linear equality and inequality constraints preserve convexity.

By contrast, cardinality constraints, integer trading decisions, or certain
nonlinear market-impact models may produce non-convex optimization problems.


% ------------------------------------------------------------
\subsection{Expected Return Maximization}
\label{subsec:expected_return_maximization}
% ------------------------------------------------------------

The simplest optimization objective maximizes expected portfolio return:

\begin{equation}
\boxed{
\mathbf{w}_t^{*}
=
\underset{\mathbf{w}\in\mathcal{W}_t}{\arg\max}
\;
\mathbf{w}^{\top}
\widehat{\boldsymbol{\mu}}_t.
}
\end{equation}

Expected portfolio return is

\begin{equation}
\boxed{
\widehat{\mu}_{p,t}
=
\mathbf{w}^{\top}
\widehat{\boldsymbol{\mu}}_t.
}
\end{equation}


\subsubsection{Need for Constraints}

Without constraints, the problem is generally ill posed.

If arbitrary leverage is allowed and at least one asset satisfies

\begin{equation}
\widehat{\mu}_{i,t}>0,
\end{equation}

the optimizer can increase expected return indefinitely by increasing the
position.

Thus,

\begin{equation}
\boxed{
\text{Expected Return Maximization}
+
\text{No Risk or Leverage Constraint}
\Rightarrow
\text{Unbounded Problem}.
}
\end{equation}

Meaningful expected-return maximization therefore requires constraints such as

\begin{equation}
\sum_i |w_i|
\leq
G_{\max},
\end{equation}

or explicit risk limits.


\subsubsection{Corner Solutions}

Even with leverage constraints, pure expected-return maximization often produces
highly concentrated portfolios.

For example, under a long-only fully invested constraint

\begin{equation}
w_i\geq0,
\qquad
\sum_iw_i=1,
\end{equation}

the solution is generally to allocate all capital to the security with the
highest expected return:

\begin{equation}
w_{i^*}=1,
\end{equation}

where

\begin{equation}
i^*
=
\underset{i}{\arg\max}
\;
\widehat{\mu}_{i,t}.
\end{equation}

This illustrates why expected return alone is usually insufficient for
realistic portfolio construction.


% ------------------------------------------------------------
\subsection{Mean--Variance Optimization}
\label{subsec:mean_variance_optimization}
% ------------------------------------------------------------

Mean--variance optimization explicitly balances expected return and portfolio
variance.

The standard objective is

\begin{equation}
\boxed{
\mathbf{w}_t^{*}
=
\underset{\mathbf{w}\in\mathcal{W}_t}{\arg\max}
\left[
\mathbf{w}^{\top}
\widehat{\boldsymbol{\mu}}_t
-
\frac{\lambda}{2}
\mathbf{w}^{\top}
\widehat{\boldsymbol{\Sigma}}_t
\mathbf{w}
\right],
}
\end{equation}

where

\begin{equation}
\lambda>0
\end{equation}

is the risk-aversion parameter.


\subsubsection{Interpretation}

The first term,

\begin{equation}
\mathbf{w}^{\top}
\widehat{\boldsymbol{\mu}}_t,
\end{equation}

rewards expected return.

The second,

\begin{equation}
\frac{\lambda}{2}
\mathbf{w}^{\top}
\widehat{\boldsymbol{\Sigma}}_t
\mathbf{w},
\end{equation}

penalizes portfolio variance.

A larger $\lambda$ produces a more risk-averse portfolio.


\subsubsection{Unconstrained Solution}

If no constraints are imposed and

\begin{equation}
\widehat{\boldsymbol{\Sigma}}_t
\end{equation}

is positive definite, the first-order condition is

\begin{equation}
\widehat{\boldsymbol{\mu}}_t
-
\lambda
\widehat{\boldsymbol{\Sigma}}_t
\mathbf{w}
=
0.
\end{equation}

Therefore,

\begin{equation}
\boxed{
\mathbf{w}_t^{*}
=
\frac{1}{\lambda}
\widehat{\boldsymbol{\Sigma}}_t^{-1}
\widehat{\boldsymbol{\mu}}_t.
}
\end{equation}

This expression illustrates the joint role of alpha and risk.

A security does not receive a large weight merely because its expected return
is high. Its weight also depends on how its risk covaries with the rest of the
portfolio.


\subsubsection{Connection with Signal Mapping}

If

\begin{equation}
\widehat{\boldsymbol{\Sigma}}_t
=
\sigma^2
\mathbf{I},
\end{equation}

then

\begin{equation}
\mathbf{w}_t^{*}
=
\frac{1}{\lambda\sigma^2}
\widehat{\boldsymbol{\mu}}_t.
\end{equation}

Thus,

\begin{equation}
\boxed{
\mathbf{w}_t^{*}
\propto
\widehat{\boldsymbol{\mu}}_t.
}
\end{equation}

Linear signal mapping is therefore a special case of mean--variance
optimization when all assets have equal variance and zero covariance.


\subsubsection{Estimation Error}

The unconstrained solution depends on

\begin{equation}
\widehat{\boldsymbol{\Sigma}}_t^{-1}
\widehat{\boldsymbol{\mu}}_t.
\end{equation}

Both expected returns and covariance matrices are estimated with error.

Expected-return estimates are usually particularly noisy.

Consequently, unconstrained mean--variance optimization may produce extreme
and unstable weights.

This motivates:

\begin{itemize}
    \item covariance shrinkage;
    \item expected-return shrinkage;
    \item position constraints;
    \item turnover penalties;
    \item weight regularization.
\end{itemize}


% ------------------------------------------------------------
\subsection{Risk-Adjusted Optimization}
\label{subsec:risk_adjusted_optimization}
% ------------------------------------------------------------

Rather than choosing an explicit risk-aversion parameter, the optimizer may
maximize a risk-adjusted objective.


\subsubsection{Maximum Sharpe Ratio}

Assuming a zero or already subtracted risk-free rate, the ex-ante portfolio
Sharpe ratio is

\begin{equation}
\boxed{
SR_t(\mathbf{w})
=
\frac{
\mathbf{w}^{\top}
\widehat{\boldsymbol{\mu}}_t
}{
\sqrt{
\mathbf{w}^{\top}
\widehat{\boldsymbol{\Sigma}}_t
\mathbf{w}
}
}.
}
\end{equation}

The optimization problem is

\begin{equation}
\boxed{
\mathbf{w}_t^{*}
=
\underset{\mathbf{w}\in\mathcal{W}_t}{\arg\max}
\;
SR_t(\mathbf{w}).
}
\end{equation}


\subsubsection{Scale Invariance}

For any

\begin{equation}
c>0,
\end{equation}

one has

\begin{equation}
SR_t(c\mathbf{w})
=
SR_t(\mathbf{w}).
\end{equation}

Therefore, maximum-Sharpe optimization determines a portfolio direction but not
its absolute leverage.

A separate volatility or gross-exposure target is required to determine scale.


\subsubsection{Maximum Information Ratio}

For benchmark-relative portfolios, expected active return may be written as

\begin{equation}
\widehat{\mu}_{p,t}^{\mathrm{active}}
=
\mathbf{w}^{\top}
\widehat{\boldsymbol{\alpha}}_t,
\end{equation}

where

\begin{equation}
\widehat{\boldsymbol{\alpha}}_t
\end{equation}

denotes expected benchmark-relative returns.

If

\begin{equation}
\widehat{\boldsymbol{\Sigma}}_t^{\mathrm{active}}
\end{equation}

describes active-return risk, the ex-ante Information Ratio is

\begin{equation}
\boxed{
IR_t(\mathbf{w})
=
\frac{
\mathbf{w}^{\top}
\widehat{\boldsymbol{\alpha}}_t
}{
\sqrt{
\mathbf{w}^{\top}
\widehat{\boldsymbol{\Sigma}}_t^{\mathrm{active}}
\mathbf{w}
}
}.
}
\end{equation}

This formulation is natural for benchmark-relative active equity portfolios.


\subsubsection{Risk Budget Formulation}

Another formulation maximizes expected return subject to a risk constraint:

\begin{equation}
\boxed{
\begin{aligned}
\max_{\mathbf{w}}
\quad&
\mathbf{w}^{\top}
\widehat{\boldsymbol{\mu}}_t
\\
\text{s.t.}
\quad&
\mathbf{w}^{\top}
\widehat{\boldsymbol{\Sigma}}_t
\mathbf{w}
\leq
\sigma_{\max}^{2},
\\
&
\mathbf{w}
\in
\mathcal{W}_t.
\end{aligned}
}
\end{equation}

This expresses the investment problem directly as:

\begin{equation}
\boxed{
\text{maximize alpha for a fixed risk budget}.
}
\end{equation}


% ------------------------------------------------------------
\subsection{Transaction-Cost-Aware Optimization}
\label{subsec:transaction_cost_aware_optimization}
% ------------------------------------------------------------

The portfolio with the highest gross expected alpha may not be optimal after
implementation costs.

Let

\begin{equation}
\mathbf{w}_{t^-}
\in
\mathbb{R}^{N_t}
\end{equation}

denote pre-rebalance portfolio weights.

The required trade is

\begin{equation}
\boxed{
\Delta\mathbf{w}_t
=
\mathbf{w}_t
-
\mathbf{w}_{t^-}.
}
\end{equation}

Let

\begin{equation}
C_t
\left(
\Delta\mathbf{w}_t
\right)
\end{equation}

denote expected transaction costs expressed as a fraction of portfolio NAV.

A cost-aware mean--variance problem is

\begin{equation}
\boxed{
\mathbf{w}_t^{*}
=
\underset{\mathbf{w}\in\mathcal{W}_t}{\arg\max}
\left[
\mathbf{w}^{\top}
\widehat{\boldsymbol{\mu}}_t
-
\frac{\lambda}{2}
\mathbf{w}^{\top}
\widehat{\boldsymbol{\Sigma}}_t
\mathbf{w}
-
C_t
\left(
\mathbf{w}
-
\mathbf{w}_{t^-}
\right)
\right].
}
\end{equation}


\subsubsection{Linear Transaction Costs}

A simple proportional cost model is

\begin{equation}
\boxed{
C_t^{\mathrm{linear}}
=
\sum_{i=1}^{N_t}
c_{i,t}
\left|
w_{i,t}
-
w_{i,t^-}
\right|,
}
\end{equation}

where

\begin{equation}
c_{i,t}
\end{equation}

is the estimated one-way transaction cost per unit of traded portfolio notional.

The optimizer will trade only when the expected improvement in the objective is
sufficient to compensate for these costs.


\subsubsection{Nonlinear Market Impact}

Large trades generally have higher marginal cost.

A stylized nonlinear cost function may take the form

\begin{equation}
\boxed{
C_{i,t}^{\mathrm{impact}}
=
\eta_{i,t}
\left|
\Delta w_{i,t}
\right|^{p},
\qquad
p>1.
}
\end{equation}

Then

\begin{equation}
C_t^{\mathrm{impact}}
=
\sum_i
C_{i,t}^{\mathrm{impact}}.
\end{equation}

The convexity of the cost function discourages concentrated large trades.


\subsubsection{Costs in Currency Notional}

If portfolio NAV is

\begin{equation}
A_t,
\end{equation}

the currency notional traded is

\begin{equation}
Q_{i,t}
=
A_t
|\Delta w_{i,t}|.
\end{equation}

Transaction costs may therefore depend on

\begin{equation}
C_{i,t}
=
C
\left(
Q_{i,t},
\mathrm{ADV}_{i,t},
\mathrm{Spread}_{i,t},
\sigma_{i,t},
\ldots
\right).
\end{equation}

This establishes the connection between portfolio optimization and the
implementation data introduced earlier in the backtester.


% ------------------------------------------------------------
\subsection{Turnover-Aware Optimization}
\label{subsec:turnover_aware_optimization}
% ------------------------------------------------------------

Turnover-aware optimization penalizes changes in portfolio weights even when a
complete transaction-cost model is unavailable.

A natural turnover measure is

\begin{equation}
\boxed{
TO_t
=
\sum_{i=1}^{N_t}
\left|
w_{i,t}
-
w_{i,t^-}
\right|.
}
\end{equation}

A turnover-penalized objective is

\begin{equation}
\boxed{
\mathbf{w}_t^{*}
=
\underset{\mathbf{w}\in\mathcal{W}_t}{\arg\max}
\left[
\mathbf{w}^{\top}
\widehat{\boldsymbol{\mu}}_t
-
\frac{\lambda}{2}
\mathbf{w}^{\top}
\widehat{\boldsymbol{\Sigma}}_t
\mathbf{w}
-
\kappa
\left\|
\mathbf{w}
-
\mathbf{w}_{t^-}
\right\|_1
\right],
}
\end{equation}

where

\begin{equation}
\kappa\geq0.
\end{equation}


\subsubsection{$L^1$ Turnover Penalty}

The $L^1$ penalty is

\begin{equation}
\left\|
\mathbf{w}
-
\mathbf{w}_{t^-}
\right\|_1
=
\sum_i
|w_i-w_{i,t^-}|.
\end{equation}

It directly corresponds to total traded notional in weight space.


\subsubsection{Quadratic Turnover Penalty}

A smoother alternative is

\begin{equation}
\boxed{
\kappa
\left\|
\mathbf{w}
-
\mathbf{w}_{t^-}
\right\|_2^2.
}
\end{equation}

The optimization becomes

\begin{equation}
\mathbf{w}_t^{*}
=
\underset{\mathbf{w}\in\mathcal{W}_t}{\arg\max}
\left[
\mathbf{w}^{\top}
\widehat{\boldsymbol{\mu}}_t
-
\frac{\lambda}{2}
\mathbf{w}^{\top}
\widehat{\boldsymbol{\Sigma}}_t
\mathbf{w}
-
\kappa
\left\|
\mathbf{w}
-
\mathbf{w}_{t^-}
\right\|_2^2
\right].
\end{equation}

Unlike the $L^1$ penalty, this penalizes large trades disproportionately.


\subsubsection{Turnover Constraint}

Instead of penalizing turnover in the objective, a hard constraint may be
imposed:

\begin{equation}
\boxed{
\sum_i
|w_{i,t}-w_{i,t^-}|
\leq
TO_{\max}.
}
\end{equation}

The distinction is

\begin{equation}
\boxed{
\begin{aligned}
\text{Penalty}
&:
\quad
\text{turnover is allowed but costly},
\\
\text{Constraint}
&:
\quad
\text{turnover may not exceed a fixed limit}.
\end{aligned}
}
\end{equation}


\subsubsection{Endogenous No-Trade Region}

An important consequence of transaction-cost or $L^1$ turnover penalties is
that small changes in expected alpha may optimally produce no trade.

Conceptually,

\begin{equation}
\boxed{
\text{Marginal Alpha Improvement}
<
\text{Marginal Trading Cost}
\Rightarrow
\Delta w_{i,t}^{*}=0.
}
\end{equation}

Thus, the optimization can generate an endogenous no-trade region rather than
requiring an arbitrary signal threshold.


% ------------------------------------------------------------
\subsection{Regularization of Portfolio Weights}
\label{subsec:portfolio_weight_regularization}
% ------------------------------------------------------------

Portfolio optimization can be highly sensitive to estimation error in expected
returns and covariance matrices.

Regularization reduces this sensitivity by penalizing extreme or unstable
portfolio weights.


\subsubsection{$L^2$ Weight Regularization}

A quadratic weight penalty is

\begin{equation}
\boxed{
\lambda_w
\|\mathbf{w}\|_2^2
=
\lambda_w
\sum_{i=1}^{N_t}
w_i^2.
}
\end{equation}

A regularized mean--variance objective becomes

\begin{equation}
\boxed{
\mathbf{w}_t^{*}
=
\underset{\mathbf{w}\in\mathcal{W}_t}{\arg\max}
\left[
\mathbf{w}^{\top}
\widehat{\boldsymbol{\mu}}_t
-
\frac{\lambda_R}{2}
\mathbf{w}^{\top}
\widehat{\boldsymbol{\Sigma}}_t
\mathbf{w}
-
\lambda_w
\mathbf{w}^{\top}\mathbf{w}
\right].
}
\end{equation}


\subsubsection{Equivalent Covariance Interpretation}

The risk and regularization terms can be combined:

\begin{equation}
\frac{\lambda_R}{2}
\mathbf{w}^{\top}
\widehat{\boldsymbol{\Sigma}}_t
\mathbf{w}
+
\lambda_w
\mathbf{w}^{\top}\mathbf{w}.
\end{equation}

This can be rewritten as

\begin{equation}
\boxed{
\frac{\lambda_R}{2}
\mathbf{w}^{\top}
\left(
\widehat{\boldsymbol{\Sigma}}_t
+
\frac{2\lambda_w}{\lambda_R}
\mathbf{I}
\right)
\mathbf{w}.
}
\end{equation}

Thus, $L^2$ weight regularization is mathematically related to adding a positive
constant to the diagonal of the covariance matrix.

This improves numerical conditioning and discourages extreme positions.


\subsubsection{$L^1$ Weight Regularization}

An alternative is

\begin{equation}
\boxed{
\lambda_w
\|\mathbf{w}\|_1
=
\lambda_w
\sum_i |w_i|.
}
\end{equation}

However, if gross exposure is already fixed through

\begin{equation}
\sum_i|w_i|
=
G,
\end{equation}

then

\begin{equation}
\|\mathbf{w}\|_1
=
G
\end{equation}

is constant and an additional $L^1$ penalty has no effect.

This illustrates the importance of considering interactions between
regularization terms and portfolio constraints.


\subsubsection{Regularization Toward a Reference Portfolio}

Rather than shrinking weights toward zero, one may shrink toward a reference
portfolio

\begin{equation}
\mathbf{w}_t^{\mathrm{ref}}.
\end{equation}

The penalty is

\begin{equation}
\boxed{
\lambda_w
\left\|
\mathbf{w}
-
\mathbf{w}_t^{\mathrm{ref}}
\right\|_2^2.
}
\end{equation}

Possible reference portfolios include:

\begin{itemize}
    \item equal weight;
    \item benchmark weights;
    \item previous portfolio weights;
    \item raw signal-mapped weights;
    \item risk-balanced weights.
\end{itemize}

This formulation allows the optimizer to deviate from a stable baseline only
when the estimated benefits justify doing so.


\subsubsection{Diversification Effect}

Since

\begin{equation}
\|\mathbf{w}\|_2^2
=
\sum_i w_i^2,
\end{equation}

for a fixed gross or net exposure, minimizing the $L^2$ norm generally
discourages concentration.

For a long-only fully invested portfolio,

\begin{equation}
\sum_i w_i=1,
\qquad
w_i\geq0,
\end{equation}

the minimum of

\begin{equation}
\sum_iw_i^2
\end{equation}

is achieved by

\begin{equation}
\boxed{
w_i
=
\frac{1}{N_t}.
}
\end{equation}

Therefore, $L^2$ regularization naturally pulls the optimizer toward more
diversified allocations.


% ------------------------------------------------------------
\subsubsection{Unified Optimization Objective}
% ------------------------------------------------------------

The different components developed in this section can be combined into a
general portfolio optimization problem:

\begin{equation}
\boxed{
\begin{aligned}
\mathbf{w}_t^{*}
=
\underset{\mathbf{w}\in\mathcal{W}_t}{\arg\max}
\Bigg[
&
\mathbf{w}^{\top}
\widehat{\boldsymbol{\mu}}_t
-
\frac{\lambda_R}{2}
\mathbf{w}^{\top}
\widehat{\boldsymbol{\Sigma}}_t
\mathbf{w}
\\
&
-
C_t
\left(
\mathbf{w}-\mathbf{w}_{t^-}
\right)
\\
&
-
\lambda_{TO}
\left\|
\mathbf{w}-\mathbf{w}_{t^-}
\right\|_1
\\
&
-
\lambda_W
\left\|
\mathbf{w}-\mathbf{w}_t^{\mathrm{ref}}
\right\|_2^2
\Bigg].
\end{aligned}
}
\end{equation}

The optimizer therefore balances four competing objectives:

\begin{equation}
\boxed{
\begin{aligned}
\text{Expected Alpha}
&\uparrow,
\\
\text{Portfolio Risk}
&\downarrow,
\\
\text{Implementation Cost and Turnover}
&\downarrow,
\\
\text{Weight Instability and Concentration}
&\downarrow.
\end{aligned}
}
\end{equation}


\subsubsection{Optimization Output}

The optimization stage produces

\begin{equation}
\boxed{
\mathbf{w}_t^{*}
\in
\mathbb{R}^{N_t}.
}
\end{equation}

However, this vector is only valid if it satisfies all desired portfolio
constraints.

Thus, the logical portfolio-construction sequence is

\begin{equation}
\boxed{
\begin{aligned}
\mathbf{s}_t^{\mathrm{final}}
&\rightarrow
\widehat{\boldsymbol{\mu}}_t
\\
+
\widehat{\boldsymbol{\Sigma}}_t
&\rightarrow
\text{Optimization Objective}
\\
+
\mathcal{W}_t
&\rightarrow
\mathbf{w}_t^{\mathrm{target}}.
\end{aligned}
}
\end{equation}

The next step is therefore to define the feasible set
$\mathcal{W}_t$ explicitly through portfolio exposure, leverage, position,
factor, and liquidity constraints.

% ============================================================
\section{Portfolio Constraints}
\label{sec:portfolio_constraints}
% ============================================================

Portfolio optimization does not operate over all possible weight vectors.
Instead, the optimizer must choose weights inside a feasible set determined by
the investment mandate, risk limits, liquidity, leverage, and implementation
constraints.

Let

\begin{equation}
\mathbf{w}_t
=
\begin{bmatrix}
w_{1,t}\\
\vdots\\
w_{N_t,t}
\end{bmatrix}
\in
\mathbb{R}^{N_t}
\end{equation}

denote a candidate portfolio-weight vector at time $t$.

The feasible set is

\begin{equation}
\boxed{
\mathcal{W}_t
=
\left\{
\mathbf{w}\in\mathbb{R}^{N_t}
:
\mathbf{w}
\text{ satisfies all portfolio constraints}
\right\}.
}
\end{equation}

The optimization problem developed previously can therefore be written as

\begin{equation}
\boxed{
\mathbf{w}_t^{\mathrm{target}}
=
\underset{\mathbf{w}\in\mathcal{W}_t}{\arg\max}
\;
\mathcal{J}_t(\mathbf{w}).
}
\end{equation}

A constraint is fundamentally different from a penalty.

A penalty enters the objective function and may be violated if the expected
benefit is sufficiently large. A hard constraint defines admissibility:

\begin{equation}
\boxed{
\mathbf{w}\notin\mathcal{W}_t
\quad\Rightarrow\quad
\mathbf{w}
\text{ cannot be selected},
}
\end{equation}

regardless of its expected alpha.

The constraints below can be combined depending on the portfolio mandate.


% ------------------------------------------------------------
\subsection{Gross Exposure}
\label{subsec:gross_exposure}
% ------------------------------------------------------------

Gross exposure measures the total absolute size of all long and short
positions.

It is defined as

\begin{equation}
\boxed{
G_t
=
\sum_{i=1}^{N_t}
|w_{i,t}|
=
\|\mathbf{w}_t\|_1.
}
\end{equation}

A maximum gross-exposure constraint is

\begin{equation}
\boxed{
\|\mathbf{w}_t\|_1
\leq
G_{\max}.
}
\end{equation}

Alternatively, the portfolio may be required to maintain an exact gross
exposure:

\begin{equation}
\boxed{
\|\mathbf{w}_t\|_1
=
G_{\mathrm{target}}.
}
\end{equation}


\subsubsection{Interpretation}

For example,

\begin{equation}
G_t=1
\end{equation}

corresponds to $100\%$ gross exposure, whereas

\begin{equation}
G_t=2
\end{equation}

corresponds to $200\%$ gross exposure.

A market-neutral portfolio that is $100\%$ long and $100\%$ short satisfies

\begin{equation}
G_t
=
1+1
=
2.
\end{equation}


% ------------------------------------------------------------
\subsection{Net Exposure}
\label{subsec:net_exposure}
% ------------------------------------------------------------

Net exposure measures the signed total portfolio exposure:

\begin{equation}
\boxed{
N_t
=
\sum_{i=1}^{N_t}
w_{i,t}
=
\mathbf{1}^{\top}\mathbf{w}_t.
}
\end{equation}

A dollar-neutral portfolio satisfies

\begin{equation}
\boxed{
\mathbf{1}^{\top}\mathbf{w}_t
=
0.
}
\end{equation}

More generally,

\begin{equation}
N_{\min}
\leq
\mathbf{1}^{\top}\mathbf{w}_t
\leq
N_{\max}.
\end{equation}


\subsubsection{Gross vs.\ Net Exposure}

Gross and net exposure measure different concepts.

For example, suppose

\begin{equation}
G_t=2
\end{equation}

and

\begin{equation}
N_t=0.
\end{equation}

The portfolio may be $100\%$ long and $100\%$ short.

By contrast,

\begin{equation}
G_t=1
\qquad
\text{and}
\qquad
N_t=1
\end{equation}

corresponds to a fully invested long-only portfolio.

Thus,

\begin{equation}
\boxed{
\text{Gross Exposure}
\neq
\text{Net Exposure}.
}
\end{equation}


% ------------------------------------------------------------
\subsection{Long and Short Budgets}
\label{subsec:long_short_budgets}
% ------------------------------------------------------------

Define the long-leg exposure as

\begin{equation}
\boxed{
G_t^{L}
=
\sum_{i=1}^{N_t}
\max(w_{i,t},0),
}
\end{equation}

and the absolute short-leg exposure as

\begin{equation}
\boxed{
G_t^{S}
=
\sum_{i=1}^{N_t}
\max(-w_{i,t},0).
}
\end{equation}

Then

\begin{equation}
\boxed{
G_t
=
G_t^{L}
+
G_t^{S},
}
\end{equation}

while

\begin{equation}
\boxed{
N_t
=
G_t^{L}
-
G_t^{S}.
}
\end{equation}

Therefore,

\begin{equation}
\boxed{
G_t^{L}
=
\frac{
G_t+N_t
}{2},
}
\end{equation}

and

\begin{equation}
\boxed{
G_t^{S}
=
\frac{
G_t-N_t
}{2}.
}
\end{equation}

This relation is particularly useful because gross and net exposure uniquely
determine long and short budgets.


\subsubsection{Explicit Leg Budgets}

The portfolio may impose directly

\begin{equation}
G_t^{L}
=
G_{\mathrm{target}}^{L},
\end{equation}

and

\begin{equation}
G_t^{S}
=
G_{\mathrm{target}}^{S}.
\end{equation}

For example,

\begin{equation}
G_{\mathrm{target}}^{L}
=
G_{\mathrm{target}}^{S}
=
1
\end{equation}

implies

\begin{equation}
G_t=2,
\qquad
N_t=0.
\end{equation}


% ------------------------------------------------------------
\subsection{Leverage Constraints}
\label{subsec:leverage_constraints}
% ------------------------------------------------------------

Leverage describes portfolio exposure relative to NAV.

In weight space, gross exposure is often used as a direct leverage measure:

\begin{equation}
\boxed{
L_t
=
\|\mathbf{w}_t\|_1.
}
\end{equation}

A leverage constraint is therefore

\begin{equation}
\boxed{
\|\mathbf{w}_t\|_1
\leq
L_{\max}.
}
\end{equation}

However, economic leverage may also depend on the instruments used.

For example, derivatives, futures, swaps, and financed positions may have
different cash requirements despite generating similar economic exposures.

Thus, the backtester should distinguish where necessary between

\begin{equation}
\boxed{
\text{Economic Exposure}
\qquad\text{and}\qquad
\text{Balance-Sheet or Margin Usage}.
}
\end{equation}

For a cash-equity backtester, gross exposure is generally a sufficient first
representation.


% ------------------------------------------------------------
\subsection{Position Limits}
\label{subsec:position_limits}
% ------------------------------------------------------------

Individual-security limits prevent the optimizer from concentrating too much
capital in a single position.

For each security,

\begin{equation}
\boxed{
w_{i,t}^{\min}
\leq
w_{i,t}
\leq
w_{i,t}^{\max}.
}
\end{equation}

A symmetric constraint is

\begin{equation}
\boxed{
|w_{i,t}|
\leq
w_{\max}.
}
\end{equation}


\subsubsection{Asymmetric Long and Short Limits}

Long and short limits need not be identical.

For example,

\begin{equation}
-w_{\max}^{S}
\leq
w_{i,t}
\leq
w_{\max}^{L}.
\end{equation}

This may reflect different liquidity, borrow, or risk constraints on the short
side.


\subsubsection{Security-Specific Limits}

Position bounds may depend on security characteristics:

\begin{equation}
\boxed{
w_{i,t}^{\max}
=
g
\left(
\mathrm{ADV}_{i,t},
\mathrm{MCap}_{i,t},
\sigma_{i,t},
\mathrm{Borrow}_{i,t},
\ldots
\right).
}
\end{equation}

Thus, liquid large-cap securities may support larger positions than illiquid
small-cap securities.


% ------------------------------------------------------------
\subsection{Beta Neutrality}
\label{subsec:portfolio_beta_neutrality}
% ------------------------------------------------------------

Let

\begin{equation}
\boldsymbol{\beta}_t
=
\begin{bmatrix}
\beta_{1,t}\\
\vdots\\
\beta_{N_t,t}
\end{bmatrix}
\in
\mathbb{R}^{N_t}
\end{equation}

denote security-level market betas.

Portfolio beta is

\begin{equation}
\boxed{
\beta_{p,t}
=
\boldsymbol{\beta}_t^{\top}
\mathbf{w}_t.
}
\end{equation}

Exact beta neutrality requires

\begin{equation}
\boxed{
\boldsymbol{\beta}_t^{\top}
\mathbf{w}_t
=
0.
}
\end{equation}

A tolerance band may instead be imposed:

\begin{equation}
\boxed{
-\beta_{\max}
\leq
\boldsymbol{\beta}_t^{\top}
\mathbf{w}_t
\leq
\beta_{\max}.
}
\end{equation}


\subsubsection{Signal vs.\ Portfolio Beta Neutrality}

As discussed previously,

\begin{equation}
\boldsymbol{\beta}_t^\top
\mathbf{s}_t^{\mathrm{neutral}}
=
0
\end{equation}

does not necessarily imply

\begin{equation}
\boldsymbol{\beta}_t^\top
\mathbf{w}_t
=
0.
\end{equation}

Therefore,

\begin{equation}
\boxed{
\text{Signal-Level Neutrality}
\neq
\text{Portfolio-Level Neutrality}.
}
\end{equation}

If portfolio beta neutrality is a mandate requirement, it should be imposed
directly on the portfolio weights.


% ------------------------------------------------------------
\subsection{Industry and Sector Neutrality}
\label{subsec:portfolio_industry_neutrality}
% ------------------------------------------------------------

Let

\begin{equation}
\mathbf{D}_t
\in
\mathbb{R}^{N_t\times G}
\end{equation}

denote the industry-exposure matrix, where $G$ is the number of industries.

Portfolio industry exposures are

\begin{equation}
\boxed{
\mathbf{e}_t^{\mathrm{industry}}
=
\mathbf{D}_t^\top
\mathbf{w}_t
\in
\mathbb{R}^{G}.
}
\end{equation}

Exact industry neutrality requires

\begin{equation}
\boxed{
\mathbf{D}_t^\top
\mathbf{w}_t
=
\mathbf{0}_{G}.
}
\end{equation}


\subsubsection{Industry Exposure Bands}

Exact neutrality may be unnecessarily restrictive.

Instead,

\begin{equation}
-\mathbf{u}_t
\leq
\mathbf{D}_t^\top
\mathbf{w}_t
\leq
\mathbf{u}_t,
\end{equation}

where

\begin{equation}
\mathbf{u}_t
\in
\mathbb{R}_{+}^{G}
\end{equation}

contains maximum absolute industry exposures.


\subsubsection{Benchmark-Relative Constraints}

For benchmark-relative portfolios, sector exposure may be constrained relative
to benchmark weights.

Let

\begin{equation}
\mathbf{w}_t^{B}
\end{equation}

denote benchmark weights.

Active weights are

\begin{equation}
\mathbf{a}_t
=
\mathbf{w}_t
-
\mathbf{w}_t^{B}.
\end{equation}

Industry active exposure is

\begin{equation}
\boxed{
\mathbf{D}_t^\top
\mathbf{a}_t.
}
\end{equation}

A benchmark-relative sector constraint is therefore

\begin{equation}
-\mathbf{u}
\leq
\mathbf{D}_t^\top
\left(
\mathbf{w}_t-\mathbf{w}_t^{B}
\right)
\leq
\mathbf{u}.
\end{equation}


% ------------------------------------------------------------
\subsection{Country Neutrality}
\label{subsec:portfolio_country_neutrality}
% ------------------------------------------------------------

Let

\begin{equation}
\mathbf{C}_t
\in
\mathbb{R}^{N_t\times J}
\end{equation}

denote country exposures.

The portfolio country-exposure vector is

\begin{equation}
\boxed{
\mathbf{e}_t^{\mathrm{country}}
=
\mathbf{C}_t^\top
\mathbf{w}_t.
}
\end{equation}

Exact country neutrality requires

\begin{equation}
\boxed{
\mathbf{C}_t^\top
\mathbf{w}_t
=
\mathbf{0}_{J}.
}
\end{equation}

Alternatively, exposure bands may be imposed:

\begin{equation}
-\mathbf{c}_{\max}
\leq
\mathbf{C}_t^\top
\mathbf{w}_t
\leq
\mathbf{c}_{\max}.
\end{equation}

These constraints are particularly important in global equity portfolios where
otherwise attractive signals may generate unintended country concentration.


% ------------------------------------------------------------
\subsection{Style Factor Neutrality}
\label{subsec:portfolio_style_factor_neutrality}
% ------------------------------------------------------------

Let

\begin{equation}
\mathbf{F}_t
\in
\mathbb{R}^{N_t\times K}
\end{equation}

denote security exposures to $K$ style factors.

Possible factors include:

\begin{itemize}
    \item size;
    \item value;
    \item momentum;
    \item quality;
    \item volatility;
    \item liquidity;
    \item leverage;
    \item growth.
\end{itemize}

Portfolio factor exposures are

\begin{equation}
\boxed{
\mathbf{f}_{p,t}
=
\mathbf{F}_t^\top
\mathbf{w}_t
\in
\mathbb{R}^{K}.
}
\end{equation}

Exact neutrality requires

\begin{equation}
\boxed{
\mathbf{F}_t^\top
\mathbf{w}_t
=
\mathbf{0}_{K}.
}
\end{equation}

More commonly,

\begin{equation}
-\mathbf{f}_{\max}
\leq
\mathbf{F}_t^\top
\mathbf{w}_t
\leq
\mathbf{f}_{\max}.
\end{equation}


\subsubsection{Joint Exposure Constraints}

Beta, industry, country, and style exposures can be combined into one matrix:

\begin{equation}
\boxed{
\mathbf{B}_t^{\mathrm{port}}
=
\begin{bmatrix}
\boldsymbol{\beta}_t
&
\mathbf{D}_t
&
\mathbf{C}_t
&
\mathbf{F}_t
\end{bmatrix}
\in
\mathbb{R}^{N_t\times K_{\mathrm{tot}}}.
}
\end{equation}

Then a generic exposure constraint is

\begin{equation}
\boxed{
\mathbf{l}_t
\leq
\left(
\mathbf{B}_t^{\mathrm{port}}
\right)^\top
\mathbf{w}_t
\leq
\mathbf{u}_t.
}
\end{equation}


% ------------------------------------------------------------
\subsection{Tracking Error Constraints}
\label{subsec:tracking_error_constraints}
% ------------------------------------------------------------

For benchmark-relative strategies, portfolio risk may be controlled through
tracking error.

Let

\begin{equation}
\mathbf{w}_t^{B}
\in
\mathbb{R}^{N_t}
\end{equation}

denote benchmark weights.

Active weights are

\begin{equation}
\boxed{
\mathbf{a}_t
=
\mathbf{w}_t
-
\mathbf{w}_t^{B}.
}
\end{equation}

Ex-ante tracking-error variance is

\begin{equation}
\boxed{
TE_t^2
=
\mathbf{a}_t^\top
\widehat{\boldsymbol{\Sigma}}_t
\mathbf{a}_t.
}
\end{equation}

Tracking error is therefore

\begin{equation}
\boxed{
TE_t
=
\sqrt{
\mathbf{a}_t^\top
\widehat{\boldsymbol{\Sigma}}_t
\mathbf{a}_t
}.
}
\end{equation}

A maximum tracking-error constraint is

\begin{equation}
\boxed{
\mathbf{a}_t^\top
\widehat{\boldsymbol{\Sigma}}_t
\mathbf{a}_t
\leq
TE_{\max}^{2}.
}
\end{equation}


\subsubsection{Active vs.\ Absolute Risk}

Tracking error measures risk relative to the benchmark, not absolute portfolio
risk.

Thus,

\begin{equation}
\boxed{
\text{Portfolio Volatility}
\neq
\text{Tracking Error}.
}
\end{equation}

A portfolio may have high absolute volatility while maintaining low tracking
error if its exposures closely resemble the benchmark.


% ------------------------------------------------------------
\subsection{Turnover Constraints}
\label{subsec:portfolio_turnover_constraints}
% ------------------------------------------------------------

Let

\begin{equation}
\mathbf{w}_{t^-}
\end{equation}

denote portfolio weights immediately before rebalancing.

The desired weight change is

\begin{equation}
\Delta\mathbf{w}_t
=
\mathbf{w}_t
-
\mathbf{w}_{t^-}.
\end{equation}

Total traded notional in weight space is

\begin{equation}
\boxed{
TO_t^{\mathrm{gross}}
=
\sum_{i=1}^{N_t}
|\Delta w_{i,t}|.
}
\end{equation}

A hard turnover constraint is

\begin{equation}
\boxed{
\sum_{i=1}^{N_t}
|w_{i,t}-w_{i,t^-}|
\leq
TO_{\max}.
}
\end{equation}

If the half-turnover convention is used,

\begin{equation}
TO_t^{1/2}
=
\frac{1}{2}
\sum_i
|w_{i,t}-w_{i,t^-}|,
\end{equation}

the corresponding limit must be defined consistently.


\subsubsection{Security-Level Trade Limits}

One may also impose

\begin{equation}
\boxed{
|w_{i,t}-w_{i,t^-}|
\leq
\Delta w_{i,t}^{\max}.
}
\end{equation}

This prevents excessively large trades in individual securities even if total
portfolio turnover remains below its overall limit.


% ------------------------------------------------------------
\subsection{Liquidity and ADV Constraints}
\label{subsec:liquidity_adv_constraints}
% ------------------------------------------------------------

Portfolio weights must be compatible with the amount of liquidity available in
the market.

Let

\begin{equation}
A_t
\end{equation}

denote portfolio NAV.

The desired currency notional trade in security $i$ is

\begin{equation}
\boxed{
Q_{i,t}
=
A_t
|w_{i,t}-w_{i,t^-}|.
}
\end{equation}

Let

\begin{equation}
\mathrm{ADV}_{i,t}
\end{equation}

denote Average Daily Volume in the same currency.

The participation rate is

\begin{equation}
\boxed{
PR_{i,t}
=
\frac{
A_t
|w_{i,t}-w_{i,t^-}|
}{
\mathrm{ADV}_{i,t}
}.
}
\end{equation}

A maximum participation-rate constraint is

\begin{equation}
\boxed{
PR_{i,t}
\leq
PR_{\max}.
}
\end{equation}

Equivalently,

\begin{equation}
\boxed{
|w_{i,t}-w_{i,t^-}|
\leq
PR_{\max}
\frac{
\mathrm{ADV}_{i,t}
}{
A_t
}.
}
\end{equation}


\subsubsection{Position Capacity Constraint}

Liquidity may also constrain total position size rather than only daily trade
size.

Suppose the maximum desired position represents no more than $H$ days of ADV:

\begin{equation}
A_t|w_{i,t}|
\leq
H
\,
\mathrm{ADV}_{i,t}.
\end{equation}

Then

\begin{equation}
\boxed{
|w_{i,t}|
\leq
H
\frac{
\mathrm{ADV}_{i,t}
}{
A_t
}.
}
\end{equation}


\subsubsection{AUM Dependence}

Because

\begin{equation}
|w_{i,t}-w_{i,t^-}|_{\max}
\propto
\frac{1}{A_t},
\end{equation}

the feasible portfolio becomes more restrictive as AUM increases.

Thus,

\begin{equation}
\boxed{
A_t\uparrow
\quad\Rightarrow\quad
\text{Liquidity Constraints Tighten}.
}
\end{equation}

This establishes the direct link between portfolio constraints and strategy
capacity.


% ------------------------------------------------------------
\subsection{Short Availability Constraints}
\label{subsec:short_availability_constraints}
% ------------------------------------------------------------

A negative desired weight is implementable only if the corresponding security
can be borrowed.

Let

\begin{equation}
H_{i,t}^{\mathrm{short}}
\in
\{0,1\}
\end{equation}

denote short availability.

If

\begin{equation}
H_{i,t}^{\mathrm{short}}=0,
\end{equation}

then the portfolio must satisfy

\begin{equation}
\boxed{
w_{i,t}
\geq
0.
}
\end{equation}


\subsubsection{Borrow Quantity Constraint}

Suppose

\begin{equation}
B_{i,t}^{\$}
\end{equation}

denotes the maximum currency notional available to borrow.

Then

\begin{equation}
A_t
\max(-w_{i,t},0)
\leq
B_{i,t}^{\$}.
\end{equation}

Equivalently,

\begin{equation}
\boxed{
w_{i,t}
\geq
-
\frac{
B_{i,t}^{\$}
}{
A_t
}.
}
\end{equation}


\subsubsection{Borrow-Cost Constraints}

If the borrow fee is

\begin{equation}
b_{i,t},
\end{equation}

the strategy may exclude short positions when

\begin{equation}
b_{i,t}>b_{\max}.
\end{equation}

This can be expressed as

\begin{equation}
b_{i,t}>b_{\max}
\quad\Rightarrow\quad
w_{i,t}\geq0.
\end{equation}

Alternatively, borrow cost may remain in the optimization objective rather than
being imposed as a hard exclusion.

Thus,

\begin{equation}
\boxed{
\begin{aligned}
\text{Borrow Availability}
&\rightarrow
\text{Hard Feasibility Constraint},
\\
\text{Borrow Cost}
&\rightarrow
\text{Constraint or Cost Penalty}.
\end{aligned}
}
\end{equation}


% ------------------------------------------------------------
\subsubsection{Combined Feasible Set}
% ------------------------------------------------------------

The complete feasible set can now be represented schematically as

\begin{equation}
\boxed{
\begin{aligned}
\mathcal{W}_t
=
\Bigg\{
\mathbf{w}\in\mathbb{R}^{N_t}
:\;
&
\|\mathbf{w}\|_1
\leq
G_{\max},
\\
&
N_{\min}
\leq
\mathbf{1}^{\top}\mathbf{w}
\leq
N_{\max},
\\
&
\mathbf{w}^{\min}
\leq
\mathbf{w}
\leq
\mathbf{w}^{\max},
\\
&
\mathbf{l}_B
\leq
\mathbf{B}_t^\top\mathbf{w}
\leq
\mathbf{u}_B,
\\
&
\left(
\mathbf{w}-\mathbf{w}^{B}
\right)^\top
\widehat{\boldsymbol{\Sigma}}_t
\left(
\mathbf{w}-\mathbf{w}^{B}
\right)
\leq
TE_{\max}^{2},
\\
&
\left\|
\mathbf{w}
-
\mathbf{w}_{t^-}
\right\|_1
\leq
TO_{\max},
\\
&
A_t
|w_i-w_{i,t^-}|
\leq
PR_{\max}
\mathrm{ADV}_{i,t},
\quad
\forall i,
\\
&
\text{short-availability constraints hold}
\Bigg\}.
\end{aligned}
}
\end{equation}

Not every strategy requires all of these constraints.

The feasible set should contain only restrictions that correspond to actual
economic, regulatory, mandate, or implementation requirements.


% ------------------------------------------------------------
\subsubsection{Constraint Interactions and Feasibility}
% ------------------------------------------------------------

Constraints cannot always be considered independently.

For example, simultaneously requiring

\begin{equation}
\mathbf{1}^{\top}\mathbf{w}=0,
\end{equation}

\begin{equation}
\boldsymbol{\beta}_t^\top\mathbf{w}=0,
\end{equation}

\begin{equation}
\mathbf{D}_t^\top\mathbf{w}=0,
\end{equation}

and very tight position limits may leave no feasible portfolio.

Thus,

\begin{equation}
\boxed{
\mathcal{W}_t
=
\varnothing
}
\end{equation}

is possible.

The backtester must therefore explicitly verify feasibility before solving the
optimization problem.


\subsubsection{Constraint Hierarchy}

In practical implementations, it is useful to distinguish between:

\begin{equation}
\boxed{
\begin{aligned}
\text{Hard Constraints}
&:
\quad
\text{must never be violated},
\\
\text{Soft Constraints}
&:
\quad
\text{may be relaxed at a cost}.
\end{aligned}
}
\end{equation}

Hard constraints may include:

\begin{itemize}
    \item legal shortability;
    \item mandate leverage limits;
    \item absolute position limits.
\end{itemize}

Soft constraints may include:

\begin{itemize}
    \item preferred factor-neutrality bands;
    \item turnover targets;
    \item benchmark-relative sector limits.
\end{itemize}

A soft constraint can be incorporated using a slack variable.

For example,

\begin{equation}
-\beta_{\max}-\xi_t
\leq
\boldsymbol{\beta}_t^\top\mathbf{w}_t
\leq
\beta_{\max}+\xi_t,
\end{equation}

with

\begin{equation}
\xi_t\geq0,
\end{equation}

and a penalty

\begin{equation}
-\lambda_{\xi}\xi_t
\end{equation}

in the optimization objective.

This preserves feasibility while penalizing departures from the preferred
constraint.


% ------------------------------------------------------------
\subsubsection{Constraint Diagnostics}
% ------------------------------------------------------------

For each optimized portfolio, the backtester should record:

\begin{itemize}
    \item gross exposure;
    \item net exposure;
    \item long and short gross budgets;
    \item maximum absolute position;
    \item portfolio beta;
    \item industry exposures;
    \item country exposures;
    \item style-factor exposures;
    \item ex-ante tracking error;
    \item portfolio turnover;
    \item maximum ADV participation rate;
    \item number of short-constrained securities.
\end{itemize}

It is also useful to identify which constraints are binding.

For constraint

\begin{equation}
g_j(\mathbf{w}_t)\leq c_j,
\end{equation}

the constraint is binding when

\begin{equation}
\boxed{
g_j(\mathbf{w}_t)
\approx
c_j.
}
\end{equation}

A large number of persistently binding constraints may indicate that the
portfolio specification is excessively restrictive or that the raw signal is
poorly aligned with the investment mandate.


% ------------------------------------------------------------
\subsubsection{Position in the Portfolio-Construction Pipeline}
% ------------------------------------------------------------

At this stage, the portfolio engine has defined:

\begin{equation}
\boxed{
\begin{aligned}
\mathbf{s}_t^{\mathrm{final}}
&\in
\mathbb{R}^{N_t},
&&
\text{signal},
\\
\widehat{\boldsymbol{\Sigma}}_t
&\in
\mathbb{R}^{N_t\times N_t},
&&
\text{risk model},
\\
\mathcal{J}_t(\mathbf{w})
&&&
\text{optimization objective},
\\
\mathcal{W}_t
&&&
\text{feasible portfolio set}.
\end{aligned}
}
\end{equation}

The target portfolio is therefore fully characterized by

\begin{equation}
\boxed{
\mathbf{w}_t^{\mathrm{target}}
=
\underset{\mathbf{w}\in\mathcal{W}_t}{\arg\max}
\;
\mathcal{J}_t(\mathbf{w}).
}
\end{equation}

The next stages of the backtester determine how these target weights are scaled,
rebalanced, converted into executable trades, and ultimately transformed into
realized portfolio returns.

% ============================================================
\section{Portfolio-Level Risk Scaling}
\label{sec:portfolio_level_risk_scaling}
% ============================================================

The portfolio optimization and constraint framework developed previously
produces the target portfolio

\begin{equation}
\boxed{
\mathbf{w}_t^{\mathrm{target}}
=
\begin{bmatrix}
w_{1,t}^{\mathrm{target}}
\\
\vdots
\\
w_{N_t,t}^{\mathrm{target}}
\end{bmatrix}
\in
\mathbb{R}^{N_t}.
}
\end{equation}

These weights determine the desired relative allocation across securities after
accounting for the signal, risk model, optimization objective, and portfolio
constraints.

However, the risk of a portfolio constructed using the same methodology may
vary substantially over time as security volatilities and correlations change.

An optional additional step therefore consists of scaling the entire portfolio
to target a desired level of ex-ante volatility.

This operation is commonly referred to as \emph{volatility targeting},
\emph{volatility scaling}, or \emph{isovol scaling}.

The general transformation is

\begin{equation}
\boxed{
\mathbf{w}_t^{\mathrm{target}}
\quad
\xrightarrow{\text{portfolio-level risk scaling}}
\quad
\mathbf{w}_t^{\mathrm{final}}.
}
\end{equation}

Importantly, volatility targeting is not required in every strategy.

Portfolio risk may instead be controlled directly inside the optimization
problem through a risk constraint or an appropriate objective function. The
separate scaling procedure developed in this section is useful when the
portfolio-construction process determines a desired portfolio composition but
the final amount of risk allocated to that portfolio is controlled separately.


% ------------------------------------------------------------
\subsection{Portfolio Direction and Portfolio Scale}
\label{subsec:portfolio_direction_scale}
% ------------------------------------------------------------

It is useful to distinguish between the \emph{direction} of a portfolio and its
\emph{scale}.

The target weight vector

\begin{equation}
\mathbf{w}_t^{\mathrm{target}}
\in
\mathbb{R}^{N_t}
\end{equation}

determines the relative positions across securities.

Consider a positive scalar

\begin{equation}
\lambda_t>0.
\end{equation}

Scaling the portfolio gives

\begin{equation}
\boxed{
\mathbf{w}_t^{\mathrm{scaled}}
=
\lambda_t
\mathbf{w}_t^{\mathrm{target}}.
}
\end{equation}

For every pair of securities $i$ and $j$ with non-zero weights,

\begin{equation}
\frac{
w_{i,t}^{\mathrm{scaled}}
}{
w_{j,t}^{\mathrm{scaled}}
}
=
\frac{
\lambda_t w_{i,t}^{\mathrm{target}}
}{
\lambda_t w_{j,t}^{\mathrm{target}}
}
=
\frac{
w_{i,t}^{\mathrm{target}}
}{
w_{j,t}^{\mathrm{target}}
}.
\end{equation}

Thus, common positive scaling does not change relative portfolio composition.

It changes the amount of capital exposure assigned to that composition.

Conceptually,

\begin{equation}
\boxed{
\underbrace{
\mathbf{w}_t^{\mathrm{target}}
}_{\text{portfolio composition}}
\quad+\quad
\underbrace{
\lambda_t
}_{\text{portfolio scale}}
\quad\longrightarrow\quad
\underbrace{
\mathbf{w}_t^{\mathrm{final}}
}_{\text{final risk-scaled portfolio}}.
}
\end{equation}

This distinction is particularly useful for long--short strategies. The
portfolio-construction engine may determine which securities should be long or
short and their relative importance, while a separate risk-scaling layer
determines how aggressively the strategy should be deployed.


% ------------------------------------------------------------
\subsection{Volatility Targeting}
\label{subsec:portfolio_volatility_targeting}
% ------------------------------------------------------------

Let

\begin{equation}
\widehat{\boldsymbol{\Sigma}}_t
\in
\mathbb{R}^{N_t\times N_t}
\end{equation}

denote the ex-ante covariance matrix available at time $t$.

The forecast volatility of the target portfolio is

\begin{equation}
\boxed{
\widehat{\sigma}_{p,t}
=
\sqrt{
\left(
\mathbf{w}_t^{\mathrm{target}}
\right)^\top
\widehat{\boldsymbol{\Sigma}}_t
\mathbf{w}_t^{\mathrm{target}}
}.
}
\end{equation}

Let

\begin{equation}
\sigma^{\mathrm{target}}>0
\end{equation}

denote the desired portfolio volatility.

The isovol scaling coefficient is

\begin{equation}
\boxed{
\lambda_t^{\mathrm{vol}}
=
\frac{
\sigma^{\mathrm{target}}
}{
\widehat{\sigma}_{p,t}
}.
}
\end{equation}

The volatility-scaled portfolio is therefore

\begin{equation}
\boxed{
\mathbf{w}_t^{\mathrm{vol}}
=
\lambda_t^{\mathrm{vol}}
\mathbf{w}_t^{\mathrm{target}}.
}
\end{equation}


\subsubsection{Why the Scaling Works}

Portfolio variance after scaling is

\begin{equation}
\begin{aligned}
\widehat{\sigma}_{p,t}^{2,\mathrm{vol}}
&=
\left(
\lambda_t^{\mathrm{vol}}
\mathbf{w}_t^{\mathrm{target}}
\right)^\top
\widehat{\boldsymbol{\Sigma}}_t
\left(
\lambda_t^{\mathrm{vol}}
\mathbf{w}_t^{\mathrm{target}}
\right)
\\
&=
\left(
\lambda_t^{\mathrm{vol}}
\right)^2
\left(
\mathbf{w}_t^{\mathrm{target}}
\right)^\top
\widehat{\boldsymbol{\Sigma}}_t
\mathbf{w}_t^{\mathrm{target}}
\\
&=
\left(
\lambda_t^{\mathrm{vol}}
\right)^2
\widehat{\sigma}_{p,t}^{2}.
\end{aligned}
\end{equation}

Hence,

\begin{equation}
\widehat{\sigma}_{p,t}^{\mathrm{vol}}
=
\lambda_t^{\mathrm{vol}}
\widehat{\sigma}_{p,t}.
\end{equation}

Substituting

\begin{equation}
\lambda_t^{\mathrm{vol}}
=
\frac{
\sigma^{\mathrm{target}}
}{
\widehat{\sigma}_{p,t}
}
\end{equation}

gives

\begin{equation}
\boxed{
\widehat{\sigma}_{p,t}^{\mathrm{vol}}
=
\sigma^{\mathrm{target}}.
}
\end{equation}

Thus, under the covariance forecast used for scaling, the resulting portfolio
has exactly the desired ex-ante volatility.


\subsubsection{Time Variation in the Scaling Coefficient}

The scaling coefficient varies through time:

\begin{equation}
\lambda_t^{\mathrm{vol}}
=
\frac{
\sigma^{\mathrm{target}}
}{
\widehat{\sigma}_{p,t}
}.
\end{equation}

Therefore,

\begin{equation}
\widehat{\sigma}_{p,t}
<
\sigma^{\mathrm{target}}
\quad\Rightarrow\quad
\lambda_t^{\mathrm{vol}}>1,
\end{equation}

so portfolio exposure is increased.

Conversely,

\begin{equation}
\widehat{\sigma}_{p,t}
>
\sigma^{\mathrm{target}}
\quad\Rightarrow\quad
\lambda_t^{\mathrm{vol}}<1,
\end{equation}

so portfolio exposure is reduced.

Hence,

\begin{equation}
\boxed{
\text{Low Forecast Risk}
\Rightarrow
\text{Higher Exposure},
}
\end{equation}

while

\begin{equation}
\boxed{
\text{High Forecast Risk}
\Rightarrow
\text{Lower Exposure}.
}
\end{equation}


\subsubsection{Point-in-Time Volatility Estimation}

The volatility forecast used in the scaling coefficient must be strictly
point-in-time.

At time $t$, the backtester must compute

\begin{equation}
\widehat{\sigma}_{p,t}
\end{equation}

using only information available at or before the portfolio decision time.

For example, if the covariance matrix is estimated from historical returns,

\begin{equation}
\widehat{\boldsymbol{\Sigma}}_t
=
f
\left(
\mathbf{r}_{t-1},
\mathbf{r}_{t-2},
\ldots
\right),
\end{equation}

and not from future realized returns.

Using future realized volatility to determine $\lambda_t^{\mathrm{vol}}$ would
introduce look-ahead bias.


\subsubsection{Forecast vs.\ Realized Volatility}

Volatility targeting guarantees the target only with respect to the risk
estimate used at time $t$.

In general,

\begin{equation}
\boxed{
\widehat{\sigma}_{p,t}^{\mathrm{vol}}
=
\sigma^{\mathrm{target}}
\quad\not\Rightarrow\quad
\sigma_{p,t\rightarrow t+1}^{\mathrm{realized}}
=
\sigma^{\mathrm{target}}.
}
\end{equation}

The realized covariance structure over the subsequent holding period may differ
from the covariance forecast.

Volatility targeting should therefore be interpreted as an \emph{ex-ante risk
control mechanism}, not as a guarantee of future realized volatility.


% ------------------------------------------------------------
\subsection{Leverage Caps and Scaling Limits}
\label{subsec:leverage_caps_scaling}
% ------------------------------------------------------------

Pure volatility targeting may imply very large scaling coefficients when
forecast volatility is unusually low.

For example,

\begin{equation}
\widehat{\sigma}_{p,t}
\rightarrow 0
\quad\Rightarrow\quad
\lambda_t^{\mathrm{vol}}
\rightarrow\infty.
\end{equation}

Such leverage is generally economically unrealistic and potentially dangerous.

A maximum scaling coefficient may therefore be imposed:

\begin{equation}
\boxed{
\lambda_t
=
\min
\left(
\lambda_t^{\mathrm{vol}},
\lambda_{\max}
\right).
}
\end{equation}

The resulting portfolio is

\begin{equation}
\mathbf{w}_t^{\mathrm{scaled}}
=
\lambda_t
\mathbf{w}_t^{\mathrm{target}}.
\end{equation}


\subsubsection{Gross Exposure Cap}

Scaling also changes gross exposure.

Before scaling,

\begin{equation}
G_t^{\mathrm{target}}
=
\left\|
\mathbf{w}_t^{\mathrm{target}}
\right\|_1.
\end{equation}

After multiplying all weights by $\lambda_t$,

\begin{equation}
\begin{aligned}
G_t^{\mathrm{scaled}}
&=
\left\|
\lambda_t
\mathbf{w}_t^{\mathrm{target}}
\right\|_1
\\
&=
\lambda_t
\left\|
\mathbf{w}_t^{\mathrm{target}}
\right\|_1.
\end{aligned}
\end{equation}

If gross exposure may not exceed $G_{\max}$, then

\begin{equation}
\lambda_t
G_t^{\mathrm{target}}
\leq
G_{\max}.
\end{equation}

Therefore,

\begin{equation}
\boxed{
\lambda_t^{\mathrm{gross}}
=
\frac{
G_{\max}
}{
\left\|
\mathbf{w}_t^{\mathrm{target}}
\right\|_1
}.
}
\end{equation}

Combining the volatility target, leverage cap, and gross-exposure constraint
gives

\begin{equation}
\boxed{
\lambda_t
=
\min
\left(
\frac{
\sigma^{\mathrm{target}}
}{
\widehat{\sigma}_{p,t}
},
\lambda_{\max},
\frac{
G_{\max}
}{
\left\|
\mathbf{w}_t^{\mathrm{target}}
\right\|_1
}
\right).
}
\end{equation}

If the gross-exposure constraint is binding, the strategy cannot attain its
desired volatility target without violating the mandate.

In that case,

\begin{equation}
\boxed{
\sigma_{p,t}^{\mathrm{final}}
<
\sigma^{\mathrm{target}}
}
\end{equation}

may be the correct portfolio outcome.


% ------------------------------------------------------------
\subsection{Interaction with Portfolio Constraints}
\label{subsec:risk_scaling_constraints}
% ------------------------------------------------------------

Portfolio-level scaling occurs after the constrained portfolio has been
constructed. It is therefore necessary to determine which constraints are
preserved by common scaling and which may be violated.


\subsubsection{Homogeneous Equality Constraints}

Consider a linear neutrality constraint

\begin{equation}
\mathbf{B}_t^\top
\mathbf{w}_t^{\mathrm{target}}
=
\mathbf{0}.
\end{equation}

After scaling,

\begin{equation}
\begin{aligned}
\mathbf{B}_t^\top
\left(
\lambda_t
\mathbf{w}_t^{\mathrm{target}}
\right)
&=
\lambda_t
\mathbf{B}_t^\top
\mathbf{w}_t^{\mathrm{target}}
\\
&=
\mathbf{0}.
\end{aligned}
\end{equation}

Therefore,

\begin{equation}
\boxed{
\mathbf{B}_t^\top
\mathbf{w}_t^{\mathrm{target}}
=
\mathbf{0}
\quad\Rightarrow\quad
\mathbf{B}_t^\top
\mathbf{w}_t^{\mathrm{scaled}}
=
\mathbf{0}.
}
\end{equation}

Consequently, common portfolio scaling preserves exact zero-exposure
constraints such as:

\begin{itemize}
    \item dollar neutrality;
    \item beta neutrality;
    \item industry neutrality;
    \item country neutrality;
    \item style-factor neutrality.
\end{itemize}

For example, if

\begin{equation}
\boldsymbol{\beta}_t^\top
\mathbf{w}_t^{\mathrm{target}}
=
0,
\end{equation}

then

\begin{equation}
\boldsymbol{\beta}_t^\top
\mathbf{w}_t^{\mathrm{scaled}}
=
\lambda_t
\boldsymbol{\beta}_t^\top
\mathbf{w}_t^{\mathrm{target}}
=
0.
\end{equation}


\subsubsection{Non-Zero Exposure Targets}

The same result does not hold for constraints requiring a specific non-zero
exposure.

Suppose

\begin{equation}
\mathbf{b}_t^\top
\mathbf{w}_t^{\mathrm{target}}
=
c,
\qquad
c\neq0.
\end{equation}

After scaling,

\begin{equation}
\mathbf{b}_t^\top
\mathbf{w}_t^{\mathrm{scaled}}
=
\lambda_t c.
\end{equation}

Therefore,

\begin{equation}
\boxed{
c\neq0
\quad\Rightarrow\quad
\text{common scaling generally changes the targeted exposure}.
}
\end{equation}


\subsubsection{Inequality Constraints}

Scaling may also violate upper or lower bounds.

For example, if

\begin{equation}
|w_{i,t}^{\mathrm{target}}|
\leq
w_{i,t}^{\max},
\end{equation}

then after scaling,

\begin{equation}
|w_{i,t}^{\mathrm{scaled}}|
=
\lambda_t
|w_{i,t}^{\mathrm{target}}|.
\end{equation}

For $\lambda_t>1$, it is possible that

\begin{equation}
|w_{i,t}^{\mathrm{scaled}}|
>
w_{i,t}^{\max}.
\end{equation}

The same issue applies to:

\begin{itemize}
    \item gross exposure limits;
    \item individual position limits;
    \item liquidity and capacity limits;
    \item leverage limits;
    \item non-zero factor-exposure bands.
\end{itemize}

Thus, volatility scaling cannot be considered independently of the feasible set.


\subsubsection{Maximum Feasible Scaling Coefficient}

More generally, let

\begin{equation}
\lambda_t^{\mathrm{feasible}}
\end{equation}

denote the largest positive scalar such that

\begin{equation}
\lambda_t^{\mathrm{feasible}}
\mathbf{w}_t^{\mathrm{target}}
\end{equation}

continues to satisfy all relevant portfolio constraints.

Then the operational scaling coefficient can be written as

\begin{equation}
\boxed{
\lambda_t
=
\min
\left(
\lambda_t^{\mathrm{vol}},
\lambda_t^{\mathrm{feasible}}
\right).
}
\end{equation}

This formulation provides a general interpretation of constrained volatility
targeting:

\begin{equation}
\boxed{
\text{take as much exposure as required by the risk target,
subject to remaining feasible}.
}
\end{equation}


% ------------------------------------------------------------
\subsection{Final Target Weights}
\label{subsec:final_target_weights}
% ------------------------------------------------------------

After portfolio-level risk scaling, the final desired weights are

\begin{equation}
\boxed{
\mathbf{w}_t^{\mathrm{final}}
=
\lambda_t
\mathbf{w}_t^{\mathrm{target}}.
}
\end{equation}

where

\begin{equation}
\mathbf{w}_t^{\mathrm{final}}
\in
\mathbb{R}^{N_t}.
\end{equation}

If no separate portfolio-level risk scaling is required, then

\begin{equation}
\lambda_t=1,
\end{equation}

and therefore

\begin{equation}
\boxed{
\mathbf{w}_t^{\mathrm{final}}
=
\mathbf{w}_t^{\mathrm{target}}.
}
\end{equation}

The distinction between the two objects is therefore:

\begin{equation}
\boxed{
\begin{aligned}
\mathbf{w}_t^{\mathrm{target}}
&:
&&
\text{portfolio obtained from mapping, optimization, and constraints},
\\
\mathbf{w}_t^{\mathrm{final}}
&:
&&
\text{portfolio after optional portfolio-level risk scaling}.
\end{aligned}
}
\end{equation}

The complete portfolio-construction pipeline can now be summarized as

\begin{equation}
\boxed{
\begin{aligned}
\mathbf{s}_t^{\mathrm{final}}
&\rightarrow
\text{Signal-to-Weight Mapping}
\\
&\rightarrow
\text{Risk Model}
\\
&\rightarrow
\text{Optimization}
\\
&\rightarrow
\text{Portfolio Constraints}
\\
&\rightarrow
\mathbf{w}_t^{\mathrm{target}}
\\
&\rightarrow
\text{Optional Isovol Scaling}
\\
&\rightarrow
\mathbf{w}_t^{\mathrm{final}}.
\end{aligned}
}
\end{equation}

At this point, the portfolio-construction problem is complete.

The vector

\begin{equation}
\boxed{
\mathbf{w}_t^{\mathrm{final}}
}
\end{equation}

represents the final desired economic exposure of the strategy at time $t$.

The remaining stages of the backtester no longer determine which portfolio
\emph{should} be held. Instead, they determine how the existing portfolio is
moved toward these desired weights, how trades are executed and costed, and how
the resulting realized portfolio performance is measured.


% ============================================================
\part{Backtest Simulation and Implementation}
\label{part:backtest_simulation_implementation}
% ============================================================

The previous parts of the backtester determined the portfolio that the strategy
\emph{wants} to hold.

At each portfolio decision date, the portfolio-construction engine produces

\begin{equation}
\boxed{
\mathbf{w}_t^{\mathrm{final}}
\in
\mathbb{R}^{N_t},
}
\end{equation}

where $\mathbf{w}_t^{\mathrm{final}}$ denotes the final desired portfolio weights
after signal construction, portfolio optimization, constraints, and any optional
portfolio-level risk scaling.

The purpose of the present part is fundamentally different.

The backtester must now simulate how this desired portfolio would actually evolve
through time.

This requires modeling:

\begin{itemize}
    \item when portfolio decisions are made;
    \item when they can be executed;
    \item the portfolio already held before each rebalance;
    \item the trades required to reach the new target;
    \item execution prices and partial execution;
    \item transaction costs and market frictions;
    \item positions, cash, and mark-to-market valuation;
    \item realized portfolio profit and loss;
    \item the resulting net return series.
\end{itemize}

The simulation layer can therefore be summarized as

\begin{equation}
\boxed{
\begin{aligned}
\mathbf{w}_t^{\mathrm{final}}
&\rightarrow
\text{Rebalance Decision}
\\
&\rightarrow
\mathbf{w}_{t^-}
\rightarrow
\Delta\mathbf{w}_t
\\
&\rightarrow
\text{Orders}
\rightarrow
\text{Execution}
\\
&\rightarrow
\mathbf{w}_t^{\mathrm{implemented}}
\\
&\rightarrow
\text{Portfolio State and P\&L}
\\
&\rightarrow
R_t^{\mathrm{gross}}
\rightarrow
R_t^{\mathrm{net}}
\rightarrow
NAV_t.
\end{aligned}
}
\end{equation}

The distinction between desired and implemented weights is essential.

The portfolio-construction engine produces

\begin{equation}
\mathbf{w}_t^{\mathrm{final}},
\end{equation}

whereas the simulation engine ultimately determines

\begin{equation}
\boxed{
\mathbf{w}_t^{\mathrm{implemented}},
}
\end{equation}

the portfolio actually held after accounting for execution, tradability,
position rounding, partial fills, and other implementation effects.


% ============================================================
\section{Backtest Timing and Rebalancing}
\label{sec:backtest_timing_rebalancing}
% ============================================================

The first responsibility of the simulation engine is to define the chronological
relationship between information, portfolio decisions, execution, and realized
returns.

A backtest must ensure that the portfolio earning a return was fully determined
before that return became observable.

The general timing structure is

\begin{equation}
\boxed{
\text{Available Information}
\rightarrow
\text{Portfolio Decision}
\rightarrow
\text{Execution}
\rightarrow
\text{Holding Period}
\rightarrow
\text{Realized Return}.
}
\end{equation}


% ------------------------------------------------------------
\subsection{Decision, Execution, and Holding Timeline}
% ------------------------------------------------------------

Let

\begin{equation}
t_d
\end{equation}

denote the portfolio decision time,

\begin{equation}
t_e
\end{equation}

the execution time, and

\begin{equation}
[t_e,t_e')
\end{equation}

the subsequent holding interval.

A valid chronology requires

\begin{equation}
\boxed{
t_d
\leq
t_e
<
t_e'.
}
\end{equation}

The final desired portfolio is constructed using only information available at
$t_d$:

\begin{equation}
\boxed{
\mathbf{w}_{t_d}^{\mathrm{final}}
=
f
\left(
\mathcal{I}_{t_d}
\right),
}
\end{equation}

where $\mathcal{I}_{t_d}$ denotes the point-in-time information set.

The portfolio return must then be generated using price changes occurring after
the simulated execution.


% ------------------------------------------------------------
\subsection{Execution Lags}
% ------------------------------------------------------------

A strategy cannot generally assume that information observed at a given instant
can be traded at a price that occurred before the information was available.

If a signal requires the closing price at date $t$, for example,

\begin{equation}
P_{i,t}^{\mathrm{close}},
\end{equation}

then the resulting portfolio decision is formed only after that close is known.

A conservative timing convention may therefore be

\begin{equation}
\boxed{
\text{Close}_t
\rightarrow
\text{Portfolio Decision}
\rightarrow
\text{Execution at Open}_{t+1}.
}
\end{equation}

Alternatively,

\begin{equation}
\boxed{
\text{Close}_t
\rightarrow
\text{Portfolio Decision}
\rightarrow
\text{Execution at Close}_{t+1}.
}
\end{equation}

The appropriate lag depends on the assumed execution mechanism.

The lag should be modeled explicitly rather than being introduced accidentally
through DataFrame alignment.


\subsubsection{Avoiding Double Lags}

Suppose weights computed from information at date $t-1$ are already stored at
date $t$.

If those weights are then shifted again before multiplying by date-$t$ returns,
the backtester may unintentionally introduce an additional period of delay.

Thus,

\begin{equation}
\boxed{
\text{The economic decision should be lagged exactly once.}
}
\end{equation}

The correct implementation depends on the timestamp convention used for both
weights and returns.


% ------------------------------------------------------------
\subsection{Rebalancing Schedule}
% ------------------------------------------------------------

Let

\begin{equation}
\mathcal{T}^{R}
=
\left\{
t_1^{R},
t_2^{R},
\ldots
\right\}
\end{equation}

denote the set of portfolio rebalance dates.

Possible schedules include:

\begin{itemize}
    \item daily;
    \item weekly;
    \item monthly;
    \item quarterly;
    \item event-driven;
    \item threshold-triggered.
\end{itemize}

The signal calculation frequency and portfolio rebalance frequency need not be
identical.

Thus,

\begin{equation}
\boxed{
f_{\mathrm{signal}}
\neq
f_{\mathrm{rebalance}}
}
\end{equation}

in general.

For example, signals may be updated daily while the portfolio is rebalanced only
weekly.


% ------------------------------------------------------------
\subsection{Holding Period}
% ------------------------------------------------------------

Following execution at time $t_e$, the implemented portfolio is held until the
next rebalance or liquidation event.

The holding interval is

\begin{equation}
\boxed{
[t_e,t_e').
}
\end{equation}

If the portfolio is represented at discrete times, the implemented weights at
$t$ earn the subsequent return:

\begin{equation}
\boxed{
R_{p,t\rightarrow t+1}^{\mathrm{gross}}
=
\left(
\mathbf{w}_t^{\mathrm{implemented}}
\right)^\top
\mathbf{R}_{t\rightarrow t+1}.
}
\end{equation}

The fundamental requirement is

\begin{equation}
\boxed{
\mathbf{w}_t^{\mathrm{implemented}}
\text{ must be determined before }
\mathbf{R}_{t\rightarrow t+1}
\text{ is observed}.
}
\end{equation}


% ------------------------------------------------------------
\subsection{Timing Integrity and Look-Ahead Prevention}
% ------------------------------------------------------------

Timing integrity must hold throughout the complete portfolio simulation.

At each date, the backtester should verify that:

\begin{itemize}
    \item all signal inputs were available before the decision;
    \item all model parameters were estimated using admissible historical data;
    \item risk estimates were available before portfolio construction;
    \item target weights were computed before the returns attributed to them;
    \item execution prices occurred no earlier than the simulated execution time;
    \item transaction-cost inputs were available at the relevant decision or
    execution time.
\end{itemize}

Thus, the simulation chronology must always satisfy

\begin{equation}
\boxed{
\mathcal{I}_t
\rightarrow
\mathbf{w}_t^{\mathrm{final}}
\rightarrow
\text{Execution}
\rightarrow
\mathbf{w}_t^{\mathrm{implemented}}
\rightarrow
\text{Future Returns}.
}
\end{equation}


% ============================================================
\section{Portfolio Implementation}
\label{sec:portfolio_implementation}
% ============================================================

The portfolio-construction engine specifies what the strategy would ideally
like to hold.

The implementation engine determines the trades required to move from the
portfolio currently held to that desired portfolio.

The central objects are therefore

\begin{equation}
\boxed{
\mathbf{w}_{t^-}
\qquad\text{and}\qquad
\mathbf{w}_t^{\mathrm{final}},
}
\end{equation}

where $\mathbf{w}_{t^-}$ denotes portfolio weights immediately before the
rebalance.


% ------------------------------------------------------------
\subsection{Target vs.\ Current Portfolio}
% ------------------------------------------------------------

The required portfolio adjustment is

\begin{equation}
\boxed{
\Delta\mathbf{w}_t
=
\mathbf{w}_t^{\mathrm{final}}
-
\mathbf{w}_{t^-}.
}
\end{equation}

For security $i$,

\begin{equation}
\Delta w_{i,t}
=
w_{i,t}^{\mathrm{final}}
-
w_{i,t^-}.
\end{equation}

If

\begin{equation}
\Delta w_{i,t}>0,
\end{equation}

the strategy must increase its exposure to security $i$.

If

\begin{equation}
\Delta w_{i,t}<0,
\end{equation}

the strategy must reduce its exposure.


% ------------------------------------------------------------
\subsection{Portfolio Drift}
% ------------------------------------------------------------

The current portfolio immediately before rebalance is generally not equal to the
weights established at the previous rebalance.

Suppose the portfolio initially holds weights

\begin{equation}
w_{i,t-1}^{\mathrm{implemented}}
\end{equation}

and securities subsequently realize returns

\begin{equation}
R_{i,t}.
\end{equation}

Ignoring intermediate cash flows, pre-rebalance weights become

\begin{equation}
\boxed{
w_{i,t^-}
=
\frac{
w_{i,t-1}^{\mathrm{implemented}}
(1+R_{i,t})
}{
1+R_{p,t}
}.
}
\end{equation}

Therefore,

\begin{equation}
\boxed{
\mathbf{w}_{t^-}
\neq
\mathbf{w}_{t-1}^{\mathrm{implemented}}
}
\end{equation}

in general.

This phenomenon is referred to as \emph{portfolio drift}.

The correct trade calculation must therefore compare the new desired portfolio
with the drifted portfolio, not mechanically with the previous target weights.


% ------------------------------------------------------------
\subsection{Required Trades}
% ------------------------------------------------------------

Let

\begin{equation}
A_t
\end{equation}

denote portfolio NAV immediately before the rebalance.

The desired currency notional trade in security $i$ is approximately

\begin{equation}
\boxed{
Q_{i,t}^{\$}
=
A_t
\Delta w_{i,t}.
}
\end{equation}

Positive values correspond to purchases and negative values to sales.

If the execution price is

\begin{equation}
P_{i,t}^{\mathrm{exec}},
\end{equation}

the corresponding desired share quantity is

\begin{equation}
\boxed{
\Delta q_{i,t}^{\mathrm{desired}}
=
\frac{
A_t\Delta w_{i,t}
}{
P_{i,t}^{\mathrm{exec}}
}.
}
\end{equation}

The exact conversion may depend on contract multipliers, currency conversion,
and instrument type.


% ------------------------------------------------------------
\subsection{Order Generation}
% ------------------------------------------------------------

The desired quantity changes must be translated into executable orders.

For each security,

\begin{equation}
\Delta q_{i,t}>0
\end{equation}

generates a buy order, while

\begin{equation}
\Delta q_{i,t}<0
\end{equation}

generates a sell order.

An order can be represented schematically as

\begin{equation}
\boxed{
\mathcal{O}_{i,t}
=
\left(
i,
\mathrm{Side}_{i,t},
|\Delta q_{i,t}|,
\mathrm{OrderType}_{i,t},
t
\right).
}
\end{equation}

The simulation may use simple market orders or more detailed assumptions about
execution algorithms.


% ------------------------------------------------------------
\subsection{Execution Price Assumptions}
% ------------------------------------------------------------

The execution price determines the price at which simulated trades are assumed
to occur.

Possible conventions include:

\begin{itemize}
    \item next open;
    \item next close;
    \item VWAP;
    \item TWAP;
    \item arrival price;
    \item mid-price plus modeled spread and slippage.
\end{itemize}

The execution-price convention must be consistent with the information and
execution timeline.

For example,

\begin{equation}
\boxed{
\text{Signal requiring Close}_t
\not\Rightarrow
\text{execution at the same Close}_t
}
\end{equation}

unless the strategy is known and executable before the closing auction under the
assumed model.


% ------------------------------------------------------------
\subsection{Partial Execution and Partial Rebalancing}
% ------------------------------------------------------------

The desired trade need not always be fully executed.

Let

\begin{equation}
\phi_{i,t}
\in
[0,1]
\end{equation}

denote the fraction of the desired trade that can be executed.

Then

\begin{equation}
\boxed{
\Delta q_{i,t}^{\mathrm{executed}}
=
\phi_{i,t}
\Delta q_{i,t}^{\mathrm{desired}}.
}
\end{equation}

The executed portfolio therefore moves only partially toward the desired target.

Partial execution may result from:

\begin{itemize}
    \item liquidity constraints;
    \item participation-rate limits;
    \item unavailable borrow;
    \item trading halts;
    \item order-size limits;
    \item deliberate partial-rebalancing rules.
\end{itemize}


% ------------------------------------------------------------
\subsection{Position Rounding}
% ------------------------------------------------------------

Real portfolios generally trade discrete quantities.

If security $i$ must be traded in lot size

\begin{equation}
L_i,
\end{equation}

the executed quantity may be rounded as

\begin{equation}
\boxed{
\Delta q_{i,t}^{\mathrm{rounded}}
=
L_i
\cdot
\operatorname{Round}
\left(
\frac{
\Delta q_{i,t}^{\mathrm{desired}}
}{
L_i
}
\right).
}
\end{equation}

For ordinary equities,

\begin{equation}
L_i=1
\end{equation}

may represent whole-share trading.

Rounding means that implemented portfolio weights need not exactly equal the
desired weights.


% ------------------------------------------------------------
\subsection{Implemented Portfolio}
% ------------------------------------------------------------

After execution, the portfolio contains quantities

\begin{equation}
\boxed{
q_{i,t}^{\mathrm{implemented}}
=
q_{i,t^-}
+
\Delta q_{i,t}^{\mathrm{executed}}.
}
\end{equation}

The resulting implemented weights are

\begin{equation}
\boxed{
w_{i,t}^{\mathrm{implemented}}
=
\frac{
q_{i,t}^{\mathrm{implemented}}
P_{i,t}^{\mathrm{exec}}
}{
NAV_t^{\mathrm{post}}
}.
}
\end{equation}

Thus,

\begin{equation}
\boxed{
\mathbf{w}_t^{\mathrm{implemented}}
\neq
\mathbf{w}_t^{\mathrm{final}}
}
\end{equation}

in general.

The difference

\begin{equation}
\boxed{
\mathbf{e}_t^{\mathrm{implementation}}
=
\mathbf{w}_t^{\mathrm{implemented}}
-
\mathbf{w}_t^{\mathrm{final}}
}
\end{equation}

can be interpreted as implementation error.


% ============================================================
\section{Transaction Costs and Market Frictions}
\label{sec:transaction_costs_market_frictions}
% ============================================================

Trading changes portfolio holdings but also generates implementation costs.

Let

\begin{equation}
\Delta\mathbf{w}_t
\end{equation}

denote the executed portfolio trade.

Total implementation cost may be represented generically as

\begin{equation}
\boxed{
TC_t
=
C_t
\left(
\Delta\mathbf{w}_t,
\mathrm{Liquidity}_t,
\mathrm{Spread}_t,
\mathrm{Volatility}_t,
A_t,
\ldots
\right).
}
\end{equation}


% ------------------------------------------------------------
\subsection{General Transaction-Cost Model}
% ------------------------------------------------------------

A useful decomposition is

\begin{equation}
\boxed{
TC_t
=
TC_t^{\mathrm{fees}}
+
TC_t^{\mathrm{spread}}
+
TC_t^{\mathrm{slippage}}
+
TC_t^{\mathrm{impact}}.
}
\end{equation}

Holding-related costs such as borrow and financing are treated separately because
they accrue through time rather than only when a trade occurs.


% ------------------------------------------------------------
\subsection{Commissions and Explicit Fees}
% ------------------------------------------------------------

Let

\begin{equation}
c_{i,t}^{\mathrm{fee}}
\end{equation}

denote explicit cost per unit of traded notional.

Then

\begin{equation}
\boxed{
TC_t^{\mathrm{fees}}
=
\sum_i
c_{i,t}^{\mathrm{fee}}
|\Delta w_{i,t}|.
}
\end{equation}

These costs may include commissions, exchange fees, taxes, and regulatory fees.


% ------------------------------------------------------------
\subsection{Bid--Ask Spread}
% ------------------------------------------------------------

Let

\begin{equation}
P_{i,t}^{\mathrm{bid}}
\qquad\text{and}\qquad
P_{i,t}^{\mathrm{ask}}
\end{equation}

denote contemporaneous bid and ask prices.

The mid-price is

\begin{equation}
P_{i,t}^{\mathrm{mid}}
=
\frac{
P_{i,t}^{\mathrm{bid}}
+
P_{i,t}^{\mathrm{ask}}
}{2}.
\end{equation}

The proportional half-spread is approximately

\begin{equation}
\boxed{
c_{i,t}^{\mathrm{spread}}
=
\frac{
P_{i,t}^{\mathrm{ask}}
-
P_{i,t}^{\mathrm{bid}}
}{
2P_{i,t}^{\mathrm{mid}}
}.
}
\end{equation}

A trade crossing the spread therefore generates approximately

\begin{equation}
TC_t^{\mathrm{spread}}
=
\sum_i
c_{i,t}^{\mathrm{spread}}
|\Delta w_{i,t}|.
\end{equation}


% ------------------------------------------------------------
\subsection{Slippage}
% ------------------------------------------------------------

Slippage captures the difference between the reference price assumed when the
portfolio decision is made and the actual execution price.

For security $i$,

\begin{equation}
\boxed{
\mathrm{Slippage}_{i,t}
=
\frac{
P_{i,t}^{\mathrm{exec}}
-
P_{i,t}^{\mathrm{reference}}
}{
P_{i,t}^{\mathrm{reference}}
}
}
\end{equation}

with sign conventions adjusted according to trade direction.

Slippage may reflect price movement during the delay between decision and
execution, order-book dynamics, and execution uncertainty.


% ------------------------------------------------------------
\subsection{Market Impact}
% ------------------------------------------------------------

Large orders may move market prices.

A generic impact model can be written as

\begin{equation}
\boxed{
TC_{i,t}^{\mathrm{impact}}
=
\eta_{i,t}
|\Delta w_{i,t}|^{p},
\qquad
p>1.
}
\end{equation}

A more implementation-oriented model may depend on trade size relative to ADV:

\begin{equation}
TC_{i,t}^{\mathrm{impact}}
=
g
\left(
\frac{
A_t|\Delta w_{i,t}|
}{
\mathrm{ADV}_{i,t}
},
\sigma_{i,t},
\ldots
\right).
\end{equation}

Market impact therefore introduces an explicit dependence of strategy costs on
portfolio AUM.


% ------------------------------------------------------------
\subsection{Short Borrowing Costs}
% ------------------------------------------------------------

Short positions require borrowing the underlying securities.

Let

\begin{equation}
b_{i,t}
\end{equation}

denote the annualized borrow rate.

For a short position with notional

\begin{equation}
A_t|w_{i,t}|,
\end{equation}

the approximate borrow cost over a time fraction $\Delta t$ is

\begin{equation}
\boxed{
BC_{i,t}
=
A_t
|w_{i,t}|
b_{i,t}
\Delta t.
}
\end{equation}

Borrow costs accrue while the short position is held and should therefore be
distinguished from one-off transaction costs.


% ------------------------------------------------------------
\subsection{Financing Costs}
% ------------------------------------------------------------

Leveraged portfolios may require external financing.

Let

\begin{equation}
F_t
\end{equation}

denote the financed amount and

\begin{equation}
r_t^{\mathrm{fin}}
\end{equation}

the applicable annualized financing rate.

Over period $\Delta t$,

\begin{equation}
\boxed{
FC_t
=
F_t
r_t^{\mathrm{fin}}
\Delta t.
}
\end{equation}

The precise financing convention depends on portfolio structure, instrument
type, collateral, and cash treatment.


% ------------------------------------------------------------
\subsection{Liquidity and Capacity Effects}
% ------------------------------------------------------------

Implementation costs depend on strategy size.

Let

\begin{equation}
A_t
\end{equation}

denote portfolio AUM.

For fixed portfolio weights,

\begin{equation}
Q_{i,t}^{\$}
=
A_t|\Delta w_{i,t}|.
\end{equation}

Therefore,

\begin{equation}
A_t\uparrow
\quad\Rightarrow\quad
Q_{i,t}^{\$}\uparrow.
\end{equation}

Larger AUM may increase participation rates, market impact, execution time, and
the number of binding liquidity constraints.

Thus,

\begin{equation}
\boxed{
\text{Backtest Performance}
=
f(\text{AUM})
}
\end{equation}

whenever transaction costs and market impact are modeled realistically.


% ============================================================
\section{Portfolio Accounting and P\&L}
\label{sec:portfolio_accounting_pnl}
% ============================================================

Once trades have been executed, the backtester must maintain the evolving state
of the simulated portfolio.

The accounting engine converts security quantities, prices, and cash flows into
portfolio NAV and profit and loss.


% ------------------------------------------------------------
\subsection{Portfolio State Representation}
% ------------------------------------------------------------

A convenient portfolio state representation is

\begin{equation}
\boxed{
\mathcal{P}_t
=
\left(
\mathbf{q}_t,
\mathbf{P}_t,
C_t,
NAV_t
\right),
}
\end{equation}

where

\begin{equation}
\mathbf{q}_t
\in
\mathbb{R}^{N_t}
\end{equation}

denotes security quantities,

\begin{equation}
\mathbf{P}_t
\in
\mathbb{R}^{N_t}
\end{equation}

denotes valuation prices,

\begin{equation}
C_t
\end{equation}

denotes cash, and

\begin{equation}
NAV_t
\end{equation}

denotes portfolio net asset value.


% ------------------------------------------------------------
\subsection{Positions and Market Values}
% ------------------------------------------------------------

The market value of security $i$ is

\begin{equation}
\boxed{
MV_{i,t}
=
q_{i,t}
P_{i,t}.
}
\end{equation}

Positive market value represents a long position and negative market value a
short position.

Portfolio security value is

\begin{equation}
\boxed{
MV_t
=
\sum_i
q_{i,t}
P_{i,t}.
}
\end{equation}


% ------------------------------------------------------------
\subsection{Cash Account}
% ------------------------------------------------------------

Trading changes the cash account.

Ignoring costs temporarily,

\begin{equation}
\boxed{
C_t^{\mathrm{post}}
=
C_t^{\mathrm{pre}}
-
\sum_i
\Delta q_{i,t}
P_{i,t}^{\mathrm{exec}}.
}
\end{equation}

Purchases reduce cash, whereas sales increase cash.

Transaction costs and other cash flows are subsequently deducted or credited
according to the accounting convention.


% ------------------------------------------------------------
\subsection{Mark-to-Market Valuation}
% ------------------------------------------------------------

At valuation time $t$,

\begin{equation}
\boxed{
NAV_t
=
C_t
+
\sum_i
q_{i,t}
P_{i,t}.
}
\end{equation}

As market prices change, portfolio NAV changes even in the absence of trading.

This mark-to-market process creates the realized P\&L of the portfolio.


% ------------------------------------------------------------
\subsection{Daily Profit and Loss}
% ------------------------------------------------------------

Ignoring external capital flows, portfolio P\&L between two valuation times is

\begin{equation}
\boxed{
PnL_{t\rightarrow t+1}
=
NAV_{t+1}
-
NAV_t.
}
\end{equation}

A useful decomposition is

\begin{equation}
\boxed{
PnL_t^{\mathrm{net}}
=
PnL_t^{\mathrm{market}}
-
TC_t
-
BC_t
-
FC_t,
}
\end{equation}

where:

\begin{itemize}
    \item $PnL_t^{\mathrm{market}}$ is the mark-to-market trading P\&L;
    \item $TC_t$ denotes transaction costs;
    \item $BC_t$ denotes short-borrowing costs;
    \item $FC_t$ denotes financing costs.
\end{itemize}


% ------------------------------------------------------------
\subsection{NAV Evolution}
% ------------------------------------------------------------

If

\begin{equation}
R_{p,t+1}^{\mathrm{net}}
\end{equation}

denotes the net portfolio return, then

\begin{equation}
\boxed{
NAV_{t+1}
=
NAV_t
\left(
1+
R_{p,t+1}^{\mathrm{net}}
\right).
}
\end{equation}

Equivalently,

\begin{equation}
\boxed{
R_{p,t+1}^{\mathrm{net}}
=
\frac{
NAV_{t+1}
-
NAV_t
}{
NAV_t
}.
}
\end{equation}


% ============================================================
\section{Backtest Return Construction}
\label{sec:backtest_return_construction}
% ============================================================

The final output of the simulation engine is the realized return series of the
strategy.

The return-construction layer summarizes the effects of implemented positions,
market returns, transaction costs, and holding costs into a single net return
series.


% ------------------------------------------------------------
\subsection{Gross Portfolio Return}
% ------------------------------------------------------------

For implemented weights

\begin{equation}
\mathbf{w}_t^{\mathrm{implemented}},
\end{equation}

the gross portfolio return over the subsequent period is

\begin{equation}
\boxed{
R_{p,t\rightarrow t+1}^{\mathrm{gross}}
=
\left(
\mathbf{w}_t^{\mathrm{implemented}}
\right)^\top
\mathbf{R}_{t\rightarrow t+1}.
}
\end{equation}

This is the return generated by the holdings before implementation and financing
costs are deducted.


% ------------------------------------------------------------
\subsection{Trading-Cost Deduction}
% ------------------------------------------------------------

Let

\begin{equation}
TC_t
\end{equation}

denote transaction costs incurred when establishing or changing the portfolio,
expressed as a fraction of NAV.

Then the post-trading-cost return is

\begin{equation}
\boxed{
R_{p,t}^{\mathrm{postTC}}
=
R_{p,t}^{\mathrm{gross}}
-
TC_t.
}
\end{equation}


% ------------------------------------------------------------
\subsection{Financing and Borrowing-Cost Deduction}
% ------------------------------------------------------------

Let

\begin{equation}
FC_t
\end{equation}

and

\begin{equation}
BC_t
\end{equation}

denote financing and borrowing costs expressed as fractions of NAV over the
relevant holding period.

The return after all modeled holding costs is therefore

\begin{equation}
\boxed{
R_{p,t}^{\mathrm{net}}
=
R_{p,t}^{\mathrm{gross}}
-
TC_t
-
FC_t
-
BC_t.
}
\end{equation}

Additional taxes or strategy-specific implementation costs can be included
analogously.


% ------------------------------------------------------------
\subsection{Net Portfolio Return}
% ------------------------------------------------------------

The net portfolio return is the economically relevant return delivered after
the modeled costs of implementing and financing the strategy.

Thus,

\begin{equation}
\boxed{
R_{p,t}^{\mathrm{net}}
}
\end{equation}

rather than

\begin{equation}
R_{p,t}^{\mathrm{gross}}
\end{equation}

should generally be used when assessing implementable strategy performance.


% ------------------------------------------------------------
\subsection{Cumulative Performance}
% ------------------------------------------------------------

Starting from normalized wealth

\begin{equation}
V_0=1,
\end{equation}

the cumulative portfolio value evolves according to

\begin{equation}
\boxed{
V_T
=
\prod_{t=1}^{T}
\left(
1+
R_{p,t}^{\mathrm{net}}
\right).
}
\end{equation}

The cumulative net return is

\begin{equation}
\boxed{
R_{0\rightarrow T}^{\mathrm{cum}}
=
\prod_{t=1}^{T}
\left(
1+
R_{p,t}^{\mathrm{net}}
\right)
-1.
}
\end{equation}

At this point, the complete economic simulation of the strategy has been
constructed:

\begin{equation}
\boxed{
\begin{aligned}
\text{Final Desired Weights}
\;(\mathbf{w}_t^{\mathrm{final}})
&\rightarrow
\text{Execution}
\\
&\rightarrow
\text{Implemented Weights}
\;(\mathbf{w}_t^{\mathrm{implemented}})
\\
&\rightarrow
\text{Gross P\&L}
\\
&\rightarrow
\text{Costs and Frictions}
\\
&\rightarrow
R_{p,t}^{\mathrm{net}}
\\
&\rightarrow
\text{Cumulative Performance}.
\end{aligned}
}
\end{equation}

The resulting return series can now be evaluated using performance, risk,
drawdown, attribution, and robustness diagnostics.

% ============================================================
\part{Performance Evaluation and Attribution}
\label{part:performance_evaluation_attribution}
% ============================================================

The backtest simulation developed in the previous part produces a realized
portfolio return series

\begin{equation}
\left\{
R_{p,t}^{\mathrm{net}}
\right\}_{t=1}^{T},
\end{equation}

together with the corresponding gross returns, portfolio holdings, exposures,
trades, transaction costs, and other implementation quantities.

The final stage of the research pipeline consists of evaluating this realized
performance and understanding its economic sources.

Three distinct questions should be separated:

\begin{equation}
\boxed{
\begin{aligned}
\text{Performance Assessment}
&:
&&
\text{How attractive was the realized strategy performance?}
\\
\text{Benchmark and Factor Adjustment}
&:
&&
\text{How much performance remains after controlling for known risks?}
\\
\text{Performance Attribution}
&:
&&
\text{Which positions, exposures, signals, and costs generated the result?}
\end{aligned}
}
\end{equation}

A high backtest return is not sufficient evidence of a successful strategy.
Performance must be evaluated jointly with risk, drawdowns, turnover,
implementation costs, capacity, benchmark exposure, factor exposure, and the
stability of the underlying sources of return.


% ============================================================
\section{Performance Assessment}
\label{sec:performance_assessment}
% ============================================================

Let

\begin{equation}
R_{p,t},
\qquad
t=1,\ldots,T,
\end{equation}

denote the portfolio return series at the evaluation frequency.

Unless otherwise stated, performance should preferably be evaluated using net
returns,

\begin{equation}
R_{p,t}
=
R_{p,t}^{\mathrm{net}}.
\end{equation}

Let

\begin{equation}
A
\end{equation}

denote the number of return observations per year. Typical conventions are

\begin{equation}
A=
\begin{cases}
252, & \text{daily returns},\\
52, & \text{weekly returns},\\
12, & \text{monthly returns}.
\end{cases}
\end{equation}


% ------------------------------------------------------------
\subsection{Annualized Return}
% ------------------------------------------------------------

The arithmetic mean periodic return is

\begin{equation}
\bar{R}_p
=
\frac{1}{T}
\sum_{t=1}^{T}
R_{p,t}.
\end{equation}

A simple arithmetic annualization is

\begin{equation}
\boxed{
\mu_p^{\mathrm{ann}}
=
A\bar{R}_p.
}
\end{equation}

For cumulative investment performance, however, geometric compounding is
generally more appropriate.

Let

\begin{equation}
V_T
=
V_0
\prod_{t=1}^{T}
(1+R_{p,t}).
\end{equation}

The compound annual growth rate is

\begin{equation}
\boxed{
CAGR
=
\left(
\frac{V_T}{V_0}
\right)^{A/T}
-1.
}
\end{equation}

Equivalently,

\begin{equation}
\boxed{
CAGR
=
\left[
\prod_{t=1}^{T}
(1+R_{p,t})
\right]^{A/T}
-1.
}
\end{equation}

The distinction between arithmetic annualized return and CAGR should always be
made explicit.


% ------------------------------------------------------------
\subsection{Annualized Volatility}
% ------------------------------------------------------------

The sample volatility of periodic returns is

\begin{equation}
\widehat{\sigma}_p
=
\sqrt{
\frac{1}{T-1}
\sum_{t=1}^{T}
\left(
R_{p,t}-\bar{R}_p
\right)^2
}.
\end{equation}

Under the standard square-root-of-time convention, annualized volatility is

\begin{equation}
\boxed{
\widehat{\sigma}_p^{\mathrm{ann}}
=
\sqrt{A}
\widehat{\sigma}_p.
}
\end{equation}

This scaling assumes that the relevant return variance aggregates approximately
linearly through time. Serial correlation, volatility clustering, and other
time-series dependencies may cause the simple square-root-of-time approximation
to be imperfect.


% ------------------------------------------------------------
\subsection{Sharpe Ratio}
% ------------------------------------------------------------

Let

\begin{equation}
R_{f,t}
\end{equation}

denote the risk-free return over period $t$ and define excess returns as

\begin{equation}
R_{p,t}^{e}
=
R_{p,t}-R_{f,t}.
\end{equation}

The annualized Sharpe ratio is

\begin{equation}
\boxed{
SR
=
\sqrt{A}
\frac{
\overline{R_p^{e}}
}{
\widehat{\sigma}(R_p^{e})
}.
}
\end{equation}

When the risk-free rate is negligible relative to the strategy return, a common
approximation is

\begin{equation}
SR
\approx
\sqrt{A}
\frac{
\bar{R}_p
}{
\widehat{\sigma}_p
}.
\end{equation}

For market-neutral or leveraged strategies, the financing convention should be
consistent with the return construction developed previously.


% ------------------------------------------------------------
\subsection{Sortino Ratio}
% ------------------------------------------------------------

The Sortino ratio penalizes downside variation rather than total volatility.

For a minimum acceptable periodic return $MAR$, define downside deviation as

\begin{equation}
\widehat{\sigma}_{\mathrm{down}}
=
\sqrt{
\frac{1}{T}
\sum_{t=1}^{T}
\min
\left(
R_{p,t}-MAR,
0
\right)^2
}.
\end{equation}

The annualized Sortino ratio can then be written as

\begin{equation}
\boxed{
\mathrm{Sortino}
=
\sqrt{A}
\frac{
\bar{R}_p-MAR
}{
\widehat{\sigma}_{\mathrm{down}}
}.
}
\end{equation}

The definition of $MAR$ and its frequency must be stated explicitly.


% ------------------------------------------------------------
\subsection{Maximum Drawdown}
% ------------------------------------------------------------

Let

\begin{equation}
V_t
=
V_0
\prod_{\tau=1}^{t}
(1+R_{p,\tau})
\end{equation}

denote the portfolio wealth process.

The running historical peak is

\begin{equation}
H_t
=
\max_{0\leq u\leq t}
V_u.
\end{equation}

The drawdown at time $t$ is

\begin{equation}
\boxed{
DD_t
=
\frac{V_t}{H_t}-1.
}
\end{equation}

Hence,

\begin{equation}
DD_t\leq0.
\end{equation}

Maximum drawdown is

\begin{equation}
\boxed{
MDD
=
-\min_t DD_t.
}
\end{equation}

Under this convention, $MDD$ is reported as a positive loss magnitude.


% ------------------------------------------------------------
\subsection{Calmar Ratio}
% ------------------------------------------------------------

The Calmar ratio compares compounded annual performance with maximum drawdown:

\begin{equation}
\boxed{
\mathrm{Calmar}
=
\frac{
CAGR
}{
MDD
}.
}
\end{equation}

Unlike the Sharpe ratio, the Calmar ratio directly penalizes the largest
historical peak-to-trough loss.


% ------------------------------------------------------------
\subsection{Hit Ratio}
% ------------------------------------------------------------

A simple portfolio-level hit ratio measures the fraction of periods with
positive returns:

\begin{equation}
\boxed{
HR
=
\frac{1}{T}
\sum_{t=1}^{T}
\mathbb{I}
\left(
R_{p,t}>0
\right).
}
\end{equation}

More generally, a benchmark or required-return threshold $c_t$ may be used:

\begin{equation}
HR(c)
=
\frac{1}{T}
\sum_{t=1}^{T}
\mathbb{I}
\left(
R_{p,t}>c_t
\right).
\end{equation}

Hit ratio should not be interpreted independently of payoff magnitude. A
strategy may have a low hit ratio and remain profitable if winning periods are
substantially larger than losing periods, or conversely.


% ------------------------------------------------------------
\subsection{Skewness and Kurtosis}
% ------------------------------------------------------------

The standardized third central moment measures return skewness:

\begin{equation}
\boxed{
\widehat{\mathrm{Skew}}
=
\frac{
\frac{1}{T}
\sum_{t=1}^{T}
(R_{p,t}-\bar{R}_p)^3
}{
\widehat{\sigma}_p^3
}.
}
\end{equation}

Positive skewness indicates a relatively longer right tail, while negative
skewness indicates a relatively longer left tail.

The standardized fourth central moment is

\begin{equation}
\boxed{
\widehat{\mathrm{Kurt}}
=
\frac{
\frac{1}{T}
\sum_{t=1}^{T}
(R_{p,t}-\bar{R}_p)^4
}{
\widehat{\sigma}_p^4
}.
}
\end{equation}

For a Gaussian distribution,

\begin{equation}
\mathrm{Kurt}=3.
\end{equation}

Excess kurtosis is therefore

\begin{equation}
\boxed{
\mathrm{ExKurt}
=
\widehat{\mathrm{Kurt}}-3.
}
\end{equation}

Large positive excess kurtosis indicates heavier tails than under the Gaussian
benchmark.


% ------------------------------------------------------------
\subsection{Value at Risk}
% ------------------------------------------------------------

For confidence level $\alpha$, Value at Risk is defined through the lower
quantile of the portfolio return distribution.

Let

\begin{equation}
q_{1-\alpha}(R_p)
\end{equation}

denote the $(1-\alpha)$ return quantile.

Then

\begin{equation}
\boxed{
VaR_{\alpha}
=
-
q_{1-\alpha}(R_p).
}
\end{equation}

For example,

\begin{equation}
VaR_{95\%}
=
-
q_{5\%}(R_p).
\end{equation}

Historical VaR can be estimated directly from the empirical distribution of
backtest returns.

VaR describes a loss threshold but does not measure the magnitude of losses
beyond that threshold.


% ------------------------------------------------------------
\subsection{Expected Shortfall}
% ------------------------------------------------------------

Expected Shortfall addresses this limitation by measuring the average loss in
the tail beyond the VaR threshold.

Under a continuous return distribution,

\begin{equation}
\boxed{
ES_{\alpha}
=
-
\mathbb{E}
\left[
R_p
\mid
R_p
\leq
q_{1-\alpha}(R_p)
\right].
}
\end{equation}

For example, $ES_{95\%}$ measures the average loss conditional on returns lying
in approximately the worst $5\%$ of the distribution.

Expected Shortfall therefore contains information about tail severity that is
not captured by VaR alone.


% ------------------------------------------------------------
\subsection{Turnover}
% ------------------------------------------------------------

Turnover measures the amount of portfolio trading required to implement the
strategy.

Using the one-sided convention developed previously,

\begin{equation}
\boxed{
TO_t
=
\frac{1}{2}
\sum_{i=1}^{N_t}
\left|
w_{i,t}^{\mathrm{final}}
-
w_{i,t^-}
\right|.
}
\end{equation}

If implemented rather than desired trades are available, realized turnover
should preferably be computed from executed changes.

Average turnover is

\begin{equation}
\boxed{
\overline{TO}
=
\frac{1}{T}
\sum_{t=1}^{T}
TO_t.
}
\end{equation}

If turnover is annualized, the exact scaling convention and rebalance frequency
must be stated explicitly.

Turnover is not itself a performance measure, but it is an essential indicator
of implementation intensity and the sensitivity of gross alpha to trading
costs.


% ------------------------------------------------------------
\subsection{Capacity}
% ------------------------------------------------------------

Capacity measures how much capital can be deployed in the strategy before
liquidity constraints and market impact materially degrade its performance.

Let

\begin{equation}
A
\end{equation}

denote strategy AUM.

Because traded dollar notional scales with AUM,

\begin{equation}
Q_{i,t}^{\$}(A)
=
A
|\Delta w_{i,t}|,
\end{equation}

transaction costs generally become functions of portfolio size:

\begin{equation}
TC_t
=
TC_t(A).
\end{equation}

Consequently,

\begin{equation}
\boxed{
R_{p,t}^{\mathrm{net}}
=
R_{p,t}^{\mathrm{gross}}
-
TC_t(A)
-
FC_t(A)
-
BC_t(A).
}
\end{equation}

A capacity curve may therefore be constructed as

\begin{equation}
\boxed{
A
\longmapsto
\mathrm{Performance}^{\mathrm{net}}(A).
}
\end{equation}

For example, one may examine

\begin{equation}
A
\longmapsto
SR^{\mathrm{net}}(A)
\end{equation}

or

\begin{equation}
A
\longmapsto
CAGR^{\mathrm{net}}(A).
\end{equation}

A practical capacity threshold may then be defined as the largest AUM satisfying
a chosen economic condition, such as

\begin{equation}
SR^{\mathrm{net}}(A)
\geq
SR_{\min}.
\end{equation}

Capacity is therefore not a single universal number. It depends on the
transaction-cost model, execution assumptions, liquidity constraints, required
performance threshold, and portfolio implementation process.


% ------------------------------------------------------------
\subsection{Gross vs.\ Net Performance}
% ------------------------------------------------------------

A strategy should generally be evaluated both before and after implementation
costs.

Let

\begin{equation}
R_{p,t}^{\mathrm{gross}}
\end{equation}

denote gross strategy return and

\begin{equation}
R_{p,t}^{\mathrm{net}}
\end{equation}

denote return after the modeled costs.

The difference

\begin{equation}
\boxed{
D_t^{\mathrm{implementation}}
=
R_{p,t}^{\mathrm{gross}}
-
R_{p,t}^{\mathrm{net}}
}
\end{equation}

measures the return drag associated with implementation.

Under the cost decomposition introduced previously,

\begin{equation}
\boxed{
D_t^{\mathrm{implementation}}
=
TC_t
+
FC_t
+
BC_t.
}
\end{equation}

Comparing gross and net performance is particularly important for high-turnover
signals, less liquid universes, leveraged strategies, and short portfolios.

A strategy with attractive gross performance but weak net performance is not
necessarily economically implementable.


% ============================================================
\section{Benchmark and Factor-Adjusted Performance}
\label{sec:benchmark_factor_adjusted_performance}
% ============================================================

Absolute portfolio performance does not reveal whether returns result from
genuine alpha or compensation for systematic risk exposures.

The strategy should therefore also be evaluated relative to an appropriate
benchmark and, where relevant, a set of systematic risk factors.


% ------------------------------------------------------------
\subsection{Benchmark Excess Returns}
% ------------------------------------------------------------

Let

\begin{equation}
R_{b,t}
\end{equation}

denote the benchmark return.

The active return is

\begin{equation}
\boxed{
R_{p,t}^{\mathrm{active}}
=
R_{p,t}
-
R_{b,t}.
}
\end{equation}

The average active return is

\begin{equation}
\overline{R^{\mathrm{active}}}
=
\frac{1}{T}
\sum_{t=1}^{T}
\left(
R_{p,t}-R_{b,t}
\right).
\end{equation}

The choice of benchmark should reflect the economic mandate of the portfolio.

For a long-only equity strategy, a broad equity index may be appropriate.

For a market-neutral Pure Alpha strategy, benchmark-relative performance may be
less informative than absolute return and factor-adjusted alpha.


% ------------------------------------------------------------
\subsection{Tracking Error}
% ------------------------------------------------------------

Tracking error is the volatility of active returns:

\begin{equation}
\boxed{
TE
=
\sqrt{A}
\,
\widehat{\sigma}
\left(
R_p-R_b
\right).
}
\end{equation}

Tracking error measures the variability of portfolio returns relative to the
benchmark rather than the absolute volatility of the portfolio.


% ------------------------------------------------------------
\subsection{Information Ratio}
% ------------------------------------------------------------

The Information Ratio compares average active return with tracking error:

\begin{equation}
\boxed{
IR
=
\frac{
A\,
\overline{R_p-R_b}
}{
TE
}.
}
\end{equation}

Equivalently,

\begin{equation}
\boxed{
IR
=
\sqrt{A}
\frac{
\overline{R_p-R_b}
}{
\widehat{\sigma}(R_p-R_b)
}.
}
\end{equation}

The Information Ratio is particularly relevant for benchmark-relative
long-only portfolios.


% ------------------------------------------------------------
\subsection{CAPM Alpha and Beta}
% ------------------------------------------------------------

A first systematic-risk adjustment is provided by the CAPM regression:

\begin{equation}
\boxed{
R_{p,t}-R_{f,t}
=
\alpha_p
+
\beta_p
\left(
R_{m,t}-R_{f,t}
\right)
+
\varepsilon_t.
}
\end{equation}

Here,

\begin{equation}
\alpha_p
\end{equation}

is the periodic CAPM alpha,

\begin{equation}
\beta_p
\end{equation}

is the portfolio market beta, and

\begin{equation}
\varepsilon_t
\end{equation}

is the regression residual.

If the regression uses daily returns, annualized arithmetic alpha is commonly
reported as

\begin{equation}
\boxed{
\alpha_p^{\mathrm{ann}}
\approx
252\alpha_p.
}
\end{equation}

More generally,

\begin{equation}
\alpha_p^{\mathrm{ann}}
\approx
A\alpha_p.
\end{equation}

The estimated beta provides an ex-post diagnostic of whether the realized
portfolio behaved as intended with respect to market exposure.


% ------------------------------------------------------------
\subsection{Multifactor Regression}
% ------------------------------------------------------------

The CAPM framework can be generalized to $K$ systematic factors.

Let

\begin{equation}
\mathbf{F}_t
=
\begin{bmatrix}
F_{1,t}\\
\vdots\\
F_{K,t}
\end{bmatrix}
\in
\mathbb{R}^{K}
\end{equation}

denote factor returns.

The multifactor regression is

\begin{equation}
\boxed{
R_{p,t}-R_{f,t}
=
\alpha_p
+
\boldsymbol{\beta}_p^\top
\mathbf{F}_t
+
\varepsilon_t,
}
\end{equation}

where

\begin{equation}
\boldsymbol{\beta}_p
=
\begin{bmatrix}
\beta_{p,1}\\
\vdots\\
\beta_{p,K}
\end{bmatrix}
\in
\mathbb{R}^{K}.
\end{equation}

The factor set may include market, size, value, momentum, quality, low-risk, or
other systematic return factors appropriate to the investment universe.


% ------------------------------------------------------------
\subsection{Style Factor Exposures}
% ------------------------------------------------------------

Two distinct notions of factor exposure should be distinguished.

First, holdings-based exposure can be computed from contemporaneous security
factor exposures.

Let

\begin{equation}
\mathbf{B}_t
\in
\mathbb{R}^{N_t\times K}
\end{equation}

denote the security-level factor exposure matrix.

Portfolio exposure is

\begin{equation}
\boxed{
\mathbf{b}_{p,t}
=
\mathbf{B}_t^\top
\mathbf{w}_t.
}
\end{equation}

Second, realized return exposure can be estimated from the time-series
regression

\begin{equation}
R_{p,t}
=
\alpha_p
+
\boldsymbol{\beta}_p^\top
\mathbf{F}_t
+
\varepsilon_t.
\end{equation}

These quantities answer related but different questions.

Holdings-based exposures describe the portfolio's contemporaneous composition,
whereas regression betas describe its realized historical co-movement with
factor returns.


% ------------------------------------------------------------
\subsection{Factor-Adjusted Alpha}
% ------------------------------------------------------------

Under the multifactor model,

\begin{equation}
R_{p,t}-R_{f,t}
=
\alpha_p
+
\boldsymbol{\beta}_p^\top
\mathbf{F}_t
+
\varepsilon_t,
\end{equation}

the intercept

\begin{equation}
\boxed{
\alpha_p
}
\end{equation}

represents average portfolio performance unexplained by the included systematic
factors.

This quantity should not automatically be interpreted as structural investment
skill.

Its interpretation depends on:

\begin{itemize}
    \item the factor model being sufficiently comprehensive;
    \item the stability of estimated factor exposures;
    \item the sample period;
    \item statistical estimation uncertainty;
    \item transaction costs and financing assumptions;
    \item the absence of backtest overfitting and data leakage.
\end{itemize}

The factor-adjusted return can also be represented period by period as

\begin{equation}
\boxed{
R_{p,t}^{\mathrm{factor\text{-}adjusted}}
=
R_{p,t}
-
\widehat{\boldsymbol{\beta}}_p^\top
\mathbf{F}_t.
}
\end{equation}

Its sample mean corresponds closely to the estimated regression alpha under the
same regression specification.


% ============================================================
\section{Performance Attribution}
\label{sec:performance_attribution}
% ============================================================

Performance attribution goes beyond evaluating whether the strategy performed
well.

Its purpose is to identify the economic sources of realized P\&L.

Attribution can be performed along several dimensions:

\begin{equation}
\boxed{
\text{Portfolio P\&L}
\rightarrow
\begin{cases}
\text{Long / Short},\\
\text{Security},\\
\text{Sector / Country},\\
\text{Factor},\\
\text{Signal / Model},\\
\text{Beta / Pure Alpha},\\
\text{Implementation Costs}.
\end{cases}
}
\end{equation}

Different attribution schemes answer different questions and need not be
mutually exclusive.


% ------------------------------------------------------------
\subsection{Attribution Framework}
% ------------------------------------------------------------

For a linear holdings-based attribution, the gross portfolio return can be
written as

\begin{equation}
\boxed{
R_{p,t}^{\mathrm{gross}}
=
\sum_{i=1}^{N_t}
w_{i,t-1}^{\mathrm{implemented}}
R_{i,t},
}
\end{equation}

under the corresponding timing convention.

Define the security contribution

\begin{equation}
\boxed{
C_{i,t}
=
w_{i,t-1}^{\mathrm{implemented}}
R_{i,t}.
}
\end{equation}

Then

\begin{equation}
\boxed{
R_{p,t}^{\mathrm{gross}}
=
\sum_{i=1}^{N_t}
C_{i,t}.
}
\end{equation}

This simple additive identity forms the basis of many holdings-based attribution
methods.


% ------------------------------------------------------------
\subsection{Long vs.\ Short Attribution}
% ------------------------------------------------------------

Partition the universe into long and short positions:

\begin{equation}
\mathcal{L}_t
=
\left\{
i:w_{i,t}>0
\right\},
\end{equation}

\begin{equation}
\mathcal{S}_t
=
\left\{
i:w_{i,t}<0
\right\}.
\end{equation}

The long contribution is

\begin{equation}
\boxed{
R_{L,t}
=
\sum_{i\in\mathcal{L}_t}
w_{i,t}R_{i,t},
}
\end{equation}

while the short contribution is

\begin{equation}
\boxed{
R_{S,t}
=
\sum_{i\in\mathcal{S}_t}
w_{i,t}R_{i,t}.
}
\end{equation}

Since short weights are negative, profitable short positions generate positive
contributions when the corresponding securities decline.

Thus,

\begin{equation}
\boxed{
R_{p,t}^{\mathrm{gross}}
=
R_{L,t}
+
R_{S,t}.
}
\end{equation}

This decomposition helps determine whether performance is driven predominantly
by long selection, short selection, or both.


% ------------------------------------------------------------
\subsection{Security-Level Attribution}
% ------------------------------------------------------------

Security-level contribution is

\begin{equation}
\boxed{
C_{i,t}
=
w_{i,t}R_{i,t}.
}
\end{equation}

Over a single period,

\begin{equation}
R_{p,t}
=
\sum_i C_{i,t}.
\end{equation}

Security contributions can be aggregated through time to identify persistent
contributors and detractors.

For multi-period attribution, simple summation of percentage contributions does
not in general reproduce compounded total return exactly. An appropriate
multi-period linking methodology should therefore be used when exact cumulative
attribution is required.


% ------------------------------------------------------------
\subsection{Sector and Industry Attribution}
% ------------------------------------------------------------

Let

\begin{equation}
\mathcal{I}_{g,t}
\end{equation}

denote the set of securities belonging to industry or sector $g$.

The sector contribution is

\begin{equation}
\boxed{
C_{g,t}^{\mathrm{sector}}
=
\sum_{i\in\mathcal{I}_{g,t}}
C_{i,t}.
}
\end{equation}

Therefore,

\begin{equation}
\boxed{
R_{p,t}^{\mathrm{gross}}
=
\sum_g
C_{g,t}^{\mathrm{sector}}.
}
\end{equation}

This decomposition can reveal whether an apparently diversified strategy
derives a disproportionate fraction of its realized P\&L from a small number of
industries.


% ------------------------------------------------------------
\subsection{Country Attribution}
% ------------------------------------------------------------

Let

\begin{equation}
\mathcal{C}_{c,t}
\end{equation}

denote the set of securities assigned to country $c$.

Country-level contribution is

\begin{equation}
\boxed{
C_{c,t}^{\mathrm{country}}
=
\sum_{i\in\mathcal{C}_{c,t}}
C_{i,t}.
}
\end{equation}

The portfolio return can then be decomposed as

\begin{equation}
R_{p,t}^{\mathrm{gross}}
=
\sum_c
C_{c,t}^{\mathrm{country}}.
\end{equation}

For international portfolios, the interpretation should be consistent with the
currency-conversion convention used in the backtest.


% ------------------------------------------------------------
\subsection{Factor Attribution}
% ------------------------------------------------------------

Suppose realized portfolio returns satisfy the factor representation

\begin{equation}
R_{p,t}-R_{f,t}
=
\alpha_p
+
\boldsymbol{\beta}_p^\top
\mathbf{F}_t
+
\varepsilon_t.
\end{equation}

The systematic factor component is

\begin{equation}
\boxed{
R_{p,t}^{\mathrm{factor}}
=
\boldsymbol{\beta}_p^\top
\mathbf{F}_t
=
\sum_{k=1}^{K}
\beta_{p,k}F_{k,t}.
}
\end{equation}

The contribution of factor $k$ is therefore

\begin{equation}
\boxed{
C_{k,t}^{\mathrm{factor}}
=
\beta_{p,k}F_{k,t}.
}
\end{equation}

This provides a return-based factor attribution.

A holdings-based risk-model attribution may instead use time-varying portfolio
exposures

\begin{equation}
\mathbf{b}_{p,t}
=
\mathbf{B}_t^\top\mathbf{w}_t
\end{equation}

and corresponding factor returns.

The chosen attribution methodology should therefore be stated explicitly.


% ------------------------------------------------------------
\subsection{Signal Attribution}
% ------------------------------------------------------------

Suppose the final signal is constructed from $K$ component signals:

\begin{equation}
\mathbf{s}_t^{\mathrm{final}}
=
\sum_{k=1}^{K}
a_{k,t}
\mathbf{s}_{k,t}.
\end{equation}

It is often desirable to determine how much portfolio performance is associated
with each component signal.

A simple diagnostic is to construct standalone portfolios

\begin{equation}
\mathbf{w}_{k,t}
=
\mathcal{G}
\left(
\mathbf{s}_{k,t}
\right)
\end{equation}

and evaluate their corresponding returns

\begin{equation}
R_{k,t}
=
\mathbf{w}_{k,t}^\top
\mathbf{R}_{t+1}.
\end{equation}

This provides a useful measure of standalone signal performance.

However,

\begin{equation}
\boxed{
R_{p,t}
\neq
\sum_{k=1}^{K}
a_{k,t}R_{k,t}
}
\end{equation}

in general when the portfolio-construction operator $\mathcal{G}$ is nonlinear
or when optimization, constraints, clipping, transaction costs, or interactions
between signals affect the final portfolio.

Therefore, standalone signal backtests should not automatically be interpreted
as an exact additive attribution of the final portfolio.

Exact signal attribution requires a methodology consistent with the actual
signal-combination and portfolio-construction process.


% ------------------------------------------------------------
\subsection{Model Attribution}
% ------------------------------------------------------------

For model-based strategies, attribution can be extended to the predictive model
itself.

Suppose

\begin{equation}
\widehat{\boldsymbol{\mu}}_t
=
f_{\widehat{\theta}_t}
\left(
\mathbf{X}_t
\right),
\end{equation}

where $\mathbf{X}_t$ denotes model inputs and
$f_{\widehat{\theta}_t}(\cdot)$ the fitted prediction function.

Model attribution may investigate:

\begin{itemize}
    \item performance by model;
    \item performance by feature group;
    \item performance by prediction bucket;
    \item performance by confidence level;
    \item performance across market regimes;
    \item marginal performance from adding or removing model components.
\end{itemize}

For linear models,

\begin{equation}
\widehat{\mu}_{i,t}
=
\widehat{\beta}_{0,t}
+
\sum_{k=1}^{K}
X_{i,k,t}
\widehat{\beta}_{k,t},
\end{equation}

so the prediction itself admits an additive decomposition.

For nonlinear machine-learning or deep-learning models, feature attribution may
require model-specific methods.

Importantly, prediction attribution and portfolio-return attribution are
different objects.

A feature may contribute strongly to a model prediction without generating an
equally large realized portfolio P\&L contribution after portfolio construction.


% ------------------------------------------------------------
\subsection{Beta vs.\ Pure Alpha}
% ------------------------------------------------------------

For strategies intended to generate stock-selection alpha, it is useful to
separate realized performance associated with market exposure from residual
performance.

Under a simple market model,

\begin{equation}
R_{p,t}-R_{f,t}
=
\alpha_p
+
\beta_p
\left(
R_{m,t}-R_{f,t}
\right)
+
\varepsilon_t.
\end{equation}

The market-related component is

\begin{equation}
\boxed{
R_{p,t}^{\beta}
=
\widehat{\beta}_p
\left(
R_{m,t}-R_{f,t}
\right),
}
\end{equation}

while the residual component is

\begin{equation}
\boxed{
R_{p,t}^{\mathrm{res}}
=
R_{p,t}-R_{f,t}
-
R_{p,t}^{\beta}.
}
\end{equation}

More generally, with time-varying portfolio beta,

\begin{equation}
\boxed{
R_{p,t}^{\beta}
=
\beta_{p,t-1}
R_{m,t},
}
\end{equation}

under the appropriate timing and return convention.

The residual return can then be interpreted as the portion of realized
performance not mechanically explained by the modeled market exposure.

This decomposition is particularly informative when comparing long-only,
long--short, and beta-neutral versions of the same underlying alpha signal.


% ------------------------------------------------------------
\subsection{Gross Alpha vs.\ Implementation Costs}
% ------------------------------------------------------------

A strategy may generate economically meaningful gross alpha while losing a
substantial fraction of that alpha through implementation.

The basic decomposition is

\begin{equation}
\boxed{
R_{p,t}^{\mathrm{net}}
=
R_{p,t}^{\mathrm{gross}}
-
TC_t
-
FC_t
-
BC_t.
}
\end{equation}

Over an evaluation sample, one may compare

\begin{equation}
\boxed{
\text{Gross Alpha}
\quad\text{vs.}\quad
\text{Implementation Drag}
\quad\text{vs.}\quad
\text{Net Alpha}.
}
\end{equation}

This decomposition is particularly important when comparing signals with
different turnover, liquidity, or shorting requirements.


% ------------------------------------------------------------
\subsection{Transaction Cost Attribution}
% ------------------------------------------------------------

Using the cost decomposition developed previously,

\begin{equation}
TC_t
=
TC_t^{\mathrm{fees}}
+
TC_t^{\mathrm{spread}}
+
TC_t^{\mathrm{slippage}}
+
TC_t^{\mathrm{impact}}.
\end{equation}

The implementation drag attributable to trading can therefore be decomposed as

\begin{equation}
\boxed{
TC_t
=
\underbrace{TC_t^{\mathrm{fees}}}_{\text{explicit fees}}
+
\underbrace{TC_t^{\mathrm{spread}}}_{\text{spread}}
+
\underbrace{TC_t^{\mathrm{slippage}}}_{\text{slippage}}
+
\underbrace{TC_t^{\mathrm{impact}}}_{\text{market impact}}.
}
\end{equation}

Costs can additionally be aggregated by:

\begin{itemize}
    \item security;
    \item sector;
    \item country;
    \item rebalance date;
    \item long and short legs;
    \item signal or strategy sleeve.
\end{itemize}

This identifies which parts of the strategy consume the largest fraction of
gross alpha.


% ------------------------------------------------------------
\subsection{Financing and Borrowing Cost Attribution}
% ------------------------------------------------------------

Holding-related costs should be separated from transaction costs.

Let

\begin{equation}
FC_t
\end{equation}

denote financing costs and

\begin{equation}
BC_t
\end{equation}

denote short-borrowing costs.

Total holding-cost drag is

\begin{equation}
\boxed{
HC_t
=
FC_t
+
BC_t.
}
\end{equation}

For a short portfolio, borrow-cost attribution can be performed security by
security:

\begin{equation}
\boxed{
BC_t
=
\sum_{i\in\mathcal{S}_t}
BC_{i,t}.
}
\end{equation}

This is particularly important because securities with strong predicted alpha
may simultaneously be expensive or difficult to borrow.

The economically relevant contribution of a short position is therefore not
only its gross price return but its return after the cost of maintaining the
position.

At the highest level, realized net strategy performance can finally be viewed
as

\begin{equation}
\boxed{
\begin{aligned}
\text{Net Strategy Performance}
&=
\text{Systematic Return Contribution}
\\
&\quad+
\text{Security-Selection / Pure Alpha Contribution}
\\
&\quad-
\text{Transaction-Cost Drag}
\\
&\quad-
\text{Financing and Borrowing Drag}.
\end{aligned}
}
\end{equation}

This final decomposition closes the backtesting pipeline.

Starting from raw point-in-time financial data, the framework has successively
constructed cleaned firm-level characteristics, transformed them into signals,
neutralized unwanted exposures, mapped those signals into portfolio positions,
controlled portfolio risk and constraints, simulated implementation, and
ultimately decomposed the resulting realized performance into its economic
sources.


% ============================================================
\part{Validation, Robustness, and Production}
% ============================================================


% ============================================================
\section{Robustness and Sensitivity Analysis}
\label{sec:robustness_sensitivity}
% ============================================================

A backtest should not be considered reliable solely because one particular
specification produces attractive historical performance.

The observed results depend on many research choices:

\begin{itemize}
    \item the historical sample;
    \item the investment universe;
    \item signal-construction parameters;
    \item neutralization choices;
    \item portfolio mapping rules;
    \item rebalance frequency;
    \item holding period;
    \item transaction-cost assumptions;
    \item model hyperparameters.
\end{itemize}

A robust strategy should remain economically meaningful under reasonable
perturbations of these choices.

The purpose of robustness analysis is therefore not to search for the
specification producing the highest Sharpe ratio. Instead, it is to assess
whether the conclusions of the research depend critically on a narrow set of
assumptions.

Conceptually, let

\begin{equation}
\boldsymbol{\theta}
\end{equation}

denote the complete set of research and implementation choices.

The resulting backtest performance can be represented abstractly as

\begin{equation}
\boxed{
\mathcal{M}
=
\mathcal{B}
\left(
\boldsymbol{\theta}
\right),
}
\end{equation}

where $\mathcal{B}(\cdot)$ denotes the complete backtesting engine and
$\mathcal{M}$ denotes one or several performance metrics.

Robustness analysis studies

\begin{equation}
\boxed{
\mathcal{B}
\left(
\boldsymbol{\theta}
+
\Delta\boldsymbol{\theta}
\right)
}
\end{equation}

for economically plausible perturbations
$\Delta\boldsymbol{\theta}$.

A robust strategy should ideally exhibit a broad region of acceptable
performance rather than a single isolated optimum.


% ------------------------------------------------------------
\subsection{Subperiod Analysis}
\label{subsec:subperiod_analysis}
% ------------------------------------------------------------

Full-sample performance may conceal substantial variation through time.

Let the complete backtest sample be

\begin{equation}
\mathcal{T}
=
\{1,\ldots,T\}.
\end{equation}

Partition it into $J$ disjoint subperiods:

\begin{equation}
\boxed{
\mathcal{T}
=
\mathcal{T}_1
\cup
\mathcal{T}_2
\cup
\cdots
\cup
\mathcal{T}_J.
}
\end{equation}

Performance statistics can then be computed separately for each subperiod.

For example,

\begin{equation}
SR_j
=
\sqrt{A}
\frac{
\overline{R}_{p,j}
}{
\sigma_{p,j}
},
\end{equation}

where $\overline{R}_{p,j}$ and $\sigma_{p,j}$ are computed only over
$\mathcal{T}_j$.

Useful subperiod diagnostics include:

\begin{itemize}
    \item annualized return;
    \item Sharpe ratio;
    \item maximum drawdown;
    \item IC and Rank IC;
    \item turnover;
    \item gross and net performance;
    \item factor exposures.
\end{itemize}

A strategy whose entire historical performance is generated by one short
episode may be substantially less reliable than one with moderate but persistent
performance across several periods.


\subsubsection{Rolling Performance}

Fixed subperiods may depend on arbitrary start and end dates.

Rolling statistics provide a complementary view.

For a window of length $L$,

\begin{equation}
\boxed{
SR_t^{(L)}
=
\sqrt{A}
\frac{
\overline{R}_{t-L+1:t}
}{
\sigma_{t-L+1:t}
}.
}
\end{equation}

Likewise, rolling IC, drawdown, turnover, and factor exposures can be examined.

The objective is to understand whether strategy quality varies smoothly,
deteriorates structurally, or is concentrated in isolated episodes.


% ------------------------------------------------------------
\subsection{Market Regime Analysis}
\label{subsec:market_regime_analysis}
% ------------------------------------------------------------

Strategy performance may depend systematically on the market environment.

Let

\begin{equation}
Z_t
\end{equation}

denote a regime variable.

Possible regime dimensions include:

\begin{itemize}
    \item bull vs.\ bear markets;
    \item high vs.\ low volatility;
    \item rising vs.\ falling interest rates;
    \item high vs.\ low cross-sectional dispersion;
    \item high vs.\ low market correlation;
    \item expansion vs.\ recession;
    \item liquid vs.\ stressed markets.
\end{itemize}

Define regime $r$ through a rule

\begin{equation}
\mathcal{R}_r
=
\left\{
t:
Z_t\in\mathcal{Z}_r
\right\}.
\end{equation}

Conditional strategy performance is then

\begin{equation}
\boxed{
\mathbb{E}
\left[
R_{p,t}
\mid
t\in\mathcal{R}_r
\right].
}
\end{equation}


\subsubsection{Ex-Ante Regime Definition}

Regime definitions should preferably be economically motivated and specified
without reference to subsequent strategy performance.

Defining regimes after observing where the strategy performed well creates
another layer of data mining.

Regime analysis should therefore answer

\begin{equation}
\boxed{
\text{When does the strategy work and why?}
}
\end{equation}

rather than

\begin{equation}
\boxed{
\text{How can the sample be partitioned to make the strategy look good?}
}
\end{equation}


% ------------------------------------------------------------
\subsection{Alternative Investment Universes}
\label{subsec:alternative_investment_universes}
% ------------------------------------------------------------

The strategy should be tested under reasonable variations of the investment
universe.

Let

\begin{equation}
\mathcal{U}_t^{(1)},
\ldots,
\mathcal{U}_t^{(M)}
\end{equation}

denote alternative admissible universes.

Examples include:

\begin{itemize}
    \item different market-capitalization thresholds;
    \item alternative liquidity requirements;
    \item developed markets vs.\ broader global universes;
    \item large- and mid-cap universes;
    \item exclusion or inclusion of smaller exchanges;
    \item different minimum price thresholds.
\end{itemize}

For each universe,

\begin{equation}
\boxed{
\mathcal{M}^{(m)}
=
\mathcal{B}
\left(
\mathcal{U}^{(m)}
\right).
}
\end{equation}

The purpose is not to require identical performance across universes.

Instead, the analysis should determine whether the strategy's apparent alpha
depends excessively on a narrow or difficult-to-trade subset of securities.


\subsubsection{Composition Effects}

Differences across universes may arise from changes in:

\begin{itemize}
    \item sector composition;
    \item country composition;
    \item liquidity;
    \item volatility;
    \item factor exposures;
    \item signal dispersion.
\end{itemize}

Therefore, universe sensitivity should be interpreted jointly with exposure and
capacity diagnostics.


% ------------------------------------------------------------
\subsection{Alternative Rebalancing Frequencies}
\label{subsec:alternative_rebalancing_frequencies}
% ------------------------------------------------------------

The rebalance frequency determines how quickly portfolio weights respond to new
signal information.

Let

\begin{equation}
f_R
\end{equation}

denote the rebalance frequency.

Alternative specifications may include

\begin{equation}
f_R
\in
\{
\text{daily},
\text{weekly},
\text{monthly},
\ldots
\}.
\end{equation}

A higher rebalance frequency may capture signal changes more quickly but
typically increases turnover and transaction costs.

A lower frequency reduces implementation intensity but may allow the signal to
decay before positions are updated.

Thus,

\begin{equation}
\boxed{
\text{Faster Rebalancing}
\quad\Longleftrightarrow\quad
\text{Faster Signal Capture + Higher Trading Cost}.
}
\end{equation}

Both gross and net performance should therefore be compared across rebalance
frequencies.


% ------------------------------------------------------------
\subsection{Alternative Holding Periods}
\label{subsec:alternative_holding_periods}
% ------------------------------------------------------------

The holding period determines how long a portfolio or signal cohort remains
active.

Let

\begin{equation}
H
\end{equation}

denote the holding horizon.

The strategy should be evaluated for nearby values

\begin{equation}
H-\Delta H,
\qquad
H,
\qquad
H+\Delta H.
\end{equation}

A robust signal should generally exhibit an economically interpretable
performance-decay profile.

For example,

\begin{equation}
\boxed{
\text{Signal Strength}
\rightarrow
\text{Peak Predictive Horizon}
\rightarrow
\text{Gradual Decay}.
}
\end{equation}

An isolated performance spike at exactly one holding horizon may indicate
sensitivity to noise or timing conventions.

Holding-period analysis is particularly useful for distinguishing short-lived
alpha from persistent cross-sectional information.


% ------------------------------------------------------------
\subsection{Signal Lag Sensitivity}
\label{subsec:signal_lag_sensitivity}
% ------------------------------------------------------------

A powerful robustness test consists of deliberately delaying the signal before
implementation.

Suppose the baseline strategy uses

\begin{equation}
\mathbf{s}_t
\end{equation}

to determine portfolio positions.

A lagged specification uses

\begin{equation}
\boxed{
\mathbf{s}_{t-L}
}
\end{equation}

for lag

\begin{equation}
L>0.
\end{equation}

Performance can then be evaluated as a function of lag:

\begin{equation}
\boxed{
L
\longmapsto
\mathcal{M}(L).
}
\end{equation}

A gradual decline in performance as the signal becomes stale is economically
plausible.

By contrast, a strategy whose performance disappears completely after a very
small additional lag may be difficult to implement in practice.

Signal-lag sensitivity therefore provides information about:

\begin{itemize}
    \item alpha decay;
    \item execution urgency;
    \item timing robustness;
    \item sensitivity to potential look-ahead effects.
\end{itemize}


% ------------------------------------------------------------
\subsection{Outlier Treatment Sensitivity}
\label{subsec:outlier_treatment_sensitivity}
% ------------------------------------------------------------

The signal should not depend excessively on one arbitrary outlier-treatment
choice.

Suppose the baseline method applies winsorization at thresholds

\begin{equation}
(q_L,q_U).
\end{equation}

Alternative specifications may include

\begin{equation}
(0.5\%,99.5\%),
\qquad
(1\%,99\%),
\qquad
(2\%,98\%),
\end{equation}

or robust MAD-based treatments.

The complete strategy should be rerun under each specification.

The objective is to test whether

\begin{equation}
\boxed{
\text{Economic Conclusion}
}
\end{equation}

is stable even if the precise cross-sectional transformation changes.

If performance is generated almost entirely by a small number of extreme
observations, the underlying signal may be less robust than aggregate
performance statistics suggest.


% ------------------------------------------------------------
\subsection{Neutralization Sensitivity}
\label{subsec:neutralization_sensitivity}
% ------------------------------------------------------------

Signal and portfolio performance should be evaluated under alternative
neutralization specifications.

Possible comparisons include:

\begin{equation}
\boxed{
\begin{aligned}
&\text{No Neutralization},
\\
&\text{Beta Neutralization},
\\
&\text{Industry Neutralization},
\\
&\text{Beta + Industry},
\\
&\text{Beta + Industry + Size},
\\
&\text{Full Style-Factor Neutralization}.
\end{aligned}
}
\end{equation}

For each specification, one should compare:

\begin{itemize}
    \item IC and Rank IC;
    \item long--short returns;
    \item portfolio Sharpe ratio;
    \item factor exposures;
    \item turnover;
    \item signal dispersion.
\end{itemize}

This analysis helps answer whether the strategy captures independent
stock-selection information or merely repackages systematic factor exposures.


\subsubsection{OLS vs.\ Weighted Neutralization}

If weighted neutralization is used, robustness can also be evaluated across
alternative weighting matrices

\begin{equation}
\mathbf{W}_t.
\end{equation}

For example:

\begin{equation}
\mathbf{W}_t
=
\mathbf{I},
\end{equation}

versus market-capitalization-, liquidity-, or inverse-variance-weighted
neutralization.

Substantial changes in results may indicate that the definition of neutrality
is economically important to the strategy.


% ------------------------------------------------------------
\subsection{Portfolio Mapping Sensitivity}
\label{subsec:portfolio_mapping_sensitivity}
% ------------------------------------------------------------

A predictive signal should ideally remain useful under several reasonable
signal-to-weight mappings.

Suppose

\begin{equation}
\mathcal{G}^{(m)}
\end{equation}

denotes mapping rule $m$.

Examples include:

\begin{itemize}
    \item equal-weighted quantiles;
    \item top/bottom $N$;
    \item rank weighting;
    \item linear score weighting;
    \item nonlinear conviction weighting.
\end{itemize}

Then

\begin{equation}
\mathbf{w}_t^{(m)}
=
\mathcal{G}^{(m)}
\left(
\mathbf{s}_t^{\mathrm{final}}
\right).
\end{equation}

Performance can be compared across mappings:

\begin{equation}
\boxed{
\mathcal{M}^{(m)}
=
\mathcal{B}
\left(
\mathcal{G}^{(m)}
\right).
}
\end{equation}

If the signal is genuinely predictive, its economic value should not generally
depend on one extremely specific mapping function.

However, differences across mappings may reveal whether alpha is concentrated
in the signal tails or distributed smoothly across the cross-section.


% ------------------------------------------------------------
\subsection{Transaction Cost Sensitivity}
\label{subsec:transaction_cost_sensitivity}
% ------------------------------------------------------------

Transaction costs are estimated rather than known perfectly.

Net performance should therefore be evaluated under alternative cost
assumptions.

Let

\begin{equation}
c
\end{equation}

denote a proportional cost parameter.

A simple sensitivity analysis evaluates

\begin{equation}
c
\in
\{
c_{\mathrm{low}},
c_{\mathrm{base}},
c_{\mathrm{high}}
\}.
\end{equation}

More generally,

\begin{equation}
\boxed{
c
\longmapsto
SR^{\mathrm{net}}(c)
}
\end{equation}

or

\begin{equation}
\boxed{
c
\longmapsto
CAGR^{\mathrm{net}}(c)
}
\end{equation}

can be examined.

For nonlinear cost models, sensitivity may also be performed with respect to:

\begin{itemize}
    \item spread assumptions;
    \item market-impact coefficients;
    \item participation-rate assumptions;
    \item borrow fees;
    \item financing spreads.
\end{itemize}


\subsubsection{Break-Even Transaction Cost}

A useful diagnostic is the transaction-cost level at which expected net
performance falls to zero.

If gross expected return is

\begin{equation}
\mu^{\mathrm{gross}}
\end{equation}

and expected turnover is

\begin{equation}
\overline{TO},
\end{equation}

then under a simple proportional cost approximation,

\begin{equation}
\mu^{\mathrm{net}}
\approx
\mu^{\mathrm{gross}}
-
c\overline{TO}.
\end{equation}

The approximate break-even cost is therefore

\begin{equation}
\boxed{
c^{BE}
=
\frac{
\mu^{\mathrm{gross}}
}{
\overline{TO}
}.
}
\end{equation}

The exact expression depends on the turnover and cost conventions used.


% ------------------------------------------------------------
\subsection{Hyperparameter and Parameter Perturbation}
\label{subsec:parameter_perturbation}
% ------------------------------------------------------------

Many strategies depend on numerical parameters.

Examples include:

\begin{itemize}
    \item lookback length;
    \item winsorization threshold;
    \item signal-decay parameter;
    \item number of selected securities;
    \item quantile threshold;
    \item neutralization specification;
    \item volatility target;
    \item optimizer risk-aversion coefficient;
    \item regularization strength;
    \item machine-learning hyperparameters.
\end{itemize}

Let

\begin{equation}
\theta
\end{equation}

denote one such parameter.

Rather than reporting performance only at the selected value

\begin{equation}
\theta^*,
\end{equation}

the backtester should evaluate a neighborhood

\begin{equation}
\boxed{
\theta
\in
[
\theta^*-\Delta,
\theta^*+\Delta
].
}
\end{equation}

A robust strategy should ideally exhibit a relatively smooth response:

\begin{equation}
\boxed{
\theta
\longmapsto
\mathcal{M}(\theta).
}
\end{equation}


\subsubsection{Parameter Plateaus vs.\ Sharp Optima}

Suppose strategy performance behaves schematically as

\begin{equation}
\mathcal{M}(\theta).
\end{equation}

A broad plateau around $\theta^*$ is generally more reassuring than an isolated
sharp maximum.

Conceptually,

\begin{equation}
\boxed{
\text{Broad Performance Plateau}
\quad>\quad
\text{Single Sharp Historical Optimum}
}
\end{equation}

from a robustness perspective.

A highly localized optimum may indicate that the chosen parameter fits
historical noise.


\subsubsection{Joint Parameter Sensitivity}

Parameters may interact.

For two parameters

\begin{equation}
\theta_1,
\qquad
\theta_2,
\end{equation}

one may examine the response surface

\begin{equation}
\boxed{
(\theta_1,\theta_2)
\longmapsto
\mathcal{M}
\left(
\theta_1,\theta_2
\right).
}
\end{equation}

This can reveal whether apparently robust one-dimensional choices become
unstable when parameters are varied jointly.


% ------------------------------------------------------------
\subsection{Placebo and Randomization Tests}
\label{subsec:placebo_randomization_tests}
% ------------------------------------------------------------

Placebo and randomization tests ask whether the observed strategy performance
could plausibly arise even when the intended economic relationship is
destroyed.

They provide a useful complement to standard performance statistics.


\subsubsection{Cross-Sectional Signal Permutation}

At each date $t$, the signal may be randomly permuted across securities.

Let

\begin{equation}
\pi_t
\end{equation}

denote a random permutation of the $N_t$ securities.

Define the randomized signal

\begin{equation}
\boxed{
s_{i,t}^{\mathrm{perm}}
=
s_{\pi_t(i),t}^{\mathrm{final}}.
}
\end{equation}

This preserves the cross-sectional signal distribution but destroys the link
between individual securities and their scores.

The complete portfolio-construction and backtesting pipeline can then be rerun
using

\begin{equation}
\mathbf{s}_t^{\mathrm{perm}}.
\end{equation}


\subsubsection{Random Portfolio Benchmark}

Another placebo is to construct random portfolios satisfying approximately the
same implementation constraints as the strategy.

For example, random portfolios may preserve:

\begin{itemize}
    \item number of holdings;
    \item gross exposure;
    \item net exposure;
    \item sector neutrality;
    \item approximate turnover.
\end{itemize}

Let

\begin{equation}
\mathcal{M}^{(b)},
\qquad
b=1,\ldots,B,
\end{equation}

denote performance from $B$ randomized portfolios.

The strategy's observed metric

\begin{equation}
\mathcal{M}^{\mathrm{obs}}
\end{equation}

can then be compared with the empirical null distribution.


\subsubsection{Empirical Randomization Probability}

For a performance statistic where larger values are better, an empirical
randomization probability can be estimated as

\begin{equation}
\boxed{
\widehat{p}
=
\frac{
1+
\sum_{b=1}^{B}
\mathbb{I}
\left(
\mathcal{M}^{(b)}
\geq
\mathcal{M}^{\mathrm{obs}}
\right)
}{
B+1
}.
}
\end{equation}

A small value indicates that few randomized strategies achieved performance as
strong as the observed strategy.


\subsubsection{Temporal Placebos}

One may also deliberately use economically implausible signal timing.

For example, the relationship between current signals and sufficiently distant
or randomly reassigned future returns may be examined.

A valid signal should generally lose predictive power when the economic timing
relationship is deliberately broken.


\subsubsection{Purpose of Randomization Tests}

Randomization tests should not replace economic reasoning or formal statistical
inference.

Their role is to test whether the observed strategy result depends on the
intended correspondence between:

\begin{equation}
\boxed{
\text{Security}
\leftrightarrow
\text{Signal}
\leftrightarrow
\text{Future Return}.
}
\end{equation}

If comparable performance frequently survives after that relationship is
randomized, the original backtest may contain less genuine information than its
headline statistics suggest.


% ------------------------------------------------------------
\subsubsection{Robustness Summary}
% ------------------------------------------------------------

A robust strategy should not be defined by the requirement that every
alternative specification performs equally well.

Instead, the objective is to establish that the investment thesis survives
reasonable changes in the research design.

The robustness analysis should therefore answer:

\begin{equation}
\boxed{
\begin{aligned}
\text{Temporal Robustness:}
&\quad
\text{Does the strategy survive across periods and regimes?}
\\
\text{Cross-Sectional Robustness:}
&\quad
\text{Does it survive reasonable universe changes?}
\\
\text{Methodological Robustness:}
&\quad
\text{Does it survive alternative construction choices?}
\\
\text{Implementation Robustness:}
&\quad
\text{Does it survive realistic costs and timing assumptions?}
\\
\text{Statistical Robustness:}
&\quad
\text{Is performance distinguishable from plausible randomized outcomes?}
\end{aligned}
}
\end{equation}

The desired result is not a single historically optimal specification but a
stable region of economically coherent specifications producing comparable
conclusions.

This robustness layer therefore provides the transition from an attractive
backtest to a strategy whose historical evidence can be subjected to formal
statistical inference.

% ============================================================
\section{Statistical Inference}
\label{sec:statistical_inference}
% ============================================================

Performance statistics such as average return, Sharpe ratio, Information
Coefficient, or long--short spread are sample estimates.

They are therefore subject to sampling uncertainty.

A backtest that produces a positive average return does not automatically imply
that the underlying strategy has positive expected return.

The purpose of statistical inference is to assess whether the observed
performance is sufficiently large relative to its sampling variability.

Let

\begin{equation}
R_t,
\qquad
t=1,\ldots,T,
\end{equation}

denote the strategy return series.

The sample mean is

\begin{equation}
\bar{R}
=
\frac{1}{T}
\sum_{t=1}^{T}
R_t.
\end{equation}

The central inferential question is whether the population mean

\begin{equation}
\mu
=
\mathbb{E}[R_t]
\end{equation}

is statistically different from a reference value, usually zero.

However, financial backtests create several complications:

\begin{itemize}
    \item returns may be heteroskedastic;
    \item returns may be serially correlated;
    \item portfolio holding periods may overlap;
    \item many strategies may have been tested before the reported one was
    selected;
    \item parameters and hyperparameters may have been chosen using the same
    historical sample;
    \item the reported Sharpe ratio may therefore be upward biased by selection.
\end{itemize}

Consequently,

\begin{equation}
\boxed{
\text{Statistical Significance}
\neq
\text{Economic Significance}
\neq
\text{Research Robustness}.
}
\end{equation}

All three dimensions should be considered jointly.


% ------------------------------------------------------------
\subsection{Mean Return Tests and $t$-Statistics}
\label{subsec:mean_return_tstats}
% ------------------------------------------------------------

Consider the null hypothesis

\begin{equation}
\boxed{
H_0:
\mu
=
\mu_0,
}
\end{equation}

against an alternative such as

\begin{equation}
H_1:
\mu
>
\mu_0.
\end{equation}

Usually,

\begin{equation}
\mu_0=0.
\end{equation}

Under the simplest independent and identically distributed assumption, the
sample variance is

\begin{equation}
\widehat{\sigma}^{2}
=
\frac{1}{T-1}
\sum_{t=1}^{T}
(R_t-\bar{R})^2.
\end{equation}

The standard error of the sample mean is

\begin{equation}
\boxed{
SE(\bar{R})
=
\frac{
\widehat{\sigma}
}{
\sqrt{T}
}.
}
\end{equation}

The corresponding $t$-statistic is

\begin{equation}
\boxed{
t
=
\frac{
\bar{R}-\mu_0
}{
SE(\bar{R})
}
=
\frac{
\sqrt{T}
(\bar{R}-\mu_0)
}{
\widehat{\sigma}
}.
}
\end{equation}


\subsubsection{Connection with the Sharpe Ratio}

If

\begin{equation}
\mu_0=0,
\end{equation}

the non-annualized sample Sharpe ratio is

\begin{equation}
\widehat{SR}
=
\frac{
\bar{R}
}{
\widehat{\sigma}
}.
\end{equation}

Therefore,

\begin{equation}
\boxed{
t
=
\sqrt{T}
\widehat{SR}.
}
\end{equation}

If the Sharpe ratio is annualized using $A$ observations per year,

\begin{equation}
\widehat{SR}_{\mathrm{ann}}
=
\sqrt{A}
\frac{
\bar{R}
}{
\widehat{\sigma}
},
\end{equation}

then

\begin{equation}
\boxed{
t
=
\sqrt{
\frac{T}{A}
}
\widehat{SR}_{\mathrm{ann}}.
}
\end{equation}

Thus, statistical evidence depends jointly on the magnitude of the Sharpe ratio
and the length of the historical sample.


\subsubsection{One-Sided vs.\ Two-Sided Tests}

A two-sided hypothesis test uses

\begin{equation}
H_1:
\mu\neq0,
\end{equation}

whereas a one-sided test may use

\begin{equation}
H_1:
\mu>0.
\end{equation}

A one-sided test is appropriate only if the direction of the alternative was
specified before observing the data.

Changing from a two-sided to a one-sided test after observing positive
performance constitutes an additional research choice and can overstate
statistical evidence.


% ------------------------------------------------------------
\subsection{Serial Dependence and Heteroskedasticity}
\label{subsec:serial_dependence_heteroskedasticity}
% ------------------------------------------------------------

The standard error

\begin{equation}
SE(\bar{R})
=
\frac{
\widehat{\sigma}
}{
\sqrt{T}
}
\end{equation}

is valid under restrictive assumptions.

Financial strategy returns may exhibit:

\begin{equation}
\boxed{
\text{Heteroskedasticity}
}
\end{equation}

and

\begin{equation}
\boxed{
\text{Serial Dependence}.
}
\end{equation}


\subsubsection{Heteroskedasticity}

Heteroskedasticity means that conditional variance changes through time:

\begin{equation}
\operatorname{Var}
\left(
R_t
\mid
\mathcal{I}_{t-1}
\right)
=
\sigma_t^2,
\end{equation}

where

\begin{equation}
\sigma_t^2
\end{equation}

is not constant.

This is common in financial returns because volatility tends to cluster.


\subsubsection{Serial Correlation}

Serial correlation means

\begin{equation}
\boxed{
\operatorname{Cov}
\left(
R_t,
R_{t-k}
\right)
\neq
0
}
\end{equation}

for some lag $k>0$.

Define the autocovariance at lag $k$ as

\begin{equation}
\gamma_k
=
\operatorname{Cov}
\left(
R_t,
R_{t-k}
\right).
\end{equation}

The variance of the sample mean is then not simply

\begin{equation}
\frac{\sigma^2}{T}.
\end{equation}

More generally,

\begin{equation}
\operatorname{Var}(\bar{R})
=
\frac{1}{T^2}
\sum_{t=1}^{T}
\sum_{s=1}^{T}
\operatorname{Cov}(R_t,R_s).
\end{equation}

For a weakly stationary process, this can be written approximately as

\begin{equation}
\boxed{
\operatorname{Var}(\bar{R})
\approx
\frac{1}{T}
\left[
\gamma_0
+
2
\sum_{k=1}^{\infty}
\gamma_k
\right].
}
\end{equation}

The term

\begin{equation}
\gamma_0
+
2
\sum_{k=1}^{\infty}
\gamma_k
\end{equation}

is the long-run variance.


\subsubsection{Overlapping Returns}

Serial dependence is especially important when holding periods overlap.

Suppose a strategy forms daily forecasts of future $h$-day returns:

\begin{equation}
R_{t\rightarrow t+h}.
\end{equation}

Adjacent observations,

\begin{equation}
R_{t\rightarrow t+h}
\end{equation}

and

\begin{equation}
R_{t+1\rightarrow t+h+1},
\end{equation}

share most of their underlying daily returns.

Therefore, the observations are mechanically dependent.

Ignoring this dependence generally understates the standard error and
overstates the corresponding $t$-statistic.


% ------------------------------------------------------------
\subsection{Newey--West Standard Errors}
\label{subsec:newey_west}
% ------------------------------------------------------------

Newey--West standard errors provide a
\emph{heteroskedasticity- and autocorrelation-consistent} estimate of the
sampling variance.

For inference on the mean return, define residuals

\begin{equation}
u_t
=
R_t-\bar{R}.
\end{equation}

The sample autocovariance at lag $k$ is

\begin{equation}
\boxed{
\widehat{\gamma}_k
=
\frac{1}{T}
\sum_{t=k+1}^{T}
u_tu_{t-k}.
}
\end{equation}

Choose a truncation lag

\begin{equation}
L\geq0.
\end{equation}

The Newey--West estimate of the long-run variance is

\begin{equation}
\boxed{
\widehat{\Omega}_{NW}
=
\widehat{\gamma}_0
+
2
\sum_{k=1}^{L}
\omega_k
\widehat{\gamma}_k,
}
\end{equation}

where the Bartlett weights are

\begin{equation}
\boxed{
\omega_k
=
1
-
\frac{k}{L+1}.
}
\end{equation}

The HAC standard error of the sample mean is then

\begin{equation}
\boxed{
SE_{NW}(\bar{R})
=
\sqrt{
\frac{
\widehat{\Omega}_{NW}
}{
T
}
}.
}
\end{equation}

The corresponding $t$-statistic is

\begin{equation}
\boxed{
t_{NW}
=
\frac{
\bar{R}-\mu_0
}{
SE_{NW}(\bar{R})
}.
}
\end{equation}


\subsubsection{Choice of Lag Length}

The lag parameter $L$ determines how much serial dependence is incorporated.

If $L=0$, the estimator allows for heteroskedasticity but ignores serial
correlation.

As $L$ increases, more autocovariances are included.

For overlapping $h$-period returns, a natural minimum diagnostic choice is
often related to

\begin{equation}
h-1,
\end{equation}

although the appropriate lag depends on the exact data-generating process and
sampling frequency.

Robustness to alternative reasonable lag choices should therefore be checked.


% ------------------------------------------------------------
\subsection{Bootstrap Methods}
\label{subsec:bootstrap_methods}
% ------------------------------------------------------------

Bootstrap methods approximate the sampling distribution of a statistic by
resampling the observed data.

Let

\begin{equation}
\widehat{\theta}
=
g
\left(
R_1,\ldots,R_T
\right)
\end{equation}

denote a statistic such as:

\begin{itemize}
    \item mean return;
    \item Sharpe ratio;
    \item maximum drawdown;
    \item Information Ratio;
    \item factor alpha.
\end{itemize}

A bootstrap generates replicated samples

\begin{equation}
\left\{
R_t^{*(b)}
\right\}_{t=1}^{T},
\qquad
b=1,\ldots,B,
\end{equation}

and corresponding statistics

\begin{equation}
\boxed{
\widehat{\theta}^{*(b)}
=
g
\left(
R_1^{*(b)},
\ldots,
R_T^{*(b)}
\right).
}
\end{equation}

The empirical distribution

\begin{equation}
\left\{
\widehat{\theta}^{*(1)},
\ldots,
\widehat{\theta}^{*(B)}
\right\}
\end{equation}

approximates the sampling distribution of the statistic.


\subsubsection{IID Bootstrap}

The simplest bootstrap independently resamples individual return observations.

This is appropriate only when serial dependence is negligible.

If returns are autocorrelated, independently resampling observations destroys
the temporal dependence structure.


\subsubsection{Block Bootstrap}

To preserve local serial dependence, consecutive blocks of observations can be
resampled.

Let the block length be

\begin{equation}
L_B.
\end{equation}

A block bootstrap resamples sequences such as

\begin{equation}
(R_t,\ldots,R_{t+L_B-1})
\end{equation}

rather than isolated observations.

Common variants include:

\begin{itemize}
    \item moving-block bootstrap;
    \item circular-block bootstrap;
    \item stationary bootstrap.
\end{itemize}

Block methods are generally more appropriate for financial strategies with
serially dependent returns.


\subsubsection{Bootstrap Confidence Intervals}

For confidence level $1-\alpha$, a simple percentile interval is

\begin{equation}
\boxed{
\left[
q_{\alpha/2}
\left(
\widehat{\theta}^{*}
\right),
\;
q_{1-\alpha/2}
\left(
\widehat{\theta}^{*}
\right)
\right].
}
\end{equation}

This allows uncertainty to be quantified for statistics whose analytical
sampling distributions may be difficult to derive.


% ------------------------------------------------------------
\subsection{Multiple Hypothesis Testing}
\label{subsec:multiple_hypothesis_testing}
% ------------------------------------------------------------

Backtesting frequently involves testing many candidate strategies.

Suppose

\begin{equation}
M
\end{equation}

hypotheses are tested:

\begin{equation}
H_{0,1},
\ldots,
H_{0,M}.
\end{equation}

Even if every null hypothesis is true, some strategies will appear significant
by chance.

If each independent hypothesis is tested at significance level $\alpha$, the
probability of obtaining at least one false positive is

\begin{equation}
\boxed{
P(\text{at least one false positive})
=
1-(1-\alpha)^M.
}
\end{equation}

As

\begin{equation}
M\rightarrow\infty,
\end{equation}

this probability approaches one.


\subsubsection{Bonferroni Correction}

A conservative adjustment is to use the individual significance threshold

\begin{equation}
\boxed{
\alpha_{\mathrm{ind}}
=
\frac{
\alpha
}{
M
}.
}
\end{equation}

Equivalently, raw $p$-values may be multiplied by $M$:

\begin{equation}
\boxed{
p_i^{\mathrm{adj}}
=
\min
\left(
Mp_i,
1
\right).
}
\end{equation}


\subsubsection{False Discovery Rate}

Rather than controlling the probability of any false positive, False Discovery
Rate methods control the expected proportion of false discoveries among the
rejected hypotheses.

This can be less conservative when many strategies or signals are evaluated.

The appropriate correction depends on the research objective and dependence
structure among tested hypotheses.


% ------------------------------------------------------------
\subsection{Data Snooping Adjustments}
\label{subsec:data_snooping_adjustments}
% ------------------------------------------------------------

Multiple hypothesis testing becomes particularly problematic when only the best
historical strategy is ultimately reported.

Suppose a researcher evaluates

\begin{equation}
M
\end{equation}

strategies and selects

\begin{equation}
m^*
=
\underset{m=1,\ldots,M}{\arg\max}
\;
\widehat{SR}_m.
\end{equation}

Then

\begin{equation}
\boxed{
\mathbb{E}
\left[
\widehat{SR}_{m^*}
\right]
}
\end{equation}

is upward biased relative to the performance of a strategy chosen independently
of the historical sample.

This phenomenon is often referred to as data snooping or selection bias.


\subsubsection{What Counts as a Trial?}

The effective number of trials is not limited to completely different trading
strategies.

Research choices may include:

\begin{itemize}
    \item different characteristics;
    \item alternative lookback windows;
    \item alternative winsorization thresholds;
    \item alternative neutralization specifications;
    \item alternative portfolio mappings;
    \item different investment universes;
    \item different model hyperparameters;
    \item different transaction-cost assumptions.
\end{itemize}

If these choices are influenced by backtest performance, they contribute to the
effective amount of data snooping.


\subsubsection{Dependent Trials}

Candidate strategies are often highly correlated.

Therefore,

\begin{equation}
M_{\mathrm{effective}}
<
M
\end{equation}

may hold when many tested strategies are only minor variations of one another.

Nevertheless, correlation between trials does not eliminate the selection
problem.

The statistical correction should ideally reflect both the number and
dependence structure of the experiments performed.


% ------------------------------------------------------------
\subsection{Deflated Sharpe Ratio}
\label{subsec:deflated_sharpe_ratio}
% ------------------------------------------------------------

The Sharpe ratio of a selected backtest is generally upward biased when the
reported strategy has been chosen after testing multiple alternatives.

Suppose the researcher evaluates $N$ candidate strategy specifications and
selects the strategy with the highest observed Sharpe ratio:

\begin{equation}
m^*
=
\arg\max_{m=1,\ldots,N}
\widehat{SR}_m.
\end{equation}

Even if none of the candidate strategies possesses genuine alpha, the maximum

\begin{equation}
\max_m \widehat{SR}_m
\end{equation}

will generally be positive simply because several noisy estimates were examined.

The Deflated Sharpe Ratio (DSR) addresses this problem by asking whether the
Sharpe ratio of the selected strategy is sufficiently large relative to a
benchmark that accounts for:

\begin{itemize}
    \item finite-sample estimation uncertainty;
    \item non-Gaussian returns;
    \item multiple strategy trials;
    \item selection of the best-performing specification.
\end{itemize}

The logic can be summarized as

\begin{equation}
\boxed{
\text{Observed Sharpe}
\rightarrow
\text{Multiple-Trial Benchmark}
\rightarrow
\text{Sampling Adjustment}
\rightarrow
\text{Deflated Sharpe Ratio}.
}
\end{equation}


% ------------------------------------------------------------
\subsubsection{Observed Sharpe Ratio}
% ------------------------------------------------------------

Let

\begin{equation}
\widehat{SR}
\end{equation}

denote the Sharpe ratio of the strategy ultimately selected by the research
process.

All Sharpe ratios used in the DSR calculation must follow the same convention.

In particular, they must use the same:

\begin{itemize}
    \item return frequency;
    \item annualization convention;
    \item risk-free-rate convention;
    \item gross or net return convention;
    \item evaluation sample.
\end{itemize}

For example, if daily returns are used,

\begin{equation}
\widehat{SR}^{\mathrm{ann}}
=
\sqrt{252}
\frac{
\bar{R}
}{
\widehat{\sigma}(R)
}.
\end{equation}

The same annualization must be applied to the Sharpe ratios of all alternative
strategy trials.


% ------------------------------------------------------------
\subsubsection{What Constitutes a Strategy Trial?}
% ------------------------------------------------------------

A trial is not restricted to a completely different investment strategy.

Any research decision that was influenced by observed historical performance
may constitute an additional trial.

Examples include testing alternative:

\begin{itemize}
    \item characteristics or signals;
    \item lookback windows;
    \item signal lags;
    \item winsorization thresholds;
    \item neutralization specifications;
    \item signal-combination weights;
    \item portfolio mapping functions;
    \item numbers of selected securities;
    \item rebalance frequencies;
    \item investment universes;
    \item volatility targets;
    \item optimization parameters;
    \item model specifications;
    \item machine-learning hyperparameters.
\end{itemize}

Consequently, if the final specification was selected from a collection of
backtests

\begin{equation}
\mathcal{S}
=
\left\{
S_1,\ldots,S_N
\right\},
\end{equation}

the complete set of economically relevant trials should ideally be retained.

This is an important implementation requirement.

A research framework should therefore store, for every trial,

\begin{equation}
\boxed{
\left(
\text{Configuration},
\text{Return Series},
\widehat{SR},
\text{Experiment ID}
\right).
}
\end{equation}

Retaining only the final successful backtest makes subsequent correction for
research selection substantially more difficult.


% ------------------------------------------------------------
\subsubsection{Multiple-Trial Benchmark Sharpe Ratio}
% ------------------------------------------------------------

The DSR does not generally test the selected Sharpe ratio against zero.

Instead, it compares the observed Sharpe ratio with a benchmark

\begin{equation}
\boxed{
SR_0
}
\end{equation}

representing the Sharpe ratio that the best-performing strategy could reasonably
achieve purely because multiple alternatives were tested.

Let

\begin{equation}
\widehat{SR}_1,
\ldots,
\widehat{SR}_N
\end{equation}

denote the Sharpe ratios of the candidate strategies.

Their sample variance is

\begin{equation}
\boxed{
\widehat{V}[SR]
=
\frac{1}{N-1}
\sum_{m=1}^{N}
\left(
\widehat{SR}_m
-
\overline{SR}
\right)^2,
}
\end{equation}

where

\begin{equation}
\overline{SR}
=
\frac{1}{N}
\sum_{m=1}^{N}
\widehat{SR}_m.
\end{equation}

Let

\begin{equation}
\sigma_{SR,\mathrm{trials}}
=
\sqrt{
\widehat{V}[SR]
}.
\end{equation}

For $N$ approximately independent trials, the expected maximum Sharpe ratio can
be approximated by

\begin{equation}
\boxed{
SR_0
\approx
\sigma_{SR,\mathrm{trials}}
\left[
(1-\gamma)
\Phi^{-1}
\left(
1-\frac{1}{N}
\right)
+
\gamma
\Phi^{-1}
\left(
1-\frac{1}{Ne}
\right)
\right],
}
\end{equation}

where

\begin{equation}
\gamma
\approx
0.57721566
\end{equation}

is the Euler--Mascheroni constant,

\begin{equation}
e
\approx
2.71828,
\end{equation}

and $\Phi^{-1}(\cdot)$ denotes the inverse standard-normal cumulative
distribution function.

Hence,

\begin{equation}
\boxed{
SR_0
=
\text{approximate Sharpe ratio expected from the best trial under repeated
search}.
}
\end{equation}

As either the number of trials or the dispersion of their Sharpe ratios
increases, the benchmark $SR_0$ increases.


% ------------------------------------------------------------
\subsubsection{Why the Raw Number of Trials May Be Misleading}
% ------------------------------------------------------------

The previous expression assumes approximately independent strategy trials.

This assumption is often unrealistic.

For example, strategies using momentum lookbacks of

\begin{equation}
240,\;241,\;242,\;\ldots,\;260
\end{equation}

days will generally produce highly correlated return series.

Treating these $21$ specifications as $21$ independent experiments exaggerates
the amount of independent information contained in the search.

Conversely, simply treating them as one trial ignores the fact that the
researcher still selected among several alternatives.

It is therefore useful to distinguish

\begin{equation}
N
=
\text{number of recorded strategy trials}
\end{equation}

from

\begin{equation}
\boxed{
N_{\mathrm{eff}}
=
\text{effective number of approximately independent trials}.
}
\end{equation}

In practice,

\begin{equation}
1
\leq
N_{\mathrm{eff}}
\leq
N.
\end{equation}


% ------------------------------------------------------------
\subsubsection{Estimating the Dependence Between Strategy Trials}
% ------------------------------------------------------------

Suppose the research process stores the return series of all $N$ candidate
strategies over a common sample of $T$ observations.

Construct the matrix

\begin{equation}
\mathbf{R}
=
\begin{bmatrix}
R_{1,1} & R_{2,1} & \cdots & R_{N,1}
\\
R_{1,2} & R_{2,2} & \cdots & R_{N,2}
\\
\vdots & \vdots & \ddots & \vdots
\\
R_{1,T} & R_{2,T} & \cdots & R_{N,T}
\end{bmatrix}
\in
\mathbb{R}^{T\times N},
\end{equation}

where column $m$ contains the return series of strategy trial $m$.

Estimate the trial correlation matrix

\begin{equation}
\boxed{
\widehat{\mathbf{C}}
=
\operatorname{Corr}
\left(
\mathbf{R}
\right)
\in
\mathbb{R}^{N\times N}.
}
\end{equation}

The element

\begin{equation}
\widehat{C}_{ij}
\end{equation}

measures the historical return correlation between strategy trials $i$ and $j$.

Highly correlated trials contain less independent information than unrelated
trials.


% ------------------------------------------------------------
\subsubsection{Effective Number of Trials from the Correlation Spectrum}
% ------------------------------------------------------------

There is no unique universally correct estimator of the effective number of
independent strategy trials.

For implementation, a transparent diagnostic can be constructed from the
eigenvalues of the strategy-return correlation matrix.

Let

\begin{equation}
\lambda_1,\ldots,\lambda_N
\end{equation}

denote the eigenvalues of

\begin{equation}
\widehat{\mathbf{C}}.
\end{equation}

Because $\widehat{\mathbf{C}}$ is a correlation matrix,

\begin{equation}
\sum_{j=1}^{N}\lambda_j
=
N.
\end{equation}

If all strategy trials are independent,

\begin{equation}
\lambda_j=1
\qquad
\forall j,
\end{equation}

and the effective number of trials should be close to $N$.

If all trials are nearly identical, one eigenvalue dominates while the
remaining eigenvalues are close to zero, implying an effective number much
closer to one.


\subsubsection{Participation-Ratio Effective Rank}

A convenient estimator is the participation ratio

\begin{equation}
\boxed{
N_{\mathrm{eff}}^{PR}
=
\frac{
\left(
\sum_{j=1}^{N}\lambda_j
\right)^2
}{
\sum_{j=1}^{N}\lambda_j^2
}.
}
\end{equation}

Since

\begin{equation}
\sum_{j=1}^{N}\lambda_j=N,
\end{equation}

this becomes

\begin{equation}
\boxed{
N_{\mathrm{eff}}^{PR}
=
\frac{
N^2
}{
\sum_{j=1}^{N}\lambda_j^2
}.
}
\end{equation}

This estimator has useful limiting behavior.

If all strategies are independent,

\begin{equation}
\lambda_1=\cdots=\lambda_N=1,
\end{equation}

and therefore

\begin{equation}
N_{\mathrm{eff}}^{PR}
=
N.
\end{equation}

If all strategies are perfectly correlated,

\begin{equation}
\lambda_1=N,
\qquad
\lambda_2=\cdots=\lambda_N=0,
\end{equation}

and

\begin{equation}
N_{\mathrm{eff}}^{PR}
=
1.
\end{equation}


\subsubsection{Entropy Effective Rank}

An alternative is the entropy effective rank.

Normalize the eigenvalues as

\begin{equation}
p_j
=
\frac{
\lambda_j
}{
\sum_{k=1}^{N}\lambda_k
}.
\end{equation}

Then define

\begin{equation}
\boxed{
N_{\mathrm{eff}}^{ER}
=
\exp
\left(
-
\sum_{j=1}^{N}
p_j\log p_j
\right).
}
\end{equation}

Again,

\begin{equation}
1
\leq
N_{\mathrm{eff}}^{ER}
\leq
N.
\end{equation}

The participation ratio and entropy effective rank are not identical estimators,
but both summarize how many independent dimensions are present in the collection
of strategy returns.


% ------------------------------------------------------------
\subsubsection{Practical Default for the Backtester}
% ------------------------------------------------------------

For a programmatic implementation, a practical workflow is:

\begin{enumerate}
    \item Store the return series of every strategy trial.
    
    \item Align all trial returns on a common evaluation sample.
    
    \item Compute the $N\times N$ strategy-return correlation matrix
    $\widehat{\mathbf{C}}$.
    
    \item Compute its eigenvalues
    $\lambda_1,\ldots,\lambda_N$.
    
    \item Compute
    $N_{\mathrm{eff}}^{PR}$ and, optionally,
    $N_{\mathrm{eff}}^{ER}$.
    
    \item Report both the literal number of trials $N$ and the estimated
    effective number $N_{\mathrm{eff}}$.
    
    \item Use the chosen $N_{\mathrm{eff}}$ in the multiple-trial benchmark
    calculation.
\end{enumerate}

Thus, the benchmark becomes

\begin{equation}
\boxed{
SR_0
\approx
\sigma_{SR,\mathrm{trials}}
\left[
(1-\gamma)
\Phi^{-1}
\left(
1-\frac{1}{N_{\mathrm{eff}}}
\right)
+
\gamma
\Phi^{-1}
\left(
1-\frac{1}{N_{\mathrm{eff}}e}
\right)
\right].
}
\end{equation}

Since $N_{\mathrm{eff}}$ need not be an integer, the formula can be evaluated
directly using its real-valued estimate.


\subsubsection{Important Interpretation of $N_{\mathrm{eff}}$}

The effective-rank procedure measures dependence among the \emph{return series}
of the recorded strategy trials.

It does not reconstruct experiments that were never logged.

Therefore,

\begin{equation}
\boxed{
N_{\mathrm{eff}}
\text{ can correct for dependence among recorded trials, but not for
unrecorded research decisions.}
}
\end{equation}

This is why systematic experiment tracking is part of statistical validation
rather than merely a software-engineering convenience.


% ------------------------------------------------------------
\subsubsection{Standard Error of the Selected Sharpe Ratio}
% ------------------------------------------------------------

Once $SR_0$ has been determined, the uncertainty of the selected Sharpe ratio
must be estimated.

Let

\begin{equation}
\widehat{\gamma}_3
\end{equation}

denote the sample skewness of the selected strategy returns and

\begin{equation}
\widehat{\gamma}_4
\end{equation}

their sample kurtosis.

Under a commonly used asymptotic approximation,

\begin{equation}
\boxed{
\widehat{\sigma}_{SR}
=
\sqrt{
\frac{
1
-
\widehat{\gamma}_3
\widehat{SR}
+
\frac{
\widehat{\gamma}_4-1
}{4}
\widehat{SR}^{2}
}{
T-1
}
}.
}
\end{equation}

The presence of skewness and kurtosis means that the uncertainty of the Sharpe
ratio is not determined solely by sample size.


% ------------------------------------------------------------
\subsubsection{Computing the Deflated Sharpe Ratio}
% ------------------------------------------------------------

The standardized statistic is

\begin{equation}
\boxed{
Z_{DSR}
=
\frac{
\widehat{SR}-SR_0
}{
\widehat{\sigma}_{SR}
}.
}
\end{equation}

The Deflated Sharpe Ratio is

\begin{equation}
\boxed{
DSR
=
\Phi
\left(
Z_{DSR}
\right),
}
\end{equation}

where $\Phi(\cdot)$ denotes the standard-normal cumulative distribution
function.

Equivalently,

\begin{equation}
\boxed{
DSR
=
\Phi
\left(
\frac{
\widehat{SR}-SR_0
}{
\widehat{\sigma}_{SR}
}
\right).
}
\end{equation}

A larger value indicates stronger evidence that the selected strategy's Sharpe
ratio exceeds the level that could reasonably be expected from the research
selection process.


% ------------------------------------------------------------
\subsubsection{Implementation Algorithm}
% ------------------------------------------------------------

The complete computation can therefore be implemented as follows.

Given:

\begin{equation}
\mathbf{R}
\in
\mathbb{R}^{T\times N},
\end{equation}

the matrix of returns from all recorded strategy trials:

\begin{enumerate}
    \item Compute the Sharpe ratio of each trial:
    
    \begin{equation}
    \widehat{SR}_1,\ldots,\widehat{SR}_N.
    \end{equation}
    
    \item Identify the selected strategy and its Sharpe ratio
    
    \begin{equation}
    \widehat{SR}.
    \end{equation}
    
    \item Compute the cross-trial Sharpe variance
    
    \begin{equation}
    \widehat{V}[SR].
    \end{equation}
    
    \item Compute the trial-return correlation matrix
    
    \begin{equation}
    \widehat{\mathbf{C}}
    =
    \operatorname{Corr}(\mathbf{R}).
    \end{equation}
    
    \item Compute the eigenvalues of $\widehat{\mathbf{C}}$.
    
    \item Estimate the effective number of trials, for example using
    
    \begin{equation}
    N_{\mathrm{eff}}
    =
    \frac{
    \left(\sum_j\lambda_j\right)^2
    }{
    \sum_j\lambda_j^2
    }.
    \end{equation}
    
    \item Compute the multiple-trial benchmark
    
    \begin{equation}
    SR_0.
    \end{equation}
    
    \item Compute skewness and kurtosis of the selected strategy:
    
    \begin{equation}
    \widehat{\gamma}_3,
    \qquad
    \widehat{\gamma}_4.
    \end{equation}
    
    \item Compute
    
    \begin{equation}
    \widehat{\sigma}_{SR}.
    \end{equation}
    
    \item Finally compute
    
    \begin{equation}
    \boxed{
    DSR
    =
    \Phi
    \left(
    \frac{
    \widehat{SR}-SR_0
    }{
    \widehat{\sigma}_{SR}
    }
    \right).
    }
    \end{equation}
\end{enumerate}


% ------------------------------------------------------------
\subsubsection{Concrete Example}
% ------------------------------------------------------------

Suppose the research process contains

\begin{equation}
N=100
\end{equation}

recorded backtests.

Because many specifications are similar, their return series are strongly
correlated.

Suppose the correlation-spectrum calculation produces

\begin{equation}
N_{\mathrm{eff}}=25.
\end{equation}

Assume further that the standard deviation of the Sharpe ratios across trials is

\begin{equation}
\sigma_{SR,\mathrm{trials}}
=
0.40.
\end{equation}

The relevant benchmark is then computed using

\begin{equation}
N_{\mathrm{eff}}=25,
\end{equation}

rather than mechanically treating all $100$ trials as independent.

Thus,

\begin{equation}
SR_0
\approx
0.40
\left[
(1-\gamma)
\Phi^{-1}
\left(
1-\frac{1}{25}
\right)
+
\gamma
\Phi^{-1}
\left(
1-\frac{1}{25e}
\right)
\right].
\end{equation}

If the selected strategy has

\begin{equation}
\widehat{SR}=1.50,
\end{equation}

the relevant question is therefore not

\begin{equation}
\boxed{
\text{Is }1.50>0?
}
\end{equation}

but rather

\begin{equation}
\boxed{
\text{Is }1.50\text{ sufficiently above }SR_0
\text{ given the sample size and return distribution?}
}
\end{equation}


% ------------------------------------------------------------
\subsubsection{Interpretation and Reporting}
% ------------------------------------------------------------

The DSR should not be reported alone.

A useful backtest report should include at least

\begin{equation}
\boxed{
\left(
\widehat{SR},
SR_0,
DSR,
N,
N_{\mathrm{eff}},
T,
\widehat{\gamma}_3,
\widehat{\gamma}_4
\right).
}
\end{equation}

This makes the statistical correction auditable.

A large raw Sharpe ratio becomes less convincing when:

\begin{itemize}
    \item many independent strategy variants were tested;
    \item Sharpe ratios vary substantially across trials;
    \item the selected strategy was chosen after extensive optimization;
    \item the historical sample is short;
    \item returns exhibit substantial skewness or kurtosis.
\end{itemize}

Conversely, a strategy whose Sharpe ratio remains convincingly above the
multiple-trial benchmark provides stronger evidence that the observed
performance is not merely the result of selecting the best outcome from a large
research search.

The DSR therefore changes the inferential question from

\begin{equation}
\boxed{
H_0:
SR=0
}
\end{equation}

to the substantially more demanding question

\begin{equation}
\boxed{
H_0:
SR\leq SR_0,
}
\end{equation}

where $SR_0$ reflects the advantage created by the research and strategy
selection process.



% ------------------------------------------------------------
\subsection{Probability of Backtest Overfitting}
\label{subsec:probability_backtest_overfitting}
% ------------------------------------------------------------

The Probability of Backtest Overfitting, or PBO, evaluates how frequently a
strategy selected as the best performer in-sample performs relatively poorly
out of sample.

The central idea is to repeatedly divide the available historical observations
into complementary training and test subsets.


\subsubsection{Candidate Strategies}

Suppose there are

\begin{equation}
M
\end{equation}

candidate strategy specifications.

For each partition $b$, compute an in-sample performance statistic

\begin{equation}
S_{m,b}^{IS}
\end{equation}

and an out-of-sample statistic

\begin{equation}
S_{m,b}^{OOS}.
\end{equation}

Select the best in-sample strategy:

\begin{equation}
\boxed{
m_b^*
=
\underset{m}{\arg\max}
\;
S_{m,b}^{IS}.
}
\end{equation}

Then evaluate the relative out-of-sample rank of this selected strategy.


\subsubsection{Out-of-Sample Relative Rank}

Let

\begin{equation}
r_b^{OOS}
\in
(0,1)
\end{equation}

denote the percentile rank of strategy $m_b^*$ among all candidate strategies
in the corresponding out-of-sample subset.

If

\begin{equation}
r_b^{OOS}
<
\frac{1}{2},
\end{equation}

the strategy selected as best in-sample performs below the out-of-sample median.

PBO can therefore be estimated as

\begin{equation}
\boxed{
PBO
=
P
\left(
r_b^{OOS}
<
\frac{1}{2}
\right).
}
\end{equation}

Empirically,

\begin{equation}
\boxed{
\widehat{PBO}
=
\frac{1}{B}
\sum_{b=1}^{B}
\mathbb{I}
\left(
r_b^{OOS}
<
\frac{1}{2}
\right).
}
\end{equation}


\subsubsection{Logit Representation}

A common transformation of the out-of-sample relative rank is

\begin{equation}
\boxed{
\lambda_b
=
\log
\left(
\frac{
r_b^{OOS}
}{
1-r_b^{OOS}
}
\right).
}
\end{equation}

Then

\begin{equation}
\lambda_b<0
\end{equation}

corresponds to an out-of-sample rank below the median.

Thus,

\begin{equation}
\boxed{
PBO
=
P(\lambda_b<0).
}
\end{equation}


\subsubsection{Interpretation}

A high PBO indicates that choosing the historically best specification frequently
leads to below-median performance on data not used for selection.

This is direct evidence that the research-selection process may be fitting
historical noise.

A low PBO does not prove that the strategy will perform well in live trading,
but it provides evidence that the selected strategy is less dependent on
in-sample optimization.


% ------------------------------------------------------------
\subsubsection{Statistical Inference Summary}
% ------------------------------------------------------------

The statistical validation process should progressively answer four questions.

First,

\begin{equation}
\boxed{
\text{Is average performance statistically distinguishable from zero?}
}
\end{equation}

This is addressed through mean-return tests and appropriate standard errors.

Second,

\begin{equation}
\boxed{
\text{Are the reported standard errors robust to dependence and
heteroskedasticity?}
}
\end{equation}

This motivates Newey--West and bootstrap procedures.

Third,

\begin{equation}
\boxed{
\text{How much strategy selection and repeated testing occurred?}
}
\end{equation}

This motivates multiple-testing and data-snooping adjustments.

Finally,

\begin{equation}
\boxed{
\text{Would the strategy still appear attractive after accounting for the
research process that produced it?}
}
\end{equation}

This motivates tools such as the Deflated Sharpe Ratio and the Probability of
Backtest Overfitting.

The complete progression can therefore be summarized as

\begin{equation}
\boxed{
\begin{aligned}
\text{Observed Performance}
&\rightarrow
\text{Sampling Uncertainty}
\\
&\rightarrow
\text{Dependence-Adjusted Inference}
\\
&\rightarrow
\text{Multiple-Testing Adjustment}
\\
&\rightarrow
\text{Backtest-Selection Adjustment}
\\
&\rightarrow
\text{Statistically Credible Evidence}.
\end{aligned}
}
\end{equation}

Statistical significance should nevertheless remain only one component of the
research decision.

A credible investment strategy should simultaneously exhibit economically
meaningful performance, robustness across reasonable specifications, realistic
implementation assumptions, and statistical evidence that is not explained
primarily by repeated historical experimentation.


% ============================================================
\section{Reproducibility and Production Architecture}
\label{sec:reproducibility_production}
% ============================================================

A backtest is scientifically useful only if its results can be reproduced.

It is operationally useful only if the same logic can be transferred into a
production environment without changing the economic behavior of the strategy.

The objective of the production architecture is therefore to ensure that every
historical result can be traced back to a precise combination of:

\begin{itemize}
    \item source code;
    \item configuration;
    \item input data;
    \item model version;
    \item software environment;
    \item random state;
    \item experiment metadata.
\end{itemize}

Conceptually, a backtest result should be reproducible from a complete experiment
state

\begin{equation}
\boxed{
\mathcal{E}
=
\left(
\mathcal{C},
\mathcal{D},
\mathcal{M},
\mathcal{H},
\mathcal{R},
\mathcal{S}
\right),
}
\end{equation}

where:

\begin{itemize}
    \item $\mathcal{C}$ denotes code version;
    \item $\mathcal{D}$ denotes data version;
    \item $\mathcal{M}$ denotes model version;
    \item $\mathcal{H}$ denotes configuration and hyperparameters;
    \item $\mathcal{R}$ denotes random state;
    \item $\mathcal{S}$ denotes software environment.
\end{itemize}

The goal is that

\begin{equation}
\boxed{
\mathcal{E}
\rightarrow
\text{Backtest}
\rightarrow
\text{Identical Result}
}
\end{equation}

whenever the experiment is rerun under the same conditions.


% ------------------------------------------------------------
\subsection{Research vs.\ Production Environments}
\label{subsec:research_vs_production}
% ------------------------------------------------------------

Research and production environments serve different purposes.

The research environment prioritizes:

\begin{itemize}
    \item flexibility;
    \item experimentation;
    \item rapid iteration;
    \item diagnostics;
    \item model comparison.
\end{itemize}

The production environment prioritizes:

\begin{itemize}
    \item robustness;
    \item determinism;
    \item observability;
    \item failure handling;
    \item reproducibility;
    \item controlled deployment.
\end{itemize}

Despite these differences, the underlying economic logic should be shared.

Ideally,

\begin{equation}
\boxed{
\text{Research Logic}
=
\text{Production Logic},
}
\end{equation}

while only the orchestration, scheduling, persistence, monitoring, and execution
layers differ.

A common failure mode is to implement one signal in research and independently
rewrite it for production.

This creates two separate code paths and increases the probability of divergence.

A preferable architecture reuses common modules for:

\begin{itemize}
    \item feature construction;
    \item signal generation;
    \item neutralization;
    \item portfolio construction;
    \item risk estimation;
    \item constraints;
    \item transaction-cost logic.
\end{itemize}


% ------------------------------------------------------------
\subsection{Configuration Management}
\label{subsec:configuration_management}
% ------------------------------------------------------------

Research parameters should be externalized from the core strategy code.

Let

\begin{equation}
\boldsymbol{\theta}
\end{equation}

denote the complete strategy configuration.

This may include:

\begin{itemize}
    \item universe filters;
    \item lookback windows;
    \item outlier thresholds;
    \item signal parameters;
    \item neutralization variables;
    \item rebalance frequency;
    \item risk-model parameters;
    \item optimization coefficients;
    \item transaction-cost assumptions;
    \item volatility target;
    \item leverage limits.
\end{itemize}

The strategy should therefore behave as

\begin{equation}
\boxed{
\text{Result}
=
f
\left(
\text{Data},
\text{Code},
\boldsymbol{\theta}
\right).
}
\end{equation}

Configurations should be stored in a machine-readable format such as YAML, TOML,
or JSON.

For example, conceptually,

\begin{equation}
\boxed{
\text{config\_id}
\rightarrow
\boldsymbol{\theta}
}
\end{equation}

should uniquely identify one set of strategy assumptions.


\subsubsection{Immutable Experiment Configurations}

Once an experiment has been run, its configuration should preferably be treated
as immutable.

A modified configuration should create a new experiment rather than silently
overwriting the previous one.

This preserves the research history and allows performance differences to be
attributed to explicit changes in assumptions.


% ------------------------------------------------------------
\subsection{Code Versioning}
\label{subsec:code_versioning}
% ------------------------------------------------------------

Every backtest should be associated with an exact version of the source code.

A version-control system such as Git provides a natural identifier:

\begin{equation}
\boxed{
\text{code\_version}
=
\text{commit hash}.
}
\end{equation}

For experiment $e$, metadata should therefore include something such as

\begin{equation}
\boxed{
\mathrm{commit}_e.
}
\end{equation}

This guarantees that the precise implementation used to generate a historical
result can be recovered.


\subsubsection{Dirty Working Trees}

A common reproducibility problem occurs when an experiment is run using local
uncommitted modifications.

The experiment may then depend on code that cannot be reconstructed from the
recorded Git commit.

A robust experiment runner should therefore either:

\begin{itemize}
    \item refuse to run with uncommitted changes;
    \item automatically record the source-code diff;
    \item explicitly mark the experiment as non-reproducible.
\end{itemize}


% ------------------------------------------------------------
\subsection{Data Versioning}
\label{subsec:data_versioning}
% ------------------------------------------------------------

Code reproducibility is insufficient if the underlying data change.

The backtester must therefore identify the exact data snapshot used by each
experiment.

Let

\begin{equation}
\mathcal{D}^{(v)}
\end{equation}

denote data version $v$.

Then an experiment should explicitly depend on

\begin{equation}
\boxed{
\mathcal{D}^{(v_e)}.
}
\end{equation}

Data versioning may be implemented through:

\begin{itemize}
    \item immutable data snapshots;
    \item dated partitions;
    \item object-storage versions;
    \item checksums or hashes;
    \item dataset manifests.
\end{itemize}


\subsubsection{Dataset Manifest}

A useful data manifest may record:

\begin{itemize}
    \item dataset name;
    \item source;
    \item extraction timestamp;
    \item minimum and maximum dates;
    \item number of observations;
    \item schema version;
    \item file or partition hashes.
\end{itemize}

Conceptually,

\begin{equation}
\boxed{
\mathrm{DataManifest}_e
=
\left(
\mathrm{source},
\mathrm{version},
\mathrm{schema},
\mathrm{hashes}
\right).
}
\end{equation}

This allows the exact research dataset to be reconstructed.


\subsubsection{Restated Data}

This is particularly important for financial datasets that may be revised.

A backtest rerun several months later against the current database may not use
the same historical values as the original experiment.

Therefore,

\begin{equation}
\boxed{
\text{Same Code}
+
\text{Same Configuration}
+
\text{Different Data Version}
\neq
\text{Same Experiment}.
}
\end{equation}


% ------------------------------------------------------------
\subsection{Model Versioning}
\label{subsec:model_versioning}
% ------------------------------------------------------------

Model-based strategies require an additional versioning layer.

A model artifact should be associated with:

\begin{itemize}
    \item model class;
    \item model parameters;
    \item hyperparameters;
    \item feature specification;
    \item training sample;
    \item training code version;
    \item random seed;
    \item fitted parameters.
\end{itemize}

A model version can therefore be represented as

\begin{equation}
\boxed{
\mathcal{M}^{(v)}
=
\left(
\text{Architecture},
\text{Features},
\text{Training Data},
\text{Parameters},
\text{Code Version}
\right).
}
\end{equation}

For every historical prediction date $t$, the backtester should be able to
identify the exact model version that produced

\begin{equation}
\widehat{\mathbf{y}}_{t+h|t}.
\end{equation}


\subsubsection{Model Artifacts}

Depending on the model, persisted artifacts may include:

\begin{itemize}
    \item regression coefficients;
    \item tree ensembles;
    \item neural-network weights;
    \item preprocessing parameters;
    \item feature encoders;
    \item calibration parameters.
\end{itemize}

Preprocessing artifacts are part of the model and must be versioned jointly with
the estimator.


% ------------------------------------------------------------
\subsection{Experiment Tracking}
\label{subsec:experiment_tracking}
% ------------------------------------------------------------

Every research run should create an experiment record.

Let

\begin{equation}
e
\end{equation}

denote an experiment identifier.

A useful experiment record is

\begin{equation}
\boxed{
\mathcal{E}_e
=
\left(
\mathrm{Config}_e,
\mathrm{Code}_e,
\mathrm{Data}_e,
\mathrm{Model}_e,
\mathrm{Metrics}_e,
\mathrm{Artifacts}_e
\right).
}
\end{equation}

The experiment metadata should include at least:

\begin{itemize}
    \item experiment ID;
    \item creation timestamp;
    \item Git commit;
    \item configuration;
    \item dataset version;
    \item model version;
    \item random seed;
    \item evaluation period;
    \item key performance metrics;
    \item output artifact locations.
\end{itemize}


\subsubsection{Experiment Outputs}

Useful persisted outputs include:

\begin{itemize}
    \item return series;
    \item NAV series;
    \item signal matrices;
    \item target and implemented weights;
    \item trade history;
    \item transaction costs;
    \item exposure histories;
    \item performance metrics;
    \item diagnostic plots.
\end{itemize}


\subsubsection{Connection with Statistical Validation}

Experiment tracking is also required for multiple-testing corrections.

If the researcher tests

\begin{equation}
N
\end{equation}

strategy variants but retains only the best one, the true research search cannot
later be reconstructed.

Therefore,

\begin{equation}
\boxed{
\text{Experiment Tracking}
\rightarrow
\text{Auditability}
+
\text{Statistical Validity}.
}
\end{equation}

This is particularly important for the Deflated Sharpe Ratio and Probability of
Backtest Overfitting.


% ------------------------------------------------------------
\subsection{Random Seeds}
\label{subsec:random_seeds}
% ------------------------------------------------------------

Stochastic procedures introduce additional sources of variability.

Examples include:

\begin{itemize}
    \item train-validation splitting;
    \item bootstrap sampling;
    \item random forests;
    \item stochastic optimization;
    \item neural-network initialization;
    \item randomized hyperparameter search;
    \item placebo and randomization tests.
\end{itemize}

Let

\begin{equation}
s_e
\end{equation}

denote the random seed associated with experiment $e$.

The seed should be treated as part of the experiment state and stored explicitly:

\begin{equation}
\boxed{
\mathrm{RandomState}_e
=
s_e.
}
\end{equation}

When multiple stochastic libraries are used, each relevant random-number
generator should be initialized explicitly.

For a stochastic learning algorithm, the fitted model may be represented as

\begin{equation}
\widehat{\boldsymbol{\theta}}^{(s)}
=
\mathcal{A}
\left(
\mathcal{D}_{\mathrm{train}},
s
\right),
\end{equation}

where $\mathcal{A}(\cdot)$ denotes the training algorithm and $s$ the random
seed.

Consequently, changing the seed may propagate through the entire investment
pipeline:

\begin{equation}
\boxed{
s
\rightarrow
\widehat{\boldsymbol{\theta}}^{(s)}
\rightarrow
\widehat{\mathbf{y}}^{(s)}
\rightarrow
\mathbf{s}^{(s)}
\rightarrow
\mathbf{w}^{(s)}
\rightarrow
R^{(s)}.
}
\end{equation}

Random seeds therefore have two distinct roles in a quantitative research
framework:

\begin{equation}
\boxed{
\text{Reproducibility}
\qquad\text{and}\qquad
\text{Robustness Analysis}.
}
\end{equation}


% ------------------------------------------------------------
\subsubsection{Fixed Seeds for Reproducibility}
% ------------------------------------------------------------

For reproducibility, the seed used by a given experiment must be recorded and
reused.

The complete model specification therefore includes the random state:

\begin{equation}
\boxed{
\mathcal{M}
=
\left(
\text{Architecture},
\text{Hyperparameters},
\text{Training Data},
\text{Seed},
\ldots
\right).
}
\end{equation}

Given identical data, code, configuration, software environment, and random
state, rerunning the experiment should produce the same result whenever the
underlying computational operations are deterministic.

Once a final strategy is deployed, its seed or seed-generation rule should
therefore be part of the frozen production configuration.


% ------------------------------------------------------------
\subsubsection{Multiple Seeds for Statistical Robustness}
% ------------------------------------------------------------

Fixing the seed is necessary for reproducibility but does not establish that the
strategy is robust to the stochasticity of the training procedure.

During research, a stochastic model should therefore also be evaluated across a
predefined collection of seeds

\begin{equation}
\mathcal{S}
=
\left\{
s_1,\ldots,s_K
\right\}.
\end{equation}

For each seed $s_k$, the complete model-based backtest should be rerun:

\begin{equation}
s_k
\rightarrow
\widehat{\boldsymbol{\theta}}^{(s_k)}
\rightarrow
\widehat{\mathbf{y}}^{(s_k)}
\rightarrow
\mathbf{s}^{(s_k)}
\rightarrow
\mathbf{w}^{(s_k)}
\rightarrow
R^{(s_k)}.
\end{equation}

This produces a distribution of strategy outcomes rather than a single
realization:

\begin{equation}
\boxed{
\left\{
\mathcal{M}^{(s_1)},
\ldots,
\mathcal{M}^{(s_K)}
\right\},
}
\end{equation}

where $\mathcal{M}^{(s_k)}$ denotes a performance metric obtained under seed
$s_k$.

For example, the mean Sharpe ratio across seeds is

\begin{equation}
\boxed{
\overline{SR}_{\mathrm{seed}}
=
\frac{1}{K}
\sum_{k=1}^{K}
SR^{(s_k)},
}
\end{equation}

and its cross-seed dispersion is

\begin{equation}
\boxed{
\sigma_{SR,\mathrm{seed}}
=
\sqrt{
\frac{1}{K-1}
\sum_{k=1}^{K}
\left(
SR^{(s_k)}
-
\overline{SR}_{\mathrm{seed}}
\right)^2
}.
}
\end{equation}

The same analysis can be performed for:

\begin{itemize}
    \item out-of-sample IC and Rank IC;
    \item annualized return;
    \item net Sharpe ratio;
    \item maximum drawdown;
    \item turnover;
    \item factor exposures;
    \item transaction costs.
\end{itemize}

The objective is not to require identical results across seeds.

Instead, the strategy should remain economically meaningful across reasonable
realizations of the stochastic training process.


% ------------------------------------------------------------
\subsubsection{Signal and Portfolio Stability Across Seeds}
% ------------------------------------------------------------

Performance metrics alone may conceal substantial instability in the underlying
predictions.

It is therefore useful to compare the actual signals and portfolios produced by
different seeds.

For two seeds $s_a$ and $s_b$, one may compute, for example,

\begin{equation}
\boxed{
\rho^{\mathrm{signal}}_{a,b}
=
\operatorname{Corr}
\left(
\mathbf{s}^{(s_a)},
\mathbf{s}^{(s_b)}
\right),
}
\end{equation}

and

\begin{equation}
\boxed{
\rho^{\mathrm{weights}}_{a,b}
=
\operatorname{Corr}
\left(
\mathbf{w}^{(s_a)},
\mathbf{w}^{(s_b)}
\right).
}
\end{equation}

Depending on the strategy, additional diagnostics may include:

\begin{itemize}
    \item rank correlation between predictions;
    \item overlap between selected securities;
    \item average absolute difference in portfolio weights;
    \item correlation between realized return series.
\end{itemize}

Ideally,

\begin{equation}
\boxed{
\text{Different Seeds}
\rightarrow
\text{Similar Economic Conclusions}.
}
\end{equation}

A strategy whose predictions, positions, or performance change dramatically
under arbitrary changes in the random seed may be insufficiently stable for
deployment.


% ------------------------------------------------------------
\subsubsection{The Seed Should Not Become a Hyperparameter}
% ------------------------------------------------------------

A critical distinction must be made between testing seed sensitivity and
optimizing the seed.

Suppose $K$ seeds are evaluated.

Selecting

\begin{equation}
\boxed{
s^*
=
\arg\max_{s\in\mathcal{S}}
SR^{(s)}
}
\end{equation}

and subsequently reporting only the performance associated with $s^*$ turns the
seed into an implicitly optimized hyperparameter.

The resulting Sharpe ratio is then subject to the same selection bias and
multiple-testing concerns discussed previously.

Seed variation should therefore be used as a robustness diagnostic rather than
as a mechanism for searching for the best historical realization.

A production seed may instead be determined by a rule independent of observed
backtest performance, for example:

\begin{itemize}
    \item a seed fixed before the final evaluation;
    \item the first seed from a predefined sequence;
    \item a deterministic seed-generation rule.
\end{itemize}

Thus,

\begin{equation}
\boxed{
\begin{aligned}
\text{Reproducibility} &:\quad \text{Fix Seeds},\\
\text{Robustness} &:\quad \text{Vary Seeds},\\
\text{Model Selection} &:\quad \text{Do Not Optimize Seeds on Backtest Performance}.
\end{aligned}
}
\end{equation}


% ------------------------------------------------------------
\subsubsection{Seed Ensembles}
% ------------------------------------------------------------

If random initialization materially affects model predictions, the variation
across seeds may itself be interpreted as a source of model uncertainty.

Rather than deploying one arbitrary stochastic realization, the strategy may
explicitly define an ensemble over a predetermined collection of seeds.

For example,

\begin{equation}
\boxed{
\widehat{y}_{i,t}^{\mathrm{ens}}
=
\frac{1}{K}
\sum_{k=1}^{K}
\widehat{y}_{i,t}^{(s_k)}.
}
\end{equation}

More generally, predictions may be combined using predetermined ensemble
weights

\begin{equation}
\widehat{y}_{i,t}^{\mathrm{ens}}
=
\sum_{k=1}^{K}
\omega_k
\widehat{y}_{i,t}^{(s_k)},
\end{equation}

with

\begin{equation}
\sum_{k=1}^{K}\omega_k=1.
\end{equation}

In this case, the ensemble itself becomes part of the strategy specification:

\begin{equation}
\boxed{
\mathcal{M}^{\mathrm{ens}}
=
\left(
\text{Model},
\text{Hyperparameters},
\{s_1,\ldots,s_K\},
\{\omega_1,\ldots,\omega_K\}
\right).
}
\end{equation}

The complete seed set and aggregation rule must then be fixed, versioned, and
reproduced in production.

An ensemble may reduce sensitivity to one particular random realization, at the
cost of additional training, storage, and inference requirements.


% ------------------------------------------------------------
\subsubsection{Interaction with Train, Validation, and Test Samples}
% ------------------------------------------------------------

Seed robustness should respect the same separation between training,
validation, and final test samples used elsewhere in the research process.

Seed sensitivity may be investigated during model development and validation.

The final strategy specification should then be frozen, including:

\begin{equation}
\boxed{
\left(
\text{Model},
\text{Hyperparameters},
\text{Seed or Ensemble Rule},
\text{Portfolio Construction}
\right),
}
\end{equation}

before evaluation on the final untouched test sample.

The final test sample should not be repeatedly inspected to identify the seed
that happens to perform best.

Otherwise,

\begin{equation}
\boxed{
\text{Seed Selection on Test Performance}
\rightarrow
\text{Test-Set Contamination}.
}
\end{equation}

For a walk-forward backtest, the same principle applies at every estimation
date: the rule determining the seed or ensemble must be defined using only
information available within the corresponding training and validation process.


% ------------------------------------------------------------
\subsubsection{Seeds Are Necessary but Not Always Sufficient}
% ------------------------------------------------------------

Setting a seed does not guarantee perfect reproducibility in every computational
environment.

Sources of nondeterminism may include:

\begin{itemize}
    \item parallel execution;
    \item hardware-specific numerical behavior;
    \item GPU operations;
    \item nondeterministic numerical kernels;
    \item changes in library implementations.
\end{itemize}

Therefore, random seeds should be combined with:

\begin{itemize}
    \item code versioning;
    \item dependency versioning;
    \item environment versioning;
    \item hardware information when relevant;
    \item deterministic computational settings when available.
\end{itemize}

The complete principle is therefore

\begin{equation}
\boxed{
\begin{aligned}
\text{Reproducibility}
&:
&
\text{Record and Reuse the Random State},
\\
\text{Robustness}
&:
&
\text{Evaluate Sensitivity Across Random States},
\\
\text{Model Selection}
&:
&
\text{Never Select a Seed from Final-Test Performance},
\\
\text{Production}
&:
&
\text{Freeze the Seed or Ensemble Rule}.
\end{aligned}
}
\end{equation}


% ------------------------------------------------------------
\subsection{Logging}
\label{subsec:logging}
% ------------------------------------------------------------

A production-grade backtester should produce structured logs rather than relying
only on interactive output.

At minimum, logs should record:

\begin{itemize}
    \item start and end of each pipeline stage;
    \item experiment ID;
    \item timestamps;
    \item data coverage;
    \item number of securities;
    \item missing-data statistics;
    \item model-fitting events;
    \item optimizer status;
    \item infeasible constraints;
    \item execution failures;
    \item transaction-cost calculations;
    \item warnings and exceptions.
\end{itemize}

A log record can be represented abstractly as

\begin{equation}
\boxed{
\mathcal{L}_k
=
\left(
t_k,
e_k,
\mathrm{Level}_k,
\mathrm{Component}_k,
\mathrm{Message}_k
\right).
}
\end{equation}


\subsubsection{Structured Logging}

Machine-readable structured logs facilitate automated monitoring and debugging.

Useful fields include:

\begin{equation}
\boxed{
\{
\text{timestamp},
\text{experiment\_id},
\text{module},
\text{event},
\text{severity},
\text{metadata}
\}.
}
\end{equation}

This makes it possible to search systematically for failures or anomalous
portfolio states.


% ------------------------------------------------------------
\subsection{Unit Tests}
\label{subsec:unit_tests}
% ------------------------------------------------------------

Unit tests verify individual components of the backtesting engine in isolation.

Examples include testing that:

\begin{itemize}
    \item winsorization produces the expected bounds;
    \item a z-score has approximately zero mean and unit variance;
    \item neutralized signals satisfy the required orthogonality condition;
    \item gross normalization produces the requested gross exposure;
    \item portfolio beta is computed correctly;
    \item turnover is calculated from drifted rather than stale weights;
    \item transaction-cost formulas match analytical examples;
    \item portfolio NAV reconciles with positions and cash.
\end{itemize}


\subsubsection{Mathematical Invariants}

Many portfolio components admit exact invariants that are particularly useful
for testing.

For example, after OLS neutralization,

\begin{equation}
\boxed{
\mathbf{B}_t^\top
\mathbf{s}_t^{\mathrm{neutral}}
\approx
\mathbf{0}.
}
\end{equation}

For a dollar-neutral portfolio,

\begin{equation}
\boxed{
\mathbf{1}^\top
\mathbf{w}_t
\approx
0.
}
\end{equation}

For a portfolio normalized to gross exposure $G$,

\begin{equation}
\boxed{
\|\mathbf{w}_t\|_1
\approx
G.
}
\end{equation}

Such identities provide powerful automatic correctness checks.


% ------------------------------------------------------------
\subsection{Integration Tests}
\label{subsec:integration_tests}
% ------------------------------------------------------------

Integration tests verify that multiple components interact correctly.

For example, a complete small synthetic backtest may test the chain

\begin{equation}
\boxed{
\text{Data}
\rightarrow
\text{Signal}
\rightarrow
\text{Weights}
\rightarrow
\text{Trades}
\rightarrow
\text{P\&L}.
}
\end{equation}

Useful integration scenarios include:

\begin{itemize}
    \item one-security portfolios;
    \item two-security dollar-neutral portfolios;
    \item known constant-return datasets;
    \item zero-transaction-cost environments;
    \item fixed-price environments;
    \item deterministic corporate-action examples.
\end{itemize}

The expected result should be analytically known whenever possible.


\subsubsection{Golden Tests}

A useful production technique is to maintain reference outputs for a small
deterministic dataset.

If

\begin{equation}
\mathcal{Y}^{\mathrm{reference}}
\end{equation}

denotes the expected result, the current implementation should satisfy

\begin{equation}
\boxed{
\mathcal{Y}^{\mathrm{current}}
\approx
\mathcal{Y}^{\mathrm{reference}}.
}
\end{equation}

Unexpected deviations can indicate that a code change has altered the economic
behavior of the strategy.


% ------------------------------------------------------------
\subsection{Reproducible Environments}
\label{subsec:reproducible_environments}
% ------------------------------------------------------------

The software environment used to run the backtest should itself be versioned.

Relevant dependencies include:

\begin{itemize}
    \item Python version;
    \item numerical libraries;
    \item optimization libraries;
    \item machine-learning libraries;
    \item operating-system dependencies;
    \item compiled numerical libraries.
\end{itemize}

A reproducible environment can be represented by

\begin{equation}
\boxed{
\mathcal{S}
=
\left(
\text{Runtime},
\text{Dependencies},
\text{System Libraries}
\right).
}
\end{equation}

Dependency versions should be pinned rather than allowed to change implicitly.


\subsubsection{Lock Files}

A package lock file records exact dependency versions.

Conceptually,

\begin{equation}
\boxed{
\text{Project Specification}
+
\text{Lock File}
\rightarrow
\text{Deterministic Dependency Set}.
}
\end{equation}

Examples include lock files generated by modern Python environment-management
tools.


\subsubsection{Containers}

For stronger reproducibility, the complete runtime may be containerized.

Then

\begin{equation}
\boxed{
\text{Container Image}
=
\text{Code Runtime}
+
\text{Dependencies}
+
\text{System Environment}.
}
\end{equation}

The container image itself should be versioned, for example using an immutable
image digest.


% ------------------------------------------------------------
\subsection{Backtest--Live Trading Parity}
\label{subsec:backtest_live_parity}
% ------------------------------------------------------------

The ultimate objective of the production architecture is to minimize differences
between the historical simulation and the live trading system.

Ideally, both systems should share the same core functions:

\begin{equation}
\boxed{
\begin{aligned}
\text{Feature Engine}_{BT}
&=
\text{Feature Engine}_{Live},
\\
\text{Signal Engine}_{BT}
&=
\text{Signal Engine}_{Live},
\\
\text{Risk Engine}_{BT}
&=
\text{Risk Engine}_{Live},
\\
\text{Portfolio Engine}_{BT}
&=
\text{Portfolio Engine}_{Live}.
\end{aligned}
}
\end{equation}

The main difference should be the source of inputs:

\begin{equation}
\boxed{
\begin{aligned}
\text{Backtest}
&:
&&
\text{Historical Point-in-Time Inputs},
\\
\text{Live}
&:
&&
\text{Current Real-Time Inputs}.
\end{aligned}
}
\end{equation}


\subsubsection{Avoiding Duplicate Logic}

A dangerous architecture implements one version of the strategy for research
and another for live trading.

This may produce differences in:

\begin{itemize}
    \item missing-value handling;
    \item feature calculations;
    \item signal normalization;
    \item neutralization;
    \item risk estimation;
    \item constraints;
    \item order sizing.
\end{itemize}

A preferable architecture uses the same deterministic transformation functions
in both environments.


\subsubsection{Shadow and Paper Trading}

Before deploying capital, the live system may run in shadow or paper-trading
mode.

At each live decision date, record:

\begin{equation}
\boxed{
\mathbf{s}_t^{\mathrm{live}},
\qquad
\mathbf{w}_t^{\mathrm{live}},
\qquad
\mathcal{O}_t^{\mathrm{live}}.
}
\end{equation}

The same timestamp can then be replayed using the backtest engine.

One should verify approximately

\begin{equation}
\boxed{
\mathbf{s}_t^{BT}
\approx
\mathbf{s}_t^{Live},
}
\end{equation}

and

\begin{equation}
\boxed{
\mathbf{w}_t^{BT}
\approx
\mathbf{w}_t^{Live}.
}
\end{equation}

Differences should be explainable by explicitly modeled live effects.


\subsubsection{Backtest--Live Reconciliation}

A useful reconciliation can decompose live deviations as

\begin{equation}
\boxed{
\begin{aligned}
\text{Live Performance}
-
\text{Backtest Expectation}
=
&
\text{Data Difference}
\\
&+
\text{Signal Difference}
\\
&+
\text{Portfolio Difference}
\\
&+
\text{Execution Difference}
\\
&+
\text{Cost Difference}.
\end{aligned}
}
\end{equation}

This transforms backtest-live divergence from an unexplained performance gap
into a set of identifiable engineering and economic components.


% ------------------------------------------------------------
\subsubsection{Production Architecture Summary}
% ------------------------------------------------------------

A reproducible quantitative research system should make every strategy result
traceable to a complete experiment state:

\begin{equation}
\boxed{
\begin{aligned}
\text{Code Version}
&+
\text{Data Version}
\\
&+
\text{Configuration}
+
\text{Model Version}
\\
&+
\text{Random State}
+
\text{Software Environment}
\\
&\rightarrow
\text{Reproducible Experiment}.
\end{aligned}
}
\end{equation}

The complete transition from research to production can therefore be represented
as

\begin{equation}
\boxed{
\begin{aligned}
\text{Research Idea}
&\rightarrow
\text{Versioned Experiment}
\\
&\rightarrow
\text{Validated Backtest}
\\
&\rightarrow
\text{Shared Production Logic}
\\
&\rightarrow
\text{Paper / Shadow Trading}
\\
&\rightarrow
\text{Live Deployment}
\\
&\rightarrow
\text{Backtest--Live Reconciliation}.
\end{aligned}
}
\end{equation}

Reproducibility is therefore not merely a software-engineering property.

It is part of the statistical integrity of the research process, because without
a complete record of code, data, configurations, and experiments, historical
results cannot be independently verified and the true extent of strategy
selection cannot be measured.


% ============================================================
\part{Conclusion}
% ============================================================

% ============================================================
\section{From Research Idea to Investment Strategy}
\label{sec:conclusion_research_to_strategy}
% ============================================================

A quantitative investment strategy is not merely a predictive signal, a
portfolio optimization problem, or a historical return series.

It is a complete decision-making system transforming information available at a
given point in time into positions that can realistically be implemented and
whose subsequent economic performance can be measured.

The purpose of this guide has been to describe this complete transformation.

At a high level, the quantitative investment process can be represented as

\begin{equation}
\boxed{
\begin{aligned}
\text{Raw Data}
\;(\mathcal{D}^{\mathrm{raw}})
&\rightarrow
\text{Point-in-Time Data}
\;(\mathcal{D}_t^{\mathrm{PIT}})
\\
&\rightarrow
\text{Investment Universe}
\;(\mathcal{U}_t)
\\
&\rightarrow
\text{Clean Characteristics}
\;(\mathbf{x}_t)
\\
&\rightarrow
\text{Raw Signal}
\;(\mathbf{s}_t^{\mathrm{raw}})
\\
&\rightarrow
\text{Neutralized Signal}
\;(\mathbf{s}_t^{\mathrm{neutral}})
\\
&\rightarrow
\text{Final Investable Signal}
\;(\mathbf{s}_t^{\mathrm{final}})
\\
&\rightarrow
\text{Portfolio Construction}
\;(\mathcal{G}_t,\mathcal{J}_t,\mathcal{W}_t)
\\
&\rightarrow
\text{Target Weights}
\;(\mathbf{w}_t^{\mathrm{target}})
\\
&\rightarrow
\text{(Risk Scaling)}
\;(\lambda_t)
\\
&\rightarrow
\text{Final Weights}
\;(\mathbf{w}_t^{\mathrm{final}})
\\
&\rightarrow
\text{Trades}
\;(\Delta\mathbf{w}_t,\mathcal{O}_t)
\\
&\rightarrow
\text{Realized Net Returns}
\;(R_{p,t}^{\mathrm{net}})
\\
&\rightarrow
\text{Performance Evaluation}
\;(\mathcal{M}_{\mathrm{perf}})
\\
&\rightarrow
\text{Statistical Validation}
\;(\mathcal{M}_{\mathrm{stat}}).
\end{aligned}
}
\end{equation}


Each transformation introduces assumptions.

Consequently, the credibility of the final backtest depends not only on the
quality of the original investment idea, but also on the correctness and
realism of every intermediate step.


% ------------------------------------------------------------
\subsection{Summary of the Complete Pipeline}
\label{subsec:summary_complete_pipeline}
% ------------------------------------------------------------

The complete research process begins with an economic hypothesis.

Suppose that information available at time $t$ is represented by

\begin{equation}
\mathcal{I}_t.
\end{equation}

The objective is ultimately to construct portfolio weights

\begin{equation}
\mathbf{w}_t
\end{equation}

using only this information:

\begin{equation}
\boxed{
\mathbf{w}_t
=
f
\left(
\mathcal{I}_t
\right).
}
\end{equation}

The complexity of a quantitative backtester lies in the construction and
validation of the function $f(\cdot)$.


\subsubsection{Data Layer}

The first stage constructs the information set that would actually have been
available to the investor.

Raw historical databases must be transformed into point-in-time datasets that
correctly account for:

\begin{itemize}
    \item publication and availability dates;
    \item reporting lags;
    \item restatements;
    \item historical universe membership;
    \item delisted securities;
    \item corporate actions;
    \item trading calendars;
    \item data revisions.
\end{itemize}

Thus,

\begin{equation}
\boxed{
\text{Raw Historical Database}
\neq
\text{Historical Investor Information Set}.
}
\end{equation}

Point-in-time construction is therefore a prerequisite for meaningful
historical simulation.


\subsubsection{Data Cleaning and Outlier Treatment}

The point-in-time dataset must subsequently be cleaned.

This includes identifying:

\begin{itemize}
    \item duplicate observations;
    \item missing observations;
    \item stale prices;
    \item invalid values;
    \item inconsistent frequencies;
    \item extreme observations.
\end{itemize}

Outlier treatment may involve winsorization, robust scaling, truncation, or
explicit observation removal.

Importantly, these transformations must themselves be point-in-time and
systematic.

The result is a set of economically interpretable firm-level characteristics

\begin{equation}
\mathbf{x}_{i,t}.
\end{equation}


\subsubsection{Raw Signal Construction}

Firm-level characteristics can be converted directly into cross-sectional
scores.

For example,

\begin{equation}
x_{i,t}
\rightarrow
z_{i,t}
\end{equation}

may represent a z-score, rank, percentile, Gaussian rank, or another monotonic
transformation.

Multiple characteristics may subsequently be combined to produce a raw score

\begin{equation}
\boxed{
s_{i,t}^{\mathrm{raw}}.
}
\end{equation}

Alternatively, the signal may be generated by a predictive model:

\begin{equation}
\mathbf{X}_t
\rightarrow
\widehat{f}
\rightarrow
\widehat{y}_{i,t+h|t}
\rightarrow
s_{i,t}^{\mathrm{raw}}.
\end{equation}

The predictive engine may be econometric, machine-learning based, or
deep-learning based.

In every case, estimation must respect the temporal information structure of
the problem through appropriate training, validation, embargo, and walk-forward
procedures.


\subsubsection{Signal Neutralization}

The raw score may contain unwanted systematic exposures.

Let

\begin{equation}
\mathbf{B}_t
\in
\mathbb{R}^{N_t\times K}
\end{equation}

denote the exposure matrix containing, for example:

\begin{itemize}
    \item market beta;
    \item industries;
    \item countries;
    \item size;
    \item volatility;
    \item style factors.
\end{itemize}

Neutralization transforms

\begin{equation}
\mathbf{s}_t^{\mathrm{raw}}
\end{equation}

into a residualized signal

\begin{equation}
\boxed{
\mathbf{s}_t^{\mathrm{neutral}}
}
\end{equation}

that is orthogonal, according to the chosen OLS or WLS metric, to the specified
exposure space.

After any required standardization, signal combination, smoothing,
turnover-aware transformation, or confidence adjustment, the signal engine
produces the final investable signal

\begin{equation}
\boxed{
\mathbf{s}_t^{\mathrm{final}}.
}
\end{equation}


\subsubsection{Signal Diagnostics}

Before capital is allocated, the signal itself should be evaluated.

Diagnostics such as

\begin{itemize}
    \item Information Coefficient;
    \item Rank Information Coefficient;
    \item quantile returns;
    \item long--short spreads;
    \item monotonicity;
    \item persistence;
    \item turnover;
    \item factor exposures;
\end{itemize}

help determine whether the signal behaves as economically intended.

This stage evaluates the predictive object before the additional effects of
portfolio construction are introduced.


\subsubsection{Portfolio Construction}

The portfolio engine converts the final signal into positions.

A mapping function may first produce raw portfolio weights:

\begin{equation}
\widetilde{\mathbf{w}}_t
=
g
\left(
\mathbf{s}_t^{\mathrm{final}}
\right).
\end{equation}

Examples include:

\begin{itemize}
    \item equal-weighted quantile portfolios;
    \item top/bottom $N$ portfolios;
    \item rank-weighted portfolios;
    \item linear mappings;
    \item nonlinear mappings;
    \item conviction-weighted mappings.
\end{itemize}

For direct long--short Pure Alpha construction, this mapping may itself be
sufficient to construct an economically meaningful portfolio, potentially
followed by exposure normalization and volatility scaling.

Alternatively, the signal or the resulting Pure Alpha can enter an optimization
engine.

The optimizer combines expected returns, the risk model, transaction costs, and
portfolio constraints:

\begin{equation}
\boxed{
\mathbf{w}_t^{\mathrm{target}}
=
\arg\max_{\mathbf{w}}
\;
\mathcal{J}
\left(
\mathbf{w};
\widehat{\boldsymbol{\mu}}_t,
\widehat{\boldsymbol{\Sigma}}_t,
\mathbf{s}_t^{\mathrm{final}},
\ldots
\right)
}
\end{equation}

subject to the required constraints.


\subsubsection{Portfolio-Level Risk Scaling}

Portfolio construction determines the relative allocation across securities.

Portfolio-level risk scaling determines the overall magnitude of those
positions.

If

\begin{equation}
\mathbf{w}_t^{\mathrm{target}}
\end{equation}

denotes the target portfolio and

\begin{equation}
\widehat{\sigma}_{p,t}
\end{equation}

its forecast volatility, a volatility-targeting multiplier may take the form

\begin{equation}
\lambda_t
=
\frac{
\sigma^{\mathrm{target}}
}{
\widehat{\sigma}_{p,t}
}.
\end{equation}

After applying any leverage or scaling limits,

\begin{equation}
\boxed{
\mathbf{w}_t^{\mathrm{final}}
=
\lambda_t
\mathbf{w}_t^{\mathrm{target}}.
}
\end{equation}

This emphasizes the distinction between

\begin{equation}
\boxed{
\text{Portfolio Direction}
\qquad\text{and}\qquad
\text{Portfolio Scale}.
}
\end{equation}


\subsubsection{Backtest Simulation}

Final target weights are not returns.

They must be translated into trades and subsequently into realized portfolio
P\&L.

At each rebalance date, the implementation engine compares current drifted
weights with target weights:

\begin{equation}
\Delta\mathbf{w}_t
=
\mathbf{w}_t^{\mathrm{final}}
-
\mathbf{w}_{t^-}.
\end{equation}

The resulting trades incur implementation costs.

The complete return construction therefore follows

\begin{equation}
\boxed{
\begin{aligned}
\text{Final Target Weights}
&\rightarrow
\text{Required Trades}
\\
&\rightarrow
\text{Execution}
\rightarrow
\text{Transaction Costs}
\\
&\rightarrow
\text{Realized Holdings}
\rightarrow
\text{Gross P\&L}
\\
&\rightarrow
\text{Financing / Borrowing Costs}
\rightarrow
\text{Net P\&L}.
\end{aligned}
}
\end{equation}


\subsubsection{Performance and Statistical Validation}

The resulting return series must finally be evaluated economically and
statistically.

Economic evaluation includes measures of:

\begin{itemize}
    \item return;
    \item volatility;
    \item Sharpe ratio;
    \item drawdown;
    \item downside risk;
    \item turnover;
    \item capacity;
    \item implementation costs.
\end{itemize}

Performance attribution then identifies the economic sources of realized
returns across securities, long and short books, sectors, countries, factors,
signals, models, and implementation costs.

Finally, statistical validation asks whether the observed historical performance
is sufficiently robust to survive:

\begin{itemize}
    \item sampling uncertainty;
    \item serial dependence;
    \item alternative specifications;
    \item multiple hypothesis testing;
    \item data snooping;
    \item model-selection effects;
    \item random-seed variation;
    \item backtest overfitting.
\end{itemize}

The final objective is therefore not to produce the highest historical Sharpe
ratio.

It is to produce the strongest evidence that a coherent investment process can
generate economically meaningful performance outside the historical sample
used to develop it.


% ------------------------------------------------------------
\subsection{Interactions Between Signal, Risk, and Portfolio Construction}
\label{subsec:interactions_signal_risk_portfolio}
% ------------------------------------------------------------

Although the pipeline is naturally presented as a sequence of stages, these
stages are not economically independent.

A useful decomposition is

\begin{equation}
\boxed{
\text{Signal}
\rightarrow
\text{Desired Direction},
}
\end{equation}

\begin{equation}
\boxed{
\text{Risk Model}
\rightarrow
\text{Risk of the Desired Direction},
}
\end{equation}

and

\begin{equation}
\boxed{
\text{Portfolio Construction}
\rightarrow
\text{Capital Allocated to the Desired Direction}.
}
\end{equation}


\subsubsection{Signal Quality Does Not Equal Portfolio Quality}

A signal with strong predictive power does not necessarily generate an
attractive portfolio.

For example, the highest-scoring securities may simultaneously be:

\begin{itemize}
    \item highly correlated;
    \item concentrated in one industry;
    \item exposed to the same systematic factor;
    \item illiquid;
    \item expensive to short;
    \item associated with excessive turnover.
\end{itemize}

Consequently,

\begin{equation}
\boxed{
\text{Strong Signal}
\nRightarrow
\text{Strong Portfolio}.
}
\end{equation}


\subsubsection{Risk Management Can Modify the Signal Expression}

Conversely, risk controls alter the way the signal is expressed.

Two securities may have identical signal values but receive different portfolio
weights because their:

\begin{itemize}
    \item volatilities differ;
    \item covariance structures differ;
    \item liquidity differs;
    \item factor exposures differ;
    \item transaction costs differ.
\end{itemize}

Thus,

\begin{equation}
s_{i,t}^{\mathrm{final}}
=
s_{j,t}^{\mathrm{final}}
\end{equation}

does not imply

\begin{equation}
w_{i,t}^{\mathrm{final}}
=
w_{j,t}^{\mathrm{final}}.
\end{equation}


\subsubsection{Signal Neutralization vs.\ Portfolio Neutrality}

A particularly important distinction is between neutralizing the signal and
constraining the portfolio.

Signal neutralization imposes a property on

\begin{equation}
\mathbf{s}_t,
\end{equation}

whereas portfolio neutrality imposes a property on

\begin{equation}
\mathbf{w}_t.
\end{equation}

For example,

\begin{equation}
\mathbf{B}_t^\top
\mathbf{s}_t^{\mathrm{neutral}}
=
\mathbf{0}
\end{equation}

does not necessarily imply

\begin{equation}
\mathbf{B}_t^\top
\mathbf{w}_t
=
\mathbf{0},
\end{equation}

because nonlinear mappings, optimization, constraints, clipping, and other
portfolio transformations may reintroduce exposures.

Therefore, exposure diagnostics should be performed at both the signal and
portfolio levels.


\subsubsection{Costs Feed Back into Portfolio Construction}

Transaction costs should not always be viewed solely as an ex-post deduction
from gross performance.

When turnover and liquidity materially affect economic performance, expected
implementation costs can enter the portfolio decision itself.

Thus,

\begin{equation}
\boxed{
\text{Signal}
+
\text{Risk}
+
\text{Costs}
+
\text{Constraints}
\rightarrow
\text{Investable Portfolio}.
}
\end{equation}

The economically optimal portfolio may therefore intentionally sacrifice some
predicted gross alpha in exchange for lower risk, lower turnover, greater
liquidity, or greater robustness.


\subsubsection{The Pipeline Is Sequential but the Research Process Is Iterative}

The backtest calculation follows a strict temporal sequence.

Research development, however, is iterative.

For example,

\begin{equation}
\boxed{
\begin{aligned}
\text{Signal Diagnostics}
&\rightarrow
\text{Signal Revision},
\\
\text{Portfolio Diagnostics}
&\rightarrow
\text{Constraint Revision},
\\
\text{Cost Analysis}
&\rightarrow
\text{Mapping Revision},
\\
\text{Robustness Analysis}
&\rightarrow
\text{Model Revision}.
\end{aligned}
}
\end{equation}

Every such revision constitutes an additional research decision.

Consequently, the complete experiment history should be recorded so that this
iterative process does not become invisible data snooping.


% ------------------------------------------------------------
\subsection{Research and Backtest Validation Checklist}
\label{subsec:research_backtest_validation_checklist}
% ------------------------------------------------------------

Before a strategy is considered ready for deployment, the complete research
pipeline should be reviewed systematically.

The following checklist summarizes the principal questions developed throughout
this guide.


\subsubsection{Data Integrity}

\begin{itemize}
    \item[$\square$] Are all inputs point-in-time?
    \item[$\square$] Are publication and reporting lags correctly modeled?
    \item[$\square$] Are restatements handled consistently?
    \item[$\square$] Is historical universe membership reconstructed?
    \item[$\square$] Are delisted securities and delisting returns included?
    \item[$\square$] Are corporate actions correctly reflected?
    \item[$\square$] Are duplicate, missing, stale, and invalid observations
    handled systematically?
    \item[$\square$] Are currencies, frequencies, and trading calendars aligned
    correctly?
\end{itemize}


\subsubsection{Signal Integrity}

\begin{itemize}
    \item[$\square$] Is every feature available at the signal formation date?
    \item[$\square$] Are outlier treatments defined ex ante?
    \item[$\square$] Is the score transformation economically interpretable?
    \item[$\square$] Is the signal direction convention explicit?
    \item[$\square$] Are neutralization exposures correctly specified?
    \item[$\square$] Is neutrality verified numerically after neutralization?
    \item[$\square$] Are IC, Rank IC, quantile returns, monotonicity, persistence,
    and turnover examined?
\end{itemize}


\subsubsection{Model Integrity}

For model-based strategies:

\begin{itemize}
    \item[$\square$] Is the prediction target defined precisely?
    \item[$\square$] Are training, validation, and test samples temporally
    separated?
    \item[$\square$] Is walk-forward estimation used when appropriate?
    \item[$\square$] Is an embargo applied when overlapping labels create
    leakage risk?
    \item[$\square$] Are preprocessing parameters estimated only from permitted
    observations?
    \item[$\square$] Are hyperparameters selected without contaminating the
    final test sample?
    \item[$\square$] Are stochastic models tested across multiple seeds?
    \item[$\square$] Is the production seed or ensemble rule fixed independently
    of final-test performance?
\end{itemize}


\subsubsection{Portfolio Construction Integrity}

\begin{itemize}
    \item[$\square$] Is the signal-to-weight mapping explicitly defined?
    \item[$\square$] Are gross and net exposure conventions explicit?
    \item[$\square$] Are long and short budgets correctly enforced?
    \item[$\square$] Are position limits applied?
    \item[$\square$] Are beta, industry, country, and style exposures controlled
    where required?
    \item[$\square$] Is the covariance or factor risk model estimated using only
    available information?
    \item[$\square$] Are turnover and liquidity constraints realistic?
    \item[$\square$] Is volatility targeting based on an ex-ante risk estimate?
    \item[$\square$] Are leverage caps applied after portfolio-level scaling?
    \item[$\square$] Are final portfolio exposures checked after all
    transformations?
\end{itemize}


\subsubsection{Backtest Implementation Integrity}

\begin{itemize}
    \item[$\square$] Is the information timestamp explicit?
    \item[$\square$] Is the signal-to-execution lag realistic?
    \item[$\square$] Are trades generated from target minus current drifted
    positions?
    \item[$\square$] Are execution-price assumptions realistic?
    \item[$\square$] Are transaction costs charged on actual traded notional?
    \item[$\square$] Are both purchases and sales treated consistently?
    \item[$\square$] Are bid--ask spreads, commissions, slippage, and market
    impact included where relevant?
    \item[$\square$] Are short-borrowing and financing costs modeled?
    \item[$\square$] Are liquidity and capacity assumptions plausible?
    \item[$\square$] Does portfolio accounting reconcile positions, cash, trades,
    costs, and P\&L?
\end{itemize}


\subsubsection{Performance Integrity}

\begin{itemize}
    \item[$\square$] Are gross and net performance reported separately?
    \item[$\square$] Are return, volatility, Sharpe, drawdown, and downside-risk
    statistics reported?
    \item[$\square$] Is turnover reported under an explicit convention?
    \item[$\square$] Is performance evaluated relative to appropriate benchmarks
    and factors?
    \item[$\square$] Are long and short contributions separated where relevant?
    \item[$\square$] Are factor, signal, and implementation-cost contributions
    understood?
\end{itemize}


\subsubsection{Robustness and Statistical Integrity}

\begin{itemize}
    \item[$\square$] Does the strategy survive meaningful subperiod analysis?
    \item[$\square$] Does it survive different market regimes?
    \item[$\square$] Is performance robust across reasonable universes?
    \item[$\square$] Is performance robust to rebalance and holding-period
    assumptions?
    \item[$\square$] Is performance robust to reasonable signal lags?
    \item[$\square$] Is performance robust to outlier-treatment assumptions?
    \item[$\square$] Is performance robust to neutralization choices?
    \item[$\square$] Is performance robust to portfolio-mapping choices?
    \item[$\square$] Is performance robust to higher transaction costs?
    \item[$\square$] Is performance stable under reasonable parameter and
    hyperparameter perturbations?
    \item[$\square$] Do placebo and randomization tests behave as expected?
    \item[$\square$] Are standard errors adjusted for serial dependence when
    required?
    \item[$\square$] Has multiple hypothesis testing been considered?
    \item[$\square$] Has the research history been recorded sufficiently to
    assess data snooping?
    \item[$\square$] Are DSR and/or PBO considered when substantial strategy
    selection has occurred?
\end{itemize}


\subsubsection{Reproducibility}

\begin{itemize}
    \item[$\square$] Is the exact code version recorded?
    \item[$\square$] Is the configuration version recorded?
    \item[$\square$] Is the exact data version recoverable?
    \item[$\square$] Are model artifacts versioned?
    \item[$\square$] Are random states recorded?
    \item[$\square$] Is the software environment reproducible?
    \item[$\square$] Are all relevant experiments tracked?
    \item[$\square$] Are key mathematical invariants covered by unit tests?
    \item[$\square$] Is the complete pipeline covered by integration tests?
\end{itemize}

No individual item proves that a strategy is valid.

The objective of the checklist is instead to accumulate evidence that the
observed historical performance survives progressively stronger attempts to
explain it through data errors, implementation assumptions, statistical noise,
or research selection.


% ------------------------------------------------------------
\subsection{From Backtesting to Live Implementation}
\label{subsec:backtesting_to_live}
% ------------------------------------------------------------

A successful historical backtest is not the end of the quantitative investment
process.

It is the beginning of the transition from a research hypothesis to an
operational investment strategy.

The final transition can be represented as

\begin{equation}
\boxed{
\begin{aligned}
\text{Research}
&\rightarrow
\text{Historical Backtest}
\\
&\rightarrow
\text{Robustness and Statistical Validation}
\\
&\rightarrow
\text{Frozen Strategy Specification}
\\
&\rightarrow
\text{Paper / Shadow Trading}
\\
&\rightarrow
\text{Controlled Live Deployment}
\\
&\rightarrow
\text{Continuous Monitoring}.
\end{aligned}
}
\end{equation}


\subsubsection{Freeze the Strategy Specification}

Before live evaluation begins, the economic specification should be frozen.

This includes:

\begin{itemize}
    \item data sources;
    \item universe definition;
    \item feature definitions;
    \item signal transformations;
    \item model specification;
    \item hyperparameters;
    \item random seed or ensemble rule;
    \item neutralization;
    \item portfolio mapping;
    \item risk model;
    \item constraints;
    \item volatility target;
    \item rebalance schedule;
    \item transaction-cost assumptions.
\end{itemize}

Changes after this point should be treated as new strategy versions rather than
silent modifications of the existing strategy.


\subsubsection{Paper and Shadow Trading}

Before capital is committed, the production system should generate decisions
using live information without necessarily executing them.

At each decision date $t$, the system should record

\begin{equation}
\boxed{
\left(
\mathcal{I}_t,
\mathbf{s}_t,
\mathbf{w}_t^{\mathrm{target}},
\mathcal{O}_t
\right),
}
\end{equation}

where $\mathcal{O}_t$ denotes the orders that would have been submitted.

This phase verifies aspects that cannot be fully tested using static historical
data, including:

\begin{itemize}
    \item actual data arrival times;
    \item production data quality;
    \item computation latency;
    \item operational failures;
    \item portfolio generation;
    \item order generation;
    \item backtest--production parity.
\end{itemize}


\subsubsection{Controlled Live Deployment}

Live deployment should preferably occur progressively.

The initial objective is not to maximize capital allocation, but to verify that
the realized investment process behaves consistently with the validated
research process.

Important quantities to monitor include:

\begin{itemize}
    \item live signal distributions;
    \item portfolio exposures;
    \item realized turnover;
    \item execution prices;
    \item slippage;
    \item transaction costs;
    \item borrow availability;
    \item realized volatility;
    \item realized P\&L.
\end{itemize}


\subsubsection{Backtest--Live Reconciliation}

Differences between live and historical behavior should be decomposed rather
than treated as one unexplained performance gap.

Conceptually,

\begin{equation}
\boxed{
\begin{aligned}
\text{Live}
-
\text{Backtest}
=
&
\text{Data Difference}
+
\text{Signal Difference}
\\
&+
\text{Portfolio Difference}
+
\text{Execution Difference}
\\
&+
\text{Cost Difference}
+
\text{Statistical Variation}.
\end{aligned}
}
\end{equation}

This decomposition helps distinguish an ordinary unfavorable realization from a
structural discrepancy between research assumptions and live implementation.


\subsubsection{Continuous Monitoring}

Deployment does not terminate model validation.

A live strategy should continuously monitor whether the assumptions underlying
the historical research remain valid.

Examples include changes in:

\begin{itemize}
    \item signal distributions;
    \item predictive relationships;
    \item IC and Rank IC;
    \item factor exposures;
    \item realized volatility;
    \item correlations;
    \item turnover;
    \item transaction costs;
    \item liquidity;
    \item capacity;
    \item borrow conditions.
\end{itemize}

However, poor short-term realized performance alone does not necessarily imply
that the strategy is invalid.

Similarly, strong short-term live performance does not validate an otherwise
weak research process.

Live monitoring should therefore distinguish between

\begin{equation}
\boxed{
\text{Normal Statistical Variation}
}
\end{equation}

and

\begin{equation}
\boxed{
\text{Structural Model or Implementation Failure}.
}
\end{equation}


% ------------------------------------------------------------
\subsubsection{Final Perspective}
% ------------------------------------------------------------

The central lesson of the complete backtesting framework is that historical
performance is an output of an entire chain of assumptions.

A reported Sharpe ratio ultimately depends on

\begin{equation}
\boxed{
\begin{aligned}
\text{Data}
&+
\text{Cleaning}
+
\text{Signal}
+
\text{Model}
\\
&+
\text{Neutralization}
+
\text{Risk}
+
\text{Portfolio Construction}
\\
&+
\text{Constraints}
+
\text{Execution}
+
\text{Costs}
+
\text{Validation}.
\end{aligned}
}
\end{equation}

A weakness at any stage can invalidate conclusions drawn from the final
performance statistics.

The objective of a rigorous backtester is therefore not to construct the most
attractive historical equity curve.

It is to construct the most faithful possible simulation of the investment
process that could actually have been executed using the information available
at each point in time.

Ultimately,

\begin{equation}
\boxed{
\begin{aligned}
\text{A Good Backtest}
&\neq
\text{A Good Historical Equity Curve},
\\[0.2cm]
\text{A Good Backtest}
&=
\text{A Credible Counterfactual Investment Experiment}.
\end{aligned}
}
\end{equation}

The transition from research to investment strategy is complete only when the
same economic logic that generated the historical evidence can be reproduced,
implemented, monitored, and challenged in the live environment.

% ============================================================
\appendix
\part{Appendices}
% ============================================================

\section{Mathematical Notation}

\section{Portfolio and Matrix Algebra}

\section{Regression and Neutralization Derivations}

\section{Risk Model Derivations}

\section{Portfolio Optimization Derivations}

\section{Performance Metric Derivations}

\section{Statistical Tests}

\section{Implementation Pseudocode}

\section{Backtesting Checklist}


% ============================================================
\end{document}
% ============================================================